\documentclass{beamer}

\usetheme{metropolis}  
\title{BNU Beamer Template}
\subtitle{A Metropolis-based template, dedicated for BNU}
\date{\today}
\author{Halve Luve}
\institute{School of Artificial Intelligence, Beijing Normal University}
\titlegraphic{
    \flushright\includegraphics[height=1cm]{BNU/BNU_name.pdf}
}
\addbibresource{ref.bib}
\setbeamertemplate{background canvas}{
    \centering
    \begin{tikzpicture}[remember picture,overlay]
        % Warning: compilation of tikz is time-consuming. 
        % \node[opacity=0.2] at (current page.east) {% \documentclass{standalone}
% \usepackage{graphicx} 
% \usepackage{pgfplots}
% \pgfplotsset{compat=1.18}
% \begin{document}
  \definecolor{c005093}{RGB}{0,80,147}
  \def \globalscale {1.2}
  \begin{tikzpicture}[y=1cm, x=1cm, yscale=\globalscale,xscale=\globalscale, every node/.append style={scale=\globalscale}, inner sep=0pt, outer sep=0pt]
    \begin{scope}[cm={ 1.3333333,-0.0,-0.0,-1.3333333,(-1.143, 2.5124065869)}]
      \begin{scope}[scale=154.0]
        % Image image16 not included. Base64 still not supported;
  
  
  
      \end{scope}
      \begin{scope}[shift={(0.3439583333, -0.7672916667)}]
        \path[fill=white,nonzero rule] (0.7408333333, -0.2232202271) -- 
    (2.6458333333, -0.2232202271).. controls (3.1397220458, -0.2232202271) and 
    (3.3866666667, 0.0237243937) .. (3.3866666667, 0.5176131062) -- (3.3866666667,
     2.4226131063).. controls (3.3866666667, 2.9165019246) and (3.1397220458, 
    3.1634464396) .. (2.6458333333, 3.1634464396) -- (0.7408333333, 
    3.1634464396).. controls (0.2469444621, 3.1634464396) and (0.0, 2.9165019246) 
    .. (0.0, 2.4226131063) -- (0.0, 0.5176131062).. controls (0.0, 0.0237243937) 
    and (0.2469444621, -0.2232202271) .. (0.7408333333, -0.2232202271) -- cycle;
  
  
  
      \end{scope}
      \begin{scope}[shift={(0.8731702173, -1.9486754588)}]
        \path[fill=c005093,even odd rule] (0.5631957875, 2.0542641054).. controls 
    (0.4048911342, 2.2125234885) and (0.317697829, 2.4230776881) .. (0.317697829, 
    2.646873675).. controls (0.317697829, 3.1090670769) and (0.6937856806, 
    3.4849283394) .. (1.1558431396, 3.4849283394).. controls (1.6179005721, 
    3.4849283394) and (1.9939882385, 3.1090670769) .. (1.9939882385, 
    2.646873675).. controls (1.9939882385, 2.1848162425) and (1.6179005721, 
    1.8088114435) .. (1.1558431396, 1.8088114435).. controls (0.9319187504, 
    1.8088114435) and (0.7214097946, 1.895959505) .. (0.5631957875, 2.0542641054) 
    -- cycle(0.3570908213, 2.646873675).. controls (0.3570908213, 2.4335321694) 
    and (0.4401824763, 2.2329640327) .. (0.5910391088, 2.0821074002).. controls 
    (0.741805095, 1.9312582554) and (0.9423732317, 1.8480760335) .. (1.1558431396,
     1.8480760335).. controls (1.5962438177, 1.8480760335) and (1.9545122729, 
    2.2064729969) .. (1.9545122729, 2.646873675).. controls (1.9545122729, 
    3.0872743796) and (1.5962438177, 3.4455805908) .. (1.1558431396, 
    3.4455805908).. controls (0.7154423821, 3.4455805908) and (0.3570908213, 
    3.0872743796) .. (0.3570908213, 2.646873675) -- cycle;
  
  
  
      \end{scope}
      \begin{scope}[shift={(0.8731702173, -1.2964583333)}]
        \path[fill=c005093,even odd rule] (0.0, 1.9993628256).. controls (0.0, 
    2.6412532635) and (0.5221856619, 3.1634464396) .. (1.1641666667, 
    3.1634464396).. controls (1.806049511, 3.1634464396) and (2.3282425019, 
    2.6412532635) .. (2.3282425019, 1.9993628256).. controls (2.3282425019, 
    1.3573894408) and (1.806049511, 0.8351131062) .. (1.1641666667, 
    0.8351131062).. controls (0.5221856619, 0.8351131062) and (0.0, 1.3573894408) 
    .. (0.0, 1.9993628256) -- cycle(0.0590403685, 1.9993628256).. controls 
    (0.0590403685, 1.3899839644) and (0.5547802119, 0.8942366069) .. 
    (1.1641666667, 0.8942366069).. controls (1.7734624752, 0.8942366069) and 
    (2.2691190012, 1.3899839644) .. (2.2691190012, 1.9993628256).. controls 
    (2.2691190012, 2.6086587665) and (1.7734624752, 3.1043078048) .. 
    (1.1641666667, 3.1043078048).. controls (0.5547802119, 3.1043078048) and 
    (0.0590403685, 2.6086587665) .. (0.0590403685, 1.9993628256) -- cycle;
  
  
  
      \end{scope}
      \begin{scope}[shift={(0.8731702173, -3.4438328063)}]
        \path[fill=c005093,even odd rule] (0.316519851, 4.787532335) -- 
    (0.2691576087, 4.7155069044).. controls (0.2612261412, 4.7034812215) and 
    (0.2581744106, 4.6923772675) .. (0.2599193377, 4.682005839).. controls 
    (0.2616189152, 4.6716722723) and (0.2668914525, 4.6636577785) .. 
    (0.2756085208, 4.657901821).. controls (0.2831018912, 4.6530673279) and 
    (0.2908596333, 4.6510656756) .. (0.2990479319, 4.6522818065).. controls 
    (0.3045848408, 4.6530673544) and (0.3103332842, 4.6556357177) .. 
    (0.3162629935, 4.6602133269).. controls (0.3124709852, 4.649796681) and 
    (0.3119422419, 4.6407774321) .. (0.3143897435, 4.6331556594).. controls 
    (0.3167842227, 4.6255263196) and (0.3218831404, 4.6190374927) .. 
    (0.3299883333, 4.613810305).. controls (0.3363486256, 4.6096179556) and 
    (0.3431017183, 4.6073594458) .. (0.350247585, 4.606966619).. controls 
    (0.3574388012, 4.6065737127) and (0.364191894, 4.6079715065) .. (0.3706881556,
     4.6110684279).. controls (0.3746085662, 4.6131533446) and (0.3790955027, 
    4.6171642427) .. (0.3842849608, 4.623131761).. controls (0.3910758096, 
    4.6309725294) and (0.3954796404, 4.6363735219) .. (0.3973529962, 4.639130639) 
    -- (0.4409685269, 4.7056566785) -- (0.3165199569, 4.787532335) -- 
    cycle(0.3397930508, 4.7165115538) -- (0.3288551758, 4.6998707438).. controls 
    (0.324844119, 4.6938956583) and (0.3210143281, 4.690322725) .. (0.317524474, 
    4.6893180492).. controls (0.3138608944, 4.6884039138) and (0.3101595588, 
    4.6891434242) .. (0.3062391217, 4.6917199367).. controls (0.3025755421, 
    4.6940239283) and (0.3005284608, 4.6970833848) .. (0.3000525548, 
    4.7009130962).. controls (0.2995313256, 4.7044406271) and (0.3011855006, 
    4.7092826608) .. (0.3051513667, 4.7150688073) -- (0.3162629935, 4.7321026558) 
    -- (0.3397930508, 4.7165115008) -- cycle(0.3887189013, 4.6843476396) -- 
    (0.3758699144, 4.6648285862).. controls (0.3715567033, 4.6581585727) and 
    (0.3672434923, 4.6542458085) .. (0.3631493298, 4.6530674338).. controls 
    (0.3590098706, 4.6518889796) and (0.3549534906, 4.6526745275) .. 
    (0.3509424338, 4.6552505108).. controls (0.3473317179, 4.6576904983) and 
    (0.3451486673, 4.6609158485) .. (0.3447558669, 4.6651383869).. controls 
    (0.3442798815, 4.6693609252) and (0.3462364752, 4.6747318346) .. 
    (0.3506402531, 4.6815226304) -- (0.3635345896, 4.7009132021) -- (0.3887189013,
     4.6843476925) -- cycle;
  
  
  
      \end{scope}
      \begin{scope}[shift={(0.8731702173, -3.4523682646)}]
        \path[fill=c005093,even odd rule] (0.2075404631, 4.6148984304) -- 
    (0.1611375735, 4.5005566527) -- (0.1905897731, 4.4885689375) -- (0.2196114163,
     4.5602089231) -- (0.2415702454, 4.5512349973) -- (0.2146032506, 4.4848675225)
     -- (0.2427486467, 4.4734537944) -- (0.2697685317, 4.539775946) -- 
    (0.2969168454, 4.5287852073) -- (0.2669811992, 4.4551433842) -- (0.298261405, 
    4.4424756898) -- (0.3455405417, 4.5588115527) -- cycle;
  
  
  
      \end{scope}
      \begin{scope}[shift={(0.8731702173, -3.5448724188)}]
        \path[fill=c005093,even odd rule] (0.1418312717, 4.5237774385) -- 
    (0.131029366, 4.4789005069) -- (0.2758278604, 4.4438584287) -- (0.2866751156, 
    4.4887807894) -- cycle;
  
  
  
      \end{scope}
      \begin{scope}[shift={(0.8731702173, -3.5054950104)}]
        \path[fill=c005093,even odd rule] (0.1105510923, 4.3034338938) -- 
    (0.1063209869, 4.257416211) -- (0.18706338, 4.2501042125).. controls 
    (0.2039234237, 4.2484045292) and (0.2169537031, 4.2487897625) .. 
    (0.2260635719, 4.2509728396).. controls (0.2352112231, 4.2531559167) and 
    (0.2432333369, 4.2585567504) .. (0.2501979375, 4.2671001198).. controls 
    (0.2571700788, 4.2755075463) and (0.2612717554, 4.2868834125) .. 
    (0.2626616381, 4.3009939062).. controls (0.2640590879, 4.3159050821) and 
    (0.2630620056, 4.3277947956) .. (0.2598289825, 4.336338165).. controls 
    (0.256565744, 4.34495702) and (0.251202534, 4.3519291083) .. (0.2438451594, 
    4.3571261335).. controls (0.2363895715, 4.362300431) and (0.2270606542, 
    4.3658280412) .. (0.2156922227, 4.3677014235) -- (0.2055852717, 
    4.3244333173).. controls (0.2121646125, 4.3237385215) and (0.2169537031, 
    4.3227791423) .. (0.2200054338, 4.3213891537).. controls (0.22314027, 
    4.3200368683) and (0.225497046, 4.317899379) .. (0.227234406, 4.3152404488).. 
    controls (0.2283297281, 4.3133671987) and (0.2287224756, 4.310798809) .. 
    (0.2285487502, 4.3074901944).. controls (0.2279369013, 4.3023838154) and 
    (0.2257991738, 4.2988637723) .. (0.2216597146, 4.296809124).. controls 
    (0.2176486577, 4.294716799) and (0.2110768575, 4.2941577873) .. (0.2020500944,
     4.294935874) -- (0.1105510923, 4.3034338408) -- cycle;
  
  
  
      \end{scope}
      \begin{scope}[shift={(0.8731702173, -3.5708576771)}]
        \path[fill=c005093,even odd rule] (0.1071066142, 4.2481850044) -- 
    (0.1095464958, 4.2021749152) -- (0.2582200777, 4.21001571) -- (0.2559085983, 
    4.2561089842) -- cycle;
  
  
  
      \end{scope}
      \begin{scope}[shift={(0.8731702173, -3.4535472479)}]
        \path[fill=c005093,even odd rule] (0.1162538156, 4.0058145796) -- 
    (0.1254467635, 3.9638080067) -- (0.2180486019, 3.9268246277) -- (0.1376083365,
     3.90899777) -- (0.1469372802, 3.8667343131) -- (0.2923778654, 3.8988907925) 
    -- (0.2831018119, 3.9411466558) -- (0.1911571192, 3.9780619046) -- 
    (0.2710308587, 3.9955716327) -- (0.2616641325, 4.037978626) -- cycle;
  
  
  
      \end{scope}
      \begin{scope}[shift={(0.8731702173, -3.4633521771)}]
        \path[fill=c005093,even odd rule] (0.2913343488, 3.7691239917) -- 
    (0.2631436031, 3.7560483363) -- (0.292905524, 3.6915163819) -- (0.350639089, 
    3.7181434929).. controls (0.353388639, 3.7343085556) and (0.3540760529, 
    3.747814741) .. (0.352731414, 3.7588432623).. controls (0.3513338848, 
    3.7697358406) and (0.3476779252, 3.7818067673) .. (0.3417104333, 
    3.7947011038).. controls (0.3343454917, 3.8108284369) and (0.3255377771, 
    3.8225896688) .. (0.3153477204, 3.8302567646).. controls (0.3052332023, 
    3.8378861044) and (0.2934719704, 3.841942405) .. (0.2801319698, 
    3.8425014167).. controls (0.2667164835, 3.8430678896) and (0.2529912496, 
    3.8400616938) .. (0.2389109977, 3.833655946).. controls (0.2242264375, 
    3.8269029062) and (0.2127295773, 3.8178836308) .. (0.2045790348, 
    3.8069004327).. controls (0.1965644881, 3.7957057267) and (0.1923796794, 
    3.782985195) .. (0.1920317523, 3.768640545).. controls (0.1919867731, 
    3.757445839) and (0.1953860869, 3.7441058383) .. (0.2026603823, 
    3.7286054883).. controls (0.2095041213, 3.7135281277) and (0.2159927894, 
    3.7031115083) .. (0.2221415738, 3.6968418358).. controls (0.2282374944, 
    3.6907308604) and (0.2353002819, 3.6864251635) .. (0.2431411296, 
    3.684279869).. controls (0.2511179202, 3.6819306335) and (0.2600011469, 
    3.6816284794) .. (0.2699419396, 3.6831014148) -- (0.2585659146, 
    3.7268605083).. controls (0.2526437458, 3.7261580396) and (0.2474543406, 
    3.7271552542) .. (0.2429673777, 3.7297762696).. controls (0.2384351181, 
    3.7325183054) and (0.234733756, 3.7370581056) .. (0.2319011004, 
    3.7431087031).. controls (0.2276709685, 3.7523469742) and (0.2275878365, 
    3.7611017985) .. (0.2313798713, 3.7693807169).. controls (0.2352096092, 
    3.7778787102) and (0.2440173502, 3.7851983552) .. (0.2578709865, 
    3.7916115113).. controls (0.2724649269, 3.7983570633) and (0.2841430267, 
    3.8004041975) .. (0.2929885238, 3.7977527815).. controls (0.3017962119, 
    3.7951392273) and (0.3085040079, 3.7888166646) .. (0.3129909444, 
    3.7790571644).. controls (0.3151287777, 3.7744418256) and (0.3165185546, 
    3.7696451415) .. (0.3171757267, 3.7648107277).. controls (0.3176969558, 
    3.7598403446) and (0.31756096, 3.7540087956) .. (0.3169111433, 3.7469384675) 
    -- (0.3041905852, 3.741061569) -- (0.2913340577, 3.7691239387) -- cycle;
  
  
  
      \end{scope}
      \begin{scope}[shift={(0.8731702173, -3.4192000979)}]
        \path[fill=c005093,even odd rule] (0.3381757323, 3.4664061302) -- 
    (0.3673257512, 3.4348539856) -- (0.4658950815, 3.4494554667) -- (0.4052835112,
     3.3936783102) -- (0.4346904406, 3.3617786619) -- (0.5442806985, 3.4626141219)
     -- (0.5148738485, 3.4945968494) -- (0.4168710442, 3.4801313642) -- 
    (0.4770444644, 3.5355534498) -- (0.4476829375, 3.567370019) -- cycle;
  
  
  
      \end{scope}
      \begin{scope}[shift={(0.8731702173, -3.4606478708)}]
        \path[fill=c005093,even odd rule] (0.53574442, 3.4124272144).. controls 
    (0.521800111, 3.3925985456) and (0.5165276002, 3.3731248419) .. (0.5198361354,
     3.3542554994).. controls (0.5230616179, 3.3353408338) and (0.5347397442, 
    3.3190473423) .. (0.5545608988, 3.3052314621).. controls (0.5748730142, 
    3.2910303485) and (0.594301289, 3.2854934131) .. (0.6129138269, 
    3.2887641923).. controls (0.6315641208, 3.2918990021) and (0.6476838869, 
    3.3033581063) .. (0.6614997671, 3.3229677265).. controls (0.6714707488, 
    3.3373048094) and (0.6772720823, 3.3506372429) .. (0.6788432575, 
    3.3631009965).. controls (0.6804447804, 3.3755647235) and (0.6784429429, 
    3.3874997073) .. (0.6729513042, 3.398921029).. controls (0.6674672325, 
    3.4104632123) and (0.6578362404, 3.4208798581) .. (0.6443300815, 
    3.430291881).. controls (0.6305594979, 3.4398776558) and (0.6176576208, 
    3.4456713694) .. (0.6055866677, 3.4476353715).. controls (0.593432609, 
    3.449637209) and (0.5812332535, 3.4477173923) .. (0.5690716806, 
    3.4422721879).. controls (0.5568269756, 3.4366521733) and (0.5457606983, 
    3.4267642973) .. (0.53574442, 3.4124272144) -- cycle(0.5736039667, 
    3.3862382269).. controls (0.58210196, 3.3985206879) and (0.5906453294, 
    3.4057572008) .. (0.5990073533, 3.4080233571).. controls (0.6073316213, 
    3.4101610581) and (0.6154821108, 3.4085067773) .. (0.6233229321, 
    3.403015165).. controls (0.6314734746, 3.3973045306) and (0.6358773848, 
    3.3902870398) .. (0.6365723394, 3.3819250158).. controls (0.6373578873, 
    3.3736007479) and (0.6329539771, 3.362489121) .. (0.6234135783, 
    3.3487639135).. controls (0.6153914646, 3.3374332117) and (0.6071578165, 
    3.3306348223) .. (0.5988788717, 3.3285424179).. controls (0.5904715246, 
    3.3264953367) and (0.582321035, 3.3282855075) .. (0.5744801873, 
    3.3337771727).. controls (0.5668508475, 3.3390874925) and (0.5626283092, 
    3.3460672273) .. (0.5617596027, 3.3547314054).. controls (0.560928811, 
    3.363320045) and (0.5648491423, 3.3737744998) .. (0.5736039667, 3.3862382269) 
    -- cycle;
  
  
  
      \end{scope}
      \begin{scope}[shift={(0.8731702173, -3.4475269188)}]
        \path[fill=c005093,even odd rule] (0.7466394558, 3.350198405) -- 
    (0.6884677408, 3.2129839273) -- (0.7591031829, 3.183048281).. controls 
    (0.772216515, 3.1774282665) and (0.7827238335, 3.1743311862) .. (0.7905646813,
     3.1736362581).. controls (0.798405529, 3.1729337365) and (0.8057704706, 
    3.1747693098) .. (0.8126142096, 3.1790371977).. controls (0.8193672758, 
    3.1833050856) and (0.8246775956, 3.1898844529) .. (0.8284242808, 
    3.1985184421).. controls (0.8315666577, 3.2061855115) and (0.832745059, 
    3.2135429125) .. (0.8318310294, 3.220386625).. controls (0.8309547294, 
    3.2272756871) and (0.8284242808, 3.2335453067) .. (0.8240733137, 
    3.2394297194).. controls (0.8214520867, 3.2431310815) and (0.8173201946, 
    3.2470514921) .. (0.8116548833, 3.2511078721).. controls (0.8173201946, 
    3.2507150717) and (0.8217618608, 3.2508887971) .. (0.8248135915, 
    3.2515837517).. controls (0.8267320852, 3.2519765521) and (0.8300861023, 
    3.2535477537) .. (0.8346183883, 3.2559423123).. controls (0.8391053512, 
    3.2582537652) and (0.842376104, 3.26030082) .. (0.8442494333, 3.2619173183) --
     (0.881504751, 3.2929785021) -- (0.8336137125, 3.3132830769) -- (0.7932236379,
     3.280915116).. controls (0.7879888831, 3.2766849842) and (0.7839853669, 
    3.2743810719) .. (0.781062065, 3.2735954975).. controls (0.7770132521, 
    3.2727645735) and (0.7730928415, 3.2732026971) .. (0.7693008331, 3.2748569779)
     -- (0.7655541479, 3.2764281531) -- (0.7890766117, 3.3321146898) -- 
    cycle(0.7545708969, 3.2504960231) -- (0.7724356429, 3.242874224).. controls 
    (0.7743919721, 3.2421717287) and (0.7777988, 3.2399509221) .. (0.7828069656, 
    3.2362948831).. controls (0.7854658694, 3.2345046329) and (0.7871201767, 
    3.2322007206) .. (0.7879510742, 3.2290130206).. controls (0.7887744575, 
    3.2260972858) and (0.7885175471, 3.222954909) .. (0.7872030971, 
    3.2199938246).. controls (0.7853750908, 3.2155068881) and (0.7825047056, 
    3.2127648523) .. (0.7787126708, 3.2115864775).. controls (0.7748753658, 
    3.2104080762) and (0.7690816258, 3.2113674025) .. (0.7613239102, 3.2147212608)
     -- (0.7427566954, 3.2225696756) -- (0.7545707646, 3.2504960231) -- cycle;
  
  
  
      \end{scope}
      \begin{scope}[shift={(0.8731702173, -3.4479727417)}]
        \path[fill=c005093,even odd rule] (0.888477051, 3.1364795506) -- 
    (0.9480839983, 3.1254963525) -- (0.987431641, 3.2104235544) -- (0.9939202827, 
    3.1170814383) -- (1.0532250494, 3.1061057808) -- (1.0803281194, 3.2526038527) 
    -- (1.0432464477, 3.259402295) -- (1.0226397452, 3.1476742567) -- 
    (1.014579849, 3.2646295092) -- (0.9811544808, 3.2708613462) -- (0.9319188563, 
    3.164360575) -- (0.9525255323, 3.2761716925) -- (0.9153986433, 3.2829247588) 
    -- cycle;
  
  
  
      \end{scope}
      \begin{scope}[shift={(0.8731702173, -3.471668825)}]
        \path[fill=c005093,even odd rule] (1.223298634, 3.2537669081) -- 
    (1.1709659556, 3.2516291806) -- (1.1625962852, 3.2758995156) -- (1.1157175954,
     3.2740714829) -- (1.1774999998, 3.1274449823) -- (1.2277477615, 3.1294089579)
     -- (1.277723505, 3.2805677446) -- (1.2296967352, 3.2785584194) -- 
    cycle(1.2149742867, 3.2212176548) -- (1.2006825269, 3.166973944) -- 
    (1.1821229056, 3.2198277721) -- cycle;
  
  
  
      \end{scope}
      \begin{scope}[shift={(0.8731702173, -3.4529805104)}]
        \path[fill=c005093,even odd rule] (1.347089926, 3.1242718079) -- 
    (1.391929234, 3.1346431571) -- (1.3663067988, 3.2440974985) -- (1.4362850688, 
    3.2603985835) -- (1.4279229654, 3.2960826467) -- (1.3130601967, 3.269198731) 
    -- cycle;
  
  
  
      \end{scope}
      \begin{scope}[shift={(0.8731702173, -3.4429115271)}]
        \path[fill=c005093,even odd rule] (1.7101021223, 3.2534644894) -- 
    (1.7499256975, 3.2765186671) -- (1.7053961373, 3.3532953).. controls 
    (1.7009998206, 3.3609623958) and (1.6956819867, 3.3674511169) .. 
    (1.6894123935, 3.3729427292).. controls (1.6831352069, 3.3783059127) and 
    (1.6759440171, 3.3821432177) .. (1.6679218504, 3.3842733781).. controls 
    (1.6599072508, 3.3863657825) and (1.6525423621, 3.3868492292) .. 
    (1.6457515398, 3.3856708544).. controls (1.6363320027, 3.3840996527) and 
    (1.6262627813, 3.3800885958) .. (1.615332394, 3.3738189498).. controls 
    (1.6091911767, 3.3701553702) and (1.6026119152, 3.3658421592) .. 
    (1.5955491013, 3.360652701).. controls (1.5886675533, 3.3554254869) and 
    (1.5834781217, 3.3501529231) .. (1.5800335642, 3.3447142275).. controls 
    (1.576513521, 3.3391772921) and (1.5741566921, 3.3329001054) .. (1.5727591894,
     3.3255351373).. controls (1.5714145769, 3.3183892706) and (1.5714148415, 
    3.3118099298) .. (1.5728888352, 3.3058424379).. controls (1.5751623998, 
    3.2962113665) and (1.5780782933, 3.2881892527) .. (1.5818249256, 3.2819196067)
     -- (1.62635438, 3.2049692219) -- (1.6660495265, 3.2279704829) -- 
    (1.6204322642, 3.3067111179).. controls (1.616466504, 3.3137739054) and 
    (1.6152428325, 3.3203532462) .. (1.6168517638, 3.3266228923).. controls 
    (1.6185136117, 3.3329001054) and (1.6226983675, 3.3379158117) .. 
    (1.6296705352, 3.3419193279).. controls (1.6365520831, 3.3458397385) and 
    (1.6430029688, 3.3470181398) .. (1.64910643, 3.3454922612).. controls 
    (1.6552400802, 3.3437095781) and (1.660474835, 3.3393888) .. (1.6646142942, 
    3.3321145046) -- (1.7101031806, 3.2534644894) -- cycle;
  
  
  
      \end{scope}
      \begin{scope}[shift={(0.8731702173, -3.4188151292)}]
        \path[fill=c005093,even odd rule] (1.817125234, 3.3015672475) -- 
    (1.8488888865, 3.3304604619) -- (1.8349447362, 3.4291582208) -- (1.8904119069,
     3.3682369558) -- (1.922439984, 3.3973870012) -- (1.8223146921, 3.5075438115) 
    -- (1.7901505398, 3.478393766) -- (1.8039665258, 3.3803909617) -- (1.7489297, 
    3.4409648287) -- (1.7169922692, 3.4117694073) -- cycle;
  
  
  
      \end{scope}
      \begin{scope}[shift={(0.8731702173, -3.496332225)}]
        \path[fill=c005093,even odd rule] (1.972942804, 3.5346469344) -- 
    (2.0015261383, 3.5709428729) -- (1.8846994469, 3.6631065612) -- (1.8560327423,
     3.6270221306) -- cycle;
  
  
  
      \end{scope}
      \begin{scope}[shift={(0.8731702173, -3.4592881771)}]
        \path[fill=c005093,even odd rule] (2.0401338738, 3.5833087835) -- 
    (2.0636185815, 3.6254135963) -- (1.9864945506, 3.7070171552) -- (2.0961678877,
     3.6834493154) -- (2.1190484194, 3.7241415971) -- (1.9620051671, 3.7488122467)
     -- (1.9376895883, 3.7052344721) -- cycle;
  
  
  
      \end{scope}
      \begin{scope}[shift={(0.8731702173, -3.4589862875)}]
        \path[fill=c005093,even odd rule] (2.1430555469, 3.7844508015) -- 
    (2.181481831, 3.9016175354) -- (2.1512892744, 3.9115582752) -- (2.1271549617, 
    3.8382261998) -- (2.1047049865, 3.8455835215) -- (2.1270641302, 3.9135600333) 
    -- (2.0982163713, 3.9230627554) -- (2.0757288517, 3.8550862435) -- 
    (2.0478930181, 3.8642338683) -- (2.0728204194, 3.9396205921) -- (2.040747125, 
    3.9502563658) -- (2.0015279375, 3.8309518779) -- cycle;
  
  
  
      \end{scope}
      \begin{scope}[shift={(0.8731702173, -3.4613802375)}]
        \path[fill=c005093,even odd rule] (2.0506790806, 3.998139414) -- 
    (2.198567141, 3.9800934283) -- (2.2077601154, 4.0564016694).. controls 
    (2.2094672071, 4.070429084) and (2.2095502863, 4.0813217152) .. (2.2081151863,
     4.0891700771).. controls (2.2064985821, 4.0967918498) and (2.2027973523, 
    4.1034920258) .. (2.1968299398, 4.1087796708).. controls (2.1908170983, 
    4.1141730963) and (2.1830517098, 4.1173986317) .. (2.1736397398, 
    4.118614736).. controls (2.1654060123, 4.1195288715) and (2.1581769606, 
    4.1186597152) .. (2.1518544773, 4.1158273506).. controls (2.1453658621, 
    4.1130853677) and (2.1400026256, 4.1089912052) .. (2.1355988213, 
    4.1032805708).. controls (2.1328567325, 4.0995792352) and (2.1301146438, 
    4.0944804233) .. (2.1278485669, 4.0879010825).. controls (2.1266701127, 
    4.0936117169) and (2.1252799919, 4.0977889321) .. (2.1238373777, 
    4.1005384821).. controls (2.1228401631, 4.1022834092) and (2.1204836517, 
    4.1050707681) .. (2.1169938769, 4.1087796708).. controls (2.1134284048, 
    4.1124356833) and (2.1105052617, 4.1148681038) .. (2.1084959365, 4.1162277446)
     -- (2.0683247781, 4.1435044077) -- (2.062009756, 4.09191206) -- 
    (2.1041902396, 4.0621048425).. controls (2.1096742054, 4.0581844054) and 
    (2.1132017098, 4.0549967054) .. (2.1146898323, 4.0523906919).. controls 
    (2.1165630823, 4.0487724354) and (2.1172584073, 4.044935051) .. (2.1167824219,
     4.0407125127) -- (2.1161704406, 4.0367015352) -- (2.056253825, 4.0440664504) 
    -- (2.0506790277, 3.998139414) -- cycle(2.1440967617, 4.0333400569) -- 
    (2.1465819929, 4.0526927667).. controls (2.1467645554, 4.054611525) and 
    (2.1479337492, 4.0585243421) .. (2.1500263388, 4.0644541044).. controls 
    (2.1509854533, 4.0674227823) and (2.1527306715, 4.069643536) .. (2.1555179775,
     4.0713506806).. controls (2.1582600663, 4.0728691244) and (2.1612738556, 
    4.073480841) .. (2.1644087448, 4.0731786869).. controls (2.1692431056, 
    4.0724762181) and (2.1726876631, 4.0706479473) .. (2.1748253906, 
    4.0672866542).. controls (2.1772198698, 4.063978066) and (2.1776580992, 
    4.0581844583) .. (2.1766989052, 4.0497317087) -- (2.174213674, 4.0297293146) 
    -- (2.1440968146, 4.033340004) -- cycle;
  
  
  
      \end{scope}
      \begin{scope}[shift={(0.8731702173, -3.4895414292)}]
        \path[fill=c005093,even odd rule] (2.1167843533, 4.2033464373) -- 
    (2.1161726367, 4.2471735023).. controls (2.10902677, 4.2475664085) and 
    (2.1033990825, 4.2491375044) .. (2.0996222612, 4.2517964081).. controls 
    (2.093254349, 4.2560189465) and (2.089590981, 4.2624774785) .. (2.0890243494, 
    4.2713153556).. controls (2.0884123681, 4.2778041031) and (2.0895910869, 
    4.2831220165) .. (2.0924689069, 4.2868309192).. controls (2.0953016154, 
    4.29066071) and (2.0987007704, 4.2927152525) .. (2.1027571504, 4.2929721365)..
     controls (2.1065794006, 4.2933272073) and (2.1101597952, 4.2917860094) .. 
    (2.1134305744, 4.2887041956).. controls (2.1167844856, 4.2856524385) and 
    (2.1203122017, 4.2779778021) .. (2.1238397325, 4.2658237433).. controls 
    (2.1298977913, 4.2460403448) and (2.136960711, 4.2320507656) .. (2.1452396294,
     4.2241646477).. controls (2.1532541231, 4.2160594283) and (2.1632779421, 
    4.2124940621) .. (2.1748276131, 4.2133627685).. controls (2.1824872213, 
    4.2139292415) and (2.189557814, 4.2166712773) .. (2.1960464292, 
    4.2216794429).. controls (2.2024973148, 4.2266120171) and (2.2072864583, 
    4.2336220467) .. (2.2104288352, 4.2428603177).. controls (2.2136921267, 
    4.252015536) and (2.2146515852, 4.2642602146) .. (2.2134731046, 
    4.2797227821).. controls (2.2120755754, 4.29876585) and (2.2073696433, 
    4.3128914779) .. (2.199619389, 4.3223034479).. controls (2.1918162444, 
    4.3318363325) and (2.180145844, 4.3367992015) .. (2.1646227398, 4.3376376396) 
    -- (2.1654082877, 4.2941958079).. controls (2.172078169, 4.2934480954) and 
    (2.1772224098, 4.2915369042) .. (2.1805308921, 4.2880092675).. controls 
    (2.1837564275, 4.284647895) and (2.1857203502, 4.2795489773) .. (2.1862340917,
     4.2730980123).. controls (2.1866344063, 4.2677876396) and (2.1858112875, 
    4.2636861217) .. (2.1836734542, 4.2608156571).. controls (2.1817094521, 
    4.2579829485) and (2.1789218285, 4.2564117733) .. (2.1757417221, 
    4.2561927512).. controls (2.1734755658, 4.2560181262) and (2.1712020542, 
    4.2568875471) .. (2.1691172433, 4.2589876244).. controls (2.16697941, 
    4.2608987098) and (2.1647962006, 4.2659975481) .. (2.1624923677, 
    4.2739743917).. controls (2.1565701725, 4.293667144) and (2.1510787454, 
    4.3077849931) .. (2.1461083094, 4.3160639379).. controls (2.1410094975, 
    4.3244787992) and (2.135042085, 4.3303556977) .. (2.12815297, 4.333890875).. 
    controls (2.1213167452, 4.3372448127) and (2.1138157813, 4.3385969658) .. 
    (2.105801314, 4.33803044).. controls (2.0962532688, 4.3373356442) and 
    (2.0875439792, 4.3341024358) .. (2.0800052856, 4.3282255637).. controls 
    (2.0723381898, 4.3223033421) and (2.0668540388, 4.3144928421) .. 
    (2.0634094813, 4.3046577769).. controls (2.0598742775, 4.2948454129) and 
    (2.0587034963, 4.2829104821) .. (2.0598742775, 4.2683921067).. controls 
    (2.0618382796, 4.2430718785) and (2.0679341738, 4.226000274) .. (2.0783962485,
     4.2169734579).. controls (2.0888960529, 4.2080524223) and (2.1016619342, 
    4.2034748131) .. (2.116784565, 4.2033464373) -- cycle;
  
  
  
      \end{scope}
      \begin{scope}[shift={(0.8731702173, -3.5462329062)}]
        \path[fill=c005093,even odd rule] (2.1922961456, 4.4696630296) -- 
    (2.1819247171, 4.5144569085) -- (2.0367335281, 4.4808199002) -- (2.047104745, 
    4.4358976454) -- cycle;
  
  
  
      \end{scope}
      \begin{scope}[shift={(0.8731702173, -3.4812086417)}]
        \path[fill=c005093,even odd rule] (2.1655796319, 4.51379844) -- 
    (2.1138135588, 4.6437387281) -- (2.0796552419, 4.6302249492) -- (2.096953515, 
    4.5866471746) -- (1.9927718167, 4.5450786458) -- (2.0097602363, 4.5022939863) 
    -- (2.1139950365, 4.5438096248) -- (2.1313837442, 4.5001562852) -- cycle;
  
  
  
      \end{scope}
      \begin{scope}[shift={(0.8731702173, -3.4566746229)}]
        \path[fill=c005093,even odd rule] (2.0892866308, 4.6693544958) -- 
    (2.0616096796, 4.7123128544) -- (2.003052599, 4.7103036085) -- (2.0290678346, 
    4.7628175529) -- (2.0015267204, 4.8055719442) -- (1.9600489173, 4.7099259158) 
    -- (1.9076256192, 4.6761982877) -- (1.9327720425, 4.6374549798) -- 
    (1.9851424769, 4.6711373112) -- (2.0892866308, 4.6693545488) -- cycle;
  
  
  
      \end{scope}
      \begin{scope}[shift={(0.8731702173, -3.4655733542)}]
        \path[fill=c005093,even odd rule] (0.9512183848, 5.0405175521).. controls 
    (0.9437249879, 5.091770069) and (0.9362317233, 5.1430149129) .. (0.9287383265,
     5.1942068931).. controls (0.9130037673, 5.1918954931) and (0.8972315579, 
    5.1893725323) .. (0.8815498625, 5.1866229823).. controls (0.887426761, 
    5.1532050754) and (0.8933110944, 5.1197869833) .. (0.8991879929, 
    5.0863161862).. controls (0.8911356373, 5.0903271637) and (0.8834610008, 
    5.0932050896) .. (0.8760129271, 5.0952144148).. controls (0.8686857679, 
    5.0970876648) and (0.8596287629, 5.098515674) .. (0.8489930156, 5.099308789)..
     controls (0.851342251, 5.0879781135) and (0.8536612446, 5.0766547142) .. 
    (0.8560105065, 5.0653240652).. controls (0.8715636937, 5.063707461) and 
    (0.8839822035, 5.0603461679) .. (0.8933110944, 5.0553834048).. controls 
    (0.9025040687, 5.050549044) and (0.9101711381, 5.043659929) .. (0.9162746258, 
    5.0349881573).. controls (0.9279074554, 5.0369974031) and (0.9395855287, 
    5.0388255417) .. (0.9512183848, 5.0405175521) -- cycle;
  
  
  
      \end{scope}
      \begin{scope}[shift={(0.8731702173, -3.4649153354)}]
        \path[fill=c005093,even odd rule] (1.0089134, 5.1688024746).. controls 
    (1.0242928883, 5.1685458288) and (1.0396723767, 5.1680544975) .. 
    (1.0549688123, 5.1674049454).. controls (1.0558374394, 5.1734252481) and 
    (1.0575371227, 5.177693189) .. (1.0602867521, 5.1803446315).. controls 
    (1.0628550625, 5.1828751065) and (1.0663373231, 5.1844464404) .. 
    (1.0705674285, 5.1847032185).. controls (1.0779323173, 5.1851413685) and 
    (1.0839451587, 5.1819159125) .. (1.088386825, 5.1752987098).. controls 
    (1.0916575777, 5.1703283002) and (1.0944902863, 5.1597378496) .. 
    (1.0966279873, 5.1434444375).. controls (1.0910079727, 5.1482787983) and 
    (1.0852520152, 5.151934705) .. (1.07932982, 5.1539440302).. controls 
    (1.0734907833, 5.1560289469) and (1.0668207962, 5.156905141) .. (1.0593273994,
     5.1563840177).. controls (1.0445974896, 5.1553793948) and (1.0325266158, 
    5.1494949027) .. (1.0235073669, 5.138866696).. controls (1.0142690694, 
    5.1283216744) and (1.010696136, 5.1155557931) .. (1.0121766915, 
    5.1007353429).. controls (1.0132265581, 5.0906737944) and (1.0166258983, 
    5.0817828154) .. (1.022283669, 5.073896724).. controls (1.0279112506, 
    5.0660483621) and (1.0351855196, 5.0603453742) .. (1.0439856671, 
    5.0569462456).. controls (1.0527857881, 5.0533807735) and (1.0635046146, 
    5.0520210269) .. (1.0763687092, 5.0528066806).. controls (1.0915669315, 
    5.053720816) and (1.10381161, 5.0568176581) .. (1.1128837756, 5.0624376727).. 
    controls (1.1220238333, 5.0680576873) and (1.1293509925, 5.0765104104) .. 
    (1.134751985, 5.0881507277).. controls (1.1401529775, 5.099647535) and 
    (1.1428570985, 5.1148535096) .. (1.1425927533, 5.1335490738).. controls 
    (1.1421622762, 5.1610901879) and (1.1355375065, 5.1812585287) .. 
    (1.1224618775, 5.1937299288).. controls (1.1093937892, 5.2061936294) and 
    (1.0911741046, 5.2118968025) .. (1.0683012458, 5.2105447552).. controls 
    (1.0549234892, 5.2097213719) and (1.0443784675, 5.2075382154) .. 
    (1.0368850971, 5.2041843042).. controls (1.0292633244, 5.2008229052) and 
    (1.0231598631, 5.1960339733) .. (1.018582254, 5.1902023979).. controls 
    (1.014057535, 5.1844011173) and (1.0107414592, 5.1772552242) .. (1.0089134, 
    5.1688024746) -- cycle(1.0966279873, 5.106929424).. controls (1.0970132206, 
    5.0986957229) and (1.0951021352, 5.0920258417) .. (1.0911364544, 
    5.0871914544).. controls (1.0870800744, 5.0821758275) and (1.0820264798, 
    5.0796527608) .. (1.0759229921, 5.07929769).. controls (1.0701292785, 
    5.0789502921) and (1.0652948119, 5.0807404629) .. (1.0612006758, 
    5.0846607148).. controls (1.0570612167, 5.08870966) and (1.0547043877, 
    5.094767904) .. (1.0540547827, 5.1032658708).. controls (1.0534428015, 
    5.1118092137) and (1.0549236744, 5.1185168775) .. (1.0586702273, 
    5.123306021).. controls (1.06237913, 5.1281026523) and (1.0673872956, 
    5.1306711215) .. (1.0735361065, 5.1310714096).. controls (1.0797150533, 
    5.1314188075) and (1.0851688567, 5.1295000492) .. (1.0895198502, 
    5.1253605635).. controls (1.0938406019, 5.121085135) and (1.0962729165, 
    5.1149893467) .. (1.0966279873, 5.106929424) -- cycle;
  
  
  
      \end{scope}
      \begin{scope}[shift={(0.8731702173, -3.4656937396)}]
        \path[fill=c005093,even odd rule] (1.1839502504, 5.1345466588).. controls 
    (1.1836936046, 5.105479666) and (1.189003845, 5.085227955) .. (1.1993298446, 
    5.0734969915).. controls (1.2097918135, 5.061516764) and (1.2255188056, 
    5.0554207375) .. (1.2469564585, 5.0542877123).. controls (1.2571918119, 
    5.0537135665) and (1.2657352871, 5.0543326915) .. (1.2726242963, 
    5.0565009519).. controls (1.2794150921, 5.0585935415) and (1.285080324, 
    5.0613884146) .. (1.28961261, 5.06518037).. controls (1.2940618433, 
    5.0687458421) and (1.2978539046, 5.0727568196) .. (1.3006411842, 
    5.0770322481).. controls (1.3034284902, 5.0813454062) and (1.3059062338, 
    5.0863988156) .. (1.3080439348, 5.092147206).. controls (1.3119265365, 
    5.1031756744) and (1.3145401435, 5.1149823352) .. (1.3156731952, 
    5.127272231).. controls (1.318332099, 5.1548056723) and (1.3153633681, 
    5.1754578831) .. (1.3061327171, 5.1891529546).. controls (1.2969397692, 
    5.2028251925) and (1.2794604152, 5.210673766) .. (1.2540117054, 
    5.2120183521).. controls (1.2397275127, 5.2128115729) and (1.2281324127, 
    5.2111947042) .. (1.2192869421, 5.2074028017).. controls (1.2103583394, 
    5.2035202) and (1.2030765033, 5.1976809781) .. (1.1973205194, 5.1898855065).. 
    controls (1.1932717065, 5.1844015406) and (1.1900009538, 5.1766817663) .. 
    (1.1876063688, 5.1669221867).. controls (1.1853402125, 5.1573440848) and 
    (1.1840786527, 5.1464817483) .. (1.1839502504, 5.1345467117) -- 
    cycle(1.2285251338, 5.1336325233).. controls (1.2292654379, 5.1530156337) and 
    (1.2316600229, 5.1662649881) .. (1.2356635392, 5.173199294).. controls 
    (1.2397275127, 5.1802091383) and (1.245256881, 5.1836082669) .. (1.2524027742,
     5.18321544).. controls (1.2569803569, 5.1829508567) and (1.2609838731, 
    5.1812136025) .. (1.2642168698, 5.1780334696).. controls (1.2674347588, 
    5.1746795583) and (1.2697008356, 5.1698377098) .. (1.270962369, 
    5.1630921312).. controls (1.2720122356, 5.1565051969) and (1.272178394, 
    5.146315246) .. (1.2713097669, 5.132582445).. controls (1.269957746, 
    5.112330734) and (1.2673064623, 5.0988623575) .. (1.2633784581, 
    5.0921470473).. controls (1.2592842956, 5.085393981) and (1.2538908702, 
    5.0821763302) .. (1.2468280563, 5.0825614577).. controls (1.2398105654, 
    5.082954364) and (1.2347496154, 5.0867008375) .. (1.2317960717, 
    5.0938467042).. controls (1.2288727169, 5.1012570219) and (1.2277773154, 
    5.1144609738) .. (1.2285251338, 5.1336325233) -- cycle;
  
  
  
      \end{scope}
      \begin{scope}[shift={(0.8731702173, -3.4586388896)}]
        \path[fill=c005093,even odd rule] (1.502477019, 5.1703672998).. controls 
    (1.456111859, 5.1799530746) and (1.4094369513, 5.1875369854) .. (1.3626034523,
     5.1931192704).. controls (1.3627780773, 5.180481765) and (1.3661764915, 
    5.1680634669) .. (1.3728011819, 5.15607554).. controls (1.379335226, 
    5.1441330421) and (1.3924939342, 5.1291916773) .. (1.4115823252, 
    5.1108964012).. controls (1.423124535, 5.0997772867) and (1.43014192, 
    5.0914152098) .. (1.4328461733, 5.0861425931).. controls (1.4355881563, 
    5.0807416006) and (1.4365021594, 5.0759449694) .. (1.4357165321, 
    5.0717147846).. controls (1.4349309842, 5.0672655513) and (1.4325363992, 
    5.0635189983) .. (1.4284875863, 5.0608223385).. controls (1.4244765029, 
    5.0583447802) and (1.4198084592, 5.0572946225) .. (1.4144981129, 
    5.0581632231).. controls (1.4090517973, 5.0590773585) and (1.4048669092, 
    5.061570051) .. (1.4019057983, 5.0653922748).. controls (1.398944714, 
    5.0694561425) and (1.3975020204, 5.0755973598) .. (1.3975473171, 
    5.084306861).. controls (1.3827721106, 5.0853567277) and (1.3680497679, 
    5.0862256194) .. (1.3532745615, 5.0868376006).. controls (1.353584124, 
    5.0748496738) and (1.3555407177, 5.0651357083) .. (1.3592496469, 
    5.0579444127).. controls (1.3627772571, 5.05076079) and (1.3684427006, 
    5.0449140804) .. (1.3760644733, 5.0400797196).. controls (1.3837315427, 
    5.0354189519) and (1.3946239623, 5.0321104431) .. (1.4088779925, 5.02975372)..
     controls (1.423690955, 5.0273592408) and (1.4353692135, 5.0269662023) .. 
    (1.444305304, 5.0286657533).. controls (1.45318861, 5.0304183533) and 
    (1.460681901, 5.0343311704) .. (1.4669213317, 5.0402986094).. controls 
    (1.473062549, 5.0462736948) and (1.4770282033, 5.0534573175) .. (1.4788109394,
     5.0620459042).. controls (1.480684454, 5.0711028033) and (1.4796795665, 
    5.0805602023) .. (1.4757667758, 5.089889199).. controls (1.4717934485, 
    5.0993465979) and (1.4631292704, 5.1104505783) .. (1.4497892698, 
    5.1232239208).. controls (1.4418578552, 5.1306342385) and (1.436328381, 
    5.1357253775) .. (1.433714774, 5.1385581125).. controls (1.4310105471, 
    5.1413453921) and (1.4278304142, 5.1450921567) .. (1.424302804, 
    5.1495413635).. controls (1.4480896923, 5.1457039792) and (1.4717935279, 
    5.1413456037) .. (1.4954595281, 5.1364280579).. controls (1.4977709281, 
    5.1477512721) and (1.50012019, 5.1590365185) .. (1.502477019, 5.1703671675) --
     cycle;
  
  
  
      \end{scope}
      \begin{scope}[shift={(0.8731702173, -3.5322052271)}]
        \path[fill=c005093,even odd rule] (1.3181497217, 4.0392482292) -- 
    (1.2864768213, 4.0449511906).. controls (1.2781148238, 4.046212724) and 
    (1.2764528965, 4.0513191029) .. (1.2768079673, 4.0539780331) -- (1.2766730298,
     4.0705283027).. controls (1.2765830715, 4.0755440354) and (1.2776776527, 
    4.0787695973) .. (1.2801174815, 4.0803407725).. controls (1.2818624085, 
    4.0815946329) and (1.2842117233, 4.0813453954) .. (1.2854278277, 4.080809085) 
    -- (1.3121003412, 4.0750606946).. controls (1.3183775279, 4.0727114592) and 
    (1.319722114, 4.0686097296) .. (1.3195105531, 4.0664795692) -- (1.3193359281, 
    4.0574149971).. controls (1.3193359281, 4.0544085896) and (1.3180669865, 
    4.0527090121) .. (1.3169792844, 4.0518780617).. controls (1.3142373015, 
    4.0501405429) and (1.3107021242, 4.0511001867) .. (1.3103092973, 4.0513189971)
     -- (1.2908809696, 4.0550203327).. controls (1.2896120279, 4.0552399369) and 
    (1.2887431362, 4.0563194369) .. (1.2887431362, 4.0578151794) -- (1.2887431362,
     4.0666909183).. controls (1.2887431362, 4.0676050538) and (1.2891360425, 
    4.0683527663) .. (1.2897855946, 4.0689644829).. controls (1.2905258987, 
    4.0694780392) and (1.2912738758, 4.0696595433) .. (1.2921877467, 4.0694783037)
     -- (1.3050367071, 4.0662225794).. controls (1.3064720717, 4.0658222648) and 
    (1.30735578, 4.0642585773) .. (1.3069554654, 4.0627780219).. controls 
    (1.3065625592, 4.0612900052) and (1.3050367071, 4.060330626) .. (1.3036392837,
     4.0607309406).. controls (1.3036392837, 4.0607309406) and (1.2979437571, 
    4.0621663052) .. (1.2943179335, 4.063080176).. controls (1.2943179335, 
    4.0616448115) and (1.2943179335, 4.0612521698) .. (1.2943179335, 
    4.0599452869).. controls (1.2978984604, 4.0593333056) and (1.3117972875, 
    4.0567198308) .. (1.3117972875, 4.0567198308).. controls (1.3126659146, 
    4.0564102683) and (1.3134440542, 4.0564102683) .. (1.3136708021, 
    4.0565848933).. controls (1.3136708021, 4.0567198308) and (1.3137528229, 
    4.0574611933) .. (1.3137528229, 4.0574611933) -- (1.3140549771, 4.066827946)..
     controls (1.3140549771, 4.066827946) and (1.3138433104, 4.0683537981) .. 
    (1.3106104196, 4.0695246058) -- (1.2836433983, 4.0755449085).. controls 
    (1.2831674129, 4.075727471) and (1.2822458692, 4.0748424398) .. (1.2822458692,
     4.0705291758) -- (1.2822458692, 4.0536692379).. controls (1.2822458692, 
    4.0527098588) and (1.2833412442, 4.0511009275) .. (1.287344681, 4.0504513754) 
    -- (1.319508701, 4.044649989).. controls (1.3199016073, 4.0445679681) and 
    (1.3212083844, 4.0442570827) .. (1.3220770115, 4.0448695931).. controls 
    (1.3229456385, 4.045436066) and (1.3234668677, 4.0474000681) .. (1.3234668677,
     4.0499230819) -- (1.3234668677, 4.0728035342).. controls (1.3234668677, 
    4.0728934925) and (1.3234668677, 4.0728934925) .. (1.3234668677, 
    4.0728934925).. controls (1.3234668677, 4.0738451987) and (1.3230816344, 
    4.0762398631) .. (1.3181490338, 4.0777657152).. controls (1.3181490338, 
    4.0777657152) and (1.2850786306, 4.0852591121) .. (1.284421485, 4.085432811)..
     controls (1.2833411913, 4.085432811) and (1.2809014154, 4.0855227694) .. 
    (1.2789374133, 4.0874346485).. controls (1.2777213883, 4.0885752673) and 
    (1.2762784038, 4.090576999) .. (1.2764521821, 4.0944596006) -- (1.2764521821, 
    4.1221215235).. controls (1.2763701613, 4.122256461) and (1.2761955363, 
    4.1258228592) .. (1.2789374133, 4.1282250113).. controls (1.2809014154, 
    4.1299699383) and (1.2836432925, 4.1305742467) .. (1.2871331731, 4.1297508633)
     -- (1.3206796675, 4.1221215235).. controls (1.3231648988, 4.1220315652) and 
    (1.3287469456, 4.1199461194) .. (1.329132179, 4.1128833583) -- (1.329132179, 
    4.0948449663).. controls (1.329132179, 4.0944144892) and (1.3289125748, 
    4.0916269979) .. (1.3263903019, 4.0895799167).. controls (1.3242072248, 
    4.0877897458) and (1.3211177117, 4.0873060875) .. (1.3170688723, 4.0879633125)
     -- (1.2935388415, 4.0944142775).. controls (1.2912726852, 4.0944592567) and 
    (1.2864761333, 4.0965822733) .. (1.2864761333, 4.1025195233) -- (1.2864761333,
     4.1110552992).. controls (1.2865660917, 4.1117048512) and (1.2869518542, 
    4.1144923954) .. (1.2895277581, 4.116237296).. controls (1.2914388435, 
    4.1176348252) and (1.2938409163, 4.1178160648) .. (1.296756704, 4.1173249981) 
    -- (1.3145383444, 4.1127095271).. controls (1.3157998777, 4.1124978604) and 
    (1.3166759131, 4.1111458396) .. (1.3165854256, 4.1098768185) -- (1.3157998777,
     4.101333555).. controls (1.3157998777, 4.1005554154) and (1.315316484, 
    4.0997698675) .. (1.3147121756, 4.0991581508).. controls (1.3139266277, 
    4.0987652446) and (1.3130580006, 4.0986821654) .. (1.3122724527, 4.0987652446)
     -- (1.3042503125, 4.1011218883).. controls (1.3028604562, 4.1015147946) and 
    (1.3018558333, 4.1030858904) .. (1.3023770625, 4.1045664458).. controls 
    (1.3027244604, 4.1060393812) and (1.3043410646, 4.1069080083) .. 
    (1.3057837581, 4.1065304479).. controls (1.3057837581, 4.1065304479) and 
    (1.3082237456, 4.1057372271) .. (1.310489849, 4.1051707542).. controls 
    (1.3107015156, 4.1065304479) and (1.3107015156, 4.1068326021) .. 
    (1.3107015156, 4.1080034627).. controls (1.3073098483, 4.1088720898) and 
    (1.2954957262, 4.1118408471) .. (1.2954957262, 4.1118408471).. controls 
    (1.292836796, 4.1124528283) and (1.2921871117, 4.1113196179) .. (1.2920511688,
     4.11053407) -- (1.2920511688, 4.1025194969).. controls (1.2920511688, 
    4.1005554948) and (1.2936222646, 4.0999890219) .. (1.2944985646, 4.0997699469)
     -- (1.3183760198, 4.0935002479).. controls (1.3211180027, 4.0930168542) and 
    (1.3222965627, 4.0935003273) .. (1.3228555215, 4.0938854813).. controls 
    (1.3234672381, 4.0944141188) and (1.3234672381, 4.0948448604) .. 
    (1.3234672381, 4.0949279396) -- (1.3234672381, 4.1128454435).. controls 
    (1.3234672381, 4.11602555) and (1.3210272506, 4.1164109156) .. (1.3199018454, 
    4.1165468585) -- (1.2860836769, 4.1244782731).. controls (1.284421829, 
    4.1246899398) and (1.2832509419, 4.124560294) .. (1.2826391194, 
    4.1241308752).. controls (1.2819895673, 4.1235191585) and (1.282034811, 
    4.1222576252) .. (1.282034811, 4.1222576252) -- (1.282034811, 4.0942859281).. 
    controls (1.282034811, 4.0934096281) and (1.2820349169, 4.0920954427) .. 
    (1.2826391194, 4.0915363781).. controls (1.2833415881, 4.0909622323) and 
    (1.2844218819, 4.0911434719) .. (1.2844218819, 4.0911434719).. controls 
    (1.2846938735, 4.0911434719) and (1.2849957631, 4.0911433925) .. 
    (1.2854265048, 4.0910085344) -- (1.319720791, 4.0832583858).. controls 
    (1.3304396175, 4.0795494831) and (1.3291326552, 4.0725396652) .. 
    (1.3291326552, 4.0721921615) -- (1.3291326552, 4.0499235581).. controls 
    (1.3291326552, 4.045345949) and (1.3278713865, 4.0420827633) .. (1.3254314519,
     4.0403378892).. controls (1.3221682663, 4.0379886538) and (1.318549904, 
    4.0391671079) .. (1.31814951, 4.0392501871) -- cycle;
  
  
  
      \end{scope}
      \begin{scope}[shift={(0.8731702173, -3.5324161)}]
        \path[fill=c005093,even odd rule] (1.3797749635, 4.0252576446) -- 
    (1.3480868231, 4.0309229558).. controls (1.3397248256, 4.0322146517) and 
    (1.3379345225, 4.0374116769) .. (1.3383273494, 4.0399799608) -- (1.3381606619,
     4.0564093423).. controls (1.3381606619, 4.061462831) and (1.3392409556, 
    4.0647714192) .. (1.341726134, 4.0663426208).. controls (1.343471061, 
    4.0676041542) and (1.3456993554, 4.0674379958) .. (1.3470062119, 4.0668109333)
     -- (1.3737616988, 4.0609416283).. controls (1.3800312919, 4.0586680637) and 
    (1.3813003923, 4.0546190127) .. (1.3812550692, 4.0525342019) -- (1.3808999983,
     4.0433790894).. controls (1.3808999983, 4.0403349258) and (1.3795929567, 
    4.0387561835) .. (1.3784676838, 4.0379705298).. controls (1.3758011865, 
    4.0361425235) and (1.3724020579, 4.0371019027) .. (1.372009231, 4.0373209777) 
    -- (1.3525807975, 4.040893911).. controls (1.3513192642, 4.0411135152) and 
    (1.3503601496, 4.0422911756) .. (1.3503601496, 4.0435452477) -- (1.3503601496,
     4.0527003602).. controls (1.3503601496, 4.0535766602) and (1.3507530558, 
    4.054490531) .. (1.3513195287, 4.054876029).. controls (1.3520598329, 
    4.0554501748) and (1.3529737038, 4.0556163331) .. (1.3536684996, 4.0554501748)
     -- (1.3666083179, 4.0520887758).. controls (1.3681794137, 4.0517035425) and 
    (1.3689649617, 4.0502607696) .. (1.36857232, 4.0487877548).. controls 
    (1.3681794137, 4.0472997381) and (1.3666083179, 4.046340359) .. (1.3652637319,
     4.0467330006).. controls (1.3652637319, 4.0467330006) and (1.359553071, 
    4.0481681006) .. (1.3558517354, 4.0490367277).. controls (1.3558517354, 
    4.047609036) and (1.3558517354, 4.0472087215) .. (1.3558517354, 
    4.0460379931).. controls (1.3594171281, 4.0453355244) and (1.3734973271, 
    4.0426764883) .. (1.3734973271, 4.0426764883).. controls (1.3742754667, 
    4.0424198425) and (1.3750684229, 4.0424198425) .. (1.3752800896, 
    4.0425018633).. controls (1.3752800896, 4.0426764883) and (1.3754547146, 
    4.0434612425) .. (1.3754547146, 4.0434612425) -- (1.3756372771, 
    4.0527825663).. controls (1.3756372771, 4.0527825663) and (1.3754547146, 
    4.0544896579) .. (1.3722380427, 4.0556151954) -- (1.3452182371, 4.061461905)..
     controls (1.3448253308, 4.0615915508) and (1.3437828725, 4.0607216008) .. 
    (1.3438736246, 4.056529225) -- (1.3438736246, 4.0396314252).. controls 
    (1.3438736246, 4.038672046) and (1.3449086746, 4.0371009502) .. (1.3490480279,
     4.0364059956) -- (1.3811213223, 4.0306121763).. controls (1.3815142285, 
    4.0306121763) and (1.3829114931, 4.0301287825) .. (1.3836896327, 
    4.0307418221).. controls (1.3846942556, 4.0314821263) and (1.385170241, 
    4.0332648887) .. (1.385170241, 4.0359236867) -- (1.385170241, 4.058668196).. 
    controls (1.385170241, 4.0588031335) and (1.385170241, 4.0588957377) .. 
    (1.385170241, 4.0589777585).. controls (1.385170241, 4.0599371377) and 
    (1.3846945202, 4.0621200296) .. (1.379776895, 4.0639027656).. controls 
    (1.379776895, 4.0639027656) and (1.3468273271, 4.0712601931) .. (1.3460492404,
     4.0714414856).. controls (1.3449086217, 4.0713515273) and (1.3425140631, 
    4.0715235065) .. (1.340512305, 4.0733978148).. controls (1.3392433633, 
    4.0744402731) and (1.3379363217, 4.0766233502) .. (1.3379363217, 4.0803699825)
     -- (1.3379363217, 4.1081225517).. controls (1.3379363217, 4.10821251) and 
    (1.33789928, 4.1117483752) .. (1.340512305, 4.1140976371).. controls 
    (1.3423858196, 4.1160616392) and (1.3452183958, 4.1164468725) .. 
    (1.3488819754, 4.1158878079) -- (1.3822998558, 4.1080016106).. controls 
    (1.3847398433, 4.1080016106) and (1.3904503983, 4.1058640419) .. 
    (1.3907526054, 4.0989749004) -- (1.3907526054, 4.0807627829).. controls 
    (1.3906705846, 4.0803698767) and (1.3904504513, 4.0776279731) .. 
    (1.3880029496, 4.075663971).. controls (1.3858275454, 4.0738812085) and 
    (1.3826851156, 4.0733147356) .. (1.3785909796, 4.0740550398) -- (1.3551138125,
     4.080369956).. controls (1.3529759792, 4.0804149352) and (1.3480887546, 
    4.0825077894) .. (1.3480887546, 4.0886035777) -- (1.3480887546, 
    4.0970637621).. controls (1.3481787129, 4.0977133142) and (1.3485721483, 
    4.1003724825) .. (1.3511479465, 4.1021173567).. controls (1.3529759527, 
    4.1035072129) and (1.3554612369, 4.1039075275) .. (1.3583769981, 4.1032955463)
     -- (1.3760225898, 4.0988010163).. controls (1.377420119, 4.0984157829) and 
    (1.378288746, 4.0972373288) .. (1.3782056669, 4.0958777673) -- (1.377420119, 
    4.0872966419).. controls (1.377420119, 4.0864732585) and (1.3768081377, 
    4.0856424669) .. (1.376196421, 4.0852572335).. controls (1.3756373565, 
    4.0846830877) and (1.3746779773, 4.0846830877) .. (1.373885021, 4.084856919) 
    -- (1.3658325596, 4.087039996).. controls (1.3643519512, 4.0875159815) and 
    (1.3634759158, 4.0890039981) .. (1.3638685575, 4.0905675004).. controls 
    (1.3643519512, 4.0919197858) and (1.3658325596, 4.0929244088) .. 
    (1.3673960883, 4.0924862588).. controls (1.3673960883, 4.0924862588) and 
    (1.3697453238, 4.0917459546) .. (1.3721021792, 4.0911718087).. controls 
    (1.3722768042, 4.0924862588) and (1.3722768042, 4.0928336567) .. 
    (1.3724043333, 4.0938837085).. controls (1.3689673429, 4.0948809231) and 
    (1.357198544, 4.0977964727) .. (1.357198544, 4.0977964727).. controls 
    (1.354456561, 4.0984158623) and (1.3536708015, 4.0972374081) .. (1.3535877488,
     4.0965426123) -- (1.3536697696, 4.0886036306).. controls (1.3536697696, 
    4.0864734702) and (1.3552408654, 4.0859521881) .. (1.3561097571, 4.0858615419)
     -- (1.3798589158, 4.0793729267).. controls (1.3826840573, 4.0789800204) and 
    (1.3839530519, 4.0795925308) .. (1.3844818217, 4.079977235).. controls 
    (1.3850482946, 4.0803701413) and (1.3851692092, 4.0809366142) .. 
    (1.3851692092, 4.0809366142) -- (1.3851692092, 4.0988012279).. controls 
    (1.3850475008, 4.1019436048) and (1.382517846, 4.1025026429) .. (1.3815131967,
     4.1025026429) -- (1.3475731875, 4.110396275).. controls (1.3460925792, 
    4.1106529208) and (1.3448235317, 4.1104862333) .. (1.3442570058, 
    4.1099959604).. controls (1.34356221, 4.1093916521) and (1.34356221, 
    4.1082131979) .. (1.34356221, 4.1082131979) -- (1.34356221, 4.0802869298).. 
    controls (1.34356221, 4.0793727944) and (1.3436442308, 4.0781038527) .. 
    (1.3442570058, 4.0774995179).. controls (1.3448234788, 4.0768499658) and 
    (1.3460471767, 4.0770161242) .. (1.3460471767, 4.0770161242).. controls 
    (1.3463038225, 4.0770161242) and (1.3466136496, 4.0770160183) .. 
    (1.3470065558, 4.0769261658) -- (1.3812554131, 4.0692666635).. controls 
    (1.3920950483, 4.0656484071) and (1.3907504623, 4.0584949469) .. 
    (1.3906673831, 4.0581927663) -- (1.390749404, 4.0359242423).. controls 
    (1.390749404, 4.0313013365) and (1.3894878706, 4.028000289) .. (1.3869195867, 
    4.0262554148).. controls (1.3837772098, 4.0240723377) and (1.3801136302, 
    4.025122469) .. (1.3797736935, 4.0252584648) -- cycle;
  
  
  
      \end{scope}
      \begin{scope}[shift={(0.8731702173, -3.5268413292)}]
        \path[fill=c005093,even odd rule] (1.4409891487, 4.0000727508) -- 
    (1.4092707135, 4.0102703481).. controls (1.399156169, 4.0131861358) and 
    (1.3998587435, 4.0200601167) .. (1.3998587435, 4.0202867058) -- (1.3998587435,
     4.0367161667).. controls (1.399676181, 4.0417243587) and (1.4007273706, 
    4.0448667092) .. (1.403295734, 4.0463396975).. controls (1.405259736, 
    4.0475181517) and (1.4075258129, 4.0471252454) .. (1.4087873463, 4.0462576767)
     -- (1.435671209, 4.0365433937).. controls (1.4429077219, 4.0328874077) and 
    (1.4428169963, 4.0280907765) .. (1.4427339171, 4.0275242506) -- (1.4425143129,
     4.0182784125).. controls (1.4425143129, 4.015445704) and (1.4411244567, 
    4.0138367462) .. (1.4399838379, 4.0130965215).. controls (1.4373325013, 
    4.0113968381) and (1.4338879438, 4.0129218965) .. (1.4333138508, 4.0130965215)
     -- (1.4137949033, 4.01968335).. controls (1.4126241221, 4.0200609104) and 
    (1.4119213887, 4.0210657979) .. (1.4119213887, 4.022334819) -- (1.4119213887, 
    4.0312255862).. controls (1.4119213887, 4.0322225362) and (1.412314295, 
    4.0330988362) .. (1.4131072513, 4.0336200654).. controls (1.4138853908, 
    4.03417913) and (1.41484477, 4.03417913) .. (1.4157133971, 4.0338846487) -- 
    (1.4283961198, 4.0288688896).. controls (1.4299672156, 4.0283024167) and 
    (1.430571524, 4.0266405687) .. (1.4300502948, 4.0251674746).. controls 
    (1.429566901, 4.0237776183) and (1.4277844031, 4.0229843975) .. (1.4264321442,
     4.0235963787).. controls (1.4264321442, 4.0235963787) and (1.4211595275, 
    4.0257265392) .. (1.4174128952, 4.027131556).. controls (1.4174128952, 
    4.0257266185) and (1.4174128952, 4.025167554) .. (1.4174128952, 
    4.0242536302).. controls (1.4205552721, 4.0232490073) and (1.4353682346, 
    4.0182785448) .. (1.4353682346, 4.0182785448).. controls (1.4356703887, 
    4.0182785448) and (1.4371507325, 4.0178933115) .. (1.4371507325, 
    4.0178933115).. controls (1.4371507325, 4.0178933115) and (1.4369390658, 
    4.0185883719) .. (1.4369390658, 4.0184976198) -- (1.4371957117, 
    4.0277811081).. controls (1.4371957117, 4.0277811081) and (1.4370210867, 
    4.0296998665) .. (1.4334037562, 4.0314372) -- (1.4064368144, 4.0412494581).. 
    controls (1.4061272519, 4.0413394165) and (1.4059155852, 4.0414161456) .. 
    (1.4059155852, 4.0414161456).. controls (1.4059155852, 4.0412494581) and 
    (1.4054321915, 4.0403813602) .. (1.4054321915, 4.0375032756).. controls 
    (1.4054321915, 4.0373207131) and (1.4054321915, 4.0198501169) .. 
    (1.4054321915, 4.0198501169).. controls (1.4054321915, 4.0196834294) and 
    (1.4056517956, 4.0171080546) .. (1.4109615598, 4.0155368529) -- (1.4431255004,
     4.0052108533).. controls (1.4439110483, 4.00482562) and (1.4449535067, 
    4.0045613012) .. (1.4455654879, 4.00494627).. controls (1.4460414733, 
    4.0053391762) and (1.4467439421, 4.0067367054) .. (1.4467439421, 4.0095238792)
     -- (1.4467439421, 4.0323060123).. controls (1.4467439421, 4.0323959706) and 
    (1.4467439421, 4.032533554) .. (1.4467439421, 4.0326232477).. controls 
    (1.4467439421, 4.0336201977) and (1.4462153046, 4.0361432908) .. 
    (1.4410786044, 4.0387191419).. controls (1.4410786044, 4.0387191419) and 
    (1.4072150069, 4.0510467937) .. (1.4070413079, 4.0510467937).. controls 
    (1.4059157704, 4.0512214187) and (1.4038612015, 4.051658775) .. (1.401942496, 
    4.0535775333).. controls (1.4008624669, 4.0547483146) and (1.3995480169, 
    4.0567576662) .. (1.3995480169, 4.0602022237).. controls (1.3995480169, 
    4.0602921821) and (1.3995480169, 4.0604588696) .. (1.3995480169, 4.0606782092)
     -- (1.3995480169, 4.0883099696).. controls (1.3995480169, 4.0885216362) and 
    (1.3995480169, 4.0885216362) .. (1.3995480169, 4.0886952029).. controls 
    (1.3995480169, 4.08956383) and (1.399857844, 4.0925324815) .. (1.4022070529, 
    4.0944058637).. controls (1.4042541342, 4.0959695512) and (1.4072151127, 
    4.0961507908) .. (1.4108333692, 4.0950101721) -- (1.4441681969, 
    4.0824255569).. controls (1.4466534281, 4.082070486) and (1.4520545, 
    4.07928318) .. (1.4523188452, 4.0725301402) -- (1.4523188452, 4.0543633458).. 
    controls (1.4522288869, 4.0539704396) and (1.4521442202, 4.0511378104) .. 
    (1.4494861367, 4.0493929098).. controls (1.4472199804, 4.0477387348) and 
    (1.4439944979, 4.0476101473) .. (1.4398550387, 4.0490455119) -- (1.4163249021,
     4.0585858842).. controls (1.4142778208, 4.0590692779) and (1.4097002381, 
    4.0617283404) .. (1.4097002381, 4.0674390012) -- (1.4097002381, 
    4.0759823706).. controls (1.4097452173, 4.0767605102) and (1.4101383881, 
    4.0792833917) .. (1.4127067515, 4.08089989).. controls (1.4146707535, 
    4.0820706712) and (1.4173295779, 4.0822066671) .. (1.4203813085, 4.0810665775)
     -- (1.4380195183, 4.074109544).. controls (1.4392053808, 4.0736184773) and 
    (1.4398549329, 4.0724021877) .. (1.4397718537, 4.0712768354) -- (1.4389786329,
     4.0628166246).. controls (1.4389786329, 4.0619479975) and (1.4384198329, 
    4.0610716975) .. (1.4376340204, 4.0606787912).. controls (1.4368029642, 
    4.0602028058) and (1.4359343371, 4.0600744829) .. (1.43506571, 4.0604591871) 
    -- (1.4270436492, 4.0639037446).. controls (1.4256537929, 4.0645609696) and 
    (1.42508732, 4.0661699008) .. (1.4255630408, 4.0675295681).. controls 
    (1.4262655096, 4.0688665077) and (1.4277836887, 4.069697564) .. (1.4293096731,
     4.0690932556).. controls (1.4293096731, 4.0690932556) and (1.4315758294, 
    4.0680886327) .. (1.433796636, 4.0671747619).. controls (1.4338865944, 
    4.068481539) and (1.4338865944, 4.068700614) .. (1.4341061985, 4.0696978285)..
     controls (1.4309713887, 4.0708762827) and (1.4183717452, 4.0758465071) .. 
    (1.4183717452, 4.0758465071).. controls (1.4166268181, 4.0764508154) and 
    (1.4158486785, 4.0761486613) .. (1.415539116, 4.0761486613) -- (1.4151917181, 
    4.0753631133) -- (1.4152816765, 4.06743916).. controls (1.4152816765, 
    4.064908685) and (1.4176309119, 4.0641155171) .. (1.418114385, 4.0639039562) 
    -- (1.441727495, 4.0542804254).. controls (1.4447792256, 4.0533210462) and 
    (1.446040759, 4.0538348671) .. (1.4462144579, 4.0539708629).. controls 
    (1.4466526079, 4.0542804254) and (1.4467430954, 4.0545825796) .. 
    (1.4467430954, 4.0545825796) -- (1.4467430954, 4.0724020554).. controls 
    (1.4465605329, 4.0761487142) and (1.4432985379, 4.0769343415) .. 
    (1.4427320121, 4.0771533635) -- (1.4088759817, 4.0897001433).. controls 
    (1.4070858108, 4.0903044517) and (1.4059148708, 4.0903044517) .. 
    (1.4056505521, 4.0899647267).. controls (1.4052576458, 4.0897001433) and 
    (1.4051671583, 4.0889225329) .. (1.4051671583, 4.0885219537) -- (1.4051671583,
     4.060678659).. controls (1.4050375125, 4.0572416685) and (1.4069120854, 
    4.0566222525) .. (1.4079998933, 4.0565467669).. controls (1.4082194975, 
    4.0565467669) and (1.4085663662, 4.056546661) .. (1.4088761933, 4.0564118294) 
    -- (1.4432081298, 4.0438197265).. controls (1.4535341294, 4.0388946135) and 
    (1.4523179985, 4.0321415473) .. (1.4522273787, 4.0318393667) -- (1.4523173371,
     4.009525546).. controls (1.4523173371, 4.0049479369) and (1.4509650517, 
    4.0018131535) .. (1.4484798733, 4.0002419519).. controls (1.4450806125, 
    3.9982025435) and (1.4414701348, 3.9998567185) .. (1.4409866881, 4.0000752644)
     -- cycle;
  
  
  
      \end{scope}
      \begin{scope}[shift={(0.8731702173, -3.5323182042)}]
        \path[fill=c005093,even odd rule] (1.3297068804, 4.2375798429) -- 
    (1.2981546829, 4.2432905037).. controls (1.2897851183, 4.2446351163) and 
    (1.2881309698, 4.2498245215) .. (1.2883499654, 4.2523928054) -- (1.2883499654,
     4.2689959917).. controls (1.2881303613, 4.2738682144) and (1.28930908, 
    4.2772295869) .. (1.2916660413, 4.2788007885).. controls (1.2936223704, 
    4.2800623219) and (1.2957526367, 4.2798884906) .. (1.2970594667, 4.2791936948)
     -- (1.3238602767, 4.2733999812).. controls (1.3301374369, 4.2711264167) and 
    (1.3313988644, 4.2669867723) .. (1.3313158117, 4.2649472581) -- (1.3310962075,
     4.2558373629).. controls (1.3310962075, 4.2528687115) and (1.3297063513, 
    4.2513430446) .. (1.3285202242, 4.2505573908).. controls (1.3259140783, 
    4.2484650658) and (1.3224697325, 4.2495601763) .. (1.3220769056, 4.2496432554)
     -- (1.3026484985, 4.2533899142).. controls (1.3012888048, 4.2536619058) and 
    (1.3004201777, 4.2548779308) .. (1.3004201777, 4.2561394642) -- (1.3004201777,
     4.265067961).. controls (1.3004201777, 4.2660273402) and (1.300684761, 
    4.266820561) .. (1.301462636, 4.2674248694).. controls (1.3021574319, 
    4.267825184) and (1.3030337319, 4.2681196652) .. (1.3038647881, 4.267825184) 
    -- (1.3166760983, 4.2646828071).. controls (1.3181490337, 4.2642899008) and 
    (1.3190253337, 4.2626735613) .. (1.3185493483, 4.2613213288).. controls 
    (1.3183747233, 4.2596671537) and (1.3166760983, 4.2587908537) .. 
    (1.3151954371, 4.2591911683).. controls (1.3151954371, 4.2591911683) and 
    (1.3096964429, 4.26061886) .. (1.3059951073, 4.2615025683).. controls 
    (1.3059951073, 4.2601505475) and (1.3059951073, 4.2595840746) .. 
    (1.3059951073, 4.2584509435).. controls (1.309530311, 4.2577863102) and 
    (1.3234668942, 4.2551348677) .. (1.3234668942, 4.2551348677).. controls 
    (1.3243431942, 4.2548782219) and (1.3251287421, 4.2548782219) .. 
    (1.3252118212, 4.254923201).. controls (1.3252118212, 4.2551348677) and 
    (1.3254314254, 4.2558373365) .. (1.3254314254, 4.2558373365) -- (1.3256060504,
     4.2651587925).. controls (1.3256060504, 4.2652487508) and (1.3254314254, 
    4.2669412904) .. (1.3222975417, 4.2679913952) -- (1.2953229533, 
    4.2738682937).. controls (1.2948395596, 4.2741778562) and (1.293842345, 
    4.2731809063) .. (1.293842345, 4.268996071) -- (1.293842345, 4.2520905983).. 
    controls (1.293842345, 4.2511312192) and (1.2950207992, 4.2495601233) .. 
    (1.2991602848, 4.2489105977) -- (1.3313166583, 4.243063994).. controls 
    (1.3315812417, 4.2428443898) and (1.3329708333, 4.2426710877) .. (1.333839725,
     4.2432915356).. controls (1.3346252729, 4.2438503356) and (1.3352372542, 
    4.2458521731) .. (1.3352372542, 4.2483826481) -- (1.3352372542, 
    4.2711271575).. controls (1.3352372542, 4.271262095) and (1.3352372542, 
    4.27130972) .. (1.3352372542, 4.2715274721).. controls (1.3352372542, 
    4.2723056117) and (1.3348443479, 4.2745717415) .. (1.32970778, 4.2762335629)..
     controls (1.32970778, 4.2762335629) and (1.2967581856, 4.2837646894) .. 
    (1.2961463367, 4.2839006323).. controls (1.2950207992, 4.2839006323) and 
    (1.2926640496, 4.2839826531) .. (1.2905716452, 4.2858495531).. controls 
    (1.289400864, 4.2868089323) and (1.287958091, 4.2890826556) .. (1.2881316577, 
    4.2927839648) -- (1.2881316577, 4.3205818571).. controls (1.2880496369, 
    4.3207564821) and (1.2879570327, 4.3242001135) .. (1.2905716452, 
    4.3266928854).. controls (1.292580891, 4.3284304042) and (1.2954135996, 
    4.3288682896) .. (1.2989336427, 4.3284301396) -- (1.3323591167, 
    4.3205817777).. controls (1.3348443479, 4.3204521319) and (1.340509765, 
    4.3184063735) .. (1.3407287871, 4.3113435067) -- (1.3407287871, 
    4.2933125758).. controls (1.3407287871, 4.29273843) and (1.3405938496, 
    4.2899512827) .. (1.3381604767, 4.2880779004).. controls (1.3358869121, 
    4.2862498942) and (1.332706594, 4.2857739087) .. (1.3287484008, 4.2864311337) 
    -- (1.3052108823, 4.2927385887).. controls (1.3030354781, 4.2929211512) and 
    (1.2980196925, 4.2950499887) .. (1.2980196925, 4.3009268608) -- (1.2980196925,
     4.3093796104).. controls (1.29815463, 4.3103011542) and (1.2987676696, 
    4.3129146819) .. (1.3010789902, 4.3146595825).. controls (1.3030353194, 
    4.3158756075) and (1.3056036033, 4.3163592658) .. (1.3084363383, 4.3157472846)
     -- (1.3262179787, 4.3111697812).. controls (1.3274795121, 4.3108297917) and 
    (1.3283933829, 4.30956085) .. (1.3281743079, 4.3082085646) -- (1.3275625912, 
    4.2997560531).. controls (1.3273959037, 4.2989705052) and (1.3270035267, 
    4.298230201) .. (1.3262632225, 4.2976258927).. controls (1.3256059975, 
    4.2972255781) and (1.3248204496, 4.2970591552) .. (1.3239518225, 4.2973160656)
     -- (1.315974979, 4.2994988781).. controls (1.3144491269, 4.2999748635) and 
    (1.3136709873, 4.3015382865) .. (1.3140562206, 4.3029358421).. controls 
    (1.3144491269, 4.3044616942) and (1.3159747144, 4.305254915) .. (1.3174629162,
     4.3048546004).. controls (1.3174629162, 4.3048546004) and (1.3199029037, 
    4.3042955358) .. (1.3222973829, 4.3035930671).. controls (1.3222973829, 
    4.3048546004) and (1.3222974887, 4.3052096713) .. (1.3224720079, 
    4.3064635317).. controls (1.3189897208, 4.3073398317) and (1.3071754929, 
    4.3103838629) .. (1.3071754929, 4.3103838629).. controls (1.3045241298, 
    4.3108296858) and (1.3038668783, 4.3096964754) .. (1.3036402892, 4.3090016796)
     -- (1.3037302475, 4.3009266227).. controls (1.3037302475, 4.2989702935) and 
    (1.3053013433, 4.298403556) .. (1.306170235, 4.2983207415) -- (1.3299192879, 
    4.2919150731).. controls (1.332706594, 4.2915147585) and (1.3340134504, 
    4.2919150731) .. (1.3346252729, 4.2923984669).. controls (1.3351465021, 
    4.2927384565) and (1.3352372542, 4.2933123377) .. (1.3352372542, 4.2935240044)
     -- (1.3352372542, 4.3111697019).. controls (1.3351472958, 4.3144857777) and 
    (1.3327065146, 4.3148785781) .. (1.331490516, 4.3148785781) -- (1.2977628615, 
    4.3228554217).. controls (1.2961462573, 4.3231120675) and (1.2948848298, 
    4.3230300467) .. (1.294318304, 4.3224095988).. controls (1.2936232435, 
    4.3219412863) and (1.2936232435, 4.3207551592) .. (1.2936232435, 4.3205815925)
     -- (1.2936232435, 4.2927384035).. controls (1.2936232435, 4.2918242681) and 
    (1.2938428477, 4.2905553265) .. (1.294318304, 4.2899509917).. controls 
    (1.2950207727, 4.2893845188) and (1.2961084748, 4.2895657583) .. 
    (1.2961084748, 4.2895657583).. controls (1.2963730581, 4.2895657583) and 
    (1.2967580269, 4.2895656525) .. (1.2969847748, 4.2894758) -- (1.3313167112, 
    4.2816350052).. controls (1.3421185904, 4.2780167487) and (1.3407287871, 
    4.2709087175) .. (1.3407287871, 4.2705687808) -- (1.3407287871, 
    4.2483831244).. controls (1.3407287871, 4.2436771394) and (1.3395503329, 
    4.2404139538) .. (1.3370198844, 4.238585921).. controls (1.3338397515, 
    4.236395171) and (1.3302289827, 4.2375812981) .. (1.32970778, 4.2375812981) --
     cycle;
  
  
  
      \end{scope}
      \begin{scope}[shift={(0.8731702173, -3.5321898813)}]
        \path[fill=c005093,even odd rule] (1.3914463694, 4.2241108579) -- 
    (1.3597732573, 4.2297309783).. controls (1.3514036927, 4.23099992) and 
    (1.3497494648, 4.2361819433) .. (1.349915676, 4.2387577944) -- (1.349915676, 
    4.2553155517).. controls (1.3498389469, 4.2603237173) and (1.3509200344, 
    4.2635491733) .. (1.3532769692, 4.2651203485).. controls (1.3552409712, 
    4.2664271256) and (1.3573786988, 4.2660723194) .. (1.3587686079, 4.265596334) 
    -- (1.3853503694, 4.2597118948).. controls (1.3917106087, 4.2574457385) and 
    (1.3929797092, 4.2534874925) .. (1.3928890894, 4.2512138485) -- (1.3925795269,
     4.242202246).. controls (1.3925795269, 4.2392335946) and (1.3913179935, 
    4.2374509644) .. (1.3901472123, 4.2367560363).. controls (1.3876167373, 
    4.23492803) and (1.3840890477, 4.2358797363) .. (1.3836886538, 4.2360911383) 
    -- (1.3642604319, 4.2398001204).. controls (1.3628629027, 4.2400197246) and 
    (1.3619035235, 4.2411976496) .. (1.3619035235, 4.2425421033) -- (1.3619035235,
     4.2514781938).. controls (1.3619035235, 4.2523468208) and (1.3622961652, 
    4.2531323688) .. (1.362998634, 4.25374435).. controls (1.3637389381, 
    4.2541372563) and (1.3645623215, 4.2543563313) .. (1.3654309485, 4.2541372563)
     -- (1.3782043175, 4.2510024465).. controls (1.3797754133, 4.2506095402) and 
    (1.380644305, 4.2490384444) .. (1.3802513988, 4.247557889).. controls 
    (1.3798584925, 4.2461225244) and (1.3783781487, 4.2451179015) .. 
    (1.3768068413, 4.2455031348).. controls (1.3768068413, 4.2455031348) and 
    (1.3712321498, 4.2468628285) .. (1.3675685702, 4.2478674515).. controls 
    (1.3675685702, 4.2463794348) and (1.3675685702, 4.2459034494) .. 
    (1.3675685702, 4.2448081538).. controls (1.3710962069, 4.2441207663) and 
    (1.38504803, 4.2414995392) .. (1.38504803, 4.2414995392).. controls 
    (1.3859621654, 4.241197385) and (1.3867400404, 4.2411976496) .. (1.3868307925,
     4.241279935).. controls (1.3868307925, 4.241279935) and (1.3868307925, 
    4.2412798292) .. (1.3868307925, 4.2414148725).. controls (1.3868307925, 
    4.2414968933) and (1.3870054175, 4.2422004204) .. (1.3870054175, 4.2422004204)
     -- (1.3872250217, 4.2516124169).. controls (1.3872250217, 4.2516124169) and 
    (1.3870054175, 4.253130596) .. (1.3839089458, 4.2543543998) -- (1.3569872477, 
    4.2603219975).. controls (1.3564207748, 4.2604119558) and (1.3554613956, 
    4.259536185) .. (1.355506904, 4.255313726) -- (1.355506904, 4.2384536823).. 
    controls (1.355506904, 4.2374490594) and (1.3566324415, 4.235877699) .. 
    (1.3607340917, 4.2353113054) -- (1.3928904652, 4.2294269721).. controls 
    (1.3931926194, 4.2294269721) and (1.3945901485, 4.2289509867) .. 
    (1.3953756965, 4.2296015971).. controls (1.396372911, 4.2302209867) and 
    (1.3968033881, 4.2321699075) .. (1.3968033881, 4.2347005148) -- (1.3968033881,
     4.2574448919).. controls (1.3968033881, 4.2576565585) and (1.3968033881, 
    4.2576565585) .. (1.3968033881, 4.2577922898).. controls (1.3968033881, 
    4.2586609169) and (1.396372911, 4.2610253129) .. (1.3914477981, 
    4.2625435979).. controls (1.3914477981, 4.2625435979) and (1.3585056915, 
    4.2701277204) .. (1.3577654667, 4.2702183667).. controls (1.3566325208, 
    4.2702183667) and (1.3542377506, 4.2702183667) .. (1.3520622935, 
    4.2721746958).. controls (1.3509216748, 4.2732700708) and (1.3495318185, 
    4.2753095056) .. (1.3497508935, 4.279063705) -- (1.3497508935, 4.3069069204)..
     controls (1.3497508935, 4.3071185871) and (1.3494942477, 4.3105252563) .. 
    (1.3520622935, 4.3129575973).. controls (1.354064131, 4.3147477681) and 
    (1.3569874065, 4.3153597494) .. (1.3605603398, 4.3145286931) -- (1.393970759, 
    4.3068163006).. controls (1.3964181548, 4.3067342798) and (1.4021287892, 
    4.3047239756) .. (1.4023856467, 4.2977065113) -- (1.4023856467, 
    4.2795395317).. controls (1.4023856467, 4.2790635463) and (1.4022110217, 
    4.2763972606) .. (1.3996814198, 4.2743501794).. controls (1.3974226719, 
    4.2725674169) and (1.3943711529, 4.272000944) .. (1.3903600696, 4.2727790835) 
    -- (1.3667832869, 4.2790647104).. controls (1.364562639, 4.2793213562) and 
    (1.3597732573, 4.2814139458) .. (1.3597732573, 4.2872983321) -- (1.3597732573,
     4.2958341079).. controls (1.3597732573, 4.2964913329) and (1.3603323219, 
    4.2991501838) .. (1.3628174208, 4.3008875967).. controls (1.3646531, 
    4.3022927987) and (1.3670929552, 4.3026777675) .. (1.3700540396, 4.3020660508)
     -- (1.3877902775, 4.2974961146).. controls (1.3890518108, 4.2973135521) and 
    (1.3899204379, 4.2959247542) .. (1.3897542796, 4.2945273573) -- (1.3890518108,
     4.2861199308).. controls (1.3890068317, 4.2852513037) and (1.3886136608, 
    4.2844657558) .. (1.3878810296, 4.2839823621).. controls (1.3872238046, 
    4.2835971287) and (1.3863551775, 4.2833780537) .. (1.3855619567, 4.2836349642)
     -- (1.377509601, 4.2859011204).. controls (1.3760215844, 4.286301435) and 
    (1.3751529573, 4.2879029579) .. (1.375636351, 4.2894363242).. controls 
    (1.3760215844, 4.2908261804) and (1.377509601, 4.2916948075) .. (1.3790732356,
     4.2912188221).. controls (1.3790732356, 4.2912188221) and (1.3815132231, 
    4.2906068408) .. (1.3837793265, 4.2898212929).. controls (1.3839539515, 
    4.2912188221) and (1.3839539515, 4.2916114637) .. (1.384088889, 
    4.2926540015).. controls (1.3805083885, 4.2935303015) and (1.3687849921, 
    4.2965820056) .. (1.3687849921, 4.2965820056).. controls (1.3660430092, 
    4.297186314) and (1.3654310808, 4.2960155327) .. (1.3652648431, 4.2953127994) 
    -- (1.3652648431, 4.2872982262).. controls (1.3652648431, 4.285251145) and 
    (1.3670020973, 4.2847299158) .. (1.3679161798, 4.2845108144) -- (1.3915368569,
     4.2782713837).. controls (1.3942788398, 4.2777578275) and (1.3956309929, 
    4.2783613421) .. (1.3961597892, 4.2786716983).. controls (1.3967262621, 
    4.2790646046) and (1.3968019329, 4.2797594004) .. (1.3968019329, 4.2797594004)
     -- (1.3968019329, 4.297586364).. controls (1.3966352454, 4.3007514421) and 
    (1.3942788662, 4.3012424823) .. (1.3931912435, 4.301287779) -- (1.3593729427, 
    4.3091663298).. controls (1.3577640115, 4.3094309131) and (1.3564194254, 
    4.3092483506) .. (1.3558528996, 4.3088718485).. controls (1.3552409183, 
    4.308252459) and (1.3552409183, 4.3071192485) .. (1.3552409183, 4.3071192485) 
    -- (1.3552409183, 4.2790644988).. controls (1.3552409183, 4.2782712779) and 
    (1.3552409183, 4.2768814217) .. (1.3560264662, 4.2762695463).. controls 
    (1.3566307746, 4.2757030733) and (1.3575901537, 4.2758390692) .. 
    (1.3575901537, 4.2758390692).. controls (1.35798306, 4.2758840483) and 
    (1.3583757017, 4.2758390692) .. (1.3585495329, 4.2758390692) -- (1.3928889306,
     4.26807376).. controls (1.4036983769, 4.264417774) and (1.4023839533, 
    4.2572793421) .. (1.4023462237, 4.2569697002) -- (1.4023832654, 
    4.2347012027).. controls (1.4023832654, 4.2300858375) and (1.4011217321, 
    4.226852735) .. (1.3985987712, 4.2250246758).. controls (1.3954563944, 
    4.2227585196) and (1.3918382437, 4.2238991383) .. (1.3914454169, 4.2241105404)
     -- cycle;
  
  
  
      \end{scope}
      \begin{scope}[shift={(0.8731702173, -3.5267058625)}]
        \path[fill=c005093,even odd rule] (1.4527151381, 4.1995921056) -- 
    (1.4209515121, 4.2097443004).. controls (1.410836941, 4.2127507344) and 
    (1.4115395156, 4.219541636) .. (1.4115395156, 4.2198059813) -- (1.4115395156, 
    4.2363186271).. controls (1.4113569531, 4.2412890367) and (1.4125367302, 
    4.2444237406) .. (1.4148480244, 4.24586654).. controls (1.416932941, 
    4.2470902379) and (1.4191612883, 4.2464782567) .. (1.420558791, 4.2457316025) 
    -- (1.4472236052, 4.2361081775).. controls (1.4547170021, 4.2325276771) and 
    (1.4544978742, 4.2275723752) .. (1.4544978742, 4.2270360383) -- (1.45419572, 
    4.2178883871).. controls (1.45419572, 4.2149273027) and (1.4529265138, 
    4.2133636946) .. (1.4517557325, 4.2126611994).. controls (1.4490137496, 
    4.2109236806) and (1.4456975679, 4.2123590452) .. (1.4449573431, 4.2126611994)
     -- (1.425476231, 4.2191498146).. controls (1.4243054498, 4.2195427208) and 
    (1.423602981, 4.2205925875) .. (1.423602981, 4.2218163119) -- (1.423602981, 
    4.2307827765).. controls (1.423602981, 4.2316590765) and (1.4240863748, 
    4.232527439) .. (1.4246980915, 4.2331470931).. controls (1.425476231, 
    4.2336606494) and (1.4265260977, 4.2337514015) .. (1.4273116456, 4.2333587598)
     -- (1.4402134962, 4.2283505942).. controls (1.441512865, 4.2277841213) and 
    (1.4422529046, 4.2260770296) .. (1.44173194, 4.2247399048).. controls 
    (1.4411276317, 4.2233347027) and (1.4394657838, 4.2225570923) .. 
    (1.4380230373, 4.2231158923).. controls (1.4380230373, 4.2231158923) and 
    (1.4328864958, 4.2252082173) .. (1.4291020017, 4.2266887462).. controls 
    (1.4291020017, 4.2252081379) and (1.4291020017, 4.2249059838) .. 
    (1.4291020017, 4.2238560377).. controls (1.4323198642, 4.2227760085) and 
    (1.447049774, 4.2178885988) .. (1.447049774, 4.2178885988).. controls 
    (1.447224399, 4.2177589529) and (1.4488325365, 4.2175033654) .. (1.4488325365,
     4.2175033654) -- (1.4486208698, 4.2180245946) -- (1.4488325365, 
    4.2273912415).. controls (1.4488325365, 4.2273912415) and (1.4487108281, 
    4.2291361685) .. (1.4450858777, 4.2310096037) -- (1.4181187506, 
    4.2408142948).. controls (1.4178165965, 4.2408963156) and (1.4176353569, 
    4.2408963156) .. (1.4176353569, 4.2408963156).. controls (1.4176353569, 
    4.2408142948) and (1.4170233756, 4.2398538573) .. (1.4170233756, 
    4.2370589577).. controls (1.4170233756, 4.236847291) and (1.4170233756, 
    4.2195417419) .. (1.4170233756, 4.219413366).. controls (1.4170233756, 
    4.2190281327) and (1.4173329381, 4.216663816) .. (1.4225979612, 4.2151002079) 
    -- (1.4547166846, 4.2046834562).. controls (1.4554569887, 4.2045088313) and 
    (1.4566354429, 4.2041243917) .. (1.4571640804, 4.2045088313).. controls 
    (1.457723145, 4.2049091458) and (1.4582896179, 4.2062915937) .. (1.4582896179,
     4.2091243023) -- (1.4582896179, 4.2319669985).. controls (1.4582896179, 
    4.2319669985) and (1.4582896179, 4.2319668927) .. (1.4582896179, 
    4.2321866027).. controls (1.4582896179, 4.2331459819) and (1.4580329721, 
    4.2357142129) .. (1.4527148206, 4.238146554).. controls (1.4527148206, 
    4.238146554) and (1.4187757375, 4.2506103604) .. (1.4187304144, 
    4.2507009802).. controls (1.4176350394, 4.2507909385) and (1.4155427144, 
    4.2510938865) .. (1.4135862794, 4.2530954594).. controls (1.4125364127, 
    4.2542284052) and (1.411237044, 4.2562755923) .. (1.411237044, 4.2596369648)..
     controls (1.411237044, 4.2598036523) and (1.411237044, 4.2600674419) .. 
    (1.411237044, 4.2602034377) -- (1.411237044, 4.2880014358).. controls 
    (1.411237044, 4.2880014358) and (1.411237044, 4.288046415) .. (1.411237044, 
    4.2883112629).. controls (1.411237044, 4.2892706421) and (1.4114037315, 
    4.2920125985) .. (1.413797946, 4.2937500115).. controls (1.4158527002, 
    4.2955401823) and (1.4187758433, 4.2957592573) .. (1.4225224756, 4.2945355594)
     -- (1.4558498421, 4.2819886737).. controls (1.4582443212, 4.2815957675) and 
    (1.4637435271, 4.2788992929) .. (1.4639097383, 4.2721839827) -- (1.4639097383,
     4.25397168).. controls (1.4639097383, 4.2535787738) and (1.4637430508, 
    4.2507915471) .. (1.461167676, 4.2490466729).. controls (1.4589015198, 
    4.2473924979) and (1.4556760373, 4.2472565021) .. (1.4515441452, 4.2484725271)
     -- (1.4280065473, 4.2581563829).. controls (1.425959466, 4.2586776121) and 
    (1.4213365337, 4.2613290546) .. (1.4213365337, 4.2669943658) -- (1.4213365337,
     4.2755376558).. controls (1.4213365337, 4.2762401246) and (1.4219559233, 
    4.27889924) .. (1.4243958579, 4.280417525).. controls (1.4264352662, 
    4.2815959792) and (1.4290111967, 4.2815957146) .. (1.4320176042, 4.2806371292)
     -- (1.4496633017, 4.2735819881).. controls (1.4507510037, 4.2731890819) and 
    (1.4515442246, 4.2720032194) .. (1.4514534725, 4.2707493854) -- (1.4506602517,
     4.262379715).. controls (1.4505385433, 4.2613750921) and (1.4501842662, 
    4.2605895442) .. (1.4493987183, 4.2602040462).. controls (1.4486205787, 
    4.2596375733) and (1.4476162204, 4.2596375733) .. (1.4468304079, 4.2600691087)
     -- (1.4388158083, 4.2633851581).. controls (1.4374182792, 4.263951631) and 
    (1.4366327312, 4.2656513144) .. (1.4373352, 4.2670864937).. controls 
    (1.4378111854, 4.2685671021) and (1.4395939479, 4.269133575) .. (1.4409913183,
     4.2685671021).. controls (1.4409913183, 4.2685671021) and (1.4432122308, 
    4.2675701521) .. (1.4454782812, 4.2666938521).. controls (1.4455682396, 
    4.2678723063) and (1.4456978854, 4.268257275) .. (1.4456978854, 4.269133575)..
     controls (1.4426461548, 4.2704025167) and (1.4300540519, 4.2753276825) .. 
    (1.4300540519, 4.2753276825).. controls (1.4281808019, 4.2760224783) and 
    (1.427395254, 4.2757129158) .. (1.4272290163, 4.2756224283) -- (1.4268739454, 
    4.2749273679) -- (1.4268739454, 4.2669959269).. controls (1.4268739454, 
    4.2643370231) and (1.4293136683, 4.2635514752) .. (1.4297971944, 4.2633852375)
     -- (1.453538786, 4.253754166).. controls (1.4564621408, 4.2527569515) and 
    (1.4576329221, 4.2533612598) .. (1.458033316, 4.253579541).. controls 
    (1.4582899619, 4.253754166) and (1.4582899619, 4.2540555265) .. (1.4582899619,
     4.2541386056) -- (1.4582899619, 4.2720033252).. controls (1.4582449827, 
    4.2757122015) and (1.4549812679, 4.2764977494) .. (1.4543316894, 4.2765808285)
     -- (1.4205588175, 4.2892712242).. controls (1.4187308112, 4.2898302887) and 
    (1.4176355685, 4.2896490492) .. (1.4173333615, 4.2894379117).. controls 
    (1.417023799, 4.2892712242) and (1.416857376, 4.2884406971) .. (1.4169327823, 
    4.2880480554) -- (1.4169327823, 4.2601142202).. controls (1.4168560531, 
    4.2566696627) and (1.4185038781, 4.2561936773) .. (1.4195540094, 
    4.2561030575).. controls (1.4199014073, 4.2561030575) and (1.4202490698, 
    4.2561031633) .. (1.4205588969, 4.2558913908) -- (1.4548908069, 4.243344611)..
     controls (1.4651714833, 4.23845736) and (1.4639100558, 4.2316134623) .. 
    (1.4639100558, 4.2313944402) -- (1.4639100558, 4.2091259427).. controls 
    (1.4639100558, 4.2045105775) and (1.4627392746, 4.2013757677) .. 
    (1.4601633177, 4.1998045925).. controls (1.4568547296, 4.1976744321) and 
    (1.4531080708, 4.1994193592) .. (1.452715244, 4.1995929258) -- cycle;
  
  
  
      \end{scope}
      \begin{scope}[shift={(0.8731702173, -3.5322052271)}]
        \path[fill=c005093,even odd rule] (0.9944876023, 4.0403358519).. controls 
    (0.9919192919, 4.042080779) and (0.99074097, 4.0453439117) .. (0.99074097, 
    4.0499215208) -- (0.99074097, 4.0727113269).. controls (0.99074097, 
    4.0725367019) and (0.9892603617, 4.0796834152) .. (0.9998960825, 4.0832185925)
     -- (1.0344924435, 4.0910064971).. controls (1.0347945977, 4.0911414346) and 
    (1.0351041602, 4.0911414346) .. (1.0353610706, 4.0911414346).. controls 
    (1.0353610706, 4.0911414346) and (1.0364487727, 4.0910064971) .. (1.037060754,
     4.0915343408).. controls (1.0377555498, 4.0920934054) and (1.0377558144, 
    4.0934983429) .. (1.0377558144, 4.094147895) -- (1.0377558144, 4.1221195921)..
     controls (1.0377558144, 4.1223471337) and (1.037755735, 4.1235171212) .. 
    (1.037151506, 4.1239928421).. controls (1.0365395248, 4.124559315) and 
    (1.035187504, 4.1246876379) .. (1.0334879265, 4.1243931567) -- (1.0003721208, 
    4.1166355469).. controls (0.998801025, 4.1165455885) and (0.9963610375, 
    4.1160235656) .. (0.9962326617, 4.1127075427) -- (0.9962326617, 4.094842929)..
     controls (0.9962326617, 4.0949249498) and (0.9963623075, 4.0944124519) .. 
    (0.9969277221, 4.0938835498).. controls (0.9974111158, 4.0935813956) and 
    (0.9986726492, 4.0930146581) .. (1.0017243269, 4.0934983165) -- (1.02481634, 
    4.0997680154).. controls (1.0261683608, 4.0999876196) and (1.0278679648, 
    4.1005535633) .. (1.0278679648, 4.1025175654) -- (1.0278679648, 
    4.1110533677).. controls (1.027738319, 4.1112359302) and (1.026954094, 
    4.1124508969) .. (1.0240760094, 4.1118389156).. controls (1.0240760094, 
    4.1118389156) and (1.0123978302, 4.1088701583) .. (1.0089985958, 
    4.1080014519).. controls (1.0090885542, 4.1068306706) and (1.0090885542, 
    4.1065285165) .. (1.0093081583, 4.1051688227).. controls (1.0114835625, 
    4.1058712915) and (1.0140142492, 4.1065285165) .. (1.0140142492, 
    4.1065285165).. controls (1.0154493492, 4.1069060769) and (1.0170584127, 
    4.1060374498) .. (1.0174512131, 4.1044737623).. controls (1.0179346069, 
    4.103083906) and (1.0170583069, 4.1015126515) .. (1.0154493756, 4.101119851) 
    -- (1.0075178287, 4.0987632073).. controls (1.0067775246, 4.0985885823) and 
    (1.0058633892, 4.0987632073) .. (1.0051685933, 4.099246601).. controls 
    (1.0045113683, 4.0997678302) and (1.0039901392, 4.1005533781) .. 
    (1.0038996517, 4.1013315177) -- (1.0032048558, 4.1098747813).. controls 
    (1.0029852517, 4.1111437229) and (1.0038996517, 4.1124960083) .. 
    (1.0053879329, 4.1127075162) -- (1.0228146083, 4.117322855).. controls 
    (1.0259493917, 4.1180253238) and (1.0283061412, 4.1176326821) .. 
    (1.0303078994, 4.1162351529).. controls (1.0328383744, 4.1144902258) and 
    (1.0333143069, 4.1117028669) .. (1.0334426827, 4.1115366292) -- (1.0334426827,
     4.102517486).. controls (1.0334426827, 4.096580236) and (1.0283061412, 
    4.0944575633) .. (1.0258662331, 4.0942838644) -- (1.002895161, 4.0879612752)..
     controls (0.9986726227, 4.0874326377) and (0.995575569, 4.0877866502) .. 
    (0.9933094921, 4.0895778794).. controls (0.9908771775, 4.0916249606) and 
    (0.9907411817, 4.0944121608) .. (0.9907411817, 4.0946690183) -- (0.9907411817,
     4.1127075427).. controls (0.9909158067, 4.1199440556) and (0.9966254356, 
    4.1220288665) .. (0.99958652, 4.1221195127) -- (1.032445574, 4.1298773342).. 
    controls (1.0360638304, 4.13057213) and (1.038934295, 4.1299672925) .. 
    (1.040852974, 4.1282228946).. controls (1.0435119042, 4.1257829071) and 
    (1.0433382052, 4.1223461019) .. (1.0433382052, 4.1219457344) -- (1.0433382052,
     4.094283785).. controls (1.0435128302, 4.0905749087) and (1.0420314281, 
    4.0885731506) .. (1.040852974, 4.0874325319).. controls (1.0388058927, 
    4.0855214465) and (1.036366011, 4.0854306944) .. (1.0352782825, 
    4.0854306944).. controls (1.0346210575, 4.0852560694) and (1.0013315, 
    4.0777635985) .. (1.0013315, 4.0777635985).. controls (0.9964517367, 
    4.0761017506) and (0.9962327146, 4.0737601087) .. (0.9962327146, 4.0732690685)
     -- (0.9962327146, 4.0499203037).. controls (0.9962327146, 4.0473972371) and 
    (0.9967539437, 4.0455163935) .. (0.9977133229, 4.044866815).. controls 
    (0.9984988708, 4.0442548338) and (0.9998964, 4.0445646608) .. (0.9998964, 
    4.0445646608) -- (1.0322265519, 4.0504491).. controls (1.036366011, 
    4.0510986521) and (1.0375369246, 4.0527075833) .. (1.0375369246, 4.0535762369)
     -- (1.0375369246, 4.0705269004).. controls (1.0377565287, 4.074840085) and 
    (1.0366682975, 4.0756332794) .. (1.0366682975, 4.0757239256) -- (1.008915649, 
    4.0695223569).. controls (1.005946971, 4.0683515756) and (1.0057731662, 
    4.0668256971) .. (1.0057731662, 4.0668256971) -- (1.0060377496, 
    4.0575873202).. controls (1.0060377496, 4.0574126952) and (1.0061673954, 
    4.0564996181) .. (1.0061673954, 4.0564996181).. controls (1.0063869996, 
    4.0564096598) and (1.0071267746, 4.0564096598) .. (1.0076480037, 
    4.056581639).. controls (1.0076480037, 4.056581639) and (1.0219397635, 
    4.059331189) .. (1.0254749408, 4.0600260906).. controls (1.0254749408, 
    4.0611137927) and (1.0254749408, 4.061725774) .. (1.0254749408, 
    4.0630778213).. controls (1.0217736317, 4.0621184421) and (1.0162367756, 
    4.0608569087) .. (1.0162367756, 4.0608569087).. controls (1.0147940027, 
    4.0603282713) and (1.013139616, 4.0612873858) .. (1.0128374354, 
    4.0627756671).. controls (1.0123992854, 4.0642562754) and (1.0133134208, 
    4.0657292108) .. (1.0148845167, 4.0662202246) -- (1.0276503979, 4.069475949)..
     controls (1.0283907021, 4.0695209281) and (1.029395325, 4.0694758431) .. 
    (1.0300070417, 4.0689621281).. controls (1.0306565937, 4.0683501469) and 
    (1.0310495, 4.0676024344) .. (1.0310495, 4.0666885635) -- (1.0310495, 
    4.057812851).. controls (1.0310495, 4.0563171615) and (1.0300901208, 
    4.0552371323) .. (1.0287833437, 4.0549348988) -- (1.0089999452, 
    4.0512260225).. controls (1.0089999452, 4.0510963767) and (1.0055629548, 
    4.0501380558) .. (1.002896484, 4.0518755746).. controls (1.0017257027, 
    4.0528801975) and (1.0003734173, 4.0544060496) .. (1.0003734173, 4.0574125365)
     -- (1.0001538131, 4.0665980231).. controls (1.0000717923, 4.068607269) and 
    (1.0013322673, 4.0727089456) .. (1.0072996798, 4.0748390796) -- (1.0350145194,
     4.0809424615).. controls (1.0356640715, 4.0814184469) and (1.0378019048, 
    4.0815922781) .. (1.03981123, 4.0803381531).. controls (1.0421604654, 
    4.0787670573) and (1.0432481146, 4.0754507963) .. (1.0430290925, 4.0705256833)
     -- (1.0430290925, 4.0536657454).. controls (1.043338655, 4.0514903413) and 
    (1.0416392362, 4.0462101046) .. (1.0332318627, 4.0449485712) -- (1.0012415417,
     4.0391624248).. controls (1.0012415417, 4.0391624248) and (0.9976308523, 
    4.0379839706) .. (0.9944884754, 4.040333206) -- cycle;
  
  
  
      \end{scope}
      \begin{scope}[shift={(0.8731702173, -3.5300147417)}]
        \path[fill=c005093,even odd rule] (0.932484456, 4.0237249927).. controls 
    (0.9300370602, 4.0253791677) and (0.9286471775, 4.028521624) .. (0.9286471775,
     4.0332729321) -- (0.9286471775, 4.0560022808).. controls (0.9286471775, 
    4.0557906142) and (0.9273856442, 4.0628459669) .. (0.9377948287, 4.0668646173)
     -- (0.9723080312, 4.0765787944).. controls (0.9726175938, 4.0765787944) and 
    (0.9728745042, 4.0765787944) .. (0.9732674104, 4.0765787944).. controls 
    (0.9732674104, 4.0765787944) and (0.9743551125, 4.0765789002) .. (0.975140925,
     4.0772812631).. controls (0.9757529063, 4.0779686506) and (0.9758357208, 
    4.0791545131) .. (0.9758357208, 4.0800231402) -- (0.9758357208, 
    4.1079948373).. controls (0.9758357208, 4.1083422352) and (0.9757537, 
    4.1093471227) .. (0.9752692479, 4.109732356).. controls (0.9746649396, 
    4.1103065019) and (0.9733959979, 4.1103062373) .. (0.9717416377, 4.1098672935)
     -- (0.938580509, 4.1003268419).. controls (0.9371906527, 4.1000993002) and 
    (0.9344410233, 4.099541294) .. (0.9342220277, 4.0960136837) -- (0.9342220277, 
    4.07814907).. controls (0.9342220277, 4.0782707783) and (0.9343569652, 
    4.0777487554) .. (0.9348337444, 4.0773558492).. controls (0.9352266506, 
    4.0771891617) and (0.9365334277, 4.0765779742) .. (0.9396681317, 4.07727912) 
    -- (0.9626769598, 4.0847270879).. controls (0.9637646619, 4.0849837337) and 
    (0.9658646333, 4.0855052275) .. (0.9658646333, 4.0879450562).. controls 
    (0.9658646333, 4.0879450562) and (0.9658646333, 4.095921794) .. (0.9658646333,
     4.0961861392).. controls (0.9658646333, 4.0961861392) and (0.9656370917, 
    4.0965713725) .. (0.9651243292, 4.0967981204).. controls (0.9645956917, 
    4.0971833538) and (0.9635908042, 4.0971833538) .. (0.9622009744, 
    4.0967981204).. controls (0.9622009744, 4.0967981204) and (0.9503491227, 
    4.0930891383) .. (0.9470330469, 4.0920920296).. controls (0.947078026, 
    4.0909966546) and (0.947078026, 4.0906870921) .. (0.9472076719, 
    4.0893424796).. controls (0.9495645802, 4.0900827837) and (0.9519136569, 
    4.0909061671) .. (0.9519136569, 4.0909061671).. controls (0.9534016735, 
    4.0913895608) and (0.9549653875, 4.0906865629) .. (0.9554488606, 
    4.0891686483).. controls (0.955924846, 4.0877711192) and (0.9551842773, 
    4.0861999969) .. (0.953794421, 4.085595715) -- (0.9456439844, 4.0828536263).. 
    controls (0.9448584365, 4.0826710638) and (0.9439895448, 4.0828086471) .. 
    (0.943287076, 4.0832388596).. controls (0.9425015281, 4.0837222533) and 
    (0.9419877073, 4.0845078013) .. (0.9418972198, 4.0853764283) -- (0.941202424, 
    4.0938743687).. controls (0.9411124656, 4.0950073146) and (0.9418972198, 
    4.0963974354) .. (0.9431587531, 4.0967978294) -- (0.9607591275, 
    4.1021534194).. controls (0.9638561812, 4.1029768027) and (0.9663413067, 
    4.1026746485) .. (0.9682977417, 4.1015038673).. controls (0.9708282167, 
    4.0997136965) and (0.9712210965, 4.0971831421) .. (0.9713494723, 4.0967978823)
     -- (0.9713494723, 4.087944871).. controls (0.9713494723, 4.0819773527) and 
    (0.9666434873, 4.0797565725) .. (0.9640752033, 4.0793184225) -- (0.9411116983,
     4.0718703752).. controls (0.9367985137, 4.070910996) and (0.9337013806, 
    4.0712586585) .. (0.9314352773, 4.072950669).. controls (0.9289952898, 
    4.0748315919) and (0.9287763735, 4.0777473002) .. (0.9287763735, 4.0779663223)
     -- (0.9286467277, 4.0961406838).. controls (0.929039634, 4.1030675548) and 
    (0.9345764106, 4.1054167108) .. (0.9374921983, 4.1058095113) -- (0.9703435529,
     4.1153122333).. controls (0.9740071325, 4.1162263688) and (0.9767944385, 
    4.1158334625) .. (0.9788415198, 4.1140432917).. controls (0.9811529198, 
    4.1120869625) and (0.9814175031, 4.1093447679) .. (0.9814175031, 
    4.1081663931).. controls (0.9814175031, 4.1079917681) and (0.9814175031, 
    4.1077737515) .. (0.9814175031, 4.1077737515) -- (0.9814175031, 
    4.0801041556).. controls (0.9814175031, 4.0800221348) and (0.9814175031, 
    4.0797567577) .. (0.9814175031, 4.0797567577).. controls (0.9814175031, 
    4.0763197673) and (0.9800654823, 4.0743104421) .. (0.9789777802, 
    4.0732604696).. controls (0.9770137781, 4.0712586321) and (0.9747476748, 
    4.0710773925) .. (0.9736146231, 4.0710773925).. controls (0.9732217169, 
    4.0709953717) and (0.9394564121, 4.061499079) .. (0.9394564121, 4.061499079)..
     controls (0.9345766223, 4.0595803206) and (0.9342216573, 4.0570499515) .. 
    (0.9342216573, 4.0565740719).. controls (0.9342216573, 4.056492051) and 
    (0.9342216573, 4.0332706302) .. (0.9342216573, 4.0332706302).. controls 
    (0.9342216573, 4.0296901033) and (0.9352262802, 4.0285193221) .. 
    (0.9356189219, 4.028338056).. controls (0.9363213906, 4.0279377415) and 
    (0.9374093573, 4.028163431) .. (0.9377945906, 4.0282084102) -- (0.9701700656, 
    4.0359661258).. controls (0.9743548479, 4.0367895092) and (0.9755332227, 
    4.0384891925) .. (0.9755332227, 4.0394031163) -- (0.9755332227, 
    4.0563991029).. controls (0.9756628685, 4.0604932654) and (0.9747476748, 
    4.0614979148) .. (0.9747476748, 4.0614979148) -- (0.9470328617, 
    4.0536119027).. controls (0.9439886981, 4.052342961) and (0.9437695702, 
    4.0506884685) .. (0.9437695702, 4.0505147431) -- (0.9440341535, 
    4.0413672242).. controls (0.9440341535, 4.0412772658) and (0.9441637994, 
    4.0403623367) .. (0.9441637994, 4.0403623367).. controls (0.944375466, 
    4.0403623367) and (0.9450324265, 4.0404919825) .. (0.9455536556, 
    4.0406718992).. controls (0.9455536556, 4.0406718992) and (0.9601549779, 
    4.0441996152) .. (0.9634635925, 4.0451136448).. controls (0.9634635925, 
    4.0462465906) and (0.9634635925, 4.046631824) .. (0.9634635925, 
    4.0481200788).. controls (0.9598906327, 4.0469871329) and (0.9543613702, 
    4.0452873438) .. (0.9543613702, 4.0452873438).. controls (0.9529185973, 
    4.0448944375) and (0.9513096131, 4.04568025) .. (0.9509168127, 4.0471608583)..
     controls (0.9504786627, 4.0486414667) and (0.9512189669, 4.05016716) .. 
    (0.9526995752, 4.0506883892) -- (0.9655107796, 4.0546163933).. controls 
    (0.9663418358, 4.0548730392) and (0.9672104629, 4.05482806) .. (0.9679960108, 
    4.0542160788).. controls (0.9686984796, 4.0536117704) and (0.9691742004, 
    4.0528262225) .. (0.9691742004, 4.0518668433) -- (0.9691742004, 
    4.0429307529).. controls (0.9691742004, 4.0416768925) and (0.968298165, 
    4.0404907654) .. (0.9669083088, 4.0402717433) -- (0.9472079894, 
    4.0353919535).. controls (0.9469058352, 4.0351723494) and (0.943635056, 
    4.0341833369) .. (0.941029016, 4.0359660994).. controls (0.939714566, 
    4.0367894827) and (0.9384530327, 4.0382248473) .. (0.9384530327, 4.0412764456)
     -- (0.9382784077, 4.0505147167).. controls (0.9381963869, 4.0509076229) and 
    (0.9380588035, 4.0559986825) .. (0.9451599556, 4.0589220373) -- (0.9728294721,
     4.0665514829).. controls (0.97361502, 4.0669517975) and (0.9758358796, 
    4.0673293579) .. (0.9777998817, 4.0660678246).. controls (0.9801491171, 
    4.06463246) and (0.9813275183, 4.0614978883) .. (0.9811538194, 4.0561800544) 
    -- (0.9811538194, 4.0394937096).. controls (0.9811538194, 4.0394937096) and 
    (0.9817202923, 4.0326424565) .. (0.9713489167, 4.0303914079) -- (0.9393208396,
     4.0229357671).. controls (0.9393208396, 4.0228537463) and (0.9357478798, 
    4.0214477504) .. (0.9324846942, 4.023721315) -- cycle;
  
  
  
      \end{scope}
      \begin{scope}[shift={(0.8731702173, -3.5267130063)}]
        \path[fill=c005093,even odd rule] (0.8713078417, 4.0001105862).. controls 
    (0.8688226104, 4.0016816821) and (0.8674704573, 4.0048165713) .. 
    (0.8674704573, 4.0095225563) -- (0.8674704573, 4.0322669598).. controls 
    (0.8674704573, 4.032010314) and (0.8661636802, 4.0387633538) .. (0.8764066535,
     4.0435523915) -- (0.910912289, 4.0562804373).. controls (0.9112218515, 
    4.0564153748) and (0.9114787619, 4.056492104) .. (0.911780916, 4.056492104).. 
    controls (0.912876291, 4.0566270415) and (0.9147495675, 4.0571945727) .. 
    (0.9146589477, 4.0603293825) -- (0.9146589477, 4.0883918052).. controls 
    (0.9146589477, 4.0887921198) and (0.914447281, 4.0896607469) .. (0.9142737144,
     4.0898345781).. controls (0.9137450769, 4.0901745677) and (0.9126571102, 
    4.0901743031) .. (0.9107385371, 4.0895699948) -- (0.8775774083, 
    4.0771062677).. controls (0.8764897062, 4.0768041135) and (0.8732189006, 
    4.0760185656) .. (0.8730526894, 4.0722719069) -- (0.8730526894, 
    4.0543618906).. controls (0.8737474852, 4.0537046656) and (0.8750090185, 
    4.0531911094) .. (0.878234554, 4.054150224) -- (0.9013264612, 4.0636528402).. 
    controls (0.9020667654, 4.0639851569) and (0.9045141348, 4.0647783777) .. 
    (0.9045141348, 4.0673995783) -- (0.9045141348, 4.0758523015).. controls 
    (0.903939989, 4.076018989) and (0.903154441, 4.076320614) .. (0.9013264612, 
    4.075717364).. controls (0.9013264612, 4.075717364) and (0.8888627342, 
    4.07083768) .. (0.8856826277, 4.0695686854).. controls (0.8858175652, 
    4.0685714708) and (0.8859022319, 4.0682693167) .. (0.8859921902, 
    4.0670456187).. controls (0.8882583465, 4.0679597542) and (0.8905244762, 
    4.0689641125) .. (0.8905244762, 4.0689641125).. controls (0.8920050846, 
    4.0695684208) and (0.8935686398, 4.0687365708) .. (0.8942258119, 
    4.0674760958).. controls (0.8948377931, 4.0660409958) and (0.894143791, 
    4.0644319058) .. (0.892654716, 4.0638654065) -- (0.8846402487, 4.0603302027)..
     controls (0.8838091925, 4.0600735569) and (0.8828500779, 4.0601635152) .. 
    (0.8821097737, 4.0605498069).. controls (0.8813694696, 4.0609879569) and 
    (0.8808934842, 4.0619018277) .. (0.8808029967, 4.0626876402) -- (0.8800174487,
     4.071147851).. controls (0.8799354279, 4.0722733885) and (0.8805915946, 
    4.0734894135) .. (0.8817623758, 4.0739805596) -- (0.8992342421, 
    4.0809375931).. controls (0.9024596717, 4.0820782119) and (0.9050280615, 
    4.0820329681) .. (0.9070751427, 4.0807709056).. controls (0.9096056177, 
    4.0792450535) and (0.9100437942, 4.0766313406) .. (0.9101344404, 4.0763216988)
     -- (0.9101344404, 4.0674006367).. controls (0.9101344404, 4.061599356) and 
    (0.9055115081, 4.0589404523) .. (0.9029809802, 4.0583663594) -- (0.8801081215,
     4.0489165275).. controls (0.8757949369, 4.0476097504) and (0.8725694015, 
    4.0477835817) .. (0.8703033246, 4.0493546775).. controls (0.8676443944, 
    4.0510921962) and (0.8674705896, 4.0538416404) .. (0.8674705896, 4.0541512823)
     -- (0.8674705896, 4.0722728594).. controls (0.8677727437, 4.0792450271) and 
    (0.8731359273, 4.0820322273) .. (0.8760970381, 4.0824703773) -- (0.9088200169,
     4.0948812935).. controls (0.9126572954, 4.0961502352) and (0.9155353271, 
    4.0958406727) .. (0.9175824083, 4.0942769852).. controls (0.9199393167, 
    4.0924942227) and (0.9202413121, 4.08966162) .. (0.9202413121, 4.0885663244)..
     controls (0.9202413121, 4.0883916994) and (0.9202413121, 4.088090339) .. 
    (0.9202413121, 4.088090339) -- (0.9202413121, 4.0604209283).. controls 
    (0.9202413121, 4.06033097) and (0.9202413121, 4.0601642825) .. (0.9202413121, 
    4.0600735304).. controls (0.9202413121, 4.056674296) and (0.9189797787, 
    4.0546196477) .. (0.91783916, 4.0534488665).. controls (0.9159280746, 
    4.0515301081) and (0.9138809669, 4.0510919581) .. (0.9127856712, 
    4.051008879).. controls (0.9126110462, 4.0509189206) and (0.878536814, 
    4.0383790465) .. (0.878536814, 4.0383790465).. controls (0.8733549494, 
    4.0360147298) and (0.8730527687, 4.0333632079) .. (0.8730527687, 
    4.032758926).. controls (0.8730527687, 4.0326689677) and (0.8730527687, 
    4.0095235088) .. (0.8730527687, 4.0095235088).. controls (0.8730527687, 
    4.006607721) and (0.8737478292, 4.0052103506) .. (0.8743521375, 
    4.0048175237).. controls (0.8747902875, 4.0045608779) and (0.8758856625, 
    4.0048175237) .. (0.8764068917, 4.0050371279) -- (0.9087823931, 
    4.0154084506).. controls (0.9143570846, 4.017100461) and (0.9143570846, 
    4.0195554769) .. (0.9143570846, 4.0197217146) -- (0.9143570846, 
    4.0365893254).. controls (0.9144470429, 4.0406381648) and (0.9136620242, 
    4.0415068713) .. (0.9136620242, 4.0415068713) -- (0.8859924283, 
    4.0313092475).. controls (0.8827669987, 4.0295717287) and (0.8825479767, 
    4.0277815315) .. (0.8825479767, 4.0276531292) -- (0.8828501308, 
    4.0183696408).. controls (0.8828501308, 4.0183696408) and (0.8829850683, 
    4.0177653325) .. (0.8829850683, 4.0176745804).. controls (0.8829850683, 
    4.0179841429) and (0.8841558496, 4.0181505658) .. (0.8841558496, 
    4.0181505658).. controls (0.8841558496, 4.0181505658) and (0.8991501046, 
    4.0231209754) .. (0.9022848879, 4.0241256512).. controls (0.9022848879, 
    4.0252210262) and (0.9022848879, 4.0255985867) .. (0.9022848879, 
    4.0270035771).. controls (0.8985382556, 4.0255986396) and (0.893356391, 
    4.0235590196) .. (0.893356391, 4.0235590196).. controls (0.8920041056, 
    4.0229925467) and (0.890304634, 4.0236489779) .. (0.8897381346, 
    4.0250396279).. controls (0.8892621492, 4.0266033154) and (0.8898201554, 
    4.0281744113) .. (0.8913092304, 4.0286578844) -- (0.9041128677, 
    4.0338851779).. controls (0.9048984156, 4.0341418237) and (0.9059030385, 
    4.0340518654) .. (0.906688851, 4.0334922717).. controls (0.9073836469, 
    4.0330541217) and (0.907904876, 4.0320947425) .. (0.907904876, 4.0310977925) 
    -- (0.907904876, 4.0222070252).. controls (0.907904876, 4.0209380835) and 
    (0.9070740844, 4.019933196) .. (0.9059863823, 4.0195555827) -- (0.8862105508, 
    4.0129687542).. controls (0.8858176446, 4.0128391083) and (0.8825015687, 
    4.0112690708) .. (0.8797142627, 4.0129687542).. controls (0.8785358085, 
    4.0138450542) and (0.8772745398, 4.0152272375) .. (0.8772745398, 4.0183697467)
     -- (0.8769723856, 4.0275626152).. controls (0.8768824273, 4.0279629298) and 
    (0.8768824273, 4.0328804227) .. (0.8838085046, 4.0363249802) -- (0.9113949419,
     4.046386661).. controls (0.912263569, 4.0469986423) and (0.9145297517, 
    4.047391284) .. (0.9164937537, 4.046212036).. controls (0.918978985, 
    4.0448221798) and (0.920021364, 4.04167975) .. (0.9199383112, 4.0365885052) --
     (0.9199383112, 4.0199324552).. controls (0.9199383112, 4.0199324552) and 
    (0.9206331071, 4.0130584744) .. (0.9104355892, 4.0101426867) -- (0.878452941, 
    3.9999450629).. controls (0.8783232952, 3.9997254588) and (0.8746156625, 
    3.9980718129) .. (0.8713070479, 4.0001117504) -- cycle;
  
  
  
      \end{scope}
      \begin{scope}[shift={(0.8731702173, -3.5348560875)}]
        \path[fill=c005093,even odd rule] (0.9789792354, 4.2408124427).. controls 
    (0.9765392479, 4.2425499615) and (0.9753608731, 4.245828096) .. (0.9753608731,
     4.2502244127) -- (0.9753608731, 4.2724023962).. controls (0.9753608731, 
    4.2721002421) and (0.9738805294, 4.2791554625) .. (0.9844708477, 4.2825999935)
     -- (1.0189764831, 4.2901386077).. controls (1.0193693894, 4.2903502744) and 
    (1.0197244602, 4.2903502744) .. (1.019981106, 4.2903502744).. controls 
    (1.0202832602, 4.2901386077) and (1.0212426394, 4.2903502744) .. 
    (1.0217712769, 4.290750589).. controls (1.0223303415, 4.2913548973) and 
    (1.0223300769, 4.2922687681) .. (1.0223300769, 4.2931829035) -- (1.0223300769,
     4.320240571).. controls (1.0223300769, 4.3204601752) and (1.0223301563, 
    4.321555021) .. (1.0217712769, 4.322113821).. controls (1.0212423748, 
    4.3226048877) and (1.019890354, 4.3227332106) .. (1.0182360731, 4.3223784044) 
    -- (0.9849465685, 4.3148850075).. controls (0.983511204, 4.31475007) and 
    (0.9809807025, 4.3142806992) .. (0.9808523002, 4.3111382694) -- (0.9808523002,
     4.2937571288).. controls (0.9808523002, 4.2938867746) and (0.980981946, 
    4.2934020579) .. (0.9815021169, 4.2930924954).. controls (0.982023346, 
    4.292714935) and (0.9833753669, 4.2921029537) .. (0.9863440448, 4.292714935) 
    -- (1.0094813546, 4.298757939).. controls (1.0105690567, 4.298757939) and 
    (1.0123517133, 4.2991582535) .. (1.0123517133, 4.3013339223) -- (1.0123517133,
     4.3095750052).. controls (1.0123517133, 4.3097866719) and (1.01152833, 
    4.3108365385) .. (1.0086957273, 4.3102700656).. controls (1.0086957273, 
    4.3102700656) and (0.9971460298, 4.3073921398) .. (0.9935730965, 
    4.3065234069).. controls (0.9937027423, 4.3054735402) and (0.9937027423, 
    4.305171386) .. (0.9937927006, 4.3039098527).. controls (0.996141936, 
    4.3045594048) and (0.9986346285, 4.3052544652) .. (0.9986346285, 
    4.3052544652).. controls (1.0001530723, 4.305609536) and (1.0016334954, 
    4.304733236) .. (1.002071619, 4.303418786).. controls (1.002509769, 
    4.3018175277) and (1.0015506544, 4.3002538137) .. (1.0001531252, 4.2998534198)
     -- (0.9921385256, 4.2977232594).. controls (0.9913074694, 4.2974211052) and 
    (0.9904388423, 4.2975883219) .. (0.9897816173, 4.2981087573).. controls 
    (0.9890868215, 4.298592151) and (0.9886939152, 4.299158624) .. (0.9886031631, 
    4.3001635115) -- (0.9878250235, 4.3083971067).. controls (0.9877350652, 
    4.3097869629) and (0.9886028985, 4.3109201733) .. (0.9900004277, 4.3113128944)
     -- (1.0075177229, 4.3156260525).. controls (1.0105694535, 4.3164116004) and 
    (1.0131377375, 4.3158456567) .. (1.0148826381, 4.3145835942).. controls 
    (1.0174509485, 4.3129217462) and (1.0179343687, 4.3103608706) .. 
    (1.0179343687, 4.3100512023) -- (1.0179343687, 4.3013342398).. controls 
    (1.0179343687, 4.2954573677) and (1.0131376317, 4.2934026135) .. 
    (1.0104787279, 4.2931835915) -- (0.9875153023, 4.2871784229).. controls 
    (0.9832927904, 4.2864381187) and (0.980067255, 4.2870487771) .. (0.9779295275,
     4.2887873542).. controls (0.9754442962, 4.2906153604) and (0.9753612171, 
    4.2934027194) .. (0.9753612171, 4.2936746317) -- (0.9753612171, 
    4.3113128415).. controls (0.9755358421, 4.3182850092) and (0.9812456562, 
    4.3202413648) .. (0.984161444, 4.3203319846) -- (1.0170581481, 4.3278252756)..
     controls (1.0207670244, 4.3284825006) and (1.0235997594, 4.3280078381) .. 
    (1.025601491, 4.3262615881).. controls (1.0282150452, 4.3239046798) and 
    (1.027912891, 4.3204603075) .. (1.027912891, 4.3200750477) -- (1.027912891, 
    4.2932742642).. controls (1.0281324952, 4.2897466275) and (1.0267344369, 
    4.2877826254) .. (1.025601491, 4.2866495737).. controls (1.023509166, 
    4.2846477362) and (1.0209861523, 4.2845646571) .. (1.0198908567, 
    4.2846477362).. controls (1.0191960608, 4.2844731112) and (0.9859517471, 
    4.2772375244) .. (0.9859517471, 4.2772375244).. controls (0.9810266342, 
    4.2756285931) and (0.9807623154, 4.2732717642) .. (0.9807623154, 4.2728865044)
     -- (0.9808522738, 4.2502250742).. controls (0.9808522738, 4.2477850867) and 
    (0.9814187467, 4.2461309381) .. (0.9822498029, 4.2453831462).. controls 
    (0.98311843, 4.2447335942) and (0.9844707154, 4.2452085212) .. (0.9844707154, 
    4.2452085212) -- (1.0169292696, 4.2508359971).. controls (1.020985676, 
    4.2514932221) and (1.0221564573, 4.2529661575) .. (1.0221564573, 4.2537442971)
     -- (1.0221564573, 4.2701360283).. controls (1.0222464156, 4.2742301908) and 
    (1.0213709094, 4.2751441146) .. (1.0213709094, 4.2751441146) -- (0.9933992123,
     4.2691388402).. controls (0.9905665037, 4.2682172965) and (0.9904381015, 
    4.2666083652) .. (0.9904381015, 4.2666083652) -- (0.9907402556, 
    4.2576722748).. controls (0.9907402556, 4.2574606081) and (0.9908751931, 
    4.2566221435) .. (0.9908751931, 4.2566221435).. controls (0.9908751931, 
    4.2566221435) and (0.9917438202, 4.2566222494) .. (0.9922198056, 
    4.256804706).. controls (0.9922198056, 4.256804706) and (1.0065568885, 
    4.259335181) .. (1.010175145, 4.2601132148).. controls (1.010175145, 
    4.2612461606) and (1.010175145, 4.2615935585) .. (1.010175145, 4.2630817604)..
     controls (1.0064285127, 4.2621223812) and (1.0008462542, 4.2607251167) .. 
    (1.0008462542, 4.2607251167).. controls (0.9992825667, 4.2603248021) and 
    (0.9978020906, 4.2612460812) .. (0.9974092637, 4.2626891187).. controls 
    (0.9970163575, 4.2642528062) and (0.9980135721, 4.2658240079) .. 
    (0.9995015887, 4.2660355687) -- (1.0122222527, 4.2692684596).. controls 
    (1.0131363881, 4.2694801262) and (1.013921936, 4.2693584179) .. (1.0146619756,
     4.2687398221).. controls (1.0154022798, 4.2682638367) and (1.0156665985, 
    4.2673952096) .. (1.0156665985, 4.2666096617) -- (1.0156665985, 
    4.2578094083).. controls (1.0156665985, 4.2565026312) and (1.0147902985, 
    4.2552410979) .. (1.0134004423, 4.2550145617) -- (0.9935718265, 4.251449169)..
     controls (0.9933072431, 4.2513142315) and (0.9900441104, 4.2503085502) .. 
    (0.9874305033, 4.2521818002).. controls (0.9863428012, 4.2530959356) and 
    (0.9850812679, 4.2546289315) .. (0.9850812679, 4.2574617723) -- (0.9848537262,
     4.2664357246).. controls (0.9846791012, 4.2683544829) and (0.9860321804, 
    4.2724032429) .. (0.9919165402, 4.2744502977) -- (1.019586136, 4.2805537854)..
     controls (1.020371684, 4.2809390187) and (1.022547141, 4.28111285) .. 
    (1.0244205233, 4.2798134812).. controls (1.0267319233, 4.2781516333) and 
    (1.027864975, 4.2750621731) .. (1.0277365992, 4.2701370602) -- (1.0277365992, 
    4.2538811925).. controls (1.0277365992, 4.2538811925) and (1.0283030721, 
    4.2469997504) .. (1.017841156, 4.245337929) -- (0.985813079, 4.2397633433).. 
    controls (0.985813079, 4.2397633433) and (0.9822477127, 4.2385470537) .. 
    (0.97897696, 4.24081321) -- cycle;
  
  
  
      \end{scope}
      \begin{scope}[shift={(0.8731702173, -3.5301041708)}]
        \path[fill=c005093,even odd rule] (0.9189787469, 4.2224260442).. controls 
    (0.9164935156, 4.2240802192) and (0.9152321146, 4.2273434577) .. 
    (0.9152321146, 4.2319588229) -- (0.9152321146, 4.2545823913).. controls 
    (0.9152321146, 4.2543998287) and (0.9138800938, 4.2614639392) .. 
    (0.9241607173, 4.2653843763) -- (0.958756999, 4.2751891731).. controls 
    (0.9591499052, 4.275318819) and (0.9594973031, 4.275318819) .. (0.9597616219, 
    4.275318819).. controls (0.9597616219, 4.275318819) and (0.9609324031, 
    4.2753187396) .. (0.9615443844, 4.2759757794).. controls (0.9622016094, 
    4.2765877606) and (0.9623299323, 4.2777659502) .. (0.9623299323, 4.2787253294)
     -- (0.9623299323, 4.3066970265).. controls (0.9623299323, 4.3070822598) and 
    (0.9622002865, 4.3080871473) .. (0.961717951, 4.308351466).. controls 
    (0.9612419656, 4.3089557744) and (0.9598899448, 4.3089557744) .. 
    (0.9581903408, 4.3085710702) -- (0.9250292121, 4.2989400781).. controls 
    (0.9236845996, 4.2987125365) and (0.92089732, 4.2981545302) .. (0.9208066738, 
    4.2946268142) -- (0.9208066738, 4.2768074442).. controls (0.9208066738, 
    4.2768074442) and (0.9208066738, 4.2763616213) .. (0.9214109821, 
    4.2759763879).. controls (0.9216755654, 4.2758017629) and (0.9229445071, 
    4.27519084) .. (0.9261169671, 4.2759763879) -- (0.9491710919, 4.2834244617).. 
    controls (0.950258794, 4.2836361283) and (0.9521775258, 4.2842176825) .. 
    (0.9521775258, 4.286604674).. controls (0.9521775258, 4.286604674) and 
    (0.9521775258, 4.2946267348) .. (0.9521775258, 4.2948836188).. controls 
    (0.9521775258, 4.2948836188) and (0.9521325467, 4.2952310167) .. 
    (0.9516110529, 4.2954122562).. controls (0.9511276592, 4.2958051625) and 
    (0.9500399571, 4.2958051625) .. (0.9487405883, 4.2954122562).. controls 
    (0.9487405883, 4.2954122562) and (0.9367980904, 4.2917939998) .. 
    (0.9335726344, 4.2907062713).. controls (0.9335726344, 4.2896108963) and 
    (0.9337022802, 4.2893010692) .. (0.9337022802, 4.2879641825).. controls 
    (0.9360136802, 4.2887423221) and (0.938408371, 4.2896109492) .. (0.938408371, 
    4.2896109492).. controls (0.9398437356, 4.2901395867) and (0.9414146727, 
    4.2893013867) .. (0.9418981458, 4.2877376992).. controls (0.9423741312, 
    4.2863930867) and (0.9417155833, 4.2848218056) .. (0.9401532187, 4.2842175502)
     -- (0.9321763752, 4.2816870752).. controls (0.9313077481, 4.281294169) and 
    (0.9304388565, 4.2815124502) .. (0.9297363877, 4.2818987419).. controls 
    (0.9290413273, 4.2823821356) and (0.9285656065, 4.283258171) .. (0.9284748544,
     4.2840363106) -- (0.9276893065, 4.2924437106).. controls (0.9276072856, 
    4.2938412398) and (0.9283917752, 4.295012021) .. (0.9296533085, 4.2954123885) 
    -- (0.947208254, 4.3008585983).. controls (0.9503506308, 4.3015914942) and 
    (0.9529190206, 4.3013798275) .. (0.9547923765, 4.3001182942).. controls 
    (0.9573228515, 4.2985471983) and (0.9577533815, 4.2958052419) .. 
    (0.9577533815, 4.2956239494) -- (0.957888319, 4.2866047004).. controls 
    (0.957888319, 4.2806825052) and (0.9530010679, 4.2783256763) .. (0.9504781071,
     4.2779328494) -- (0.9276052219, 4.2705680665).. controls (0.9232844967, 
    4.2694803644) and (0.9201496075, 4.2699182498) .. (0.9179741504, 
    4.2716557685).. controls (0.9153981671, 4.2736197706) and (0.9152321675, 
    4.2763617535) .. (0.9152321675, 4.2765883427) -- (0.9152321675, 
    4.2948457569).. controls (0.915398855, 4.3018556806) and (0.9211089337, 
    4.3042502656) .. (0.9240700446, 4.3044239646) -- (0.9568383465, 
    4.3139266867).. controls (0.9604112798, 4.3148408221) and (0.9632893379, 
    4.3144479158) .. (0.9653363927, 4.3126651533).. controls (0.967602549, 
    4.3107011513) and (0.967912376, 4.307959089) .. (0.967912376, 4.3067882813).. 
    controls (0.967912376, 4.3066983229) and (0.967912376, 4.3064784542) .. 
    (0.967912376, 4.3063879667) -- (0.967912376, 4.2788545254).. controls 
    (0.967912376, 4.2786349213) and (0.967912376, 4.2785521067) .. (0.967912376, 
    4.2783256233).. controls (0.967912376, 4.2750170087) and (0.9666431698, 
    4.2729247367) .. (0.9654649802, 4.2718823313).. controls (0.9634631427, 
    4.2700921604) and (0.9613330881, 4.2696992542) .. (0.9601547398, 
    4.2696992542).. controls (0.959671346, 4.2696092958) and (0.9259889352, 
    4.2600682621) .. (0.9259889352, 4.2600682621).. controls (0.9210260927, 
    4.2582402558) and (0.9208070706, 4.2557549981) .. (0.9208070706, 4.2551960394)
     -- (0.9208070706, 4.2319605163).. controls (0.9208070706, 4.2283045038) and 
    (0.9216756977, 4.227345151) .. (0.9221969269, 4.2269145946).. controls 
    (0.9227710727, 4.22651428) and (0.923941854, 4.2267346779) .. (0.9242440081, 
    4.2268696154) -- (0.9567101294, 4.2346272252).. controls (0.9609326677, 
    4.2353675294) and (0.9620280427, 4.2371502919) .. (0.9620280427, 4.2381095387)
     -- (0.9620280427, 4.2549696883).. controls (0.962118001, 4.2591091475) and 
    (0.9613329823, 4.2600230713) .. (0.9613329823, 4.2600230713) -- (0.933572846, 
    4.2523183252).. controls (0.9304380627, 4.2510038752) and (0.930302014, 
    4.2492211921) .. (0.930302014, 4.2492211921) -- (0.9305216181, 4.240118864).. 
    controls (0.9305216181, 4.2398468723) and (0.9306962431, 4.239114241) .. 
    (0.9306962431, 4.239114241).. controls (0.930825889, 4.239114241) and 
    (0.9315272994, 4.239114241) .. (0.9319577765, 4.2393338452).. controls 
    (0.9319577765, 4.2393338452) and (0.9466045277, 4.242906805) .. (0.9500415181,
     4.2438208081).. controls (0.9500415181, 4.2447801873) and (0.9500415181, 
    4.2453844956) .. (0.9500415181, 4.2467440571).. controls (0.9462948594, 
    4.2455581946) and (0.94093919, 4.2440398302) .. (0.94093919, 4.2440398302).. 
    controls (0.9393225858, 4.2434204406) and (0.9377968131, 4.2444250635) .. 
    (0.9374115533, 4.2457847573).. controls (0.9368829158, 4.2473484448) and 
    (0.9376681992, 4.2489195406) .. (0.9392850679, 4.2492670444) -- (0.9520056525,
     4.2532327252).. controls (0.9529197879, 4.2535348794) and (0.9537958233, 
    4.2534443919) .. (0.9544453754, 4.2528324106).. controls (0.9551856796, 
    4.2523188544) and (0.9556238296, 4.2514499627) .. (0.9556238296, 4.2506570065)
     -- (0.9556238296, 4.2416378369).. controls (0.9556238296, 4.2402479806) and 
    (0.9547930379, 4.2391147702) .. (0.9534859962, 4.2389335042) -- (0.9337025977,
     4.2340084177).. controls (0.9333551998, 4.2338787719) and (0.9300390181, 
    4.2328450448) .. (0.9273876815, 4.2346278073).. controls (0.9263378148, 
    4.2353681115) and (0.924947694, 4.236931799) .. (0.924947694, 4.2398475337) --
     (0.9247280898, 4.2491311279).. controls (0.924646069, 4.2496145217) and 
    (0.9245534648, 4.2544943644) .. (0.931700284, 4.2575385544) -- (0.959369774, 
    4.2652131908).. controls (0.9601553219, 4.2656135054) and (0.9623308848, 
    4.2659910658) .. (0.9642042406, 4.2647750408).. controls (0.9667347156, 
    4.263203945) and (0.9676941212, 4.2600238385) .. (0.9676035015, 4.2549703498) 
    -- (0.9676035015, 4.2381102002).. controls (0.9676035015, 4.2381102002) and 
    (0.9682982973, 4.231221191) .. (0.9578892185, 4.2290986242) -- (0.9259443, 
    4.2215599042).. controls (0.9259443, 4.2215599042) and (0.922288314, 
    4.2201549667) .. (0.9189796994, 4.2224285312) -- cycle;
  
  
  
      \end{scope}
      \begin{scope}[shift={(0.8731702173, -3.5267058625)}]
        \path[fill=c005093,even odd rule] (0.8597123448, 4.1984585512).. controls 
    (0.8570534146, 4.1999767304) and (0.8557919077, 4.2031646421) .. 
    (0.8557919077, 4.2077422513) -- (0.8557919077, 4.2304790877).. controls 
    (0.8557919077, 4.2302224419) and (0.8545756181, 4.2370583227) .. 
    (0.8647279717, 4.2418928158) -- (0.8992337129, 4.2545378883).. controls 
    (0.8994083379, 4.254749555) and (0.8998078588, 4.254749555) .. (0.9001478483, 
    4.254749555).. controls (0.901197715, 4.2548395133) and (0.9030710973, 
    4.2553237008) .. (0.9028520752, 4.2586775592) -- (0.9028520752, 4.286520854)..
     controls (0.902981721, 4.2871251623) and (0.902852181, 4.2879107102) .. 
    (0.9025954294, 4.288046706).. controls (0.9020667919, 4.288394104) and 
    (0.9010243335, 4.2884396123) .. (0.8991508719, 4.2879117685) -- (0.8658991233,
     4.2753271269).. controls (0.8648114213, 4.2751525019) and (0.8616764792, 
    4.2742772602) .. (0.8614574571, 4.2705305221) -- (0.8614574571, 
    4.2525751827).. controls (0.8620694383, 4.2519632015) and (0.86333838, 
    4.2514874806) .. (0.8666015921, 4.2525751827) -- (0.8896028531, 
    4.2619872585).. controls (0.8905243969, 4.2622148002) and (0.8928358498, 
    4.2629918815) .. (0.8928358498, 4.2656508381) -- (0.8928358498, 
    4.2741487785).. controls (0.892261704, 4.2743683827) and (0.8915213998, 
    4.2746623348) .. (0.8896028531, 4.2740588202).. controls (0.8896028531, 
    4.2740588202) and (0.8771844492, 4.26908849) .. (0.8741402856, 4.2677892271)..
     controls (0.8742223065, 4.2669129271) and (0.8742223065, 4.2666107729) .. 
    (0.8743969315, 4.2652965875).. controls (0.8765800085, 4.2662107229) and 
    (0.8787177625, 4.2671698375) .. (0.8787177625, 4.2671698375).. controls 
    (0.8803266938, 4.2677892271) and (0.8818904077, 4.2670878167) .. 
    (0.8825022567, 4.2656969021).. controls (0.8830234858, 4.2643372083) and 
    (0.8823726108, 4.2625922813) .. (0.8809764046, 4.2619880258) -- (0.8730525306,
     4.2586794113).. controls (0.8721762306, 4.2583320133) and (0.8711716077, 
    4.2584598071) .. (0.8704313035, 4.2588460988).. controls (0.8696909994, 
    4.2592464133) and (0.8691245265, 4.2601150404) .. (0.8691245265, 4.2609913404)
     -- (0.8683389785, 4.2694893867).. controls (0.8683389785, 4.2706223325) and 
    (0.8689131244, 4.2717933783) .. (0.8700839056, 4.2722767721) -- (0.8876917413,
     4.2792489663).. controls (0.8907812279, 4.2803289954) and (0.8933571583, 
    4.2802385079) .. (0.8954419956, 4.2790293621).. controls (0.8979724706, 
    4.2774958371) and (0.8984559702, 4.2749276325) .. (0.8985390494, 4.2745348321)
     -- (0.8985390494, 4.2656516319).. controls (0.8985390494, 4.259940971) and 
    (0.8938330644, 4.2572819615) .. (0.8913101035, 4.2566247894) -- (0.868512704, 
    4.2472581425).. controls (0.8641163873, 4.2458230425) and (0.8608003115, 
    4.2458682863) .. (0.8585416958, 4.2474698092).. controls (0.8560943, 
    4.2493128967) and (0.8557921458, 4.2521757942) .. (0.8557921458, 4.2524930031)
     -- (0.8557921458, 4.2706221473).. controls (0.8560943, 4.2774961281) and 
    (0.8615481033, 4.280238111) .. (0.8644185679, 4.2807291248) -- (0.8972776219, 
    4.2931852848).. controls (0.9009336079, 4.294363739) and (0.9039400419, 
    4.2941443994) .. (0.9059870967, 4.2926185473).. controls (0.9082984967, 
    4.2907450327) and (0.9085628154, 4.2879125623) .. (0.9085628154, 
    4.2867870777).. controls (0.9085628154, 4.2865224944) and (0.9085628154, 
    4.2863941715) .. (0.9085628154, 4.286303684) -- (0.9085628154, 4.2587625698)..
     controls (0.9085628154, 4.2585429656) and (0.9085628154, 4.2584604156) .. 
    (0.9085628154, 4.2582865844).. controls (0.9085628154, 4.2549326731) and 
    (0.9073012821, 4.2528403746) .. (0.9061230925, 4.2517979692).. controls 
    (0.9042495779, 4.2498339671) and (0.9022026554, 4.2493127379) .. 
    (0.9011073333, 4.2493127379).. controls (0.9009327083, 4.2492677588) and 
    (0.866858476, 4.236720635) .. (0.866858476, 4.236720635).. controls 
    (0.8616766115, 4.2341446517) and (0.8613744308, 4.2316143619) .. 
    (0.8613744308, 4.2310931592) -- (0.8614564517, 4.2077443944).. controls 
    (0.8614564517, 4.2049116858) and (0.8619398454, 4.2035519392) .. 
    (0.8626727413, 4.2031666794).. controls (0.863156135, 4.2027663648) and 
    (0.8643345892, 4.2030370335) .. (0.8647272308, 4.2033413044) -- (0.8971028381,
     4.2137127065).. controls (0.9025944504, 4.2153293106) and (0.9026775296, 
    4.2179353506) .. (0.9026775296, 4.2179353506) -- (0.9026775296, 
    4.2349313373).. controls (0.9027674879, 4.2389424206) and (0.9019827338, 
    4.2398563444) .. (0.9018919817, 4.2398563444) -- (0.8743960848, 
    4.2295227777).. controls (0.8710874702, 4.2278230944) and (0.8709592002, 
    4.225995141) .. (0.8709592002, 4.2259120883) -- (0.8711708669, 4.2166662502)..
     controls (0.8711708669, 4.2165842294) and (0.8713058044, 4.2161071856) .. 
    (0.8713058044, 4.2159788627).. controls (0.871342846, 4.2161085085) and 
    (0.8725673377, 4.216364096) .. (0.8725673377, 4.216364096).. controls 
    (0.8725673377, 4.216364096) and (0.8874710523, 4.2213723675) .. (0.8906889148,
     4.2224676631).. controls (0.8906889148, 4.223472286) and (0.8906889148, 
    4.2239480069) .. (0.8906889148, 4.2253002658).. controls (0.8868590975, 
    4.2238574929) and (0.8816772329, 4.2218178729) .. (0.8816772329, 
    4.2218178729).. controls (0.8803249475, 4.2212058917) and (0.878534856, 
    4.2218179788) .. (0.8781420292, 4.2233363167).. controls (0.8774472333, 
    4.2247412542) and (0.8782319875, 4.2265240167) .. (0.8796300458, 4.2269545731)
     -- (0.8923959535, 4.2319703058).. controls (0.8932645806, 4.2323555392) and 
    (0.8942239598, 4.2323555392) .. (0.8950095077, 4.2317877433).. controls 
    (0.895832891, 4.2312286788) and (0.8963162848, 4.2303978871) .. (0.8963162848,
     4.2295215871) -- (0.8963162848, 4.2204117977).. controls (0.8963162848, 
    4.2191957727) and (0.8954854931, 4.2181607227) .. (0.8943522827, 4.217760461) 
    -- (0.8745311281, 4.2111810938).. controls (0.8742215656, 4.2110461563) and 
    (0.8709127658, 4.2095266542) .. (0.8681707829, 4.2114006979).. controls 
    (0.8669923288, 4.2119597625) and (0.8655950642, 4.2135308583) .. 
    (0.8655950642, 4.2165826683) -- (0.8653384183, 4.2259115856).. controls 
    (0.8653384183, 4.2261682315) and (0.8651188142, 4.2310934502) .. 
    (0.8721745373, 4.2346210869) -- (0.8997156515, 4.2446072556).. controls 
    (0.9005842785, 4.2452189723) and (0.9028504613, 4.2457325285) .. 
    (0.9048522194, 4.2445622765).. controls (0.9072996152, 4.2431647473) and 
    (0.9084327198, 4.2398941533) .. (0.9082967769, 4.2348029085) -- (0.9082967769,
     4.2180258381).. controls (0.9082967769, 4.2180258381) and (0.908818006, 
    4.2114012535) .. (0.8987563252, 4.2083495229) -- (0.8667736506, 
    4.1980612529).. controls (0.8666387131, 4.1980612529) and (0.8629362663, 
    4.1962784904) .. (0.8597108367, 4.1984615675) -- cycle;
  
  
  
      \end{scope}
      \begin{scope}[shift={(0.8731702173, -3.5863580271)}]
        \path[fill=c005093,even odd rule] (0.9074754044, 3.9753959365).. controls 
    (0.8960163531, 3.9734470156) and (0.8855089817, 3.9795806129) .. 
    (0.8838094042, 3.9892192779).. controls (0.8822835521, 3.9988428088) and 
    (0.8903055865, 4.0080433237) .. (0.901681585, 4.009916706).. controls 
    (0.9131406362, 4.0118807081) and (0.9237689223, 4.0058149765) .. 
    (0.9253325304, 3.9961838785).. controls (0.9269036262, 3.9865603742) and 
    (0.9189797523, 3.9773599385) .. (0.9074754044, 3.9753959365) -- 
    cycle(0.9067729356, 3.9987597296).. controls (0.9030716265, 3.99814034) and 
    (0.9004126169, 3.9949223452) .. (0.9009338196, 3.9916515925).. controls 
    (0.9014550487, 3.9884790531) and (0.9050279556, 3.9864244048) .. 
    (0.9087821815, 3.9869456075).. controls (0.9124834906, 3.9875575887) and 
    (0.9152331465, 3.9907375629) .. (0.9146590535, 3.9939252629).. controls 
    (0.9140547452, 3.9972716071) and (0.9106102142, 3.9993640115) .. 
    (0.9067729356, 3.9987597296) -- cycle;
  
  
  
      \end{scope}
      \begin{scope}[shift={(0.8731702173, -3.5881934417)}]
        \path[fill=c005093,even odd rule] (0.9614155058, 3.9987211798).. controls 
    (0.9498657025, 3.9979356319) and (0.9399325298, 4.005081525) .. (0.9392829513,
     4.0147050294).. controls (0.9385804825, 4.0243361273) and (0.9474258208, 
    4.0328793908) .. (0.9589755977, 4.0337480973).. controls (0.9705706977, 
    4.0345336452) and (0.9804586531, 4.0273877785) .. (0.9811535813, 
    4.0176283048).. controls (0.9818483771, 4.0080500971) and (0.9730030388, 
    3.9995898862) .. (0.9614155058, 3.9987211798) -- cycle(0.9629866017, 
    4.0219415687).. controls (0.9590585975, 4.0217669438) and (0.9560974867, 
    4.018844515) .. (0.9564525575, 4.0155887906).. controls (0.9566271825, 
    4.0122651477) and (0.9599347652, 4.0098781298) .. (0.9637645825, 
    4.0102180665).. controls (0.967601861, 4.0103926915) and (0.9705706183, 
    4.013315226) .. (0.9702155475, 4.0165784117).. controls (0.9700409225, 
    4.0198038412) and (0.9668238802, 4.0223343692) .. (0.9629866017, 4.0219415687)
     -- cycle;
  
  
  
      \end{scope}
      \begin{scope}[shift={(0.8731702173, -3.5862749479)}]
        \path[fill=c005093,even odd rule] (1.018891499, 4.0131792831).. controls 
    (1.0075155269, 4.0111322019) and (0.9969250763, 4.0173263094) .. 
    (0.9952254987, 4.0269045171).. controls (0.9935713238, 4.0363240542) and 
    (1.0015480879, 4.0458116421) .. (1.0130071392, 4.0476849979).. controls 
    (1.0244661904, 4.049649) and (1.0351019642, 4.0435153498) .. (1.0366655723, 
    4.033967331).. controls (1.0383197473, 4.024464609) and (1.0304337087, 
    4.0151432852) .. (1.018891499, 4.0131792831) -- cycle(1.018105951, 
    4.0363996456).. controls (1.0142686725, 4.0358408456) and (1.0117456058, 
    4.0326151515) .. (1.0123498877, 4.029397289).. controls (1.0129163606, 
    4.0262096154) and (1.0164516967, 4.0239888352) .. (1.0201530323, 
    4.0246459808).. controls (1.0238996646, 4.0253410412) and (1.0265585948, 
    4.0285210154) .. (1.0259467458, 4.03182963).. controls (1.0253802729, 
    4.0349719804) and (1.0219357683, 4.0371021408) .. (1.018105951, 4.0363996456) 
    -- cycle;
  
  
  
      \end{scope}
      \begin{scope}[shift={(0.8731702173, -3.5863580271)}]
        \path[fill=c005093,even odd rule] (1.4121516821, 3.9753952221).. controls 
    (1.4006397406, 3.9773592242) and (1.3925799238, 3.9865596598) .. 
    (1.3942795013, 3.9961831906).. controls (1.3958431888, 4.0058142621) and 
    (1.4064410213, 4.0118798879) .. (1.4178169933, 4.0100368004).. controls 
    (1.4294045263, 4.0080426358) and (1.4373360467, 3.9988420944) .. 
    (1.4357648715, 3.9891279437).. controls (1.4341106965, 3.9795798985) and 
    (1.423482331, 3.9734463542) .. (1.4121516821, 3.9753952221) -- 
    cycle(1.4128012342, 3.9987590152).. controls (1.4089639556, 3.9993633235) and 
    (1.4055723942, 3.9971802465) .. (1.4050436244, 3.9939245485).. controls 
    (1.4044771515, 3.990736875) and (1.4070454619, 3.9875567421) .. (1.4108373379,
     3.9869449196).. controls (1.4145839702, 3.9864236904) and (1.418119174, 
    3.9884784446) .. (1.4187309965, 3.9916509046).. controls (1.4191614735, 
    3.9949216573) and (1.4166309985, 3.9980565465) .. (1.4128012342, 3.9987590152)
     -- cycle;
  
  
  
      \end{scope}
      \begin{scope}[shift={(0.8731702173, -3.5880116729)}]
        \path[fill=c005093,even odd rule] (1.36073261, 4.0025351485).. controls 
    (1.34913751, 4.0033283694) and (1.340337389, 4.0118716065) .. (1.3409038883, 
    4.0215857835).. controls (1.3417272717, 4.0311715583) and (1.351713361, 
    4.0383627481) .. (1.3632555708, 4.0374864746).. controls (1.3748431038, 
    4.036708335) and (1.3835676333, 4.0281651244) .. (1.3829104877, 
    4.0184509738).. controls (1.3822985065, 4.0088198758) and (1.3722748198, 
    4.0015833629) .. (1.36073261, 4.0025351485) -- cycle(1.3592520017, 
    4.0258989417).. controls (1.3554600463, 4.0260735667) and (1.3520608119, 
    4.0236327854) .. (1.3518795458, 4.0204074087).. controls (1.3517128583, 
    4.0170534975) and (1.3546216346, 4.0141377098) .. (1.358375834, 
    4.0139111206).. controls (1.362296271, 4.0136089665) and (1.3655218065, 
    4.0160941977) .. (1.3657408285, 4.0193195744).. controls (1.3660429827, 
    4.0226735121) and (1.363081819, 4.0255061148) .. (1.3592520017, 4.0258989417) 
    -- cycle;
  
  
  
      \end{scope}
      \begin{scope}[shift={(0.8731702173, -3.5862297042)}]
        \path[fill=c005093,even odd rule] (1.3026511179, 4.015098359).. controls 
    (1.2911844996, 4.0171454402) and (1.2833438106, 4.0264742517) .. (1.284998065,
     4.0359693802).. controls (1.2865691608, 4.0456004781) and (1.2971520179, 
    4.0516963723) .. (1.3085733396, 4.0497399373).. controls (1.3201155494, 
    4.0476399394) and (1.3279563177, 4.0384092883) .. (1.3263927096, 
    4.0288612431).. controls (1.3247308617, 4.0193207915) and (1.3140573848, 
    4.0131343569) .. (1.3026511179, 4.015098359) -- cycle(1.3034745012, 
    4.0384092883).. controls (1.299727869, 4.0392326717) and (1.2961472627, 
    4.0370572675) .. (1.29562606, 4.033831785).. controls (1.2951125038, 
    4.0305231704) and (1.2977638933, 4.0273429581) .. (1.3014652025, 
    4.0266933796).. controls (1.3052118348, 4.02607399) and (1.308739445, 
    4.0281287442) .. (1.309351294, 4.031399444).. controls (1.3099177669, 
    4.0347004915) and (1.3073118856, 4.0377596833) .. (1.3034745012, 4.0384092883)
     -- cycle;
  
  
  
      \end{scope}
      \begin{scope}[shift={(0.8731702173, -3.5863201917)}]
        \path[fill=c005093,even odd rule] (1.4257463177, 4.1785917031).. controls 
    (1.4142796994, 4.1804649531) and (1.4063028558, 4.1897863033) .. 
    (1.4079193542, 4.1992890254).. controls (1.4095735292, 4.208784154) and 
    (1.4200808742, 4.2151066373) .. (1.4315777608, 4.2131049056).. controls 
    (1.4430443792, 4.2111861473) and (1.4509683325, 4.2019554963) .. 
    (1.4493971573, 4.1923621544).. controls (1.4478260615, 4.1829501844) and 
    (1.4371978813, 4.1766352681) .. (1.4257463177, 4.1785917031) -- 
    cycle(1.4265241927, 4.2018649558).. controls (1.4226869142, 4.2025597517) and 
    (1.419159304, 4.2003692663) .. (1.418728721, 4.1971588915).. controls 
    (1.4181170044, 4.1938805452) and (1.4207681294, 4.1907079) .. (1.4246055138, 
    4.190088484).. controls (1.4283068229, 4.1893936881) and (1.4318420267, 
    4.191622009) .. (1.4323179063, 4.194885221).. controls (1.4328843792, 
    4.1981182177) and (1.430270825, 4.2011624871) .. (1.4265241927, 4.2018649558) 
    -- cycle;
  
  
  
      \end{scope}
      \begin{scope}[shift={(0.8731702173, -3.5882312771)}]
        \path[fill=c005093,even odd rule] (1.3707042531, 4.1991538762).. controls 
    (1.3591620433, 4.2000301763) and (1.3503618958, 4.2085582792) .. 
    (1.3509208546, 4.218151621).. controls (1.3517140754, 4.227820475) and 
    (1.361601925, 4.2349285063) .. (1.3732801042, 4.2341807144).. controls 
    (1.3848676371, 4.233266579) and (1.3937131871, 4.2248593642) .. (1.3929804235,
     4.2152282662).. controls (1.3921872027, 4.2055141156) and (1.3822993531, 
    4.1983682225) .. (1.3707042531, 4.1991538762) -- cycle(1.3692688885, 
    4.2225101023).. controls (1.36543161, 4.2226847273) and (1.3620852394, 
    4.2203270252) .. (1.3619039733, 4.2170184635).. controls (1.3616843692, 
    4.2137477108) and (1.3645629829, 4.2108243825) .. (1.3684002615, 
    4.2105221754).. controls (1.3722754019, 4.2103105088) and (1.3756368008, 
    4.2127885963) .. (1.3757198535, 4.2160214871).. controls (1.3760220077, 
    4.2193225081) and (1.3731062994, 4.2222004604) .. (1.3692688885, 4.2225101023)
     -- cycle;
  
  
  
      \end{scope}
      \begin{scope}[shift={(0.8731702173, -3.5863201917)}]
        \path[fill=c005093,even odd rule] (1.3107037646, 4.2117458733).. controls 
    (1.2993353596, 4.2137098754) and (1.2914414894, 4.2231218454) .. 
    (1.2930579877, 4.2325338154).. controls (1.2946290835, 4.2421724804) and 
    (1.3053027721, 4.2482608075) .. (1.3166787442, 4.2463874517).. controls 
    (1.3281756573, 4.2443403704) and (1.3361977181, 4.2349737235) .. 
    (1.3345358967, 4.2254785685).. controls (1.33288913, 4.21598344) and 
    (1.3222988646, 4.2097894383) .. (1.3107037646, 4.2117458733) -- 
    cycle(1.3116176354, 4.2350567763).. controls (1.30782568, 4.2358044888) and 
    (1.3043434723, 4.2336290846) .. (1.3037316498, 4.2303961144).. controls 
    (1.3032556644, 4.227223575) and (1.3057861394, 4.223990446) .. (1.3095328246, 
    4.2233408675).. controls (1.3133701031, 4.2227288862) and (1.316844929, 
    4.22468548) .. (1.317464345, 4.2280544196).. controls (1.3179855742, 
    4.2312648208) and (1.3154095908, 4.2344071977) .. (1.3116176354, 4.2350567763)
     -- cycle;
  
  
  
      \end{scope}
      \begin{scope}[shift={(0.8731702173, -3.5863201917)}]
        \path[fill=c005093,even odd rule] (1.4363295717, 4.3745820454).. controls 
    (1.4248629533, 4.3763269725) and (1.4169314329, 4.3857465625) .. 
    (1.4185026346, 4.3953247702).. controls (1.4200814033, 4.4048728154) and 
    (1.4307548802, 4.4110516829) .. (1.4420402325, 4.4090952215).. controls 
    (1.4535371192, 4.4072219715) and (1.46155918, 4.3979384831) .. (1.4599955719, 
    4.388398005).. controls (1.4582884802, 4.3789028765) and (1.4477508669, 
    4.3725802873) .. (1.4363295717, 4.3745820454) -- cycle(1.4370243675, 
    4.3979383773).. controls (1.4333154648, 4.3985125231) and (1.4296594523, 
    4.396329446) .. (1.4293119485, 4.3932398535).. controls (1.4286999673, 
    4.3898331579) and (1.4313667027, 4.3867361042) .. (1.435067906, 
    4.3861318223).. controls (1.4388145648, 4.3854822702) and (1.4423422808, 
    4.3875290869) .. (1.4429088067, 4.3908378073).. controls (1.4435207879, 
    4.3940632369) and (1.4408540525, 4.3973340954) .. (1.4370243675, 4.3979383773)
     -- cycle;
  
  
  
      \end{scope}
      \begin{scope}[shift={(0.8731702173, -3.5881556063)}]
        \path[fill=c005093,even odd rule] (1.3830843719, 4.3995547698).. controls 
    (1.3714892719, 4.4005593927) and (1.3628175267, 4.4089667398) .. 
    (1.3634746988, 4.418688484).. controls (1.3640790071, 4.4284102017) and 
    (1.3740651229, 4.4354125583) .. (1.3856073327, 4.4346344981).. controls 
    (1.3972477558, 4.433894194) and (1.4060477975, 4.4252226075) .. (1.4053906254,
     4.4156745887).. controls (1.4046579942, 4.4060133283) and (1.394762551, 
    4.3987767096) .. (1.3830843719, 4.3995547698) -- cycle(1.381732351, 
    4.4229185894).. controls (1.377811914, 4.4233038227) and (1.374450515, 
    4.4207355123) .. (1.374276816, 4.4175931883).. controls (1.3740651494, 
    4.4141561979) and (1.377026366, 4.4113234894) .. (1.3807352688, 4.411013821)..
     controls (1.3846934619, 4.4107651127) and (1.3878810296, 4.4132873856) .. 
    (1.3881832367, 4.416512921).. controls (1.3883949033, 4.4197761067) and 
    (1.3854412538, 4.4226995673) .. (1.381732351, 4.4229185894) -- cycle;
  
  
  
      \end{scope}
      \begin{scope}[shift={(0.8731702173, -3.5863654354)}]
        \path[fill=c005093,even odd rule] (1.3231649517, 4.4122763598).. controls 
    (1.3116606038, 4.4142403619) and (1.3038650263, 4.4235691733) .. 
    (1.3055117135, 4.4330643283).. controls (1.3071735615, 4.4426953998) and 
    (1.3177639592, 4.4487459973) .. (1.3291323642, 4.4467895623).. controls 
    (1.3405914154, 4.444870804) and (1.348568259, 4.4354210515) .. (1.3469140046, 
    4.4260090815).. controls (1.3452143213, 4.4164233067) and (1.3346240029, 
    4.4103199248) .. (1.3231649517, 4.4122763598) -- cycle(1.3240412517, 
    4.43550421).. controls (1.3202492962, 4.4362897579) and (1.3166762571, 
    4.43406911) .. (1.3161928104, 4.4309266008).. controls (1.3155808292, 
    4.4276709029) and (1.3182398917, 4.4244378798) .. (1.3220772496, 
    4.4238336244).. controls (1.3258239083, 4.423297314) and (1.3293061954, 
    4.4252160723) .. (1.3299180444, 4.4285396094).. controls (1.3304392735, 
    4.4317953073) and (1.327787884, 4.4350283304) .. (1.3240412517, 4.43550421) --
     cycle;
  
  
  
      \end{scope}
      \begin{scope}[shift={(0.8731702173, -3.5863127833)}]
        \path[fill=c005093,even odd rule] (0.8906905552, 4.177625286).. controls 
    (0.8792389652, 4.1756234485) and (0.8687316996, 4.18193855) .. (0.8671680915, 
    4.1914412721).. controls (0.8655062435, 4.2009817238) and (0.8734376846, 
    4.2103936938) .. (0.884942059, 4.2122292142).. controls (0.8964010837, 
    4.2141027288) and (0.9069915344, 4.2079009219) .. (0.9085627096, 
    4.198413334).. controls (0.9101338054, 4.1889106383) and (0.9022024702, 
    4.1795439915) .. (0.8906905552, 4.177625286) -- cycle(0.8899124156, 
    4.2009817238).. controls (0.8862111065, 4.2003623342) and (0.8836351496, 
    4.1971518271) .. (0.8841563523, 4.1938734808).. controls (0.8847228252, 
    4.1907009415) and (0.8881296796, 4.1885632404) .. (0.89200482, 4.1891674958)..
     controls (0.8957061292, 4.1898699646) and (0.8983650594, 4.1928688315) .. 
    (0.8978891533, 4.1962830942).. controls (0.8972774367, 4.1994934954) and 
    (0.8938328527, 4.2016313023) .. (0.8899124156, 4.2009817238) -- cycle;
  
  
  
      \end{scope}
      \begin{scope}[shift={(0.8731702173, -3.5873174063)}]
        \path[fill=c005093,even odd rule] (0.9472987415, 4.197453796).. controls 
    (0.9357489646, 4.1960562669) and (0.9255136113, 4.20277163) .. (0.9243352365, 
    4.2124404575).. controls (0.9232853698, 4.2220715554) and (0.931700231, 
    4.2310076458) .. (0.943287764, 4.2322616121).. controls (0.9547921383, 
    4.2335305538) and (0.9650804877, 4.2269060221) .. (0.9661228931, 
    4.2172371946).. controls (0.9672105952, 4.2076438527) and (0.9587577927, 
    4.1987153294) .. (0.9472987415, 4.197453796) -- cycle(0.9478199706, 
    4.2208100221).. controls (0.9438617775, 4.22028112) and (0.9411121217, 
    4.2172371946) .. (0.9415049485, 4.21392858).. controls (0.9418978548, 
    4.2107862296) and (0.9452516073, 4.2084370471) .. (0.9490813981, 
    4.2088223069).. controls (0.9529186767, 4.2093057006) and (0.9557060885, 
    4.2123574048) .. (0.9553585848, 4.215711316).. controls (0.9549656785, 
    4.2188083698) and (0.9515288733, 4.2212859281) .. (0.9478199706, 4.2208100221)
     -- cycle;
  
  
  
      \end{scope}
      \begin{scope}[shift={(0.8731702173, -3.5863201917)}]
        \path[fill=c005093,even odd rule] (1.008604446, 4.2108768494).. controls 
    (0.997236041, 4.2087845244) and (0.9866002938, 4.2149784731) .. (0.9849460394,
     4.2245189248).. controls (0.9832918644, 4.2340971325) and (0.9912609556, 
    4.2433430765) .. (1.002727574, 4.2453825906).. controls (1.0141866252, 
    4.2472558406) and (1.0248149113, 4.2411675135) .. (1.0263407898, 
    4.2316648179).. controls (1.0279949648, 4.2220714496) and (1.0201088204, 
    4.2128408515) .. (1.008604446, 4.2108768494) -- cycle(1.0078188981, 
    4.2340971325).. controls (1.0041175625, 4.2335380679) and (1.0014662258, 
    4.230267421) .. (1.0020704812, 4.2270495585).. controls (1.0025086313, 
    4.2238165619) and (1.0060361621, 4.2217316188) .. (1.0098734406, 
    4.2223434677).. controls (1.0136200729, 4.2230308552) and (1.01623368, 
    4.2262639048) .. (1.0156671542, 4.2293911475).. controls (1.0150551729, 
    4.2327601135) and (1.0115277744, 4.2347996012) .. (1.0078188981, 4.2340971325)
     -- cycle;
  
  
  
      \end{scope}
      \begin{scope}[shift={(0.8731702173, -3.5863201917)}]
        \path[fill=c005093,even odd rule] (0.8762698904, 4.3708354925).. controls 
    (0.8647276806, 4.3688791633) and (0.854220415, 4.3751561383) .. (0.8525284046,
     4.3846512669).. controls (0.8510025525, 4.3941917185) and (0.8590171256, 
    4.4035206358) .. (0.8704308538, 4.4054770708).. controls (0.881889905, 
    4.4073505854) and (0.8924350325, 4.4011714798) .. (0.8940515308, 
    4.391623461).. controls (0.8957057058, 4.382120739) and (0.8876911856, 
    4.3727994946) .. (0.8762698904, 4.3708354925) -- cycle(0.8755295863, 
    4.3941917185).. controls (0.871699769, 4.3936329185) and (0.8691239444, 
    4.3904450863) .. (0.8696451471, 4.3871364717).. controls (0.8700832971, 
    4.383865719) and (0.873746956, 4.3817354792) .. (0.8774482917, 4.3824304073)..
     controls (0.8812855702, 4.3829892073) and (0.883808531, 4.3861318223) .. 
    (0.8833779481, 4.3894932213).. controls (0.8827659669, 4.3927110838) and 
    (0.879238489, 4.3947658115) .. (0.8755295863, 4.3941917185) -- cycle;
  
  
  
      \end{scope}
      \begin{scope}[shift={(0.8731702173, -3.5882312771)}]
        \path[fill=c005093,even odd rule] (0.93457776, 4.3917449842).. controls 
    (0.9228996073, 4.3910120883) and (0.9130041642, 4.398150626) .. (0.9122639394,
     4.407781724).. controls (0.9116143873, 4.417458145) and (0.920414376, 
    4.4260391381) .. (0.9319566123, 4.4268247919).. controls (0.943634765, 
    4.4275725044) and (0.9535756106, 4.4204191235) .. (0.9542705123, 
    4.4108409158).. controls (0.9549655727, 4.4011191981) and (0.9461199698, 
    4.3925306115) .. (0.93457776, 4.3917449842) -- cycle(0.9359223725, 
    4.415018131).. controls (0.9321757137, 4.4149281727) and (0.929259926, 
    4.4119665062) .. (0.9294336515, 4.4086578917).. controls (0.9297358056, 
    4.4053418158) and (0.9329613675, 4.4029472573) .. (0.936798646, 
    4.4031662794).. controls (0.9406284369, 4.4035591856) and (0.9435894419, 
    4.4063011685) .. (0.9432872612, 4.4096625675).. controls (0.9431126362, 
    4.4129635885) and (0.9397143279, 4.4153277994) .. (0.9359223725, 4.415018131) 
    -- cycle;
  
  
  
      \end{scope}
      \begin{scope}[shift={(0.8731702173, -3.58628315)}]
        \path[fill=c005093,even odd rule] (0.9948786035, 4.4052208219).. controls 
    (0.9832910706, 4.403271901) and (0.9727007523, 4.4094509008) .. (0.9710464979,
     4.418998946).. controls (0.9694298937, 4.4284562392) and (0.9774068167, 
    4.4379136381) .. (0.9888658679, 4.4397416708).. controls (1.000279596, 
    4.4417810792) and (1.0109607723, 4.4356400735) .. (1.0125244069, 
    4.4260014085).. controls (1.0141862548, 4.4165138206) and (1.0061639558, 
    4.4073131998) .. (0.9948786035, 4.4052208219) -- cycle(0.9939647327, 
    4.428456345).. controls (0.9901727773, 4.4278747908) and (0.9875136619, 
    4.4247549035) .. (0.9881255108, 4.4214312606).. controls (0.9886919838, 
    4.4182964508) and (0.9922197527, 4.4161210202) .. (0.9960570313, 
    4.4168159219).. controls (0.9998036635, 4.4175107177) and (1.0024172706, 
    4.4206532004) .. (1.0018507448, 4.4238710629).. controls (1.0012387635, 
    4.4270964925) and (0.997711365, 4.4292797283) .. (0.9939647327, 4.428456345) 
    -- cycle;
  
  
  
      \end{scope}
      \begin{scope}[shift={(0.8731702173, -3.614556525)}]
        \path[fill=c005093,even odd rule] (0.8862944238, 3.9662101323).. controls 
    (0.8862944238, 3.9689974913) and (0.8840282675, 3.971263621) .. (0.8811956119,
     3.971263621).. controls (0.878453629, 3.971263621) and (0.8761874463, 
    3.9689974648) .. (0.8761874463, 3.9662101323).. controls (0.8761874463, 
    3.963294371) and (0.8784536025, 3.9610282677) .. (0.8811956119, 
    3.9610282677).. controls (0.8840283204, 3.9610282677) and (0.8862944238, 
    3.963294424) .. (0.8862944238, 3.9662101323) -- cycle;
  
  
  
      \end{scope}
      \begin{scope}[shift={(0.8731702173, -3.6145943604)}]
        \path[fill=c005093,even odd rule] (0.8695159246, 3.9825414063).. controls 
    (0.8695159246, 3.985419385) and (0.8672497683, 3.9875646267) .. (0.864507759, 
    3.9875646267).. controls (0.861675024, 3.9875646267) and (0.8594089471, 
    3.985419385) .. (0.8594089471, 3.9825414063).. controls (0.8594089471, 
    3.9797087506) and (0.8616751033, 3.9773670029) .. (0.864507759, 
    3.9773670029).. controls (0.8672497419, 3.9773670029) and (0.8695159246, 
    3.97970883) .. (0.8695159246, 3.9825414063) -- cycle;
  
  
  
      \end{scope}
      \begin{scope}[shift={(0.8731702173, -3.6145943604)}]
        \path[fill=c005093,even odd rule] (0.8667739417, 4.0150073158).. controls 
    (0.8667739417, 4.0177492988) and (0.8645077854, 4.0201061277) .. 
    (0.8616751298, 4.0201061277).. controls (0.8589255798, 4.0201061277) and 
    (0.8565763179, 4.0177492194) .. (0.8565763179, 4.0150073158).. controls 
    (0.8565763179, 4.0121746073) and (0.8589255533, 4.0099085304) .. 
    (0.8616751298, 4.0099085304).. controls (0.8645078648, 4.0099085304) and 
    (0.8667739417, 4.0121746867) .. (0.8667739417, 4.0150073158) -- cycle;
  
  
  
      \end{scope}
      \begin{scope}[shift={(0.8731702173, -3.6145943604)}]
        \path[fill=c005093,even odd rule] (0.8649002683, 4.046424179).. controls 
    (0.8649002683, 4.0491284058) and (0.8626719475, 4.051394615) .. (0.8598845356,
     4.051394615).. controls (0.8570518006, 4.051394615) and (0.8547026446, 
    4.0491284588) .. (0.8547026446, 4.046424179).. controls (0.8547026446, 
    4.0435537144) and (0.85705188, 4.0411969912) .. (0.8598845356, 4.0411969912)..
     controls (0.8626718152, 4.0411969912) and (0.8649002683, 4.043553635) .. 
    (0.8649002683, 4.046424179) -- cycle;
  
  
  
      \end{scope}
      \begin{scope}[shift={(0.8731702173, -3.6146316667)}]
        \path[fill=c005093,even odd rule] (1.061768604, 4.0379254183).. controls 
    (1.061768604, 4.0407581269) and (1.0595024477, 4.0429411775) .. (1.0567528713,
     4.0429411775).. controls (1.0539201363, 4.0429411775) and (1.0516616265, 
    4.0407581004) .. (1.0516616265, 4.0379254183).. controls (1.0516616265, 
    4.0350473867) and (1.0539201098, 4.0327812833) .. (1.0567528713, 
    4.0327812833).. controls (1.0595024213, 4.0327812833) and (1.061768604, 
    4.0350474396) .. (1.061768604, 4.0379254183) -- cycle;
  
  
  
      \end{scope}
      \begin{scope}[shift={(0.8731702173, -3.6146396042)}]
        \path[fill=c005093,even odd rule] (1.0588905723, 4.0686019508).. controls 
    (1.0588905723, 4.0713439337) and (1.0567530035, 4.0736630067) .. 
    (1.0539202421, 4.0736630067).. controls (1.0510951006, 4.0736630067) and 
    (1.0488289973, 4.0713439338) .. (1.0488289973, 4.0686019508).. controls 
    (1.0488289973, 4.0657692423) and (1.0510951535, 4.063510706) .. (1.0539202421,
     4.063510706).. controls (1.0567529506, 4.063510706) and (1.0588905723, 
    4.0657691894) .. (1.0588905723, 4.0686019508) -- cycle;
  
  
  
      \end{scope}
      \begin{scope}[shift={(0.8731702173, -3.6147679271)}]
        \path[fill=c005093,even odd rule] (1.056191981, 4.1001084283).. controls 
    (1.056191981, 4.1028957079) and (1.0539184165, 4.1050334354) .. (1.0510931956,
     4.1050334354).. controls (1.0483436456, 4.1050334354) and (1.0460774365, 
    4.1028956021) .. (1.0460774365, 4.1001084283).. controls (1.0460774365, 
    4.0971926406) and (1.0483435927, 4.0950096165) .. (1.0510931956, 
    4.0950096165).. controls (1.0539183371, 4.0950096165) and (1.056191981, 
    4.0971926935) .. (1.056191981, 4.1001084283) -- cycle;
  
  
  
      \end{scope}
      \begin{scope}[shift={(0.8731702173, -3.6146316667)}]
        \path[fill=c005093,even odd rule] (1.2598960388, 4.0379254183).. controls 
    (1.2598960388, 4.0407581269) and (1.262162195, 4.0429411775) .. (1.2649948506,
     4.0429411775).. controls (1.2678275592, 4.0429411775) and (1.2700936625, 
    4.0407581004) .. (1.2700936625, 4.0379254183).. controls (1.2700936625, 
    4.0350473867) and (1.2678275063, 4.0327812833) .. (1.2649948506, 
    4.0327812833).. controls (1.2621621421, 4.0327812833) and (1.2598960388, 
    4.0350474396) .. (1.2598960388, 4.0379254183) -- cycle;
  
  
  
      \end{scope}
      \begin{scope}[shift={(0.8731702173, -3.6146396042)}]
        \path[fill=c005093,even odd rule] (1.2628579433, 4.0686019508).. controls 
    (1.2628579433, 4.0713439337) and (1.2649955121, 4.0736630067) .. 
    (1.2678282735, 4.0736630067).. controls (1.2705778235, 4.0736630067) and 
    (1.2728364392, 4.0713439338) .. (1.2728364392, 4.0686019508).. controls 
    (1.2728364392, 4.0657692423) and (1.2705776912, 4.063510706) .. (1.2678282735,
     4.063510706).. controls (1.264995565, 4.063510706) and (1.2628579433, 
    4.0657691894) .. (1.2628579433, 4.0686019508) -- cycle;
  
  
  
      \end{scope}
      \begin{scope}[shift={(0.8731702173, -3.6147679271)}]
        \path[fill=c005093,even odd rule] (1.2673501715, 4.1001084283).. controls 
    (1.2673501715, 4.1028957079) and (1.2696994069, 4.1050334354) .. 
    (1.2724489833, 4.1050334354).. controls (1.2751909663, 4.1050334354) and 
    (1.2775403075, 4.1028956021) .. (1.2775403075, 4.1001084283).. controls 
    (1.2775403075, 4.0971926406) and (1.2751910721, 4.0950096165) .. 
    (1.2724489833, 4.0950096165).. controls (1.2696994333, 4.0950096165) and 
    (1.2673501715, 4.0971926935) .. (1.2673501715, 4.1001084283) -- cycle;
  
  
  
      \end{scope}
      \begin{scope}[shift={(0.8731702173, -3.6145943604)}]
        \path[fill=c005093,even odd rule] (1.4478260615, 3.9825414063).. controls 
    (1.4478260615, 3.985419385) and (1.4501827052, 3.9875646267) .. (1.4529247675,
     3.9875646267).. controls (1.4557575025, 3.9875646267) and (1.4580312523, 
    3.985419385) .. (1.4580312523, 3.9825414063).. controls (1.4580312523, 
    3.9797087506) and (1.4557574231, 3.9773670029) .. (1.4529247675, 
    3.9773670029).. controls (1.450182811, 3.9773670029) and (1.4478260615, 
    3.97970883) .. (1.4478260615, 3.9825414063) -- cycle;
  
  
  
      \end{scope}
      \begin{scope}[shift={(0.8731702173, -3.6146848479)}]
        \path[fill=c005093,even odd rule] (1.4506585848, 4.0142291498).. controls 
    (1.4506585848, 4.0169711327) and (1.4529244765, 4.0192373154) .. 
    (1.4557572908, 4.0192373154).. controls (1.4586353225, 4.0192373154) and 
    (1.46077305, 4.0169711592) .. (1.46077305, 4.0142291498).. controls 
    (1.46077305, 4.0113964412) and (1.4586352167, 4.0091303379) .. (1.4557572908, 
    4.0091303379).. controls (1.4529245823, 4.0091303379) and (1.4506585848, 
    4.0113964942) .. (1.4506585848, 4.0142291498) -- cycle;
  
  
  
      \end{scope}
      \begin{scope}[shift={(0.8731702173, -3.6145943604)}]
        \path[fill=c005093,even odd rule] (1.4525322581, 4.046424179).. controls 
    (1.4525322581, 4.0491284058) and (1.4548889019, 4.051394615) .. (1.4576309642,
     4.051394615).. controls (1.4603805142, 4.051394615) and (1.4626014002, 
    4.0491284588) .. (1.4626014002, 4.046424179).. controls (1.4626014002, 
    4.0435537144) and (1.4603804877, 4.0411969912) .. (1.4576309642, 
    4.0411969912).. controls (1.4548889813, 4.0411969912) and (1.4525322581, 
    4.043553635) .. (1.4525322581, 4.046424179) -- cycle;
  
  
  
      \end{scope}
      \begin{scope}[shift={(0.8731702173, -3.6146017688)}]
        \path[fill=c005093,even odd rule] (0.9186248931, 3.9791119829).. controls 
    (0.9186248931, 3.981906909) and (0.9163209015, 3.9842032277) .. (0.9135260813,
     3.9842032277).. controls (0.9107387752, 3.9842032277) and (0.9084348365, 
    3.981906909) .. (0.9084348365, 3.9791119829).. controls (0.9084348365, 
    3.9762793273) and (0.9107388281, 3.974013171) .. (0.9135260813, 3.974013171)..
     controls (0.9163209544, 3.974013171) and (0.9186248931, 3.9762793273) .. 
    (0.9186248931, 3.9791119829) -- cycle;
  
  
  
      \end{scope}
      \begin{scope}[shift={(0.8731702173, -3.6146769104)}]
        \path[fill=c005093,even odd rule] (0.954271015, 3.9883342202).. controls 
    (0.954271015, 3.9912575221) and (0.9520048587, 3.9934330321) .. (0.9491722296,
     3.9934330321).. controls (0.9464226796, 3.9934330321) and (0.9441640375, 
    3.9912576279) .. (0.9441640375, 3.9883342202).. controls (0.9441640375, 
    3.9855921844) and (0.9464227854, 3.983318461) .. (0.9491722296, 3.983318461)..
     controls (0.9520049381, 3.983318461) and (0.954271015, 3.9855922902) .. 
    (0.954271015, 3.9883342202) -- cycle;
  
  
  
      \end{scope}
      \begin{scope}[shift={(0.8731702173, -3.6146017688)}]
        \path[fill=c005093,even odd rule] (0.9896049815, 3.9969765967).. controls 
    (0.9896049815, 3.9997261996) and (0.9872105023, 4.0020679208) .. (0.98446844, 
    4.0020679208).. controls (0.98171889, 4.0020679208) and (0.9793696281, 
    3.9997260938) .. (0.9793696281, 3.9969765967).. controls (0.9793696281, 
    3.9942345608) and (0.9817188635, 3.9918777848) .. (0.98446844, 3.9918777848)..
     controls (0.9872104229, 3.9918777848) and (0.9896049815, 3.9942346931) .. 
    (0.9896049815, 3.9969765967) -- cycle;
  
  
  
      \end{scope}
      \begin{scope}[shift={(0.8731702173, -3.6146840542)}]
        \path[fill=c005093,even odd rule] (1.0244665873, 4.0057764267).. controls 
    (1.0244665873, 4.0085184096) and (1.0222456748, 4.0107845923) .. 
    (1.0194961777, 4.0107845923).. controls (1.0166634692, 4.0107845923) and 
    (1.01426899, 4.008518436) .. (1.01426899, 4.0057764267).. controls 
    (1.01426899, 4.0029437181) and (1.0166634692, 4.0006776148) .. (1.0194961777, 
    4.0006776148).. controls (1.0222457277, 4.0006776148) and (1.0244665873, 
    4.002943771) .. (1.0244665873, 4.0057764267) -- cycle;
  
  
  
      \end{scope}
      \begin{scope}[shift={(0.8731702173, -3.6147679271)}]
        \path[fill=c005093,even odd rule] (1.0641164635, 4.0127417417).. controls 
    (1.0641164635, 4.0155744502) and (1.0617672281, 4.0177574744) .. 
    (1.0590252188, 4.0177574744).. controls (1.0561924838, 4.0177574744) and 
    (1.0539189456, 4.0155743973) .. (1.0539189456, 4.0127417417).. controls 
    (1.0539189456, 4.0099090331) and (1.0561925102, 4.007733576) .. (1.0590252188,
     4.007733576).. controls (1.0617672017, 4.007733576) and (1.0641164635, 
    4.0099089802) .. (1.0641164635, 4.0127417417) -- cycle;
  
  
  
      \end{scope}
      \begin{scope}[shift={(0.8731702173, -3.6156061271)}]
        \path[fill=c005093,even odd rule] (0.8550059629, 4.190534201).. controls 
    (0.8550059629, 4.1931855377) and (0.8529136379, 4.1951041373) .. 
    (0.8503905977, 4.1951041373).. controls (0.8478222873, 4.1951041373) and 
    (0.8456846127, 4.1931856435) .. (0.8456846127, 4.190534201).. controls 
    (0.8456846127, 4.1879658906) and (0.8478221815, 4.1859187565) .. 
    (0.8503905977, 4.1859187565).. controls (0.8529136644, 4.1859187565) and 
    (0.8550059629, 4.1879658377) .. (0.8550059629, 4.190534201) -- cycle;
  
  
  
      \end{scope}
      \begin{scope}[shift={(0.8731702173, -3.6154632521)}]
        \path[fill=c005093,even odd rule] (0.8534422225, 4.2194963129).. controls 
    (0.8534422225, 4.2220193796) and (0.8513874683, 4.2241192452) .. 
    (0.8488268838, 4.2241192452).. controls (0.8462964088, 4.2241192452) and 
    (0.8441133317, 4.2220192473) .. (0.8441133317, 4.2194963129).. controls 
    (0.8441133317, 4.2169732463) and (0.8462964088, 4.2147903279) .. 
    (0.8488268838, 4.2147903279).. controls (0.8513875213, 4.2147903279) and 
    (0.8534422225, 4.216973405) .. (0.8534422225, 4.2194963129) -- cycle;
  
  
  
      \end{scope}
      \begin{scope}[shift={(0.8731702173, -3.6153719708)}]
        \path[fill=c005093,even odd rule] (0.8524358533, 4.2460778892).. controls 
    (0.8524358533, 4.2486461996) and (0.8504718513, 4.2507837683) .. 
    (0.8478205146, 4.2507837683).. controls (0.8452900396, 4.2507837683) and 
    (0.8432429054, 4.2486461996) .. (0.8432429054, 4.2460778892).. controls 
    (0.8432429054, 4.2435474142) and (0.8452899867, 4.2413642313) .. 
    (0.8478205146, 4.2413642313).. controls (0.8504718513, 4.2413642313) and 
    (0.8524358533, 4.2435473083) .. (0.8524358533, 4.2460778892) -- cycle;
  
  
  
      \end{scope}
      \begin{scope}[shift={(0.8731702173, -3.6153793792)}]
        \path[fill=c005093,even odd rule] (1.4624283098, 4.1923544815).. controls 
    (1.4624283098, 4.1948775481) and (1.4646113869, 4.1970605723) .. 
    (1.4671342948, 4.1970605723).. controls (1.4697102781, 4.1970605723) and 
    (1.4717496335, 4.1948774952) .. (1.4717496335, 4.1923544815).. controls 
    (1.4717496335, 4.189786171) and (1.4697102252, 4.1876484965) .. (1.4671342948,
     4.1876484965).. controls (1.4646112281, 4.1876484965) and (1.4624283098, 
    4.1897860652) .. (1.4624283098, 4.1923544815) -- cycle;
  
  
  
      \end{scope}
      \begin{scope}[shift={(0.8731702173, -3.6154246229)}]
        \path[fill=c005093,even odd rule] (1.4641360365, 4.2214665856).. controls 
    (1.4641360365, 4.2240801398) and (1.4661754448, 4.2261726765) .. 
    (1.4687060785, 4.2261726765).. controls (1.471274389, 4.2261726765) and 
    (1.4734120635, 4.2240800869) .. (1.4734120635, 4.2214665856).. controls 
    (1.4734120635, 4.2189813544) and (1.4712742302, 4.2168060296) .. 
    (1.4687060785, 4.2168060296).. controls (1.4661756035, 4.2168060296) and 
    (1.4641360365, 4.2189814337) .. (1.4641360365, 4.2214665856) -- cycle;
  
  
  
      \end{scope}
      \begin{scope}[shift={(0.8731702173, -3.6154703958)}]
        \path[fill=c005093,even odd rule] (1.4650876104, 4.2483513479).. controls 
    (1.4650876104, 4.2508818229) and (1.4670442042, 4.2529667131) .. 
    (1.4695746792, 4.2529667131).. controls (1.4721506625, 4.2529667131) and 
    (1.4742806642, 4.2508817965) .. (1.4742806642, 4.2483513479).. controls 
    (1.4742806642, 4.245647121) and (1.4721505038, 4.2436453629) .. (1.4695746792,
     4.2436453629).. controls (1.4670442042, 4.2436453629) and (1.4650876104, 
    4.2456472004) .. (1.4650876104, 4.2483513479) -- cycle;
  
  
  
      \end{scope}
      \begin{scope}[shift={(0.8731702173, -3.6154254167)}]
        \path[fill=c005093,even odd rule] (1.2683130431, 4.2409413213).. controls 
    (1.2683130431, 4.2434717963) and (1.2703524515, 4.2456021154) .. 
    (1.2729736256, 4.2456021154).. controls (1.2755871798, 4.2456021154) and 
    (1.2776797165, 4.243471955) .. (1.2776797165, 4.2409413213).. controls 
    (1.2776797165, 4.2382824175) and (1.2755871269, 4.2362353363) .. 
    (1.2729736256, 4.2362353363).. controls (1.2703523985, 4.2362353363) and 
    (1.2683130431, 4.2382824175) .. (1.2683130431, 4.2409413213) -- cycle;
  
  
  
      \end{scope}
      \begin{scope}[shift={(0.8731702173, -3.6155156396)}]
        \path[fill=c005093,even odd rule] (1.2699210748, 4.2699629908).. controls 
    (1.2699210748, 4.272621921) and (1.2721041519, 4.2746689758) .. (1.274498684, 
    4.2746689758).. controls (1.2770669944, 4.2746689758) and (1.2791140227, 
    4.2726218946) .. (1.2791140227, 4.2699629908).. controls (1.2791140227, 
    4.2674399242) and (1.2770669415, 4.2653929752) .. (1.274498684, 
    4.2653929752).. controls (1.2721042048, 4.2653929752) and (1.2699210748, 
    4.2674400565) .. (1.2699210748, 4.2699629908) -- cycle;
  
  
  
      \end{scope}
      \begin{scope}[shift={(0.8731702173, -3.6154703958)}]
        \path[fill=c005093,even odd rule] (1.2707951258, 4.2968090975).. controls 
    (1.2707951258, 4.2993774079) and (1.27297053, 4.3014245685) .. (1.2755011108, 
    4.3014245685).. controls (1.2780694213, 4.3014245685) and (1.280116476, 
    4.2993774873) .. (1.280116476, 4.2968090975).. controls (1.280116476, 
    4.2942333788) and (1.2780693948, 4.2921031125) .. (1.2755011108, 
    4.2921031125).. controls (1.2729706358, 4.2921031125) and (1.2707951258, 
    4.2942332729) .. (1.2707951258, 4.2968090975) -- cycle;
  
  
  
      \end{scope}
      \begin{scope}[shift={(0.8731702173, -3.6155156396)}]
        \path[fill=c005093,even odd rule] (1.0532263458, 4.2420368285).. controls 
    (1.0532263458, 4.2446881652) and (1.0511792646, 4.2466068706) .. 
    (1.0485204667, 4.2466068706).. controls (1.0459521563, 4.2466068706) and 
    (1.0439050221, 4.2446881123) .. (1.0439050221, 4.2420368285).. controls 
    (1.0439050221, 4.2393779248) and (1.0459521033, 4.2373308435) .. 
    (1.0485204667, 4.2373308435).. controls (1.0511793704, 4.2373308435) and 
    (1.0532263458, 4.2393779248) .. (1.0532263458, 4.2420368285) -- cycle;
  
  
  
      \end{scope}
      \begin{scope}[shift={(0.8731702173, -3.6154632521)}]
        \path[fill=c005093,even odd rule] (1.0526569096, 4.2701294402).. controls 
    (1.0526569096, 4.2726977506) and (1.0504812408, 4.274752346) .. (1.0479508188,
     4.274752346).. controls (1.0453825083, 4.274752346) and (1.0432448338, 
    4.2726975919) .. (1.0432448338, 4.2701294402).. controls (1.0432448338, 
    4.2675688027) and (1.0453824025, 4.2654234552) .. (1.0479508188, 
    4.2654234552).. controls (1.0504812938, 4.2654234552) and (1.0526569096, 
    4.2675686969) .. (1.0526569096, 4.2701294402) -- cycle;
  
  
  
      \end{scope}
      \begin{scope}[shift={(0.8731702173, -3.6154703958)}]
        \path[fill=c005093,even odd rule] (1.0516626319, 4.2977681592).. controls 
    (1.0516626319, 4.3003364696) and (1.0496155506, 4.3023836038) .. 
    (1.0470396202, 4.3023836038).. controls (1.044426066, 4.3023836038) and 
    (1.0423337146, 4.3003365225) .. (1.0423337146, 4.2977681592).. controls 
    (1.0423337146, 4.2952376842) and (1.0444260396, 4.2930621742) .. 
    (1.0470396202, 4.2930621742).. controls (1.0496156035, 4.2930621742) and 
    (1.0516626319, 4.2952375783) .. (1.0516626319, 4.2977681592) -- cycle;
  
  
  
      \end{scope}
      \begin{scope}[shift={(0.8731702173, -3.6155981896)}]
        \path[fill=c005093,even odd rule] (0.8368748873, 4.4170647096).. controls 
    (0.8368748873, 4.4197236133) and (0.8348730498, 4.4216801542) .. 
    (0.8322141725, 4.4216801542).. controls (0.8296911058, 4.4216801542) and 
    (0.8275082404, 4.4197235604) .. (0.8275082404, 4.4170647096).. controls 
    (0.8275082404, 4.4144511554) and (0.8296913175, 4.4124872062) .. 
    (0.8322141725, 4.4124872062).. controls (0.8348730763, 4.4124872062) and 
    (0.8368748873, 4.4144512083) .. (0.8368748873, 4.4170647096) -- cycle;
  
  
  
      \end{scope}
      \begin{scope}[shift={(0.8731702173, -3.6155981896)}]
        \path[fill=c005093,even odd rule] (0.8401106621, 4.395407426).. controls 
    (0.8401106621, 4.3980663563) and (0.837927585, 4.4001134375) .. (0.8354046771,
     4.4001134375).. controls (0.8327457733, 4.4001134375) and (0.8306986392, 
    4.3980663563) .. (0.8306986392, 4.395407426).. controls (0.8306986392, 
    4.3928391156) and (0.8327457204, 4.3909203573) .. (0.8354046771, 
    4.3909203573).. controls (0.8379277438, 4.3909203573) and (0.8401106621, 
    4.3928391156) .. (0.8401106621, 4.395407426) -- cycle;
  
  
  
      \end{scope}
      \begin{scope}[shift={(0.8731702173, -3.6154254167)}]
        \path[fill=c005093,even odd rule] (1.0354890231, 4.4794219213).. controls 
    (1.0354890231, 4.4820431483) and (1.0334419419, 4.4841280121) .. 
    (1.0307830381, 4.4841280121).. controls (1.0282147277, 4.4841280121) and 
    (1.0261676729, 4.4820430954) .. (1.0261676729, 4.4794219213).. controls 
    (1.0261676729, 4.4768459379) and (1.0282147542, 4.4747613387) .. 
    (1.0307830381, 4.4747613387).. controls (1.0334419419, 4.4747613387) and 
    (1.0354890231, 4.4768459908) .. (1.0354890231, 4.4794219213) -- cycle;
  
  
  
      \end{scope}
      \begin{scope}[shift={(0.8731702173, -3.6156368188)}]
        \path[fill=c005093,even odd rule] (1.0387148496, 4.4571530004).. controls 
    (1.0387148496, 4.4598043371) and (1.0364486933, 4.4617230425) .. 
    (1.0338803829, 4.4617230425).. controls (1.0313573163, 4.4617230425) and 
    (1.0293027738, 4.4598042842) .. (1.0293027738, 4.4571530004).. controls 
    (1.0293027738, 4.4544865296) and (1.0313575279, 4.45256793) .. (1.0338803829, 
    4.45256793).. controls (1.0364486933, 4.45256793) and (1.0387148496, 
    4.4544864237) .. (1.0387148496, 4.4571530004) -- cycle;
  
  
  
      \end{scope}
      \begin{scope}[shift={(0.8731702173, -3.6155084958)}]
        \path[fill=c005093,even odd rule] (1.2908830333, 4.4818169031).. controls 
    (1.2908830333, 4.4843852135) and (1.2888359521, 4.4865228881) .. 
    (1.2863054242, 4.4865228881).. controls (1.2837371137, 4.4865228881) and 
    (1.281554116, 4.4843850548) .. (1.281554116, 4.4818169031).. controls 
    (1.281554116, 4.4792938365) and (1.2837371931, 4.4772393998) .. (1.2863054242,
     4.4772393998).. controls (1.2888358992, 4.4772393998) and (1.2908830333, 
    4.4792938894) .. (1.2908830333, 4.4818169031) -- cycle;
  
  
  
      \end{scope}
      \begin{scope}[shift={(0.8731702173, -3.6155989833)}]
        \path[fill=c005093,even odd rule] (1.2885244846, 4.4585961173).. controls 
    (1.2885244846, 4.4611644277) and (1.2863414075, 4.463083186) .. (1.2838183938,
     4.463083186).. controls (1.2812879188, 4.463083186) and (1.2791124088, 
    4.4611644277) .. (1.2791124088, 4.4585961173).. controls (1.2791124088, 
    4.456020134) and (1.2812878129, 4.4538901323) .. (1.2838183938, 
    4.4538901323).. controls (1.2863414604, 4.4538901323) and (1.2885244846, 
    4.4560202927) .. (1.2885244846, 4.4585961173) -- cycle;
  
  
  
      \end{scope}
      \begin{scope}[shift={(0.8731702173, -3.6154703958)}]
        \path[fill=c005093,even odd rule] (1.4792065973, 4.4161969027).. controls 
    (1.4792065973, 4.4188558329) and (1.481163191, 4.4209028877) .. (1.483693666, 
    4.4209028877).. controls (1.4863525698, 4.4209028877) and (1.4883995452, 
    4.4188558065) .. (1.4883995452, 4.4161969027).. controls (1.4883995452, 
    4.4136362652) and (1.486352464, 4.4115814581) .. (1.483693666, 4.4115814581)..
     controls (1.481163191, 4.4115814581) and (1.4792065973, 4.4136362123) .. 
    (1.4792065973, 4.4161969027) -- cycle;
  
  
  
      \end{scope}
      \begin{scope}[shift={(0.8731702173, -3.6154703958)}]
        \path[fill=c005093,even odd rule] (1.4769382715, 4.3937163417).. controls 
    (1.4769382715, 4.3962394083) and (1.4791213485, 4.3983317069) .. 
    (1.4816442565, 4.3983317069).. controls (1.4842202398, 4.3983317069) and 
    (1.486259701, 4.3962393819) .. (1.486259701, 4.3937163417).. controls 
    (1.486259701, 4.3910574379) and (1.4842200281, 4.3890103567) .. (1.4816442565,
     4.3890103567).. controls (1.4791211898, 4.3890103567) and (1.4769382715, 
    4.3910574379) .. (1.4769382715, 4.3937163417) -- cycle;
  
  
  
      \end{scope}
      \begin{scope}[shift={(0.8731702173, -3.6154246229)}]
        \path[fill=c005093,even odd rule] (0.8346167744, 4.4403754273).. controls 
    (0.8346167744, 4.4429437377) and (0.8325696931, 4.4449983331) .. (0.829993895,
     4.4449983331).. controls (0.8273349913, 4.4449983331) and (0.8252878571, 
    4.442943579) .. (0.8252878571, 4.4403754273).. controls (0.8252878571, 
    4.4377692815) and (0.8273349383, 4.4356315804) .. (0.829993895, 
    4.4356315804).. controls (0.8325698783, 4.4356315804) and (0.8346167744, 
    4.4377694137) .. (0.8346167744, 4.4403754273) -- cycle;
  
  
  
      \end{scope}
      \begin{scope}[shift={(0.8731702173, -3.6154246229)}]
        \path[fill=c005093,even odd rule] (1.4904903092, 4.4403754273).. controls 
    (1.4904903092, 4.4429437377) and (1.4883979842, 4.4449983331) .. 
    (1.4858674033, 4.4449983331).. controls (1.4833067658, 4.4449983331) and 
    (1.4811614183, 4.442943579) .. (1.4811614183, 4.4403754273).. controls 
    (1.4811614183, 4.4377692815) and (1.48330666, 4.4356315804) .. (1.4858674033, 
    4.4356315804).. controls (1.4883978783, 4.4356315804) and (1.4904903092, 
    4.4377694137) .. (1.4904903092, 4.4403754273) -- cycle;
  
  
  
      \end{scope}
      \begin{scope}[shift={(0.8731702173, -3.6154703958)}]
        \path[fill=c005093,even odd rule] (0.8664646967, 4.4524472271).. controls 
    (0.8664646967, 4.4551136979) and (0.8644176154, 4.4571532121) .. 
    (0.8617586058, 4.4571532121).. controls (0.8592355392, 4.4571532121) and 
    (0.8571885638, 4.4551138037) .. (0.8571885638, 4.4524472271).. controls 
    (0.8571885638, 4.4500149125) and (0.859235645, 4.4478318619) .. (0.8617586058,
     4.4478318619).. controls (0.8644175096, 4.4478318619) and (0.8664646967, 
    4.450014939) .. (0.8664646967, 4.4524472271) -- cycle;
  
  
  
      \end{scope}
      \begin{scope}[shift={(0.8731702173, -3.6155002938)}]
        \path[fill=c005093,even odd rule] (0.8992331044, 4.4643741675).. controls 
    (0.8992331044, 4.4669046425) and (0.8971407794, 4.4689594496) .. 
    (0.8946177656, 4.4689594496).. controls (0.8920494552, 4.4689594496) and 
    (0.8899117542, 4.4669046954) .. (0.8899117542, 4.4643741675).. controls 
    (0.8899117542, 4.4617228308) and (0.8920495875, 4.4596681825) .. 
    (0.8946177656, 4.4596681825).. controls (0.8971408323, 4.4596681825) and 
    (0.8992331044, 4.4617229367) .. (0.8992331044, 4.4643741675) -- cycle;
  
  
  
      \end{scope}
      \begin{scope}[shift={(0.8731702173, -3.6155907813)}]
        \path[fill=c005093,even odd rule] (0.9320928463, 4.4768000327).. controls 
    (0.9320928463, 4.4794589629) and (0.9300380921, 4.4815059383) .. 
    (0.9274775075, 4.4815059383).. controls (0.9248185773, 4.4815059383) and 
    (0.9227715225, 4.4794588571) .. (0.9227715225, 4.4768000327).. controls 
    (0.9227715225, 4.4742695577) and (0.9248186038, 4.4723055027) .. 
    (0.9274775075, 4.4723055027).. controls (0.930038145, 4.4723055027) and 
    (0.9320928463, 4.4742695048) .. (0.9320928463, 4.4768000327) -- cycle;
  
  
  
      \end{scope}
      \begin{scope}[shift={(0.8731702173, -3.6154703958)}]
        \path[fill=c005093,even odd rule] (0.9658657181, 4.4874891994).. controls 
    (0.9658657181, 4.490012266) and (0.963682641, 4.4920138919) .. (0.9611597331, 
    4.4920138919).. controls (0.9585837498, 4.4920138919) and (0.9564537481, 
    4.4900120544) .. (0.9564537481, 4.4874891994).. controls (0.9564537481, 
    4.4848302692) and (0.9585839085, 4.4826924623) .. (0.9611597331, 
    4.4826924623).. controls (0.9636827998, 4.4826924623) and (0.9658657181, 
    4.4848302956) .. (0.9658657181, 4.4874891994) -- cycle;
  
  
  
      \end{scope}
      \begin{scope}[shift={(0.8731702173, -3.6154703958)}]
        \path[fill=c005093,even odd rule] (0.9986323267, 4.4969010635).. controls 
    (0.9986323267, 4.4995524002) and (0.9964492496, 4.5016070485) .. 
    (0.9939186688, 4.5016070485).. controls (0.9913051146, 4.5016070485) and 
    (0.9892580863, 4.4995522944) .. (0.9892580863, 4.4969010635).. controls 
    (0.9892580863, 4.494461076) and (0.9913051675, 4.4922856983) .. (0.9939186688,
     4.4922856983).. controls (0.9964491438, 4.4922856983) and (0.9986323267, 
    4.4944611025) .. (0.9986323267, 4.4969010635) -- cycle;
  
  
  
      \end{scope}
      \begin{scope}[shift={(0.8731702173, -3.6153793792)}]
        \path[fill=c005093,even odd rule] (1.0313551996, 4.5044395719).. controls 
    (1.0313551996, 4.5070078823) and (1.0293004454, 4.5091456627) .. 
    (1.0267322938, 4.5091456627).. controls (1.0242018187, 4.5091456627) and 
    (1.0221546846, 4.5070078294) .. (1.0221546846, 4.5044395719).. controls 
    (1.0221546846, 4.501901424) and (1.0242017658, 4.4997335869) .. (1.0267322938,
     4.4997335869).. controls (1.0293006042, 4.4997335869) and (1.0313551996, 
    4.5019015827) .. (1.0313551996, 4.5044395719) -- cycle;
  
  
  
      \end{scope}
      \begin{scope}[shift={(0.8731702173, -3.6153801729)}]
        \path[fill=c005093,even odd rule] (1.0641187654, 4.5099693371).. controls 
    (1.0641187654, 4.5125376475) and (1.0620716842, 4.5146754279) .. 
    (1.0595034267, 4.5146754279).. controls (1.0569729517, 4.5146754279) and 
    (1.0549258175, 4.5125375946) .. (1.0549258175, 4.5099693371).. controls 
    (1.0549258175, 4.5073933538) and (1.0569728988, 4.5052633521) .. 
    (1.0595034267, 4.5052633521).. controls (1.0620717371, 4.5052633521) and 
    (1.0641187654, 4.5073935125) .. (1.0641187654, 4.5099693371) -- cycle;
  
  
  
      \end{scope}
      \begin{scope}[shift={(0.8731702173, -3.6155537396)}]
        \path[fill=c005093,even odd rule] (1.096974274, 4.515407874).. controls 
    (1.096974274, 4.5178931052) and (1.0949271927, 4.5199400806) .. (1.0923513681,
     4.5199400806).. controls (1.089737814, 4.5199400806) and (1.0876453831, 
    4.5178929994) .. (1.0876453831, 4.515407874).. controls (1.0876453831, 
    4.5127112142) and (1.0897377081, 4.510701889) .. (1.0923513681, 4.510701889)..
     controls (1.0949273515, 4.510701889) and (1.096974274, 4.5127111348) .. 
    (1.096974274, 4.515407874) -- cycle;
  
  
  
      \end{scope}
      \begin{scope}[shift={(0.8731702173, -3.6154254167)}]
        \path[fill=c005093,even odd rule] (1.1315738894, 4.5198114931).. controls 
    (1.1315738894, 4.52251572) and (1.1295268081, 4.524517584) .. (1.1269962802, 
    4.524517584).. controls (1.12433735, 4.524517584) and (1.1222902952, 
    4.5225157465) .. (1.1222902952, 4.5198114931).. controls (1.1222902952, 
    4.5172884265) and (1.1243373765, 4.5151509106) .. (1.1269962802, 
    4.5151509106).. controls (1.1295267552, 4.5151509106) and (1.1315738894, 
    4.5172884794) .. (1.1315738894, 4.5198114931) -- cycle;
  
  
  
      \end{scope}
      \begin{scope}[shift={(0.8731702173, -3.6156061271)}]
        \path[fill=c005093,even odd rule] (1.1644676302, 4.5201735225).. controls 
    (1.1644676302, 4.5228248592) and (1.162420549, 4.5248720462) .. (1.1597616452,
     4.5248720462).. controls (1.1572385785, 4.5248720462) and (1.1551463065, 
    4.522824965) .. (1.1551463065, 4.5201735225).. controls (1.1551463065, 
    4.5176504558) and (1.1572386315, 4.5156866654) .. (1.1597616452, 
    4.5156866654).. controls (1.1624205754, 4.5156866654) and (1.1644676302, 
    4.5176506675) .. (1.1644676302, 4.5201735225) -- cycle;
  
  
  
      \end{scope}
      \begin{scope}[shift={(0.8731702173, -3.6155608833)}]
        \path[fill=c005093,even odd rule] (1.197273794, 4.51934326).. controls 
    (1.197273794, 4.5219115704) and (1.1951436335, 4.5238680848) .. (1.192575376, 
    4.5238680848).. controls (1.190044901, 4.5238680848) and (1.187869391, 
    4.521911491) .. (1.187869391, 4.51934326).. controls (1.187869391, 
    4.5168201933) and (1.1900447952, 4.5146373808) .. (1.192575376, 
    4.5146373808).. controls (1.1951436865, 4.5146373808) and (1.197273794, 
    4.5168201933) .. (1.197273794, 4.51934326) -- cycle;
  
  
  
      \end{scope}
      \begin{scope}[shift={(0.8731702173, -3.6155989833)}]
        \path[fill=c005093,even odd rule] (1.229091739, 4.5164583756).. controls 
    (1.229091739, 4.5189814423) and (1.2270446577, 4.5209453385) .. (1.224385754, 
    4.5209453385).. controls (1.2218100352, 4.5209453385) and (1.2197628481, 
    4.5189813365) .. (1.2197628481, 4.5164583756).. controls (1.2197628481, 
    4.513807039) and (1.2218099294, 4.5117523906) .. (1.224385754, 4.5117523906)..
     controls (1.2270446842, 4.5117523906) and (1.229091739, 4.5138071448) .. 
    (1.229091739, 4.5164583756) -- cycle;
  
  
  
      \end{scope}
      \begin{scope}[shift={(0.8731702173, -3.6155077021)}]
        \path[fill=c005093,even odd rule] (1.2609395554, 4.5115323631).. controls 
    (1.2609395554, 4.5141912669) and (1.2589001471, 4.5162383481) .. 
    (1.2563695398, 4.5162383481).. controls (1.2537106096, 4.5162383481) and 
    (1.2516635548, 4.5141912669) .. (1.2516635548, 4.5115323631).. controls 
    (1.2516635548, 4.5090092965) and (1.253710636, 4.506954754) .. (1.2563695398, 
    4.506954754).. controls (1.2589000148, 4.506954754) and (1.2609395554, 
    4.5090095081) .. (1.2609395554, 4.5115323631) -- cycle;
  
  
  
      \end{scope}
      \begin{scope}[shift={(0.8731702173, -3.6155537396)}]
        \path[fill=c005093,even odd rule] (1.2938431123, 4.5066989019).. controls 
    (1.2938431123, 4.509320129) and (1.2916677081, 4.5113141348) .. (1.2891371273,
     4.5113141348).. controls (1.2865688169, 4.5113141348) and (1.2844235752, 
    4.5093199702) .. (1.2844235752, 4.5066989019).. controls (1.2844235752, 
    4.5041229185) and (1.2865688169, 4.5020759696) .. (1.2891371273, 
    4.5020759696).. controls (1.2916676023, 4.5020759696) and (1.2938431123, 
    4.5041230508) .. (1.2938431123, 4.5066989019) -- cycle;
  
  
  
      \end{scope}
      \begin{scope}[shift={(0.8731702173, -3.6155537396)}]
        \path[fill=c005093,even odd rule] (1.3266022331, 4.5012897338).. controls 
    (1.3266022331, 4.5038202087) and (1.3244265644, 4.5058220992) .. 
    (1.3218961688, 4.5058220992).. controls (1.3193355313, 4.5058220992) and 
    (1.3172354804, 4.5038202617) .. (1.3172354804, 4.5012897338).. controls 
    (1.3172354804, 4.4986383706) and (1.3193354783, 4.4965837487) .. 
    (1.3218961688, 4.4965837487).. controls (1.3244266438, 4.4965837487) and 
    (1.3266022331, 4.4986385029) .. (1.3266022331, 4.5012897338) -- cycle;
  
  
  
      \end{scope}
      \begin{scope}[shift={(0.8731702173, -3.6155537396)}]
        \path[fill=c005093,even odd rule] (1.3593733925, 4.4916669967).. controls 
    (1.3593733925, 4.49424298) and (1.3572884758, 4.4961993885) .. (1.3546220844, 
    4.4961993885).. controls (1.3520614469, 4.4961993885) and (1.3500066133, 
    4.4942427948) .. (1.3500066133, 4.4916669967).. controls (1.3500066133, 
    4.4891365217) and (1.3520613675, 4.4869610117) .. (1.3546220844, 
    4.4869610117).. controls (1.3572885552, 4.4869610117) and (1.3593733925, 
    4.4891364158) .. (1.3593733925, 4.4916669967) -- cycle;
  
  
  
      \end{scope}
      \begin{scope}[shift={(0.8731702173, -3.6155156396)}]
        \path[fill=c005093,even odd rule] (1.3930987188, 4.4834324754).. controls 
    (1.3930987188, 4.4859629504) and (1.3909233146, 4.4880025175) .. 
    (1.3883927338, 4.4880025175).. controls (1.3858244233, 4.4880025175) and 
    (1.3837773685, 4.4859631092) .. (1.3837773685, 4.4834324754).. controls 
    (1.3837773685, 4.480864165) and (1.3858244498, 4.4787264904) .. (1.3883927338,
     4.4787264904).. controls (1.3909232088, 4.4787264904) and (1.3930987188, 
    4.4808643237) .. (1.3930987188, 4.4834324754) -- cycle;
  
  
  
      \end{scope}
      \begin{scope}[shift={(0.8731702173, -3.615591575)}]
        \path[fill=c005093,even odd rule] (1.4258678408, 4.472306905).. controls 
    (1.4258678408, 4.4748828883) and (1.4238284325, 4.4768014615) .. 
    (1.4211619352, 4.4768014615).. controls (1.4185936248, 4.4768014615) and 
    (1.4166296492, 4.4748827031) .. (1.4166296492, 4.472306905).. controls 
    (1.4166296492, 4.4697462675) and (1.4185936513, 4.46760092) .. (1.4211619352, 
    4.46760092).. controls (1.4238284325, 4.46760092) and (1.4258678408, 
    4.4697461617) .. (1.4258678408, 4.472306905) -- cycle;
  
  
  
      \end{scope}
      \begin{scope}[shift={(0.8731702173, -3.6153801729)}]
        \path[fill=c005093,even odd rule] (1.4595965008, 4.4576746).. controls 
    (1.4595965008, 4.460205075) and (1.4574210967, 4.4623806908) .. (1.4548905158,
     4.4623806908).. controls (1.4523222054, 4.4623806908) and (1.4502750448, 
    4.4602050221) .. (1.4502750448, 4.4576746).. controls (1.4502750448, 
    4.4550988813) and (1.452322126, 4.452968615) .. (1.4548905158, 4.452968615).. 
    controls (1.4574209908, 4.452968615) and (1.4595965008, 4.4550987754) .. 
    (1.4595965008, 4.4576746) -- cycle;
  
  
  
      \end{scope}
      \begin{scope}[shift={(0.8731702173, -3.614556525)}]
        \path[fill=c005093,even odd rule] (1.0973672067, 4.0181495869).. controls 
    (1.0973672067, 4.0208538402) and (1.0951010504, 4.0232483987) .. 
    (1.0922683948, 4.0232483987).. controls (1.0895188448, 4.0232483987) and 
    (1.0871695829, 4.0208539196) .. (1.0871695829, 4.0181495869).. controls 
    (1.0871695829, 4.0151431794) and (1.0895188183, 4.0130130454) .. 
    (1.0922683948, 4.0130130454).. controls (1.0951011033, 4.0130130454) and 
    (1.0973672067, 4.0151432058) .. (1.0973672067, 4.0181495869) -- cycle;
  
  
  
      \end{scope}
      \begin{scope}[shift={(0.8731702173, -3.6146769104)}]
        \path[fill=c005093,even odd rule] (1.1313092267, 4.0213675552).. controls 
    (1.1313092267, 4.0240264854) and (1.1291338225, 4.0262926681) .. (1.126301061,
     4.0262926681).. controls (1.1234683525, 4.0262926681) and (1.1212022492, 
    4.0240265119) .. (1.1212022492, 4.0213675552).. controls (1.1212022492, 
    4.018451794) and (1.1234684054, 4.0161781235) .. (1.126301061, 4.0161781235)..
     controls (1.1291337696, 4.0161781235) and (1.1313092267, 4.0184516881) .. 
    (1.1313092267, 4.0213675552) -- cycle;
  
  
  
      \end{scope}
      \begin{scope}[shift={(0.8731702173, -3.6145943604)}]
        \path[fill=c005093,even odd rule] (1.1664336696, 4.021889684).. controls 
    (1.1664336696, 4.024639234) and (1.1641675133, 4.0269960629) .. (1.1613348577,
     4.0269960629).. controls (1.1585853077, 4.0269960629) and (1.1563265862, 
    4.0246391546) .. (1.1563265862, 4.021889684).. controls (1.1563265862, 
    4.0190116523) and (1.1585853342, 4.0167984392) .. (1.1613348577, 
    4.0167984392).. controls (1.1641675662, 4.0167984392) and (1.1664336696, 
    4.0190116788) .. (1.1664336696, 4.021889684) -- cycle;
  
  
  
      \end{scope}
      \begin{scope}[shift={(0.8731702173, -3.6146769104)}]
        \path[fill=c005093,even odd rule] (1.201250164, 4.0211938563).. controls 
    (1.201250164, 4.0240265913) and (1.1989840077, 4.0262926681) .. (1.1961513521,
     4.0262926681).. controls (1.1934093692, 4.0262926681) and (1.1910525402, 
    4.0240265119) .. (1.1910525402, 4.0211938563).. controls (1.1910525402, 
    4.0184518733) and (1.1934094485, 4.0161781235) .. (1.1961513521, 
    4.0161781235).. controls (1.1989840606, 4.0161781235) and (1.201250164, 
    4.0184519527) .. (1.201250164, 4.0211938563) -- cycle;
  
  
  
      \end{scope}
      \begin{scope}[shift={(0.8731702173, -3.614556525)}]
        \path[fill=c005093,even odd rule] (1.2364573421, 4.0181495869).. controls 
    (1.2364573421, 4.0209823219) and (1.2342364296, 4.0232483987) .. 
    (1.2314491765, 4.0232483987).. controls (1.2286996265, 4.0232483987) and 
    (1.2263429033, 4.0209822425) .. (1.2263429033, 4.0181495869).. controls 
    (1.2263429033, 4.0153623073) and (1.2286995471, 4.0130130454) .. 
    (1.2314491765, 4.0130130454).. controls (1.2342364825, 4.0130130454) and 
    (1.2364573421, 4.0153622808) .. (1.2364573421, 4.0181495869) -- cycle;
  
  
  
      \end{scope}
      \begin{scope}[shift={(0.8731702173, -3.6145935667)}]
        \path[fill=c005093,even odd rule] (1.2681312744, 4.0143950965).. controls 
    (1.2681312744, 4.017182376) and (1.2658651181, 4.019539205) .. (1.2630324625, 
    4.019539205).. controls (1.2602904796, 4.019539205) and (1.2579412177, 
    4.0171822967) .. (1.2579412177, 4.0143950965).. controls (1.2579412177, 
    4.0116077904) and (1.2602904531, 4.0093415813) .. (1.2630324625, 
    4.0093415813).. controls (1.265865171, 4.0093415813) and (1.2681312744, 
    4.0116077375) .. (1.2681312744, 4.0143950965) -- cycle;
  
  
  
      \end{scope}
      \begin{scope}[shift={(0.8731702173, -3.6146848479)}]
        \path[fill=c005093,even odd rule] (1.3087387571, 4.0075140513).. controls 
    (1.3087387571, 4.0103920829) and (1.3064726008, 4.0126581598) .. 
    (1.3037306708, 4.0126581598).. controls (1.3008979623, 4.0126581598) and 
    (1.2987678283, 4.0103920035) .. (1.2987678283, 4.0075140513).. controls 
    (1.2987678283, 4.0047267452) and (1.3008979888, 4.0025511823) .. 
    (1.3037306708, 4.0025511823).. controls (1.3064726538, 4.0025511823) and 
    (1.3087387571, 4.004726851) .. (1.3087387571, 4.0075140513) -- cycle;
  
  
  
      \end{scope}
      \begin{scope}[shift={(0.8731702173, -3.614600975)}]
        \path[fill=c005093,even odd rule] (1.3464767806, 3.9989325819).. controls 
    (1.3464767806, 4.0015991056) and (1.3442558681, 4.0039406946) .. 
    (1.3413855358, 4.0039406946).. controls (1.3385528273, 4.0039406946) and 
    (1.3364076385, 4.0015991321) .. (1.3364076385, 3.9989325819).. controls 
    (1.3364076385, 3.9959790381) and (1.3385528802, 3.9937506644) .. 
    (1.3413855358, 3.9937506644).. controls (1.3442560004, 3.9937506644) and 
    (1.3464767806, 3.9959789852) .. (1.3464767806, 3.9989325819) -- cycle;
  
  
  
      \end{scope}
      \begin{scope}[shift={(0.8731702173, -3.6146396042)}]
        \path[fill=c005093,even odd rule] (1.3811207931, 3.9901706402).. controls 
    (1.3811207931, 3.9930033223) and (1.3787715577, 3.9952693992) .. 
    (1.3760219813, 3.9952693992).. controls (1.3731061935, 3.9952693992) and 
    (1.3709307365, 3.9930032429) .. (1.3709307365, 3.9901706402).. controls 
    (1.3709307365, 3.9872926615) and (1.3731061406, 3.9851171515) .. 
    (1.3760219813, 3.9851171515).. controls (1.3787715313, 3.9851171515) and 
    (1.3811207931, 3.9872925556) .. (1.3811207931, 3.9901706402) -- cycle;
  
  
  
      \end{scope}
      \begin{scope}[shift={(0.8731702173, -3.6145943604)}]
        \path[fill=c005093,even odd rule] (1.4136676388, 3.9801090388).. controls 
    (1.4136676388, 3.9828510746) and (1.4114014825, 3.98520031) .. (1.4087047963, 
    3.98520031).. controls (1.4058267646, 3.98520031) and (1.4036059844, 
    3.9828510746) .. (1.4036059844, 3.9801090388).. controls (1.4036059844, 
    3.9773670029) and (1.4058266323, 3.9750027127) .. (1.4087047963, 
    3.9750027127).. controls (1.411401456, 3.9750027127) and (1.4136676388, 
    3.9773670294) .. (1.4136676388, 3.9801090388) -- cycle;
  
  
  
      \end{scope}
      \begin{scope}[shift={(0.8731702173, -3.6146922563)}]
        \path[fill=c005093,even odd rule] (1.4439156256, 3.9671385023).. controls 
    (1.4439156256, 3.970008914) and (1.4415134735, 3.9722675296) .. (1.4388168137,
     3.9722675296).. controls (1.4359387821, 3.9722675296) and (1.4337180019, 
    3.9700090462) .. (1.4337180019, 3.9671385023).. controls (1.4337180019, 
    3.9644267083) and (1.4359386498, 3.9621680662) .. (1.4388168137, 
    3.9621680662).. controls (1.4415134735, 3.9621680662) and (1.4439156256, 
    3.9644268142) .. (1.4439156256, 3.9671385023) -- cycle;
  
  
  
      \end{scope}
      \begin{scope}[shift={(0.8731702173, -3.4588127208)}]
        \path[fill=c005093,even odd rule] (1.458244824, 3.9086884192).. controls 
    (1.3903589583, 3.9441533808) and (1.2701858169, 3.9573196825) .. 
    (1.2689242835, 3.9575312169).. controls (1.2665599669, 3.9576979044) and 
    (1.2649132002, 3.9598804523) .. (1.2649962794, 3.962101259).. controls 
    (1.2653891856, 3.9643674152) and (1.2673529231, 3.9661501777) .. 
    (1.2697022644, 3.9659839665).. controls (1.2750126106, 3.9654627373) and 
    (1.3858266723, 3.9530366869) .. (1.4563714681, 3.9188860429).. controls 
    (1.4570286931, 3.9294235769) and (1.4621726429, 4.0084210165) .. 
    (1.4625654698, 4.0136557712).. controls (1.4464381896, 4.0204918902) and 
    (1.3581571558, 4.0552243) .. (1.2713565981, 4.0628460727).. controls 
    (1.2690904419, 4.0632011435) and (1.2674360552, 4.0650291498) .. 
    (1.2674360552, 4.0669855319).. controls (1.2674360552, 4.0671680944) and 
    (1.2674360552, 4.067287686) .. (1.2674360552, 4.0673784381).. controls 
    (1.267518076, 4.0696520027) and (1.2694831365, 4.0716161106) .. (1.2718778008,
     4.0712989017).. controls (1.3710137098, 4.0601494923) and (1.467656794, 
    4.0206658802) .. (1.4686161467, 4.0202806204) -- (1.4713656967, 4.0191021663) 
    -- (1.4639101617, 3.9056367679) -- (1.458244824, 3.9086884985) -- cycle;
  
  
  
      \end{scope}
      \begin{scope}[shift={(0.8731702173, -3.4588127208)}]
        \path[fill=c005093,even odd rule] (0.8626718152, 3.9086884192).. controls 
    (0.9305199513, 3.9441533808) and (1.0515692604, 3.9582715475) .. 
    (1.0528307938, 3.9584905696).. controls (1.0551877021, 3.9587022363) and 
    (1.0568418771, 3.9607567258) .. (1.0567512573, 3.9629850202).. controls 
    (1.0564491031, 3.965372091) and (1.0544020219, 3.9671094246) .. (1.0521357863,
     3.9669432133).. controls (1.0467801169, 3.9663767404) and (0.9350068083, 
    3.9530366869) .. (0.8645073621, 3.9188860429).. controls (0.8637218142, 
    3.9294235769) and (0.8586230023, 4.0084210165) .. (0.8582301754, 
    4.0136557712).. controls (0.8745312604, 4.0204918902) and (0.9626764306, 
    4.0552243) .. (1.0494769883, 4.0628460727).. controls (1.0516600654, 
    4.0632011435) and (1.0534426427, 4.0650291498) .. (1.0534426427, 
    4.0669855319).. controls (1.0534426427, 4.0671680944) and (1.0534426427, 
    4.067287686) .. (1.0534426427, 4.0673784381).. controls (1.0533129969, 
    4.0696520027) and (1.051266974, 4.0716161106) .. (1.0490010029, 
    4.0712989017).. controls (0.9498651204, 4.0601494923) and (0.8532597394, 
    4.0206658802) .. (0.8522626571, 4.0202806204) -- (0.8495131071, 4.0191021663) 
    -- (0.8568402663, 3.9056367679) -- (0.8626718152, 3.9086884985) -- cycle;
  
  
  
      \end{scope}
      \begin{scope}[shift={(0.8731702173, -3.4590839188)}]
        \path[fill=c005093,even odd rule] (1.4721517473, 4.1071705279).. controls 
    (1.4041752354, 4.1429075077) and (1.2711001375, 4.1551521862) .. 
    (1.2697026612, 4.1551521862).. controls (1.267436505, 4.1554921758) and 
    (1.2657824094, 4.1574937488) .. (1.2660014315, 4.1597675515).. controls 
    (1.2661760565, 4.161988464) and (1.2681315919, 4.1637710677) .. (1.270488315, 
    4.16360483).. controls (1.2758967688, 4.1630761925) and (1.3996807054, 
    4.15157158) .. (1.470316121, 4.1173680458).. controls (1.470973346, 
    4.1280491163) and (1.4762004808, 4.2069106923) .. (1.4764573383, 
    4.2122360933).. controls (1.4602997633, 4.2191176412) and (1.3673958502, 
    4.25675806) .. (1.2815546187, 4.261464045).. controls (1.2793339708, 
    4.26163867) and (1.2775436413, 4.2633828033) .. (1.2775436413, 4.2656565002)..
     controls (1.2775436413, 4.2656565002) and (1.2775436413, 4.2658681669) .. 
    (1.2775436413, 4.2660038981).. controls (1.2777262038, 4.2683078898) and 
    (1.2797190454, 4.2699620913) .. (1.282037986, 4.269871445).. controls 
    (1.3774196692, 4.2646065277) and (1.4815561502, 4.2192837731) .. 
    (1.4825607996, 4.2188985133) -- (1.4853028883, 4.2176295717) -- (1.4778548146,
     4.1042547931) -- (1.4721517473, 4.1071705808) -- cycle;
  
  
  
      \end{scope}
      \begin{scope}[shift={(0.8731702173, -3.4590839188)}]
        \path[fill=c005093,even odd rule] (0.8489942327, 4.1071705279).. controls 
    (0.9169707446, 4.1429075077) and (1.0500004135, 4.1551521862) .. 
    (1.0513978898, 4.1551521862).. controls (1.0537092898, 4.1554921758) and 
    (1.0554920258, 4.1574937488) .. (1.0551898452, 4.1597675515).. controls 
    (1.054970241, 4.161988464) and (1.052923689, 4.1637710677) .. (1.050612236, 
    4.16360483).. controls (1.0452944285, 4.1630761925) and (0.9214122785, 
    4.15157158) .. (0.8508750763, 4.1173680458).. controls (0.8502178513, 
    4.1280491163) and (0.8451191452, 4.2069106923) .. (0.8446810217, 
    4.2122360933).. controls (0.8608007877, 4.2191176412) and (0.9538783998, 
    4.25675806) .. (1.0395836619, 4.261464045).. controls (1.0418572265, 
    4.26163867) and (1.0435116925, 4.2633828033) .. (1.0435116925, 4.2656565002)..
     controls (1.0435116925, 4.2656565002) and (1.0435116925, 4.2658681669) .. 
    (1.0435116925, 4.2660038981).. controls (1.0434217342, 4.2683078898) and 
    (1.0414572029, 4.2699620913) .. (1.0391985079, 4.269871445).. controls 
    (0.9436355852, 4.2646065277) and (0.8395822098, 4.2192837731) .. 
    (0.8386228571, 4.2188985133) -- (0.8358808212, 4.2176295717) -- (0.8433288421,
     4.1042547931) -- (0.8489942327, 4.1071705808) -- cycle;
  
  
  
      \end{scope}
      \begin{scope}[shift={(0.8731702173, -3.4857795833)}]
        \path[fill=c005093,even odd rule] (1.0924439194, 4.0174321956) -- 
    (1.1274783246, 4.0162614144).. controls (1.1436434402, 4.0129075031) and 
    (1.1403273908, 3.9929504056) .. (1.1403273908, 3.9929504056).. controls 
    (1.1403273908, 3.9775255942) and (1.1319652875, 3.9792251717) .. 
    (1.1245550756, 3.9762111706).. controls (1.1171070019, 3.9732953829) and 
    (1.0954050038, 3.9779182623) .. (1.0954050038, 3.9779182623).. controls 
    (1.0829487644, 3.9772157935) and (1.086295135, 3.9671542185) .. (1.086295135, 
    3.9671542185).. controls (1.0859099017, 3.9513365802) and (1.0738388956, 
    3.9508152717) .. (1.0738388956, 3.9508152717).. controls (1.0609823681, 
    3.9500220508) and (1.0604990008, 3.9745644304) .. (1.0604990008, 3.9745644304)
     -- (1.0609823946, 4.0627549238).. controls (1.0609823946, 4.0627549238) and 
    (1.0621605842, 4.0723934829) .. (1.0629463967, 4.0793958396).. controls 
    (1.0638150237, 4.0865039767) and (1.0729627279, 4.0882488508) .. (1.077933164,
     4.0885434644).. controls (1.0829488967, 4.089034531) and (1.0870807888, 
    4.0861111498) .. (1.0870807888, 4.0861111498) -- (1.1074759304, 
    4.0857182435).. controls (1.1282639783, 4.0835428394) and (1.131965314, 
    4.0764800783) .. (1.1345335979, 4.0698100912).. controls (1.1346235562, 
    4.0695153454) and (1.1347532021, 4.0690245433) .. (1.134843425, 
    4.0687223892).. controls (1.1367997542, 4.0614858762) and (1.1345335979, 
    4.0461970077) .. (1.1345335979, 4.0461970077).. controls (1.1307037806, 
    4.0270254581) and (1.0924439458, 4.0337333071) .. (1.0924439458, 
    4.0337333071).. controls (1.0812869694, 4.034518855) and (1.0824653442, 
    4.0261568575) .. (1.0824653442, 4.0261568575).. controls (1.0820271942, 
    4.0178325896) and (1.0924439458, 4.0174322221) .. (1.0924439458, 4.0174322221)
     -- cycle(1.0916583715, 4.0504270867) -- (1.1055120077, 4.0504270867).. 
    controls (1.1104824173, 4.0504270867) and (1.1146670938, 4.0543475237) .. 
    (1.1146670938, 4.0593103927).. controls (1.1146670938, 4.0641069975) and 
    (1.1104824173, 4.0681558369) .. (1.1055120077, 4.0681558369) -- (1.0916583715,
     4.0681558369).. controls (1.0866048827, 4.0681558369) and (1.0824653177, 
    4.0641069975) .. (1.0824653177, 4.0593103927).. controls (1.0824653177, 
    4.0543475237) and (1.0866048827, 4.0504270867) .. (1.0916583715, 4.0504270867)
     -- cycle(1.0837721742, 3.9988272512).. controls (1.0837721742, 3.9936982769) 
    and (1.0883421104, 3.9895511448) .. (1.093924369, 3.9895511448) -- 
    (1.1093492863, 3.9895511448).. controls (1.1149315448, 3.9895511448) and 
    (1.1195468835, 3.9936982769) .. (1.1195468835, 3.9988272512).. controls 
    (1.1195468835, 4.0037145023) and (1.1149315448, 4.0079370406) .. 
    (1.1093492863, 4.0079370406) -- (1.093924369, 4.0079370406).. controls 
    (1.0883421104, 4.0079370406) and (1.0837721742, 4.0037145023) .. 
    (1.0837721742, 3.9988272512) -- cycle;
  
  
  
      \end{scope}
      \begin{scope}[shift={(0.8731702173, -3.4974723146)}]
        \path[fill=c005093,even odd rule] (1.1730948725, 3.98050131) -- 
    (1.2511405265, 3.9812868579).. controls (1.2511405265, 3.9812868579) and 
    (1.2649110835, 3.982155485) .. (1.2649110835, 3.9925646431).. controls 
    (1.2649110835, 3.9925646431) and (1.2628563294, 4.0067732444) .. 
    (1.2511405265, 4.0071736383) -- (1.2304432042, 4.0083444196).. controls 
    (1.2304432042, 4.0083444196) and (1.2230404535, 4.0070828863) .. 
    (1.2230404535, 4.0165857142).. controls (1.2230404535, 4.0165857142) and 
    (1.2208499681, 4.0257785562) .. (1.233404315, 4.0270400896).. controls 
    (1.233404315, 4.0270400896) and (1.26577979, 4.027954225) .. (1.2686955777, 
    4.0416793531).. controls (1.2686955777, 4.0416793531) and (1.2799884971, 
    4.0578898183) .. (1.2721023792, 4.0795464669).. controls (1.2721023792, 
    4.0795464669) and (1.2624636083, 4.1039527977) .. (1.2537088104, 4.079848621) 
    -- (1.2508761813, 4.0640310092) -- (1.2487007771, 4.0578897919).. controls 
    (1.2487007771, 4.0578897919) and (1.2445613179, 4.0495201215) .. 
    (1.2346281452, 4.0495201215).. controls (1.2346281452, 4.0495201215) and 
    (1.2230406123, 4.047389961) .. (1.2232974698, 4.0636456435) -- (1.2213787115, 
    4.0902199438).. controls (1.2213787115, 4.0902199438) and (1.2163252227, 
    4.1061282283) .. (1.2080463044, 4.1061282283).. controls (1.2080463044, 
    4.1061282283) and (1.19379238, 4.1078202387) .. (1.19379238, 4.0857330867) -- 
    (1.1925760904, 4.0599367408).. controls (1.1925760904, 4.0599367408) and 
    (1.1950536488, 4.0499129483) .. (1.1800670667, 4.0495201215).. controls 
    (1.1800670667, 4.0495201215) and (1.165125596, 4.047389961) .. (1.1659112498, 
    4.068691671) -- (1.1646420435, 4.0807173275).. controls (1.1646420435, 
    4.0807173275) and (1.161809335, 4.09249372) .. (1.1539685402, 4.0919196271).. 
    controls (1.1539685402, 4.0919196271) and (1.1426831879, 4.0944953458) .. 
    (1.1430760148, 4.0749311017) -- (1.1426831085, 4.0491349675).. controls 
    (1.1426831085, 4.0491349675) and (1.1464298467, 4.031308004) .. (1.1629877098,
     4.0308321244) -- (1.1751493356, 4.0304392181).. controls (1.1751493356, 
    4.0304392181) and (1.1918355746, 4.0320933931) .. (1.1922208344, 
    4.0165856083).. controls (1.1922208344, 4.0165856083) and (1.1904383365, 
    4.0063049054) .. (1.1775816502, 4.0063049054) -- (1.1581080258, 
    4.0059045908).. controls (1.1581080258, 4.0059045908) and (1.1484768221, 
    4.0067732179) .. (1.148869649, 3.9925646167).. controls (1.148869649, 
    3.9925646167) and (1.1476914594, 3.9805012042) .. (1.1662963244, 3.9805012042)
     -- (1.1730946873, 3.9805012042) -- cycle;
  
  
  
      \end{scope}
      \begin{scope}[shift={(0.8731702173, -3.4325480625)}]
        \path[fill=c005093,even odd rule] (1.0891265206, 4.1074284967) -- 
    (1.1294335425, 4.1083047967).. controls (1.1294335425, 4.1083047967) and 
    (1.1427736225, 4.0825086625) .. (1.1577603104, 4.082078106).. controls 
    (1.1577603104, 4.082078106) and (1.1723163096, 4.0799479456) .. (1.1839491392,
     4.1083047967) -- (1.2417281596, 4.1066053779).. controls (1.2417281596, 
    4.1066053779) and (1.2671693023, 4.1062201446) .. (1.2693145704, 4.1253839948)
     -- (1.2775407308, 4.1573290192).. controls (1.2775407308, 4.1573290192) and 
    (1.2858649723, 4.209789994) .. (1.3008969569, 4.2384641598).. controls 
    (1.3008969569, 4.2384641598) and (1.3133682246, 4.2505275723) .. 
    (1.3034729402, 4.2596450346).. controls (1.3034729402, 4.2596450346) and 
    (1.2905257665, 4.2721918144) .. (1.2742322485, 4.2434270023).. controls 
    (1.2742322485, 4.2434270023) and (1.2588980833, 4.2251317262) .. 
    (1.2414639467, 4.1452580396).. controls (1.2414639467, 4.1452580396) and 
    (1.2426424008, 4.1282167298) .. (1.2285242342, 4.1290854362) -- (1.1927948479,
     4.1299617362).. controls (1.1927948479, 4.1299617362) and (1.2031209269, 
    4.1565434183) .. (1.2114376013, 4.1972885637).. controls (1.2114376013, 
    4.1972885637) and (1.2173295281, 4.219338145) .. (1.2230703515, 
    4.2314015575).. controls (1.2230703515, 4.2314015575) and (1.2301785944, 
    4.2484428937) .. (1.2168461873, 4.2554981406).. controls (1.2168461873, 
    4.2554981406) and (1.1981126554, 4.2580738594) .. (1.1927948479, 4.2364172108)
     -- (1.1744920048, 4.1672169217) -- (1.1661225196, 4.1395021879).. controls 
    (1.1661225196, 4.1395021879) and (1.158583826, 4.1215468485) .. (1.1523972856,
     4.136065224) -- (1.1453419329, 4.1649130887).. controls (1.1453419329, 
    4.1649130887) and (1.1328782323, 4.2288331942) .. (1.12028605, 4.2434271346)..
     controls (1.12028605, 4.2434271346) and (1.1057752681, 4.2626062512) .. 
    (1.09623479, 4.2554981406).. controls (1.09623479, 4.2554981406) and 
    (1.081678685, 4.2488734502) .. (1.0940969567, 4.2269221617).. controls 
    (1.0940969567, 4.2269221617) and (1.1078221906, 4.198157376) .. (1.1108286246,
     4.1748010442) -- (1.1191529719, 4.1299617362) -- (1.0900028206, 
    4.1290854362).. controls (1.0900028206, 4.1290854362) and (1.0829476531, 
    4.1233445071) .. (1.0783247473, 4.1429012106) -- (1.0616761585, 4.1959817867) 
    -- (1.0475580977, 4.2359790344) -- (1.0455110165, 4.2505275723).. controls 
    (1.0455110165, 4.2505275723) and (1.0428974623, 4.2534509271) .. 
    (1.0388032733, 4.256812326).. controls (1.0347091373, 4.2602039933) and 
    (1.0267323996, 4.2622057515) .. (1.023423785, 4.2604230154).. controls 
    (1.0201077092, 4.2588593279) and (1.0146613935, 4.2558982965) .. 
    (1.0142685667, 4.2534508212).. controls (1.0137925813, 4.2509203463) and 
    (1.0163608917, 4.2388569867) .. (1.0167082896, 4.2372027323).. controls 
    (1.0171011958, 4.2355408844) and (1.0234237056, 4.2196855165) .. 
    (1.0234237056, 4.2196855165) -- (1.0475580183, 4.1474410637) -- (1.0541827881,
     4.12116905).. controls (1.0541827881, 4.12116905) and (1.0616760792, 
    4.1095663829) .. (1.0728330292, 4.1083048496) -- (1.0841637046, 4.1066054308) 
    -- (1.0891264413, 4.1074288142) -- cycle;
  
  
  
      \end{scope}
      \begin{scope}[shift={(0.8731702173, -3.5647397167)}]
        \path[fill=c005093,even odd rule] (1.1699624969, 4.5179765548).. controls 
    (1.168957874, 4.518935934) and (1.169048626, 4.5208545071) .. (1.1701741635, 
    4.5218969125).. controls (1.2031237315, 4.5509185821) and (1.2453040298, 
    4.5697878452) .. (1.2923640915, 4.5764125092).. controls (1.2938446998, 
    4.5766321133) and (1.2953252817, 4.5755441467) .. (1.2954989806, 
    4.5740937008).. controls (1.2956286265, 4.5724469342) and (1.2946303535, 
    4.57113259) .. (1.293194989, 4.5708379765).. controls (1.2465655102, 
    4.5642586092) and (1.2064322931, 4.5464316721) .. (1.173754611, 
    4.5175838073).. controls (1.1727952319, 4.5165791844) and (1.1709672256, 
    4.5168059323) .. (1.1699625498, 4.5179767135) -- cycle;
  
  
  
      \end{scope}
      \begin{scope}[shift={(0.8731702173, -3.5940327958)}]
        \path[fill=c005093,even odd rule] (1.1335741129, 4.5021204725).. controls 
    (1.1321842567, 4.5029136933) and (1.1317008629, 4.5045752767) .. 
    (1.1323580879, 4.5058671312).. controls (1.1327054858, 4.5064412771) and 
    (1.1379779967, 4.5166387421) .. (1.1500036796, 4.530802126).. controls 
    (1.1511366254, 4.5319729073) and (1.152791065, 4.5321467385) .. (1.1540600596,
     4.5311873594).. controls (1.1552308408, 4.5300996573) and (1.1553215929, 
    4.5284452706) .. (1.1542796638, 4.5271837373).. controls (1.1426845638, 
    4.5135415825) and (1.1374120529, 4.5035253306) .. (1.1372836771, 
    4.5032987415).. controls (1.1366716958, 4.5020372081) and (1.1349270333, 
    4.501387656) .. (1.1335747744, 4.5021202873) -- cycle;
  
  
  
      \end{scope}
      \begin{scope}[shift={(0.8731702173, -3.5582656271)}]
        \path[fill=c005093,even odd rule] (1.1606390565, 4.5200009348).. controls 
    (1.1595436815, 4.5212246327) and (1.1597627565, 4.5229619398) .. 
    (1.1609412106, 4.5240572354).. controls (1.1891772265, 4.5487505862) and 
    (1.2353610146, 4.5795700202) .. (1.2944996494, 4.5850540919).. controls 
    (1.296146416, 4.5852657585) and (1.297370114, 4.5840571419) .. (1.2975438129, 
    4.5824781085).. controls (1.2977634171, 4.5809599294) and (1.2965844338, 
    4.5796908025) .. (1.2951114983, 4.5794792681).. controls (1.2373250167, 
    4.5740707879) and (1.1922214958, 4.544044469) .. (1.1645595729, 
    4.5197363515).. controls (1.1634718708, 4.5187394015) and (1.1617267585, 
    4.5188600515) .. (1.16063903, 4.5200009348) -- cycle;
  
  
  
      \end{scope}
      \begin{scope}[shift={(0.8731702173, -3.5859278146)}]
        \path[fill=c005093,even odd rule] (1.114666591, 4.4951938131).. controls 
    (1.1132690619, 4.4960248694) and (1.1128838285, 4.4977621235) .. 
    (1.1137524556, 4.4991216321).. controls (1.1141453619, 4.4996351883) and 
    (1.1232929337, 4.5137608956) .. (1.1405080483, 4.5321090619).. controls 
    (1.1417240733, 4.5332798431) and (1.1433784071, 4.5332798431) .. 
    (1.1445567554, 4.5321910827).. controls (1.1456897013, 4.5311033806) and 
    (1.1456899658, 4.5293584535) .. (1.1446467138, 4.5281876723).. controls 
    (1.1276507271, 4.5103229527) and (1.1184576469, 4.4961973248) .. 
    (1.1183670271, 4.4960236258).. controls (1.1175888875, 4.4948076008) and 
    (1.1157987167, 4.4944146946) .. (1.1146656915, 4.4951928342) -- cycle;
  
  
  
      \end{scope}
      \begin{scope}[shift={(0.8731702173, -3.5545940042)}]
        \path[fill=c005093,even odd rule] (1.1494809688, 4.525091571).. controls 
    (1.1483480229, 4.5263454315) and (1.1485668333, 4.5281358404) .. 
    (1.1496555937, 4.5291404898).. controls (1.1925761963, 4.5660029542) and 
    (1.2421214892, 4.5878787306) .. (1.2965843015, 4.5939293281).. controls 
    (1.298019666, 4.5941118906) and (1.2995454123, 4.592969949) .. (1.2997266783, 
    4.5914970135).. controls (1.299938345, 4.5899259177) and (1.2987672992, 
    4.5885284415) .. (1.2972866908, 4.5884907119).. controls (1.2438209608, 
    4.5823872242) and (1.1954464756, 4.5609874331) .. (1.1534021467, 
    4.5248727871).. controls (1.1521861217, 4.5239134079) and (1.1503958185, 
    4.5239134079) .. (1.149481789, 4.5250923913) -- cycle;
  
  
  
      \end{scope}
      \begin{scope}[shift={(0.8731702173, -3.5748166375)}]
        \path[fill=c005093,even odd rule] (1.0957970369, 4.4825266479).. controls 
    (1.0944902598, 4.4832140354) and (1.0940521098, 4.4849137188) .. 
    (1.0948376577, 4.486273386).. controls (1.0953136431, 4.4872250923) and 
    (1.1065611335, 4.5066609342) .. (1.1298268456, 4.5305837919).. controls 
    (1.1309145477, 4.5317545731) and (1.1327047715, 4.5317545731) .. 
    (1.1337471769, 4.5307055002).. controls (1.134925631, 4.529625471) and 
    (1.134925631, 4.5278726858) .. (1.1338371352, 4.5266565815).. controls 
    (1.111221134, 4.5033455727) and (1.0998073, 4.4836982229) .. (1.0997166538, 
    4.483524524).. controls (1.0989311058, 4.4821722385) and (1.0971483433, 
    4.4817795969) .. (1.0957962167, 4.4825273094) -- cycle;
  
  
  
      \end{scope}
      \begin{scope}[shift={(0.8731702173, -3.5564220104)}]
        \path[fill=c005093,even odd rule] (1.2011580096, 4.5647786477).. controls 
    (1.1984160267, 4.5806340421) and (1.2022081408, 4.5887089931) .. 
    (1.2024272158, 4.5890111738).. controls (1.2185017117, 4.6251334133) and 
    (1.2594203708, 4.6280567417) .. (1.2611577838, 4.6281322537).. controls 
    (1.2732287633, 4.6305796496) and (1.2787278633, 4.6220816562) .. 
    (1.2798080248, 4.6171490821).. controls (1.2799905873, 4.6111891308) and 
    (1.2785464915, 4.606739924) .. (1.2754117081, 4.6036504373).. controls 
    (1.2699578254, 4.5982041217) and (1.262162539, 4.5985063023) .. (1.2618603583,
     4.5985063023).. controls (1.2616857333, 4.5985063023) and (1.230179785, 
    4.596368469) .. (1.2296963119, 4.5767211721).. controls (1.2296963119, 
    4.57519532) and (1.2283064556, 4.573979295) .. (1.2268637092, 4.573979295).. 
    controls (1.2252547779, 4.573979295) and (1.2239932446, 4.5752783992) .. 
    (1.2240762973, 4.5768495479).. controls (1.2246882785, 4.6015579535) and 
    (1.2601983252, 4.6040356971) .. (1.2616864477, 4.6040356971).. controls 
    (1.262079354, 4.6042103221) and (1.2678278767, 4.6040356971) .. (1.2713555133,
     4.6075331448).. controls (1.2734102675, 4.6096180615) and (1.274324059, 
    4.6126697921) .. (1.2742334392, 4.6163711277).. controls (1.2741434808, 
    4.6167714423) and (1.2718767954, 4.6246122106) .. (1.2618603583, 
    4.6226481821).. controls (1.2611578896, 4.6226481821) and (1.2222029415, 
    4.6195964515) .. (1.2075259219, 4.5866544506).. controls (1.2074359635, 
    4.5864348465) and (1.2041645229, 4.5796747952) .. (1.206732701, 
    4.5657758623).. controls (1.2069152635, 4.564174604) and (1.2059545615, 
    4.5628147515) .. (1.2043834656, 4.5625201379).. controls (1.2028197781, 
    4.5622105754) and (1.2014676779, 4.5632149338) .. (1.201158036, 4.5647786213) 
    -- cycle;
  
  
  
      \end{scope}
      \begin{scope}[shift={(0.8731702173, -3.57524685)}]
        \path[fill=c005093,even odd rule] (1.209483574, 4.5887027225).. controls 
    (1.209483574, 4.5891787079) and (1.2090906677, 4.6332322827) .. (1.2670509012,
     4.6355965729).. controls (1.2685238367, 4.6355965729) and (1.2699214717, 
    4.6344107104) .. (1.2699214717, 4.6328016998).. controls (1.2699214717, 
    4.6312681748) and (1.2686977738, 4.6300143144) .. (1.2672624621, 
    4.6298859121).. controls (1.2149751863, 4.6279671538) and (1.2149751863, 
    4.5904476231) .. (1.2149751863, 4.5887933688).. controls (1.2150572071, 
    4.5873053521) and (1.2137967321, 4.5860513858) .. (1.2123161767, 
    4.5860513858).. controls (1.2108810767, 4.5860513858) and (1.209483574, 
    4.5872221671) .. (1.209483574, 4.5887027225) -- cycle;
  
  
  
      \end{scope}
      \begin{scope}[shift={(0.8731702173, -3.5804816313)}]
        \path[fill=c005093,even odd rule] (1.2544937498, 4.5358566202).. controls 
    (1.2267411806, 4.5412576127) and (1.2126231198, 4.5618643946) .. 
    (1.2120112708, 4.5627784242).. controls (1.2111349708, 4.5639190429) and 
    (1.211437125, 4.5657016731) .. (1.2127894104, 4.5665703796).. controls 
    (1.2140583521, 4.5674466796) and (1.215757956, 4.5670992817) .. (1.2167173088,
     4.5657925046).. controls (1.2167173088, 4.5655729004) and (1.230178224, 
    4.5462734513) .. (1.2556722569, 4.5413483383).. controls (1.276369685, 
    4.5374203342) and (1.2991592794, 4.5436975738) .. (1.3234673969, 
    4.5603008923).. controls (1.3240415427, 4.5607768777) and (1.342601111, 
    4.5764659285) .. (1.3568097917, 4.5747210279).. controls (1.3614251569, 
    4.5742450425) and (1.3649528729, 4.5721979612) .. (1.36756648, 4.5684060058)..
     controls (1.3686541821, 4.5674466267) and (1.3708371269, 4.5640929535) .. 
    (1.3702252779, 4.5596059642).. controls (1.3695680529, 4.5538122506) and 
    (1.3645599402, 4.5482374533) .. (1.3558504654, 4.5430404546).. controls 
    (1.3554575592, 4.5427383004) and (1.3411130944, 4.532238549) .. (1.3177643296,
     4.5489248937).. controls (1.3173714233, 4.5491894771) and (1.3173714233, 
    4.5495744458) .. (1.3171901837, 4.5499747604).. controls (1.2950197673, 
    4.5370275867) and (1.2739296181, 4.532155364) .. (1.2544937498, 4.5358566996) 
    -- cycle(1.3526703854, 4.5476556875).. controls (1.3619010894, 4.5531095702) 
    and (1.3642578125, 4.5574680515) .. (1.3646506394, 4.5602176015).. controls 
    (1.3650358727, 4.5631333892) and (1.3634721852, 4.5648330725) .. 
    (1.3632531102, 4.5649236923).. controls (1.361425104, 4.5675298381) and 
    (1.3591589742, 4.56884405) .. (1.3562356194, 4.5692292833).. controls 
    (1.3443384446, 4.5705889771) and (1.3271761144, 4.5559420935) .. 
    (1.3270024154, 4.5558967704).. controls (1.3252122446, 4.5546807454) and 
    (1.3236863396, 4.5538949329) .. (1.3220774083, 4.5528526069).. controls 
    (1.3411129092, 4.5397769515) and (1.35218678, 4.5474441531) .. (1.3526702531, 
    4.5476556875) -- cycle;
  
  
  
      \end{scope}
      \begin{scope}[shift={(0.8731702173, -3.5424784688)}]
        \path[fill=c005093,even odd rule] (1.3074302867, 4.5055275915).. controls 
    (1.3061687533, 4.5066076206) and (1.3061687533, 4.5084358121) .. 
    (1.3071736408, 4.5095311077).. controls (1.3082235075, 4.5106111369) and 
    (1.3100063494, 4.5108378848) .. (1.311101645, 4.5097427744).. controls 
    (1.319901766, 4.5019019796) and (1.3432580979, 4.4856990815) .. (1.3687823198,
     4.5074766181).. controls (1.369175226, 4.5078695244) and (1.3834744735, 
    4.5231733419) .. (1.3709275615, 4.5397690404).. controls (1.3705724906, 
    4.5404262654) and (1.3619008248, 4.55240644) .. (1.3375927337, 4.5368910088)..
     controls (1.3346240558, 4.5347156046) and (1.2604232475, 4.483123204) .. 
    (1.2246862942, 4.5297526827).. controls (1.2236816713, 4.531014216) and 
    (1.2239912338, 4.5326684704) .. (1.2252149317, 4.5338015221).. controls 
    (1.2263404692, 4.5345418262) and (1.2281307194, 4.5343227513) .. 
    (1.2290900456, 4.5331898054).. controls (1.2614655471, 4.4908734848) and 
    (1.3336192477, 4.5411288135) .. (1.3344123627, 4.541597126).. controls 
    (1.3637361335, 4.5602473671) and (1.3750668883, 4.5436063719) .. 
    (1.3756334142, 4.5430020635).. controls (1.3908316365, 4.5225992754) and 
    (1.3728007056, 4.5036543944) .. (1.3724910373, 4.5033447525).. controls 
    (1.3426460902, 4.4779337723) and (1.3151048702, 4.4986387675) .. 
    (1.3074302073, 4.5055278296) -- cycle;
  
  
  
      \end{scope}
      \begin{scope}[shift={(0.8731702173, -3.5050565958)}]
        \path[fill=c005093,even odd rule] (1.2935390267, 4.4637322619).. controls 
    (1.2923605725, 4.4647368848) and (1.2921867413, 4.4665272408) .. 
    (1.2932366079, 4.4677055892).. controls (1.2942335579, 4.4688840433) and 
    (1.295978776, 4.4688840433) .. (1.2971495838, 4.4679251933).. controls 
    (1.3446477954, 4.4259261875) and (1.3791609715, 4.4634003685) .. (1.380505584,
     4.4651378079).. controls (1.4019432104, 4.49027677) and (1.3771062438, 
    4.5124094304) .. (1.3760185152, 4.5134140798).. controls (1.3532740058, 
    4.5304553896) and (1.3275610567, 4.5081037335) .. (1.326390249, 
    4.5071443808).. controls (1.2707037652, 4.4653492365) and (1.2367570885, 
    4.496901434) .. (1.2364549079, 4.4972489377).. controls (1.2353595329, 
    4.4982988044) and (1.2352764538, 4.5000816727) .. (1.2363649496, 
    4.501252454).. controls (1.2373243288, 4.502347829) and (1.2391976581, 
    4.5023930727) .. (1.2402929538, 4.501252454).. controls (1.2415544871, 
    4.5000816727) and (1.2712633854, 4.4728123385) .. (1.3228557596, 
    4.5114499719).. controls (1.3229457179, 4.5114499719) and (1.3529727248, 
    4.5377221973) .. (1.3795922423, 4.5176441852).. controls (1.3800303923, 
    4.5173346227) and (1.4097469635, 4.4908358081) .. (1.3847363773, 
    4.4614364723).. controls (1.3842982273, 4.4609454056) and (1.34513193, 
    4.4179871) .. (1.2935395823, 4.463732791) -- cycle;
  
  
  
      \end{scope}
      \begin{scope}[shift={(0.8731702173, -3.5748682313)}]
        \path[fill=c005093,even odd rule] (1.315582549, 4.4796556542).. controls 
    (1.3118812133, 4.4804337938) and (1.3091390452, 4.4824881775) .. 
    (1.3074319006, 4.4854115323).. controls (1.3042517942, 4.4914092396) and 
    (1.3071752548, 4.4992651685) .. (1.3073551715, 4.4995597821).. controls 
    (1.3116155452, 4.5153396644) and (1.3065619506, 4.5230067337) .. 
    (1.3065619506, 4.5232182946).. controls (1.3056933235, 4.5244041571) and 
    (1.3059954777, 4.5261415435) .. (1.3073551715, 4.52701025).. controls 
    (1.3086167048, 4.5279243854) and (1.3103086888, 4.5275013167) .. 
    (1.3112680415, 4.526232375).. controls (1.3114955831, 4.5259302208) and 
    (1.3176737098, 4.5165182244) .. (1.3126655706, 4.4978225544).. controls 
    (1.3120160185, 4.4958133085) and (1.3107015685, 4.4911147319) .. 
    (1.3124830081, 4.4880629748).. controls (1.313276229, 4.486491879) and 
    (1.314711329, 4.4855324998) .. (1.3168414894, 4.48501135).. controls 
    (1.3169764269, 4.4849663708) and (1.3217741694, 4.48413505) .. (1.3264801544, 
    4.4876626867).. controls (1.332100169, 4.4918097394) and (1.3354919421, 
    4.5007382363) .. (1.3364059717, 4.5133833352).. controls (1.3365885342, 
    4.5149845935) and (1.3378939883, 4.5161252917) .. (1.33946519, 4.5159064019)..
     controls (1.34090029, 4.5159064019) and (1.3421165267, 4.5145996248) .. 
    (1.341995665, 4.5129831529).. controls (1.34090029, 4.4982985662) and 
    (1.3367986398, 4.4882822085) .. (1.3297056369, 4.4829568075).. controls 
    (1.3228543838, 4.4780015321) and (1.3159729417, 4.4794368702) .. 
    (1.3155801148, 4.4796558658) -- cycle;
  
  
  
      \end{scope}
      \begin{scope}[shift={(0.8731702173, -3.5857920833)}]
        \path[fill=c005093,even odd rule] (1.3138453213, 4.4937969454).. controls 
    (1.3126668671, 4.49479416) and (1.3126668671, 4.496455955) .. (1.3136706962, 
    4.4976267627).. controls (1.3187619146, 4.5036018481) and (1.3224708173, 
    4.5224711112) .. (1.3233395238, 4.5292241775).. controls (1.3236869217, 
    4.5307500296) and (1.3250013717, 4.5317546525) .. (1.3266028152, 
    4.5315355775).. controls (1.328173911, 4.531353015) and (1.3291332902, 
    4.5298814025) .. (1.3289142152, 4.5284385237).. controls (1.328521309, 
    4.5256512177) and (1.3249560221, 4.5022571298) .. (1.3177648323, 
    4.4940159675).. controls (1.3168054531, 4.4928830217) and (1.3150152823, 
    4.4927544342) .. (1.3138444746, 4.4937963633) -- cycle;
  
  
  
      \end{scope}
      \begin{scope}[shift={(0.8731702173, -3.5064009438)}]
        \path[fill=c005093,even odd rule] (1.4126274558, 4.4163625319).. controls 
    (1.3825180577, 4.4308808808) and (1.3890066729, 4.4790287769) .. 
    (1.3892635569, 4.481204234).. controls (1.3894910985, 4.4826922506) and 
    (1.3908346527, 4.4837725444) .. (1.3924059338, 4.4835611423).. controls 
    (1.3938939504, 4.4833865173) and (1.3949742442, 4.4818617235) .. 
    (1.3948533296, 4.480426359).. controls (1.3946257879, 4.4799429652) and 
    (1.3886138196, 4.4341064956) .. (1.4146743783, 4.4216049331).. controls 
    (1.4176354892, 4.4208572206) and (1.4281805902, 4.4191953727) .. (1.431927249,
     4.4233424519).. controls (1.4328866281, 4.424649229) and (1.4333171052, 
    4.4261750546) .. (1.4333171052, 4.4280484369).. controls (1.4333171052, 
    4.4307979869) and (1.4323201552, 4.4338497175) .. (1.431368449, 
    4.4365917004).. controls (1.4291853719, 4.439024015) and (1.4163287915, 
    4.4549322731) .. (1.4158528854, 4.486001024).. controls (1.4158528854, 
    4.4862206281) and (1.4153316562, 4.5102260623) .. (1.4014251827, 
    4.5196833554).. controls (1.3938940562, 4.5245631452) and (1.3835679508, 
    4.5245631452) .. (1.370575454, 4.5194637512).. controls (1.3691779248, 
    4.5188594429) and (1.3675237233, 4.5196833554) .. (1.3670025206, 
    4.5210348471).. controls (1.3663077248, 4.5223794596) and (1.367092479, 
    4.5241696304) .. (1.3684452935, 4.5246531035).. controls (1.3832960121, 
    4.5305374369) and (1.3954575321, 4.5303562767) .. (1.4046580471, 
    4.5241697098).. controls (1.4210346442, 4.5131940258) and (1.421427471, 
    4.487095711) .. (1.421427471, 4.4861363583).. controls (1.4219561085, 
    4.4551961156) and (1.4357646333, 4.4400280558) .. (1.4358930356, 4.4398090602)
     -- (1.4362859419, 4.4394235623) -- (1.4364605669, 4.4389401685).. controls 
    (1.4398598013, 4.4303967992) and (1.4397766427, 4.4239005904) .. 
    (1.4360676606, 4.4197309687).. controls (1.4294052406, 4.4121469521) and 
    (1.4144185527, 4.4158935844) .. (1.4126283554, 4.4163618969) -- cycle;
  
  
  
      \end{scope}
      \begin{scope}[shift={(0.8731702173, -3.5333752146)}]
        \path[fill=c005093,even odd rule] (1.4328850142, 4.4469626527).. controls 
    (1.42360142, 4.448352509) and (1.4162440983, 4.4526733135) .. (1.4111452865, 
    4.4597285604).. controls (1.4042562773, 4.4692765792) and (1.4027378865, 
    4.4819064117) .. (1.4027378865, 4.4915828327).. controls (1.4027378865, 
    4.4989477215) and (1.4036065135, 4.5044772221) .. (1.403863424, 
    4.5057462167).. controls (1.403296951, 4.5162006979) and (1.4005549417, 
    4.523693989) .. (1.3956298288, 4.528188519).. controls (1.390750039, 
    4.5324035696) and (1.3856060098, 4.532101389) .. (1.3854397721, 4.532101389)..
     controls (1.3839062471, 4.5320246598) and (1.3825164173, 4.533196764) .. 
    (1.3825164173, 4.534631864).. controls (1.3822968131, 4.5362407952) and 
    (1.3834757965, 4.5375098956) .. (1.3850471569, 4.5376005419).. controls 
    (1.385349311, 4.5378122085) and (1.3924951248, 4.5382048502) .. (1.3991575448,
     4.5324035167).. controls (1.4053894083, 4.5270025242) and (1.4087886427, 
    4.5179908423) .. (1.4093627356, 4.5057461638) -- (1.409484444, 4.5054363367) 
    -- (1.4093627356, 4.5051267742).. controls (1.4093627356, 4.5048780658) and 
    (1.4046567506, 4.4783335577) .. (1.4156701906, 4.4630824452).. controls 
    (1.4197643531, 4.4571980854) and (1.4259583548, 4.4536628023) .. 
    (1.4337160704, 4.4524466979).. controls (1.435151435, 4.4523699688) and 
    (1.4362843808, 4.4508756021) .. (1.4360653058, 4.4493572112).. controls 
    (1.4358906808, 4.44774828) and (1.4345016183, 4.4467889008) .. (1.4328851729, 
    4.4469627321) -- cycle;
  
  
  
      \end{scope}
      \begin{scope}[shift={(0.8731702173, -3.5034852354)}]
        \path[fill=c005093,even odd rule] (1.2921904983, 4.4876174694).. controls 
    (1.29070989, 4.4876174694) and (1.2894030865, 4.4888866756) .. (1.2894030865, 
    4.4904501779).. controls (1.2894030865, 4.4914095571) and (1.2895327323, 
    4.5177268675) .. (1.2983317156, 4.5476625667).. controls (1.2988151094, 
    4.5485311937) and (1.3193387856, 4.5867912137) .. (1.3012928, 4.6038249565).. 
    controls (1.3001220188, 4.6049203315) and (1.3001220188, 4.6066575592) .. 
    (1.3010811333, 4.6077981779).. controls (1.3021611625, 4.6089237154) and 
    (1.3039138419, 4.6089237154) .. (1.3050393265, 4.6079278237).. controls 
    (1.3262654185, 4.5878347833) and (1.3045255056, 4.5472710098) .. 
    (1.3036417973, 4.5453976275).. controls (1.2951135885, 4.5169424838) and 
    (1.2950229423, 4.49067047) .. (1.2950229423, 4.4904514479).. controls 
    (1.2950229423, 4.4888877604) and (1.2938444881, 4.4876187394) .. 
    (1.2921903396, 4.4876187394) -- cycle;
  
  
  
      \end{scope}
      \begin{scope}[shift={(0.8731702173, -3.4967777833)}]
        \path[fill=c005093,even odd rule] (1.3046463144, 4.4892262948).. controls 
    (1.3036416915, 4.4930635733) and (1.3032564581, 4.4971653029) .. 
    (1.3032564581, 4.5014709998).. controls (1.3032564581, 4.5245704477) and 
    (1.3163698696, 4.5508425673) .. (1.3177673458, 4.5535921173).. controls 
    (1.3198144271, 4.5622109988) and (1.3206831335, 4.5696666131) .. 
    (1.3206831335, 4.5760268525).. controls (1.3206831335, 4.600689935) and 
    (1.3081814652, 4.6089311238) .. (1.3074336733, 4.6093239506).. controls 
    (1.3065650463, 4.6098904235) and (1.30617214, 4.6107590506) .. (1.30617214, 
    4.6116276777).. controls (1.30617214, 4.6122393944) and (1.30617214, 
    4.6127230527) .. (1.3065650463, 4.6132442819).. controls (1.3073582671, 
    4.6145058152) and (1.3091410296, 4.6148987215) .. (1.3103116785, 
    4.614112909).. controls (1.3114070535, 4.6135011923) and (1.3344158023, 
    4.5989071196) .. (1.323168206, 4.5520207569) -- (1.3229486019, 4.5516278506) 
    -- (1.3228586435, 4.5513256965).. controls (1.3227766227, 4.551016134) and 
    (1.3036417708, 4.5139420294) .. (1.3100095773, 4.4907141792).. controls 
    (1.310492971, 4.4892261625) and (1.3095338565, 4.4876625544) .. (1.308181571, 
    4.4872621604).. controls (1.3065649669, 4.4868846) and (1.3050391942, 
    4.487662475) .. (1.3046463673, 4.4892261625) -- cycle;
  
  
  
      \end{scope}
      \begin{scope}[shift={(0.8731702173, -3.4974731083)}]
        \path[fill=c005093,even odd rule] (1.31658741, 4.4987297313).. controls 
    (1.3165053892, 4.4996438667) and (1.3146310808, 4.5206055077) .. (1.325606659,
     4.5481843779).. controls (1.3259162215, 4.5485846925) and (1.3463039813, 
    4.5950707406) .. (1.3151068546, 4.6180342456).. controls (1.3138453213, 
    4.6188197935) and (1.3136717546, 4.620610229) .. (1.3145403817, 
    4.6219547092).. controls (1.3154090088, 4.6231331633) and (1.3171917448, 
    4.6233519738) .. (1.3183776602, 4.6225664258).. controls (1.3532761225, 
    4.5968533708) and (1.3310982448, 4.5464016419) .. (1.3308338996, 
    4.545835116).. controls (1.3202510425, 4.5199482563) and (1.3221697215, 
    4.4995605758) .. (1.3221697215, 4.499250934).. controls (1.3222993673, 
    4.4978155694) and (1.3212103423, 4.496418199) .. (1.3196392465, 
    4.4961992033).. controls (1.3181512298, 4.4960325158) and (1.3166781621, 
    4.4972038262) .. (1.3165875158, 4.4987296783) -- cycle;
  
  
  
      \end{scope}
      \begin{scope}[shift={(0.8731702173, -3.6150700812)}]
        \path[fill=c005093,even odd rule] (1.3293065658, 4.6650027879).. controls 
    (1.3277883867, 4.6653880213) and (1.3267834992, 4.6668760379) .. 
    (1.3270025742, 4.6683490263).. controls (1.3273954804, 4.6699201221) and 
    (1.3287475013, 4.6708795012) .. (1.3304395646, 4.6704868596).. controls 
    (1.3329700396, 4.6700108742) and (1.3405088125, 4.6692179179) .. (1.343870185,
     4.6717105575).. controls (1.3451317183, 4.6725791846) and (1.3469143485, 
    4.6722770304) .. (1.3477377583, 4.6709702533).. controls (1.3487423813, 
    4.66970872) and (1.348349475, 4.6679185227) .. (1.3470882063, 4.6670498163).. 
    controls (1.3409922856, 4.6628574404) and (1.3305303431, 4.6648214954) .. 
    (1.3293066452, 4.665002735) -- cycle;
  
  
  
      \end{scope}
      \begin{scope}[shift={(0.8731702173, -3.6178265104)}]
        \path[fill=c005093,even odd rule] (1.3317032675, 4.683259911).. controls 
    (1.3314916008, 4.6848386798) and (1.3327078904, 4.6861832658) .. 
    (1.3341506633, 4.6864022615) -- (1.3419990252, 4.6874068844).. controls 
    (1.3434719606, 4.6876635302) and (1.3449073252, 4.686583501) .. (1.3451339144,
     4.6850124052).. controls (1.3453006019, 4.6835696323) and (1.3442576144, 
    4.6820891562) .. (1.3426486831, 4.6818776219) -- (1.3349288294, 
    4.6809182427).. controls (1.3333651419, 4.6806986385) and (1.3319677185, 
    4.6816963823) .. (1.3317033733, 4.6832598052) -- cycle;
  
  
  
      \end{scope}
      \begin{scope}[shift={(0.8731702173, -3.6184313479)}]
        \path[fill=c005093,even odd rule] (1.3327536633, 4.6992522538).. controls 
    (1.3327086842, 4.7007328621) and (1.3338413654, 4.7021302854) .. 
    (1.3354126729, 4.7021302854) -- (1.3432610348, 4.7026967583).. controls 
    (1.3448247223, 4.7027787792) and (1.3460936375, 4.7017373792) .. 
    (1.3462144727, 4.7001284479).. controls (1.346304431, 4.6985949229) and 
    (1.3451344435, 4.6972049873) .. (1.3435631096, 4.6970312883) -- (1.335805394, 
    4.696646055).. controls (1.3342342981, 4.6964185133) and (1.3328896063, 
    4.6976808404) .. (1.3327536633, 4.6992522008) -- cycle;
  
  
  
      \end{scope}
      \begin{scope}[shift={(0.8731702173, -3.350136175)}]
        \path[fill=c005093,even odd rule] (1.4723622763, 4.3133517206).. controls 
    (1.452533581, 4.3177103077) and (1.4382040652, 4.3429851333) .. (1.4376375394,
     4.3439897827) -- (1.437547581, 4.344375016) -- (1.4374179352, 4.3446395994)..
     controls (1.4305742492, 4.3702620081) and (1.3915363277, 4.3754438992) .. 
    (1.3911359338, 4.375534519).. controls (1.3501416038, 4.3787977046) and 
    (1.3411146819, 4.4081216606) .. (1.3407294221, 4.4094359783).. controls 
    (1.3341500548, 4.456971946) and (1.3775919658, 4.4722683552) .. (1.3779847927,
     4.4723590015).. controls (1.4152853806, 4.4903067738) and (1.4238286442, 
    4.5128774519) .. (1.4238664002, 4.5130511508).. controls (1.4449565494, 
    4.5585023077) and (1.4741067006, 4.574493645) .. (1.475285049, 4.5752792723) 
    -- (1.4810787625, 4.5783234358).. controls (1.4890102035, 4.5836866988) and 
    (1.4966772729, 4.5858619442) .. (1.5037401133, 4.5844721408).. controls 
    (1.5121475133, 4.5828179658) and (1.5164607773, 4.576623779) .. (1.5166797994,
     4.57645008).. controls (1.519905229, 4.5707469863) and (1.5206908827, 
    4.5656481744) .. (1.5188628765, 4.5611612115).. controls (1.5155014775, 
    4.5527084883) and (1.505001779, 4.5497399163) .. (1.5046089521, 4.54964927).. 
    controls (1.4727546798, 4.5392779737) and (1.4490960615, 4.4893927442) .. 
    (1.4488317163, 4.4889545942).. controls (1.4345022004, 4.4604994504) and 
    (1.4038640325, 4.4473407423) .. (1.4025118794, 4.4467742165) -- (1.4022097252,
     4.4466842581) -- (1.4019451419, 4.4465546123).. controls (1.382682946, 
    4.4433367498) and (1.3792458498, 4.4236441033) .. (1.3791627971, 
    4.4227678033).. controls (1.3775086221, 4.4180995744) and (1.3778107763, 
    4.4144887792) .. (1.3797747783, 4.4113539958).. controls (1.3850472892, 
    4.4029391081) and (1.4010765942, 4.40093735) .. (1.4012502931, 4.40093735).. 
    controls (1.4458780402, 4.3994114979) and (1.4641354544, 4.3716059327) .. 
    (1.4648682179, 4.370435125).. controls (1.4727543094, 4.3533107625) and 
    (1.4803836492, 4.3495640244) .. (1.4803836492, 4.3495640244).. controls 
    (1.4887910756, 4.3465122673) and (1.494328011, 4.3419799813) .. (1.4965487913,
     4.3357481442).. controls (1.4998119769, 4.3275069554) and (1.4955063329, 
    4.3198398596) .. (1.4952872579, 4.3194546263) -- (1.4950226746, 4.3189712325) 
    -- (1.4945014454, 4.3185785908).. controls (1.4875292777, 4.3135628581) and 
    (1.4799906635, 4.3116895817) .. (1.4723613238, 4.3133514031) -- 
    cycle(1.4735404658, 4.3188431477).. controls (1.4793795025, 4.3174911269) and 
    (1.4850901633, 4.3188432535) .. (1.4905817756, 4.3227636906).. controls 
    (1.491057761, 4.3235492385) and (1.492062384, 4.3259438235) .. (1.492062384, 
    4.329380814).. controls (1.492062384, 4.3308312598) and (1.4918427798, 
    4.3324324387) .. (1.4913599152, 4.3338752381).. controls (1.4896223965, 
    4.3384981704) and (1.4853017506, 4.3418520817) .. (1.4782465038, 
    4.3443750425).. controls (1.4777252746, 4.3445496675) and (1.4687060521, 
    4.3486958735) .. (1.4599965508, 4.3676936183).. controls (1.4599965508, 
    4.3676936183) and (1.442517091, 4.3939202296) .. (1.4007295402, 4.395363029)..
     controls (1.3998609131, 4.3954926748) and (1.3817317954, 4.3976217769) .. 
    (1.375069349, 4.4082951479).. controls (1.3721006975, 4.4130087265) and 
    (1.3717078706, 4.4182359935) .. (1.3736718198, 4.4242035119).. controls 
    (1.3736718198, 4.4240209494) and (1.377811279, 4.4481262902) .. (1.4009485094,
     4.45212978) -- (1.4003367927, 4.4519630925).. controls (1.4006389469, 
    4.4520398217) and (1.4303629, 4.4649858577) .. (1.4439144615, 4.4913863531).. 
    controls (1.4447830885, 4.4934411073) and (1.4688419685, 4.5437718158) .. 
    (1.5029548829, 4.5549665219).. controls (1.5031295079, 4.5550564802) and 
    (1.5114075002, 4.5574969969) .. (1.5137643292, 4.5630793348).. controls 
    (1.5148067875, 4.5660404192) and (1.5141948063, 4.5695680558) .. 
    (1.5117550833, 4.573269365).. controls (1.5115804583, 4.5737527588) and 
    (1.502645294, 4.5863450204) .. (1.4839118944, 4.5734969067) -- (1.477944376, 
    4.5701731579).. controls (1.4776422219, 4.5700435121) and (1.4492702102, 
    4.554355546) .. (1.4290109321, 4.5109137144).. controls (1.4287913279, 
    4.5101281665) and (1.4199011427, 4.4863867071) .. (1.3801153235, 
    4.4671623467).. controls (1.3782042381, 4.4666032821) and (1.3403369921, 
    4.4531348262) .. (1.3462138642, 4.4105692946).. controls (1.3464334683, 
    4.4098668258) and (1.354545805, 4.3840328298) .. (1.3917104235, 
    4.3810188552).. controls (1.393712261, 4.3808521677) and (1.4350690702, 
    4.3753535175) .. (1.442736166, 4.3461204194) -- (1.4425165619, 4.3468228881)..
     controls (1.4426462077, 4.3465133256) and (1.4562492835, 4.3225070977) .. 
    (1.4735398838, 4.3188435181) -- cycle;
  
  
  
      \end{scope}
      \begin{scope}[shift={(0.8731702173, -3.3587695292)}]
        \path[fill=c005093,even odd rule] (1.4894027923, 4.3254298704).. controls 
    (1.4882243381, 4.3261246663) and (1.459814491, 4.3418065733) .. (1.4474414367,
     4.3681239896).. controls (1.4474414367, 4.3679123229) and (1.4331043538, 
    4.3900978735) .. (1.3943685071, 4.3935801871).. controls (1.3941938821, 
    4.3935801871) and (1.3541899933, 4.3951060392) .. (1.3546205498, 
    4.4314019777).. controls (1.3548480915, 4.4318324548) and (1.3572870471, 
    4.4611638721) .. (1.3895341463, 4.4722302023).. controls (1.3897537504, 
    4.4722302023) and (1.4286099031, 4.4887427423) .. (1.4389812258, 4.520121135) 
    -- (1.4390711842, 4.5203407392) -- (1.4392907883, 4.5206352204).. controls 
    (1.4407713967, 4.5231733683) and (1.4773316275, 4.5811260083) .. 
    (1.5088385019, 4.5903567123).. controls (1.5102736019, 4.5907496185) and 
    (1.5117541837, 4.5898807269) .. (1.5122829535, 4.5884003831).. controls 
    (1.5123649744, 4.5880908206) and (1.5123649744, 4.5879169894) .. 
    (1.5123649744, 4.5876148352).. controls (1.5123649744, 4.586436381) and 
    (1.5115794265, 4.5852581915) .. (1.5102726494, 4.5848729581).. controls 
    (1.4809942017, 4.5764202085) and (1.4446906167, 4.5187620498) .. 
    (1.4440787677, 4.5176743213).. controls (1.4324912348, 4.4844301133) and 
    (1.3931888888, 4.4675701754) .. (1.3915346344, 4.4670036496).. controls 
    (1.3628151454, 4.4571988527) and (1.3603299406, 4.4321881606) .. 
    (1.3603299406, 4.4311835113).. controls (1.3603299406, 4.4309639071) and 
    (1.3603299406, 4.4306622821) .. (1.3603299406, 4.4304810425).. controls 
    (1.3603299406, 4.4004698577) and (1.3933249375, 4.3992007573) .. 
    (1.3946242269, 4.3991101375).. controls (1.4364571008, 4.3953634788) and 
    (1.4517535365, 4.3717428281) .. (1.4523200358, 4.3706626667).. controls 
    (1.4639075688, 4.345954261) and (1.4918034627, 4.3305217767) .. (1.4920603202,
     4.330438724).. controls (1.4930649431, 4.329872251) and (1.4935409285, 
    4.3288676281) .. (1.4935409285, 4.3278630052).. controls (1.4935409285, 
    4.3274777719) and (1.4934509702, 4.326903626) .. (1.4931027785, 
    4.3265260656).. controls (1.4925815494, 4.3251663719) and (1.4907990515, 
    4.3247282219) .. (1.4894015488, 4.3254306906) -- cycle;
  
  
  
      \end{scope}
      \begin{scope}[shift={(0.8731702173, -3.5274914104)}]
        \path[fill=c005093,even odd rule] (1.4394641698, 4.562520429).. controls 
    (1.4315780783, 4.5858767873) and (1.4313666233, 4.6029558267) .. 
    (1.4385955427, 4.6134933871).. controls (1.445348609, 4.6229053571) and 
    (1.4557575819, 4.6232604279) .. (1.4562485956, 4.6232604279).. controls 
    (1.4701852846, 4.624129055) and (1.4758960248, 4.6104114675) .. (1.4768553775,
     4.6031296315).. controls (1.4764550629, 4.5987257477) and (1.4756769233, 
    4.5935438567) .. (1.4717489985, 4.592410805).. controls (1.4678361285, 
    4.5913609383) and (1.4643010306, 4.5960668175) .. (1.4633417044, 
    4.5974642937).. controls (1.4625635648, 4.5988165792) and (1.4599953338, 
    4.600123224) .. (1.4585072113, 4.6005537804).. controls (1.4554554806, 
    4.6009996033) and (1.4530157048, 4.6003791554) .. (1.4517541715, 
    4.5988996054).. controls (1.4488308167, 4.5955910173) and (1.4500092444, 
    4.588218535) .. (1.4505381465, 4.5856502775).. controls (1.4524569048, 
    4.5802945817) and (1.4554555865, 4.5766385692) .. (1.4598141735, 
    4.5747198902).. controls (1.468614321, 4.5709279348) and (1.4799902931, 
    4.5764195735) .. (1.4802018275, 4.5764195735).. controls (1.505914909, 
    4.5952889425) and (1.4891378385, 4.6209869427) .. (1.4883976137, 
    4.6220369417).. controls (1.4721493925, 4.6504089004) and (1.4375455702, 
    4.635195544) .. (1.4360650148, 4.6347272315).. controls (1.4075268183, 
    4.6216062529) and (1.4150580506, 4.5758530477) .. (1.4151939935, 
    4.5754602208).. controls (1.415328931, 4.5739722042) and (1.414415854, 
    4.5725369719) .. (1.4127995144, 4.5722347912).. controls (1.4113567415, 
    4.5719326371) and (1.4098760538, 4.5730203392) .. (1.4097023548, 
    4.5745009475).. controls (1.4091811256, 4.5766387808) and (1.4013403573, 
    4.624914841) .. (1.4337158323, 4.6398108827).. controls (1.4341087385, 
    4.639900841) and (1.4742797117, 4.6576831429) .. (1.4931037575, 
    4.6249978937).. controls (1.4932334033, 4.6248232688) and (1.5135443281, 
    4.5938839521) .. (1.4829514569, 4.5717135621).. controls (1.4819545069, 
    4.5712754121) and (1.4689164752, 4.5648697438) .. (1.4576311229, 
    4.5695757287).. controls (1.4518448706, 4.57206096) and (1.447614871, 
    4.5768501035) .. (1.445303471, 4.5841319396).. controls (1.4449182377, 
    4.5853858) and (1.4420400737, 4.5964596973) .. (1.4474410398, 4.6025631585).. 
    controls (1.4501829962, 4.6057055354) and (1.454496366, 4.6067404796) .. 
    (1.4598142, 4.6058717731).. controls (1.4613022167, 4.6054788669) and 
    (1.4660384965, 4.6038321002) .. (1.4680478217, 4.6005085367).. controls 
    (1.4687050467, 4.5995944013) and (1.4695736737, 4.5987257742) .. 
    (1.4703592217, 4.5982500533).. controls (1.4705788258, 4.5987260387) and 
    (1.4709712029, 4.5996854179) .. (1.4711447696, 4.6029560383).. controls 
    (1.4711447696, 4.6029560383) and (1.4688407779, 4.6185999777) .. 
    (1.4564601829, 4.617640625).. controls (1.4559767892, 4.617640625) and 
    (1.4482265612, 4.6173307979) .. (1.4432108285, 4.6101849842).. controls 
    (1.4370242881, 4.601377296) and (1.4375455967, 4.5853859587) .. (1.444782189, 
    4.5644846427).. controls (1.4453034181, 4.5629135469) and (1.4445176056, 
    4.561387589) .. (1.4431277494, 4.5608588192).. controls (1.4415111452, 
    4.5603828337) and (1.4399853725, 4.561115465) .. (1.4394641698, 4.5625204025) 
    -- cycle;
  
  
  
      \end{scope}
      \begin{scope}[shift={(0.8731702173, -3.5517163958)}]
        \path[fill=c005093,even odd rule] (1.4238294379, 4.594284399).. controls 
    (1.4236945004, 4.5947603844) and (1.4109653433, 4.642258596) .. (1.4436883221,
     4.656157529).. controls (1.4441717158, 4.6563771331) and (1.4713652998, 
    4.6650406233) .. (1.4849168348, 4.6402944881).. controls (1.4853020681, 
    4.6391237069) and (1.49040088, 4.6234269831) .. (1.4830359119, 4.6138487754)..
     controls (1.4793345763, 4.6089236625) and (1.4733216554, 4.6065668335) .. 
    (1.4650880338, 4.6068690142).. controls (1.4586823919, 4.6070436392) and 
    (1.4466568148, 4.6107970183) .. (1.4462186648, 4.623124776).. controls 
    (1.4460440398, 4.6247337073) and (1.4472235523, 4.6260027283) .. 
    (1.4488324835, 4.626048025).. controls (1.4504035794, 4.6261829625) and 
    (1.4517557325, 4.6249150792) .. (1.4517557325, 4.6233513652).. controls 
    (1.4521486387, 4.6130631217) and (1.4646954185, 4.6124437056) .. 
    (1.4652619444, 4.6124437056).. controls (1.4715844277, 4.6122770181) and 
    (1.4760714171, 4.6138489077) .. (1.4786396746, 4.6171497965).. controls 
    (1.4840859902, 4.6243863094) and (1.4797726204, 4.6379378444) .. 
    (1.4797726204, 4.6381191369).. controls (1.4687062902, 4.6581592871) and 
    (1.4466568148, 4.6512777392) .. (1.4456974621, 4.6509226683).. controls 
    (1.4176351452, 4.6390330606) and (1.4291018694, 4.5962862365) .. 
    (1.4291849221, 4.595893436).. controls (1.4295701554, 4.5942845048) and 
    (1.4287920158, 4.592841679) .. (1.4272285929, 4.5924564456).. controls 
    (1.4256574971, 4.591965379) and (1.4242147242, 4.592841679) .. (1.4238294644, 
    4.5942844519) -- cycle;
  
  
  
      \end{scope}
      \begin{scope}[shift={(0.8731702173, -3.5439656917)}]
        \path[fill=c005093,even odd rule] (1.3630824804, 4.7047428606).. controls 
    (1.3628628763, 4.7059213148) and (1.3585503002, 4.7314153742) .. 
    (1.3684457169, 4.743788455).. controls (1.3718449513, 4.7481016131) and 
    (1.3764149933, 4.7502846373) .. (1.3819066321, 4.7505414948).. controls 
    (1.3820362779, 4.7505414948) and (1.3909257752, 4.751931351) .. (1.3960245871,
     4.7461377169).. controls (1.3995521973, 4.7417866969) and (1.3999450242, 
    4.7355548333) .. (1.397248285, 4.727019031).. controls (1.3970286808, 
    4.7262258102) and (1.3919303717, 4.7142002596) .. (1.4023848529, 
    4.7077039715).. controls (1.4029060821, 4.7073639819) and (1.4115399654, 
    4.7015249981) .. (1.4250915004, 4.7142002596).. controls (1.4262699546, 
    4.7153332054) and (1.4280072617, 4.7152424533) .. (1.4291025838, 
    4.7140706137).. controls (1.43005429, 4.7129829117) and (1.43005429, 
    4.711200255) .. (1.4290126254, 4.7102409023).. controls (1.4117521877, 
    4.6941589452) and (1.399681314, 4.7027853673) .. (1.3991601113, 
    4.7031706271).. controls (1.387882326, 4.7102410081) and (1.3895365804, 
    4.723180509) .. (1.3919311654, 4.729057381).. controls (1.3939705738, 
    4.7351608423) and (1.3938951675, 4.739995309) .. (1.3917115613, 4.742389894)..
     controls (1.3887957735, 4.7458420185) and (1.3825185075, 4.7449656127) .. 
    (1.3825185075, 4.7449656127).. controls (1.3782053229, 4.7446634585) and 
    (1.3751536981, 4.7431754419) .. (1.3728044363, 4.7403427069).. controls 
    (1.3645632475, 4.7299713842) and (1.3685743308, 4.7060938496) .. 
    (1.3685743308, 4.7059201506).. controls (1.3756371448, 4.6778578073) and 
    (1.403782435, 4.6837421671) .. (1.4048701635, 4.6840443477).. controls 
    (1.4064412594, 4.684437254) and (1.4080049733, 4.6835683623) .. (1.4084053673,
     4.681959431).. controls (1.4087075215, 4.6805166581) and (1.4076121465, 
    4.6790360763) .. (1.4062222902, 4.6786432494).. controls (1.405829384, 
    4.6786432494) and (1.3714899598, 4.6711045558) .. (1.3630825598, 4.7047415642)
     -- cycle;
  
  
  
      \end{scope}
      \begin{scope}[shift={(0.8731702173, -3.5622163854)}]
        \path[fill=c005093,even odd rule] (1.3817308165, 4.7185894325).. controls 
    (1.3817308165, 4.7186793908) and (1.3662229258, 4.7443931338) .. 
    (1.3891864308, 4.764002701).. controls (1.3903572121, 4.7649620802) and 
    (1.3920945456, 4.764788249) .. (1.3930992215, 4.7636023865).. controls 
    (1.3941038444, 4.7624316052) and (1.3939678485, 4.7606413021) .. 
    (1.3927970673, 4.7596819494).. controls (1.3734971419, 4.7431694094) and 
    (1.3860440275, 4.7222982823) .. (1.3866105533, 4.7214220088).. controls 
    (1.3938923894, 4.7087089912) and (1.4173316681, 4.7132790333) .. 
    (1.4176338752, 4.7133696796).. controls (1.4191218919, 4.7136718338) and 
    (1.4205571242, 4.7126219671) .. (1.4208593048, 4.7111413587).. controls 
    (1.421161459, 4.709622915) and (1.420156836, 4.7081726808) .. (1.4187291444, 
    4.7079158233).. controls (1.4175054465, 4.7075684254) and (1.390923685, 
    4.7023865344) .. (1.3817307371, 4.7185894325) -- cycle;
  
  
  
      \end{scope}
      \begin{scope}[shift={(0.8731702173, -3.4658374083)}]
        \path[fill=c005093,even odd rule] (1.3050379242, 4.5190333006).. controls 
    (1.3039123867, 4.5194262069) and (1.2801178783, 4.5270553615) .. 
    (1.2733648121, 4.5479189215).. controls (1.2686965831, 4.5623843802) and 
    (1.2733648121, 4.57929742) .. (1.2873468242, 4.598166789) -- (1.2875743658, 
    4.5983863931) -- (1.2879595992, 4.598643039).. controls (1.2946295863, 
    4.603787174) and (1.3116632231, 4.6193405729) .. (1.3095330892, 
    4.6297496252).. controls (1.30857371, 4.6344556102) and (1.3036412681, 
    4.637983141) .. (1.2946295863, 4.6402568115).. controls (1.2930584904, 
    4.6409063635) and (1.284212861, 4.6454765379) .. (1.2833441281, 
    4.6527206179).. controls (1.282822899, 4.6566863781) and (1.2846964135, 
    4.6604707665) .. (1.2888358727, 4.663922891).. controls (1.2904069685, 
    4.6648370265) and (1.3154099613, 4.6775120763) .. (1.3338412067, 
    4.6601308298).. controls (1.3344153525, 4.6593831173) and (1.3460481292, 
    4.6463602727) .. (1.3448244577, 4.6277024646).. controls (1.3438727515, 
    4.6139772306) and (1.3362886554, 4.6009997621) .. (1.3222990762, 
    4.5890572642).. controls (1.3218987617, 4.58875511) and (1.3074332235, 
    4.5788597463) .. (1.3059073715, 4.5657010381).. controls (1.305212311, 
    4.5595975769) and (1.3074332235, 4.55384916) .. (1.3124867123, 4.5481385256)..
     controls (1.3136725748, 4.5469677444) and (1.3134460915, 4.5452227379) .. 
    (1.3124867123, 4.5442180621).. controls (1.3112706873, 4.5430472808) and 
    (1.3094424429, 4.5432134392) .. (1.3083471473, 4.5444376663).. controls 
    (1.3021606069, 4.5511907325) and (1.2995470262, 4.5586388063) .. 
    (1.3004232998, 4.566404089).. controls (1.3021608185, 4.5819951381) and 
    (1.31815205, 4.5930614683) .. (1.3188545188, 4.593544915).. controls 
    (1.3314014044, 4.6041806888) and (1.3383282754, 4.6158963858) .. 
    (1.3391666871, 4.6279673919).. controls (1.3403374683, 4.6442684769) and 
    (1.329920849, 4.6561657575) .. (1.3297093146, 4.6563394565).. controls 
    (1.314450635, 4.6706689723) and (1.2927563098, 4.6597311238) .. (1.2918800363,
     4.6592628113).. controls (1.2900065217, 4.6576009633) and (1.2888358727, 
    4.6557729306) .. (1.2888358727, 4.6538995748).. controls (1.2888358727, 
    4.6537646373) and (1.2888358727, 4.6536720331) .. (1.2889178935, 
    4.6535066685).. controls (1.2893107998, 4.6497600098) and (1.2945001785, 
    4.6463986373) .. (1.2965849894, 4.6454770406).. controls (1.30717544, 
    4.6429164031) and (1.3136716223, 4.6379837496) .. (1.3150162083, 
    4.6308907467).. controls (1.3151987708, 4.6299766113) and (1.3154091146, 
    4.6289719883) .. (1.3154091146, 4.6279673125).. controls (1.3154091146, 
    4.6132449698) and (1.2952329156, 4.5972083094) .. (1.2916675229, 
    4.5945494056).. controls (1.2820364515, 4.5814813173) and (1.2772398202, 
    4.5697502744) .. (1.2772398202, 4.5593336285).. controls (1.2772398202, 
    4.5559872844) and (1.277723214, 4.5527619342) .. (1.2787278369, 
    4.5497177706).. controls (1.2845140892, 4.5316718115) and (1.306397221, 
    4.5245258125) .. (1.3065634588, 4.5243974367).. controls (1.3077794838, 
    4.5239971221) and (1.3085727046, 4.5228186679) .. (1.3085727046, 
    4.5216478867).. controls (1.3085727046, 4.5214732617) and (1.3085727046, 
    4.5211644929) .. (1.3084377671, 4.5208623387).. controls (1.3079617817, 
    4.5194272387) and (1.3063983588, 4.5185586117) .. (1.3050386121, 4.5190343325)
     -- cycle;
  
  
  
      \end{scope}
      \begin{scope}[shift={(0.8731702173, -3.4894512063)}]
        \path[fill=c005093,even odd rule] (1.3079627871, 4.5552851331).. controls 
    (1.307313235, 4.5555047373) and (1.2905286769, 4.5611619258) .. (1.285429865, 
    4.5754536856).. controls (1.2823857015, 4.583740171) and (1.2840400088, 
    4.5931522469) .. (1.290317116, 4.6032667915).. controls (1.2905287827, 
    4.6035234373) and (1.2958918075, 4.6125881417) .. (1.3086199592, 
    4.6237450917).. controls (1.3087416675, 4.6237450917) and (1.3285241665, 
    4.6431734194) .. (1.324732211, 4.6611287588).. controls (1.3226928027, 
    4.6708429358) and (1.3138471469, 4.678638381) .. (1.2985507113, 
    4.6843490419).. controls (1.2970626946, 4.6849610231) and (1.296284555, 
    4.686532119) .. (1.2969870238, 4.6878767579).. controls (1.2973722571, 
    4.6892742871) and (1.2990264321, 4.690059835) .. (1.3004693373, 
    4.6896669288).. controls (1.3177675046, 4.6832991223) and (1.3277913235, 
    4.6740683125) .. (1.3302312317, 4.6622617575).. controls (1.3346275483, 
    4.6413000106) and (1.3132805417, 4.6205120421) .. (1.3124874267, 
    4.6196433356).. controls (1.3002880713, 4.6090075883) and (1.2952344767, 
    4.6003888127) .. (1.2951136679, 4.6003434896).. controls (1.2897051877, 
    4.5917548235) and (1.2883530346, 4.5839970021) .. (1.2907098635, 
    4.5772439358).. controls (1.2948417556, 4.5656564029) and (1.3097001471, 
    4.5606859669) .. (1.30983609, 4.5604669448).. controls (1.31127119, 
    4.5599909594) and (1.3121022463, 4.5584201281) .. (1.3116185879, 
    4.5569394404).. controls (1.3111050317, 4.5555040758) and (1.3095336712, 
    4.5547185279) .. (1.3079625754, 4.5552850008) -- cycle;
  
  
  
      \end{scope}
      \begin{scope}[shift={(0.8731702173, -3.6006571688)}]
        \path[fill=c005093,even odd rule] (1.2282165767, 4.6784253385).. controls 
    (1.2273479496, 4.6797397885) and (1.2277482642, 4.6814770692) .. 
    (1.2289190454, 4.6823533427).. controls (1.256195735, 4.7009961225) and 
    (1.2744079319, 4.6996514306) .. (1.2751935856, 4.6996514306).. controls 
    (1.2768101898, 4.6994768056) and (1.2778448165, 4.6980803348) .. 
    (1.2777240606, 4.6964714035).. controls (1.2775414981, 4.6949455515) and 
    (1.2762812877, 4.6938578494) .. (1.2746724358, 4.6939409285).. controls 
    (1.2746724358, 4.6939409285) and (1.2574574271, 4.6951569535) .. 
    (1.2321446338, 4.6776474106).. controls (1.2308378567, 4.6768618627) and 
    (1.2290929031, 4.6772545044) .. (1.2282166296, 4.6784252856) -- cycle;
  
  
  
      \end{scope}
      \begin{scope}[shift={(0.8731702173, -3.5986934313)}]
        \path[fill=c005093,even odd rule] (1.2243851454, 4.6871889206).. controls 
    (1.2235995975, 4.6885335331) and (1.2242105204, 4.6902330842) .. 
    (1.2255635996, 4.6910187379).. controls (1.2267420538, 4.6915852108) and 
    (1.2559750725, 4.7074935219) .. (1.2765817485, 4.7111042112).. controls 
    (1.2781152735, 4.7114137737) and (1.2795051033, 4.7104546592) .. 
    (1.2798073104, 4.708838055).. controls (1.2801168729, 4.7073122029) and 
    (1.2791125146, 4.7059298608) .. (1.2775411542, 4.7056201925).. controls 
    (1.2578484019, 4.7022209581) and (1.2285245517, 4.6863202671) .. 
    (1.2282148833, 4.6860105988).. controls (1.2268628625, 4.6853158029) and 
    (1.2252161487, 4.6858280362) .. (1.2243851719, 4.6871887883) -- cycle;
  
  
  
      \end{scope}
      \begin{scope}[shift={(0.8731702173, -3.5951657417)}]
        \path[fill=c005093,even odd rule] (1.218245939, 4.6945449458).. controls 
    (1.2174529827, 4.6959046396) and (1.2178530327, 4.6974683006) .. 
    (1.2192053181, 4.6983370071).. controls (1.2205499306, 4.6992133071) and 
    (1.25401324, 4.7190418965) .. (1.2846060848, 4.7222220029).. controls 
    (1.2860866931, 4.7223886904) and (1.2875746569, 4.7214287821) .. 
    (1.2876577096, 4.7198198508).. controls (1.2878773137, 4.7181656758) and 
    (1.2867435742, 4.7168966019) .. (1.2850819908, 4.7167303642).. controls 
    (1.2558488662, 4.7138070094) and (1.2224233923, 4.6937668327) .. 
    (1.2220381325, 4.6935931338).. controls (1.2207691908, 4.6927620775) and 
    (1.2191148042, 4.6931549838) .. (1.2182460713, 4.69454484) -- cycle;
  
  
  
      \end{scope}
      \begin{scope}[shift={(0.8731702173, -3.5204278292)}]
        \path[fill=c005093,even odd rule] (1.1718345563, 4.4249510127).. controls 
    (1.1713133271, 4.425774396) and (1.1286496085, 4.5059124544) .. (1.0159697581,
     4.5216847431).. controls (1.0144891498, 4.5219043473) and (1.0134014477, 
    4.5233844265) .. (1.0137036019, 4.5248649819).. controls (1.0137935602, 
    4.526390834) and (1.0152746977, 4.5273954569) .. (1.0168459523, 
    4.5272216256).. controls (1.1323584319, 4.5111396685) and (1.1763214398, 
    4.4283501677) .. (1.1767142667, 4.4275192173).. controls (1.1774998146, 
    4.4261746048) and (1.1769338708, 4.4244297306) .. (1.1755358125, 
    4.4237272619).. controls (1.1741459563, 4.423032466) and (1.172491649, 
    4.4235605744) .. (1.1718344769, 4.4249509598) -- cycle;
  
  
  
      \end{scope}
      \begin{scope}[shift={(0.8731702173, -3.5063104563)}]
        \path[fill=c005093,even odd rule] (1.1906584963, 4.4116566527).. controls 
    (1.1900467796, 4.4126612756) and (1.1242078629, 4.5124921921) .. 
    (1.0140964023, 4.5226898158).. controls (1.0125253065, 4.5229014825) and 
    (1.0113544194, 4.5241627512) .. (1.0114828481, 4.5258623817).. controls 
    (1.011612494, 4.5273959067) and (1.0131370231, 4.5283931212) .. (1.0145799019,
     4.5281814546).. controls (1.1274787744, 4.5178101319) and (1.1946695531, 
    4.4157131385) .. (1.1954476133, 4.4146178429).. controls (1.1963617488, 
    4.4134018179) and (1.1958479279, 4.4116567585) .. (1.1945789863, 
    4.4108333487).. controls (1.1933174529, 4.409911805) and (1.1915348227, 
    4.4103121196) .. (1.1906585227, 4.4116567321) -- cycle;
  
  
  
      \end{scope}
      \begin{scope}[shift={(0.8731702173, -3.4922909792)}]
        \path[fill=c005093,even odd rule] (1.2094832565, 4.3962017581).. controls 
    (1.2088789481, 4.3972516248) and (1.1446491215, 4.5064417269) .. 
    (1.0119593892, 4.5212546894).. controls (1.0103504579, 4.5214372519) and 
    (1.0093081319, 4.5227882144) .. (1.0094818308, 4.5243970662).. controls 
    (1.0095638517, 4.5258398392) and (1.0110907621, 4.5269198683) .. 
    (1.0125259944, 4.5268367892).. controls (1.148259996, 4.5116309998) and 
    (1.213622795, 4.4001221423) .. (1.2142724, 4.3990344137).. controls 
    (1.2150579479, 4.3976368846) and (1.2145896354, 4.3958920633) .. 
    (1.2133130208, 4.3950686535).. controls (1.2118324125, 4.3944191015) and 
    (1.2102688573, 4.3948940285) .. (1.2094832035, 4.3962015994) -- cycle;
  
  
  
      \end{scope}
      \begin{scope}[shift={(0.8731702173, -3.5565585354)}]
        \path[fill=c005093,even odd rule] (1.104733921, 4.5578152642).. controls 
    (1.1032533127, 4.5580348683) and (1.1022489544, 4.5596056996) .. 
    (1.1025963523, 4.5612069315).. controls (1.1048170002, 4.5749397325) and 
    (1.1016822169, 4.5817758779) .. (1.1016822169, 4.5818664977).. controls 
    (1.0868692544, 4.6149445737) and (1.0480427879, 4.6179055787) .. 
    (1.0475593148, 4.6179055787).. controls (1.0372786383, 4.61995266) and 
    (1.0351031813, 4.6120665421) .. (1.0347935394, 4.6111978356).. controls 
    (1.0347935394, 4.6078817598) and (1.035752654, 4.604920649) .. (1.0377545179, 
    4.602964214).. controls (1.0412443985, 4.5993383904) and (1.0471665938, 
    4.5994668721) .. (1.0471665938, 4.5994668721).. controls (1.0489115208, 
    4.5993372263) and (1.0845125048, 4.5968985617) .. (1.0850337075, 
    4.5721522413).. controls (1.0850337075, 4.5706263892) and (1.0838174179, 
    4.5694026912) .. (1.0823746979, 4.5692364535).. controls (1.0807657667, 
    4.5692364535) and (1.0795042333, 4.570490314) .. (1.0794589102, 
    4.5719785423).. controls (1.0790207602, 4.591588136) and (1.0473478863, 
    4.5938543188) .. (1.0471665938, 4.5939373979).. controls (1.0470369479, 
    4.5938553771) and (1.0391066975, 4.5937627729) .. (1.0338341867, 
    4.5989531306).. controls (1.030571001, 4.6020879404) and (1.0291280958, 
    4.6064918242) .. (1.0293017948, 4.6118852496).. controls (1.0304347406, 
    4.6172862685) and (1.0359717819, 4.6259201519) .. (1.0483448627, 
    4.6234878373).. controls (1.0496971481, 4.6233528998) and (1.0906158867, 
    4.6204285131) .. (1.1067810288, 4.5842232208).. controls (1.1069106746, 
    4.5839586375) and (1.1107014658, 4.5761104344) .. (1.10795181, 4.5602097433)..
     controls (1.107777185, 4.5586460558) and (1.1062973704, 4.5576869413) .. 
    (1.1047337358, 4.5578152642) -- cycle;
  
  
  
      \end{scope}
      \begin{scope}[shift={(0.8731702173, -3.5741131104)}]
        \path[fill=c005093,even odd rule] (1.0968842627, 4.5800759035).. controls 
    (1.0953131669, 4.5800759035) and (1.0940968508, 4.5813826806) .. 
    (1.0940968508, 4.582863289).. controls (1.0940968508, 4.584389141) and 
    (1.0934396258, 4.6223922242) .. (1.0407671696, 4.6250435873).. controls 
    (1.0391960738, 4.6250435873) and (1.0380176196, 4.6264334435) .. 
    (1.0381459425, 4.6279668362).. controls (1.0381909217, 4.6294397717) and 
    (1.0395810425, 4.6307542217) .. (1.0410691915, 4.6306258458).. controls 
    (1.0992409065, 4.627702491) and (1.099716786, 4.5832561158) .. (1.099716786, 
    4.582863289).. controls (1.099716786, 4.5812921931) and (1.0983721735, 
    4.5800759035) .. (1.0968841833, 4.5800759035) -- cycle;
  
  
  
      \end{scope}
      \begin{scope}[shift={(0.8731702173, -3.5804363875)}]
        \path[fill=c005093,even odd rule] (0.9919196358, 4.5451326208).. controls 
    (0.9917450108, 4.5447397146) and (0.9917450108, 4.5442639938) .. 
    (0.9914362421, 4.5440449188).. controls (0.9679968575, 4.5274038971) and 
    (0.9536597746, 4.5379866748) .. (0.9530479256, 4.5383795017).. controls 
    (0.9445498794, 4.5433952344) and (0.9395418196, 4.54896982) .. (0.9388468915, 
    4.5548467185).. controls (0.9384087415, 4.5591674437) and (0.9404558227, 
    4.5627329158) .. (0.9418078965, 4.5638735346).. controls (0.9441648048, 
    4.567310525) and (0.9478207379, 4.5693576062) .. (0.952179325, 4.5699694287)..
     controls (0.9666447837, 4.5716236037) and (0.9850834904, 4.555896691) .. 
    (0.9858162275, 4.5552395454).. controls (1.0099581073, 4.5389913242) and 
    (1.0328386654, 4.532578221) .. (1.0534453415, 4.5365515483).. controls 
    (1.0790224535, 4.5414313381) and (1.0923547813, 4.5607311577) .. (1.092445401,
     4.5609501798).. controls (1.0933595365, 4.562347709) and (1.0951044371, 
    4.5626046194) .. (1.0962752183, 4.5617357277).. controls (1.0976727475, 
    4.5609501798) and (1.0979749017, 4.5591674173) .. (1.0971062746, 
    4.5578076971).. controls (1.0965019662, 4.5570295575) and (1.082467011, 
    4.5364153937) .. (1.0545783931, 4.531105021).. controls (1.0351878215, 
    4.5272677425) and (1.0140976723, 4.5321927231) .. (0.991919821, 4.5451324356) 
    -- cycle(0.9530025231, 4.5643872496).. controls (0.9500414388, 4.56404726) and
     (0.9478206585, 4.5627330746) .. (0.946121081, 4.5603460038).. controls 
    (0.9457358477, 4.5600364413) and (0.9440361644, 4.558208435) .. (0.9444213977,
     4.5553756735).. controls (0.9448595477, 4.5527167433) and (0.9472087831, 
    4.5484034) .. (0.9562355198, 4.5429949198).. controls (0.9569303156, 
    4.5425266073) and (0.9679967517, 4.5348066213) .. (0.9871304394, 
    4.548010679).. controls (0.9855593435, 4.5489700581) and (0.9839881683, 
    4.549755606) .. (0.9824245602, 4.5508432817).. controls (0.9820242456, 
    4.5510999275) and (0.9647713221, 4.5657468904) .. (0.9530024967, 4.5643872496)
     -- cycle;
  
  
  
      \end{scope}
      \begin{scope}[shift={(0.8731702173, -3.5425226542)}]
        \path[fill=c005093,even odd rule] (0.93670628, 4.4985012106).. controls 
    (0.9364041258, 4.4989015252) and (0.9183657073, 4.5178540263) .. 
    (0.9336998725, 4.53833992).. controls (0.9340927788, 4.53889872) and 
    (0.9453327285, 4.5554944715) .. (0.974490341, 4.5369347179).. controls 
    (0.9754420473, 4.5363756533) and (1.0475582035, 4.4861279181) .. 
    (1.0799790017, 4.5283913485).. controls (1.0809836246, 4.5297057985) and 
    (1.0827662813, 4.5298796298) .. (1.083944656, 4.5289202506).. controls 
    (1.0850776019, 4.5280061152) and (1.0854252644, 4.5261327594) .. 
    (1.0845113935, 4.5249543846).. controls (1.0487366577, 4.4784230927) and 
    (0.9745282029, 4.5300531965) .. (0.9714387163, 4.5322286535).. controls 
    (0.9472968365, 4.5475704123) and (0.9385797681, 4.5355826706) .. 
    (0.9382775875, 4.535106791).. controls (0.9257307019, 4.5183298265) and 
    (0.9399320271, 4.5031165494) .. (0.9405437438, 4.5024291354).. controls 
    (0.9658639983, 4.4810368056) and (0.9891750071, 4.4970281165) .. 
    (0.9980128842, 4.5049897729).. controls (0.9991913383, 4.5060396396) and 
    (1.000936239, 4.5059491521) .. (1.0019408883, 4.5047781062).. controls 
    (1.0028925946, 4.503607325) and (1.0028925946, 4.501816916) .. (1.0017212842, 
    4.5007291875).. controls (0.9939635421, 4.4937948552) and (0.9664224279, 
    4.4732259087) .. (0.9367058567, 4.4985008667) -- cycle;
  
  
  
      \end{scope}
      \begin{scope}[shift={(0.8731702173, -3.5050864937)}]
        \path[fill=c005093,even odd rule] (0.9244609929, 4.4565328965).. controls 
    (0.8993220308, 4.4860833358) and (0.9290386021, 4.5126574244) .. (0.929431429,
     4.5128386904).. controls (0.9560962431, 4.5328788671) and (0.9861679117, 
    4.5067806581) .. (0.9863416106, 4.506478451).. controls (1.0377980154, 
    4.4679541917) and (1.0675599098, 4.4952761779) .. (1.0688289044, 
    4.4965301442).. controls (1.0699166065, 4.4975347671) and (1.071699369, 
    4.4975347671) .. (1.0728324206, 4.4964084358).. controls (1.0738370435, 
    4.49527549) and (1.0738370435, 4.49353051) .. (1.0727503998, 4.4925711573).. 
    controls (1.0723574935, 4.4921330073) and (1.0385392985, 4.460460001) .. 
    (0.9828074652, 4.5020284504).. controls (0.981636684, 4.5033352275) and 
    (0.9558329827, 4.5255962373) .. (0.9328694512, 4.5084340129).. controls 
    (0.9319553158, 4.507648465) and (0.9072998269, 4.4853798617) .. (0.9286468335,
     4.4602861963).. controls (0.9300366898, 4.4587151004) and (0.964466866, 
    4.421157655) .. (1.0119575106, 4.4631189312).. controls (1.0131359648, 
    4.4641613896) and (1.0148808654, 4.4641158813) .. (1.0158402181, 
    4.4628913896).. controls (1.0169279202, 4.4617206083) and (1.016844841, 
    4.4600585488) .. (1.0156655931, 4.4590237104).. controls (0.9639901663, 
    4.4131722919) and (0.9248163019, 4.4561458904) .. (0.924461231, 4.4565310708) 
    -- cycle;
  
  
  
      \end{scope}
      \begin{scope}[shift={(0.8731702173, -3.5749441667)}]
        \path[fill=c005093,even odd rule] (0.9793711363, 4.4783260171).. controls 
    (0.9723083223, 4.483432396) and (0.9681688631, 4.493629914) .. (0.9671264577, 
    4.5081785312).. controls (0.9669968119, 4.5097496271) and (0.968168916, 
    4.511094319) .. (0.9696495244, 4.511275585).. controls (0.9712206202, 
    4.511275585) and (0.9726181758, 4.5102257183) .. (0.972701255, 4.5086620308)..
     controls (0.9737058779, 4.495888556) and (0.977014519, 4.4871337317) .. 
    (0.9826345335, 4.4829489494).. controls (0.9873405185, 4.4793306929) and 
    (0.9921371498, 4.4802900456) .. (0.9923108487, 4.480380639).. controls 
    (0.9944862529, 4.4809018681) and (0.9959668612, 4.4817781681) .. 
    (0.9966240069, 4.4832511035).. controls (0.9982781819, 4.4864010475) and 
    (0.997145236, 4.4913184875) .. (0.9965340485, 4.492844366).. controls 
    (0.9915258829, 4.511887434) and (0.997629159, 4.5211258108) .. (0.9978860694, 
    4.5216016904).. controls (0.9986716173, 4.5226893925) and (1.0004617881, 
    4.5230822988) .. (1.0017233479, 4.5222964863).. controls (1.0029848813, 
    4.5214201863) and (1.0033777875, 4.5197281758) .. (1.0025088958, 
    4.5183760492).. controls (1.0025088958, 4.5181643825) and (0.9976292119, 
    4.5106257948) .. (1.0019424229, 4.4944531915).. controls (1.0019424229, 
    4.4944531915) and (1.0047751315, 4.4867030429) .. (1.0016328604, 
    4.4806902015).. controls (1.0001522521, 4.4777215235) and (0.9973197023, 
    4.4756745481) .. (0.993489885, 4.4747982746).. controls (0.9931803225, 
    4.4747982746) and (0.9863440183, 4.4732345871) .. (0.9793718242, 4.4783259113)
     -- cycle;
  
  
  
      \end{scope}
      \begin{scope}[shift={(0.8731702173, -3.5856629667)}]
        \path[fill=c005093,even odd rule] (0.9913065698, 4.4890447377).. controls 
    (0.9841607031, 4.4972859) and (0.9805878492, 4.520861254) .. (0.9802403454, 
    4.523474861).. controls (0.9799757621, 4.5249554694) and (0.9810258933, 
    4.5263075696) .. (0.9825062371, 4.5267003965).. controls (0.9840773329, 
    4.5269200006) and (0.9854674273, 4.5257788527) .. (0.9857317725, 
    4.524260409).. controls (0.9867289871, 4.5175073427) and (0.9904378633, 
    4.4986832969) .. (0.9955744313, 4.4927535081).. controls (0.9965338104, 
    4.4916205623) and (0.9964507312, 4.4898377469) .. (0.9952270333, 
    4.488795315).. controls (0.9940940875, 4.4877906921) and (0.9923112456, 
    4.4878737713) .. (0.9913065698, 4.4890440233) -- cycle;
  
  
  
      \end{scope}
      \begin{scope}[shift={(0.8731702173, -3.506492225)}]
        \path[fill=c005093,even odd rule] (0.8730518163, 4.4149804015).. controls 
    (0.8692976169, 4.4191198606) and (0.8692145377, 4.4256539048) .. 
    (0.8727422538, 4.4341972742) -- (0.8729618579, 4.4348089908) -- (0.8732185038,
     4.4351942242).. controls (0.8734381079, 4.4351942242) and (0.8871703533, 
    4.4505359565) .. (0.887691556, 4.4812949596).. controls (0.887691556, 
    4.4823903346) and (0.8880844623, 4.5085791898) .. (0.9045138438, 
    4.5195549002).. controls (0.9136614685, 4.5256583615) and (0.9258155273, 
    4.5258320604) .. (0.9406284898, 4.5199478065).. controls (0.9419881835, 
    4.5193360898) and (0.9427663231, 4.5177647294) .. (0.9421995856, 
    4.5162917146).. controls (0.9415876044, 4.5148489417) and (0.9400694252, 
    4.5141993896) .. (0.9385813292, 4.5146298667).. controls (0.9256038871, 
    4.5199477006) and (0.9152326702, 4.51994778) .. (0.9076863831, 4.5148494708)..
     controls (0.8937496942, 4.5054375008) and (0.8932661417, 4.4815221308) .. 
    (0.8932661417, 4.4812955417).. controls (0.8927901563, 4.4503174371) and 
    (0.8799337081, 4.434197724) .. (0.8777960071, 4.4317578158).. controls 
    (0.8767990571, 4.4290989121) and (0.8757941696, 4.4260018583) .. 
    (0.8757941696, 4.4232598754).. controls (0.8757941696, 4.4214697046) and 
    (0.876270155, 4.4199059377) .. (0.8772747779, 4.4186822663).. controls 
    (0.8809761135, 4.4145805631) and (0.8915212146, 4.4161591996) .. 
    (0.8949204754, 4.4171564142).. controls (0.920497561, 4.4294464423) and 
    (0.91444699, 4.4751996475) .. (0.9143563438, 4.4755849073).. controls 
    (0.9141367396, 4.4771560031) and (0.9152326438, 4.4785535852) .. 
    (0.9167129875, 4.4788632271).. controls (0.9182311667, 4.4791198729) and 
    (0.9197572304, 4.4779946) .. (0.9198478767, 4.4763705875).. controls 
    (0.9202407829, 4.4744595021) and (0.9266009429, 4.426175716) .. (0.8968843717,
     4.4118386331).. controls (0.8948372904, 4.4112795685) and (0.8797145802, 
    4.4075329362) .. (0.8730521338, 4.41498101) -- cycle;
  
  
  
      \end{scope}
      \begin{scope}[shift={(0.8731702173, -3.5331564042)}]
        \path[fill=c005093,even odd rule] (0.8730527423, 4.444296076).. controls 
    (0.8728781173, 4.4457840927) and (0.8739213694, 4.4473553737) .. 
    (0.8755303006, 4.4475290727).. controls (0.8831596404, 4.4485260227) and 
    (0.8893462867, 4.452152005) .. (0.8934404227, 4.4579382308).. controls 
    (0.90451432, 4.4734083654) and (0.8998082292, 4.4998163485) .. (0.8997176094, 
    4.5001185292) -- (0.899627651, 4.5004206833) -- (0.8997176094, 4.5008209979)..
     controls (0.900329326, 4.5130657029) and (0.9037286663, 4.5219036858) .. 
    (0.910043609, 4.5273952981).. controls (0.9165776267, 4.5332796579) and 
    (0.9238594627, 4.532675376) .. (0.9241616698, 4.532675376).. controls 
    (0.9256044427, 4.5325854177) and (0.9268205735, 4.5312249302) .. 
    (0.9267299802, 4.5297067246).. controls (0.9266003344, 4.5281808725) and 
    (0.925332451, 4.5270101706) .. (0.9237688958, 4.5270931704).. controls 
    (0.9235492917, 4.5271751913) and (0.9183679033, 4.5273953246) .. (0.913525949,
     4.5230820871).. controls (0.9086461592, 4.518768929) and (0.9059041763, 
    4.5111923735) .. (0.9052923538, 4.500730431).. controls (0.9055119579, 
    4.4994162456) and (0.9063877288, 4.4940225821) .. (0.9063877288, 
    4.4866123438).. controls (0.9063877288, 4.4769812723) and (0.9048542038, 
    4.4642153646) .. (0.8979727352, 4.4548033946).. controls (0.8928361938, 
    4.4475291256) and (0.8855996808, 4.4432537235) .. (0.8762707635, 
    4.4418637879).. controls (0.8747523198, 4.4416441838) and (0.8732191123, 
    4.4428231671) .. (0.873052901, 4.4442961025) -- cycle;
  
  
  
      \end{scope}
      \begin{scope}[shift={(0.8731702173, -3.5032963229)}]
        \path[fill=c005093,even odd rule] (1.0140978046, 4.4853363377).. controls 
    (1.0140978046, 4.4855929835) and (1.0140157837, 4.5120465015) .. 
    (1.0053883033, 4.541068171).. controls (1.0046405908, 4.5422466252) and 
    (0.9828100317, 4.5828482606) .. (1.0039907742, 4.6029866242).. controls 
    (1.0051692283, 4.6040364908) and (1.0069518585, 4.6038629242) .. 
    (1.0080395871, 4.6028569783).. controls (1.0090894537, 4.6017240325) and 
    (1.0089989663, 4.5998429773) .. (1.0078199829, 4.598845895).. controls 
    (0.9898722371, 4.5816760771) and (1.0102671142, 4.5435445387) .. 
    (1.0104788073, 4.5432423581).. controls (1.0197246454, 4.5128307794) and 
    (1.0197625073, 4.4865134425) .. (1.0197625073, 4.4853350677).. controls 
    (1.0197625073, 4.4838092156) and (1.018500974, 4.4825401946) .. (1.0169297723,
     4.4825401946).. controls (1.0154039202, 4.4825401946) and (1.0140971696, 
    4.4838091363) .. (1.0140971696, 4.4853350677) -- cycle;
  
  
  
      \end{scope}
      \begin{scope}[shift={(0.8731702173, -3.4966872958)}]
        \path[fill=c005093,even odd rule] (1.0010666521, 4.4823000587).. controls 
    (0.9995029646, 4.482692965) and (0.9986721729, 4.4842188171) .. (0.9990648146,
     4.485835236).. controls (1.0055611292, 4.5089648729) and (0.986427415, 
    4.546137085) .. (0.9863443623, 4.5464392921) -- (0.9860345352, 4.5468396067) 
    -- (0.9859525144, 4.5470062942).. controls (0.9746671621, 4.5940210592) and 
    (0.9977136404, 4.6086301602) .. (0.9986729931, 4.6092344156).. controls 
    (1.0000705223, 4.6100125552) and (1.0017248296, 4.609619649) .. (1.002593536, 
    4.6083202802).. controls (1.0028956902, 4.6078368865) and (1.0028956902, 
    4.6072249052) .. (1.0028956902, 4.606831999).. controls (1.0028956902, 
    4.6058726198) and (1.0025104569, 4.604958749) .. (1.0016341569, 4.604392276)..
     controls (1.000939361, 4.6040901219) and (0.9883923696, 4.5958943356) .. 
    (0.9883923696, 4.5711556352).. controls (0.9883923696, 4.5647878288) and 
    (0.9892609967, 4.5572491352) .. (0.9914365331, 4.5486680098).. controls 
    (0.9927888185, 4.5459637829) and (1.005947315, 4.5195633933) .. (1.005947315, 
    4.4966300508).. controls (1.005947315, 4.492324354) and (1.0055620817, 
    4.4881320046) .. (1.0045119504, 4.4842191346).. controls (1.0041190442, 
    4.4827763617) and (1.0025101129, 4.4819077346) .. (1.0010673929, 4.4823003763)
     -- cycle;
  
  
  
      \end{scope}
      \begin{scope}[shift={(0.8731702173, -3.4974723146)}]
        \path[fill=c005093,even odd rule] (0.9894796219, 4.4914090279).. controls 
    (0.9879990135, 4.4914910488) and (0.9867300983, 4.4929724508) .. 
    (0.9869565552, 4.494415356).. controls (0.9869565552, 4.4946270227) and 
    (0.9888679052, 4.51524116) .. (0.9783226719, 4.5412563956).. controls 
    (0.977975274, 4.5416493019) and (0.9558350465, 4.5920178987) .. (0.9907410494,
     4.6176403075).. controls (0.9919195035, 4.6186449304) and (0.9937020544, 
    4.6182898596) .. (0.9945783279, 4.6170283262).. controls (0.9955755425, 
    4.615857545) and (0.9952281446, 4.6141201056) .. (0.9939666113, 
    4.6132363708).. controls (0.9628072406, 4.5902728394) and (0.9832099494, 
    4.5438775169) .. (0.98351213, 4.5433864767).. controls (0.9945784602, 
    4.5158529296) and (0.9926144581, 4.4948082094) .. (0.992531379, 4.493977259)..
     controls (0.9924414206, 4.4923683277) and (0.9909147748, 4.4913258958) .. 
    (0.9894795425, 4.4914089485) -- cycle;
  
  
  
      \end{scope}
      \begin{scope}[shift={(0.8731702173, -3.6149788)}]
        \path[fill=c005093,even odd rule] (0.962027831, 4.662207809).. controls 
    (0.9607588894, 4.663076436) and (0.9604112269, 4.6648213631) .. (0.9614161144,
     4.6660905165).. controls (0.9622016623, 4.6672991331) and (0.9639918331, 
    4.6676994477) .. (0.9652533929, 4.6667853123).. controls (0.9686979504, 
    4.6643379165) and (0.9761459977, 4.6651311373) .. (0.9788049279, 
    4.6656899373).. controls (0.9803760238, 4.6659095415) and (0.9817206098, 
    4.6651308727) .. (0.9820227904, 4.6633407019).. controls (0.982423105, 
    4.6619884165) and (0.9814184821, 4.6605080727) .. (0.9798473863, 
    4.6602058921).. controls (0.9787144404, 4.6600233296) and (0.9680861544, 
    4.6578868192) .. (0.9620279898, 4.6622077296) -- cycle;
  
  
  
      \end{scope}
      \begin{scope}[shift={(0.8731702173, -3.6177217354)}]
        \path[fill=c005093,even odd rule] (0.9742736473, 4.6758808671) -- 
    (0.9664253913, 4.6770214858).. controls (0.9648617038, 4.6771511317) and 
    (0.9637663817, 4.678539665) .. (0.9639930767, 4.6801109725).. controls 
    (0.9642047433, 4.6815915808) and (0.9656398433, 4.682777364) .. (0.967211045, 
    4.6825959392) -- (0.9750518133, 4.6815082371).. controls (0.97653983, 
    4.6812806954) and (0.9775822883, 4.6798540621) .. (0.9774087217, 
    4.6782828075).. controls (0.9771065675, 4.6768852783) and (0.9757542821, 
    4.6757070888) .. (0.9742738325, 4.6758806554) -- cycle;
  
  
  
      \end{scope}
      \begin{scope}[shift={(0.8731702173, -3.61843955)}]
        \path[fill=c005093,even odd rule] (0.9733964477, 4.6917216035) -- 
    (0.9655103563, 4.6923259119).. controls (0.9639921771, 4.6925084744) and 
    (0.9628968021, 4.6937760931) .. (0.9629872896, 4.6952491608).. controls 
    (0.9630772479, 4.6968202567) and (0.9643771458, 4.6979912496) .. (0.965948374,
     4.6979912496) -- (0.9737891688, 4.697432185).. controls (0.9753602646, 
    4.6973025392) and (0.9764027229, 4.6959515767) .. (0.9763196438, 
    4.6944710742).. controls (0.9762746646, 4.6929073867) and (0.9749676229, 
    4.6916384715) .. (0.9733963948, 4.6917215242) -- cycle;
  
  
  
      \end{scope}
      \begin{scope}[shift={(0.8731702173, -3.3501888271)}]
        \path[fill=c005093,even odd rule] (0.8146597298, 4.313796935) -- 
    (0.814176336, 4.3141972496) -- (0.8137834298, 4.3147108058).. controls 
    (0.8137384506, 4.3149753892) and (0.8094324627, 4.3228160252) .. 
    (0.8125218965, 4.3310949435).. controls (0.8147880527, 4.3371531081) and 
    (0.8203627177, 4.3419044163) .. (0.8289060342, 4.344820204).. controls 
    (0.8287763883, 4.344820204) and (0.8363616221, 4.3485215131) .. (0.8440287179,
     4.3652985042).. controls (0.8449880971, 4.3668696) and (0.863336131, 
    4.3946298156) .. (0.9081678454, 4.3961934237).. controls (0.9080778871, 
    4.3961934237) and (0.9241592885, 4.3981574258) .. (0.9294317994, 
    4.4066100696).. controls (0.9314336369, 4.4097448794) and (0.9316979556, 
    4.4133632417) .. (0.9301268598, 4.4175857535).. controls (0.9299072556, 
    4.4189002035) and (0.9263800952, 4.4385929294) .. (0.9072992977, 4.441901544) 
    -- (0.9069894706, 4.441901544) -- (0.9065965644, 4.4421211481).. controls 
    (0.9052897873, 4.4426045419) and (0.8747498327, 4.4558464879) .. 
    (0.8601937542, 4.4840825037).. controls (0.8600191292, 4.484694485) and 
    (0.8363615162, 4.5345040702) .. (0.804507191, 4.544913149).. controls 
    (0.8041142848, 4.545087774) and (0.7935692896, 4.5480102027) .. (0.7903891831,
     4.556281554).. controls (0.7885159331, 4.5608591631) and (0.789210729, 
    4.5659201131) .. (0.7925193435, 4.5714948046).. controls (0.7926489894, 
    4.5718877108) and (0.796968524, 4.5780289546) .. (0.805375871, 4.5796378858)..
     controls (0.8124386321, 4.5810883317) and (0.8201887806, 4.5789428254) .. 
    (0.8283393231, 4.5733228637) -- (0.8338309619, 4.570407076).. controls 
    (0.835009416, 4.5698406031) and (0.8641141648, 4.5537207312) .. (0.8851136677,
     4.5085264583).. controls (0.8852962302, 4.5080129021) and (0.8939591119, 
    4.4854799006) .. (0.9314334252, 4.467615181).. controls (0.9316080502, 
    4.4675252227) and (0.9750490352, 4.4523187719) .. (0.9685603142, 
    4.4050396617).. controls (0.9680843287, 4.4034685658) and (0.9590576979, 
    4.3740615571) .. (0.9178366465, 4.3707001846).. controls (0.9175800006, 
    4.3707001846) and (0.8785343004, 4.3653898119) .. (0.8717812606, 4.339895779) 
    -- (0.8716992398, 4.3395031373) -- (0.8714796356, 4.3393364498).. controls 
    (0.8709584065, 4.3381127519) and (0.8564928683, 4.3128378469) .. 
    (0.8368757075, 4.30853215).. controls (0.829125506, 4.3069533812) and 
    (0.8216321621, 4.3087438167) .. (0.8146600473, 4.3137971996) -- 
    cycle(0.8356592327, 4.3140165392).. controls (0.8529120504, 4.3178538177) and 
    (0.8664635854, 4.3416330331) .. (0.8665995283, 4.3420334271) -- (0.8662521304,
     4.3413764667).. controls (0.874002279, 4.3704887825) and (0.9155329723, 
    4.3761086912) .. (0.9173157083, 4.3763202521).. controls (0.9545785667, 
    4.3792436069) and (0.9626761131, 4.4051303608) .. (0.9630765071, 
    4.4062105223).. controls (0.9687796008, 4.4483909) and (0.9309124606, 
    4.4619499227) .. (0.9293866085, 4.4623427231).. controls (0.8892608527, 
    4.4815595958) and (0.8803246829, 4.5054748865) .. (0.8800150146, 
    4.5064342392).. controls (0.8598842181, 4.5496191869) and (0.8314366415, 
    4.5653915021) .. (0.8311722698, 4.5655198779) -- (0.8253785563, 
    4.5686162438).. controls (0.806418541, 4.5815558504) and (0.7973992921, 
    4.56900915) .. (0.7971424875, 4.5684416188).. controls (0.7949670833, 
    4.56469496) and (0.7944004517, 4.561257996) .. (0.7954806396, 4.558425261).. 
    controls (0.7978372833, 4.5527599233) and (0.8060257406, 4.5504032002) .. 
    (0.80616171, 4.5501916658).. controls (0.8403652442, 4.5389817992) and 
    (0.864333425, 4.4886962021) .. (0.8653380744, 4.4866566879).. controls 
    (0.878715831, 4.4601655727) and (0.908560884, 4.4473922037) .. (0.908779906, 
    4.4472185048) -- (0.9081679248, 4.4473931298).. controls (0.9313051288, 
    4.4434726927) and (0.93553534, 4.4191571933) .. (0.93553534, 4.418945659).. 
    controls (0.9363208879, 4.4168078256) and (0.9367061213, 4.414632395) .. 
    (0.9367061213, 4.4124116148).. controls (0.9367061213, 4.4093523171) and 
    (0.9359205733, 4.4064366352) .. (0.9341378108, 4.4036417621).. controls 
    (0.9273847446, 4.3928851796) and (0.909346326, 4.3907020496) .. (0.9085606988,
     4.3906189969).. controls (0.8665994225, 4.3891383885) and (0.8491275563, 
    4.3628663219) .. (0.8489538573, 4.3626472998).. controls (0.8403652177, 
    4.3439516562) and (0.831391239, 4.3398047094) .. (0.830953089, 4.3395931485)..
     controls (0.8239431652, 4.337153161) and (0.8194940377, 4.3336256567) .. 
    (0.8178397304, 4.3290933706).. controls (0.8172280138, 4.3275598456) and 
    (0.8169258596, 4.3260870425) .. (0.8169258596, 4.32459892).. controls 
    (0.8169258596, 4.3211543625) and (0.8180966408, 4.3187220479) .. 
    (0.8184895471, 4.3181555221).. controls (0.8240264825, 4.314197329) and 
    (0.8296086617, 4.3127470419) .. (0.8356592327, 4.3140160629) -- cycle;
  
  
  
      \end{scope}
      \begin{scope}[shift={(0.8731702173, -3.358784875)}]
        \path[fill=c005093,even odd rule] (0.815967724, 4.3217439071).. controls 
    (0.8156578969, 4.3222122196) and (0.8155748177, 4.3226125342) .. 
    (0.8155748177, 4.3230885196).. controls (0.8155748177, 4.3241762217) and 
    (0.8160508031, 4.3250448488) .. (0.817055426, 4.3255811592).. controls 
    (0.8173200094, 4.3257478467) and (0.8451178223, 4.3410814827) .. (0.8565769, 
    4.3655784333).. controls (0.8574455271, 4.3669230458) and (0.8726587777, 
    4.3905362352) .. (0.9143557088, 4.3943281906).. controls (0.9159268046, 
    4.394418149) and (0.9487406148, 4.3957783719) .. (0.9487406148, 
    4.4259709021).. controls (0.9487406148, 4.4259709021) and (0.9487406148, 
    4.4262733208) .. (0.9487406148, 4.4265828833).. controls (0.9487406148, 
    4.4274970187) and (0.9464215419, 4.4523714504) .. (0.9178002663, 
    4.4620931681).. controls (0.9159267517, 4.4628334723) and (0.8765792148, 
    4.4795197377) .. (0.8651201635, 4.5129074556).. controls (0.8645084469, 
    4.5138969973) and (0.8281217563, 4.5716381294) .. (0.798797959, 
    4.5801814988).. controls (0.7975290173, 4.5806045675) and (0.7967508777, 
    4.5815261112) .. (0.7967508777, 4.5827498092).. controls (0.7967508777, 
    4.5831427154) and (0.796840836, 4.5833996258) .. (0.796840836, 4.5835279487)..
     controls (0.7972337423, 4.5850990446) and (0.7988879173, 4.5859676717) .. 
    (0.8003609321, 4.5854464425).. controls (0.831830024, 4.5764347606) and 
    (0.8683374969, 4.5183915537) .. (0.8699086985, 4.5158157027) -- (0.870043636, 
    4.5156490152) -- (0.8700806777, 4.5152487006).. controls (0.8805879433, 
    4.4839609277) and (0.9194521658, 4.467493605) .. (0.9197543465, 
    4.4674029852).. controls (0.9517824235, 4.4564725979) and (0.9543582746, 
    4.4270506402) .. (0.9543582746, 4.4267484596).. controls (0.9549625829, 
    4.3902788221) and (0.9150483615, 4.3887076469) .. (0.9146555346, 
    4.3887076469).. controls (0.8760934133, 4.3853159796) and (0.8616732513, 
    4.3632212604) .. (0.8615448754, 4.3629190798).. controls (0.8494285727, 
    4.3371153785) and (0.8208904027, 4.321342984) .. (0.8197120015, 
    4.3206480558).. controls (0.8183144723, 4.3199606683) and (0.8166603237, 
    4.3203382287) .. (0.8159653427, 4.3217434308) -- cycle;
  
  
  
      \end{scope}
      \begin{scope}[shift={(0.8731702173, -3.5274914104)}]
        \path[fill=c005093,even odd rule] (0.8661235429, 4.5560700725).. controls 
    (0.8645069387, 4.5565460579) and (0.8638121429, 4.5582076413) .. 
    (0.8643333721, 4.5595976033).. controls (0.8715623179, 4.580687726) and 
    (0.8720910877, 4.5964601471) .. (0.8659875471, 4.6053962375).. controls 
    (0.8608887352, 4.6125798602) and (0.8531310196, 4.6128442054) .. 
    (0.8527457598, 4.6128442054).. controls (0.8403273294, 4.6137128325) and 
    (0.8379252831, 4.5982503444) .. (0.837842151, 4.5974647171).. controls 
    (0.8381443052, 4.594677411) and (0.8385369469, 4.5938012433) .. (0.8387107781,
     4.5934159835).. controls (0.8395418344, 4.5938012169) and (0.8404557052, 
    4.5946775169) .. (0.8411960094, 4.595757546).. controls (0.8430240156, 
    4.5990358658) and (0.8478207263, 4.6007732788) .. (0.8488254021, 
    4.6009998679).. controls (0.8547022742, 4.6019592471) and (0.8589247067, 
    4.6009099096) .. (0.8616742567, 4.5977215481).. controls (0.8670373873, 
    4.5915425748) and (0.8642047317, 4.5805593238) .. (0.8639404129, 
    4.5794262721).. controls (0.8615459338, 4.5720613833) and (0.8572326698, 
    4.5672269167) .. (0.8514767123, 4.5647870085).. controls (0.84019136, 
    4.5600810235) and (0.8271157046, 4.5664414481) .. (0.826639825, 4.566569771)..
     controls (0.7955710742, 4.5891027196) and (0.8158756225, 4.6199900721) .. 
    (0.8160493744, 4.620382899).. controls (0.8347903146, 4.652841559) and 
    (0.8750067431, 4.6351053475) .. (0.8755279458, 4.6348937867).. controls 
    (0.9079034473, 4.6199901779) and (0.8998058744, 4.5717140119) .. 
    (0.8994054804, 4.5697953065).. controls (0.8992308554, 4.568231619) and 
    (0.8977510408, 4.5672267315) .. (0.8963159144, 4.5674383981).. controls 
    (0.894835306, 4.5676130231) and (0.893747604, 4.5691909981) .. (0.8940497581, 
    4.5706639335).. controls (0.8940497581, 4.571139919) and (0.9015884517, 
    4.616855421) .. (0.8730502552, 4.6297950013).. controls (0.8716979698, 
    4.6304897971) and (0.8370564708, 4.6456126396) .. (0.8207553594, 
    4.6174596765).. controls (0.820279374, 4.6166815369) and (0.814787841, 
    4.6082288667) .. (0.814787841, 4.5974647435).. controls (0.814787841, 
    4.5892235548) and (0.8180586467, 4.5796831031) .. (0.8293818079, 
    4.5713211056).. controls (0.8293818079, 4.5713211056) and (0.8405839752, 
    4.5662677227) .. (0.8494294723, 4.5699690848).. controls (0.8536595777, 
    4.5717140119) and (0.8567038471, 4.5754606177) .. (0.8585393675, 
    4.5810806323).. controls (0.8592341633, 4.5833920323) and (0.8604128821, 
    4.5907570533) .. (0.8574439925, 4.5941109646).. controls (0.8560919717, 
    4.5956820604) and (0.8536596042, 4.5962032896) .. (0.8502150467, 
    4.5956820604).. controls (0.8491273446, 4.5952891542) and (0.84655152, 
    4.5938842167) .. (0.8459018885, 4.5927588115).. controls (0.8448065135, 
    4.5912782031) and (0.8413696025, 4.5866553238) .. (0.83727544, 4.5876599996)..
     controls (0.8334381085, 4.5887098663) and (0.8325694021, 4.5938843754) .. 
    (0.8322143842, 4.597721654).. controls (0.8334380821, 4.6057059588) and 
    (0.8390128265, 4.6194237315) .. (0.8530477552, 4.6184190821).. controls 
    (0.853304401, 4.6184190821) and (0.8637212585, 4.6181994779) .. (0.8704290017,
     4.6086218523).. controls (0.87783924, 4.5980314017) and (0.877574789, 
    4.5810807381) .. (0.8696434538, 4.557815026).. controls (0.8691222246, 
    4.5562439302) and (0.8674680496, 4.5555414615) .. (0.8661233842, 4.556070099) 
    -- cycle;
  
  
  
      \end{scope}
      \begin{scope}[shift={(0.8731702173, -3.5515952167)}]
        \path[fill=c005093,even odd rule] (0.8818895081, 4.5874022425).. controls 
    (0.8803710644, 4.5877951488) and (0.879495029, 4.5893662446) .. (0.879933179, 
    4.5909297733).. controls (0.8800151998, 4.5913300879) and (0.8913922302, 
    4.6340693979) .. (0.8632468342, 4.6460950544).. controls (0.8625443654, 
    4.6462299919) and (0.8403664877, 4.6533315673) .. (0.8292095377, 
    4.6327248912).. controls (0.8292095377, 4.6325502662) and (0.8277366023, 
    4.6277997783) .. (0.8277366023, 4.6224366477).. controls (0.8277366023, 
    4.6187806352) and (0.8284239898, 4.6150717325) .. (0.8305692579, 
    4.6121559448).. controls (0.8330468163, 4.6089758383) and (0.8376622079, 
    4.6073214781) .. (0.8438563154, 4.6076160917).. controls (0.84441538, 
    4.6076160917) and (0.8569697004, 4.6082354813) .. (0.8574456065, 
    4.6183878083).. controls (0.8574456065, 4.6199514958) and (0.8587523835, 
    4.6211751144) .. (0.8601951565, 4.6211751144).. controls (0.861758844, 
    4.6211751144) and (0.863065621, 4.6197852581) .. (0.8629372188, 
    4.6182141094).. controls (0.8624612333, 4.6058334615) and (0.8503903596, 
    4.6021319406) .. (0.8441131729, 4.6019129185).. controls (0.8358795777, 
    4.601746231) and (0.8297836835, 4.603914756) .. (0.8260747808, 4.6089757325)..
     controls (0.818845835, 4.6185237777) and (0.8239446204, 4.6342884994) .. 
    (0.8240276996, 4.6349003483).. controls (0.8377529335, 4.6602131152) and 
    (0.8649843529, 4.6513751323) .. (0.865203375, 4.6513751323).. controls 
    (0.8981453758, 4.6373477177) and (0.8855079763, 4.5898420448) .. 
    (0.8852964154, 4.5893661387).. controls (0.8849413446, 4.5879307742) and 
    (0.8834231654, 4.5870169033) .. (0.8818897198, 4.5874021367) -- cycle;
  
  
  
      \end{scope}
      \begin{scope}[shift={(0.8731702173, -3.5440035271)}]
        \path[fill=c005093,even odd rule] (0.9028508581, 4.6738021152).. controls 
    (0.9014154935, 4.6741950215) and (0.9005394581, 4.6756377944) .. 
    (0.9008037769, 4.6772466727).. controls (0.9011966831, 4.6786820373) and 
    (0.9025939477, 4.6796789873) .. (0.904112365, 4.679293754).. controls 
    (0.9052908192, 4.6790741498) and (0.9335720788, 4.6729788377) .. 
    (0.9405442465, 4.7012526096).. controls (0.9405442465, 4.7013346304) and 
    (0.9447744577, 4.7251831402) .. (0.9363217346, 4.735592219).. controls 
    (0.9339648262, 4.7385079802) and (0.9309132544, 4.7400337529) .. (0.926902171,
     4.7401697223).. controls (0.9266000169, 4.7401697223) and (0.9203228302, 
    4.7410838577) .. (0.9173163963, 4.7376392473).. controls (0.9152314796, 
    4.7351992598) and (0.9151407275, 4.7302667915) .. (0.9171867504, 
    4.7239140133).. controls (0.9196267379, 4.718422401) and (0.9212809923, 
    4.7054827944) .. (0.9097312154, 4.6982916046).. controls (0.9094290613, 
    4.6979894504) and (0.897494104, 4.689272435) .. (0.8801959102, 4.7053544185)..
     controls (0.8791005352, 4.7064873644) and (0.8790552915, 4.7081416981) .. 
    (0.8800133477, 4.7093200729).. controls (0.8811918019, 4.7104456104) and 
    (0.8829820256, 4.7104456104) .. (0.884024431, 4.7094867604).. controls 
    (0.897575966, 4.6967208527) and (0.9061569856, 4.7026958323) .. (0.9065951092,
     4.7029073931).. controls (0.9171024013, 4.7094867604) and (0.9121772883, 
    4.72155774) .. (0.91186762, 4.7220411867).. controls (0.9091709602, 
    4.730758255) and (0.9096014637, 4.7371638175) .. (0.9131368262, 
    4.7412957096).. controls (0.918190315, 4.7471347462) and (0.9271188383, 
    4.7457902396) .. (0.9276023115, 4.7457902396).. controls (0.932784176, 
    4.7454806771) and (0.937271139, 4.7432145208) .. (0.9407156965, 
    4.7391202525).. controls (0.9505204933, 4.7267471717) and (0.9462072294, 
    4.7012531387) .. (0.9460335304, 4.7002484894).. controls (0.9375808073, 
    4.6663471358) and (0.9032414096, 4.6738026708) .. (0.9028485827, 4.6738026708)
     -- cycle;
  
  
  
      \end{scope}
      \begin{scope}[shift={(0.8731702173, -3.5622616292)}]
        \path[fill=c005093,even odd rule] (0.89038684, 4.7030442092).. controls 
    (0.8889514754, 4.7034371154) and (0.8879471171, 4.7047817279) .. 
    (0.8882566796, 4.706345336).. controls (0.8885588338, 4.7079619402) and 
    (0.8900016067, 4.7088832194) .. (0.8915198652, 4.7084905777).. controls 
    (0.8917844485, 4.7084905777) and (0.9153142413, 4.7039130744) .. 
    (0.9225431871, 4.7166336325).. controls (0.9230718246, 4.7176382554) and 
    (0.9356188425, 4.7384641123) .. (0.9163188906, 4.7548935731).. controls 
    (0.9152311885, 4.755898196) and (0.915049949, 4.7577715254) .. (0.9160992865, 
    4.7587308517).. controls (0.9171039094, 4.759992385) and (0.9188412429, 
    4.7600828725) .. (0.9200196177, 4.7591690017).. controls (0.9429831227, 
    4.7395140848) and (0.9276036344, 4.7140200519) .. (0.9273846123, 
    4.713808491).. controls (0.9181916644, 4.6975527292) and (0.8915645798, 
    4.7028253458) .. (0.890386205, 4.7030443679) -- cycle;
  
  
  
      \end{scope}
      \begin{scope}[shift={(0.8731702173, -3.4658826521)}]
        \path[fill=c005093,even odd rule] (1.0006818156, 4.5161181215).. controls 
    (1.0005468781, 4.5164202756) and (1.0005468781, 4.516684859) .. (1.0005468781,
     4.5169944215).. controls (1.0005468781, 4.5182030381) and (1.001332426, 
    4.5192529048) .. (1.0025108802, 4.5196457581).. controls (1.0027304844, 
    4.5197357165) and (1.0245981117, 4.5268369479) .. (1.0304371483, 
    4.5449659863).. controls (1.0313586921, 4.5479648531) and (1.0319630004, 
    4.5510997423) .. (1.0319630004, 4.5544611413).. controls (1.0319630004, 
    4.5649609456) and (1.0270378875, 4.5767750677) .. (1.0174521127, 
    4.589759971).. controls (1.0140151223, 4.5925472771) and (0.9937936002, 
    4.6085840169) .. (0.9937936002, 4.6232157133).. controls (0.9937936002, 
    4.6242203362) and (0.993923246, 4.6250889633) .. (0.994096019, 4.6260936656)..
     controls (0.9955766273, 4.6331187235) and (1.0019443808, 4.6380287023) .. 
    (1.0131390869, 4.6408614108).. controls (1.0145818598, 4.6416998754) and 
    (1.0197637508, 4.6450990833) .. (1.02028498, 4.6487097727).. controls 
    (1.02028498, 4.6488447102) and (1.02028498, 4.6490193352) .. (1.02028498, 
    4.6492384102).. controls (1.02028498, 4.650975929) and (1.0191065258, 
    4.6528492054) .. (1.0166666177, 4.6548962867).. controls (1.0163644635, 
    4.6550259325) and (0.9947077621, 4.665886999) .. (0.9791923044, 
    4.6512780302).. controls (0.9791923044, 4.6512780302) and (0.9687832521, 
    4.6396074181) .. (0.9699616004, 4.6233063331).. controls (0.9708302275, 
    4.6110616546) and (0.9777117754, 4.5995194183) .. (0.9905683029, 
    4.5886193258).. controls (0.9909157008, 4.5883171717) and (1.0070430869, 
    4.5772055977) .. (1.0088258229, 4.5615163617).. controls (1.00969445, 
    4.553894589) and (1.0070430604, 4.5464389481) .. (1.0007659002, 
    4.5396481523).. controls (0.9996781981, 4.5385152065) and (0.9978878685, 
    4.538386619) .. (0.9967548433, 4.5393915065).. controls (0.9956218975, 
    4.5405622877) and (0.9955763892, 4.5422997271) .. (0.9965352392, 
    4.5433950227).. controls (1.001634051, 4.5489697142) and (1.0039001279, 
    4.554854074) .. (1.0032052262, 4.5609122121).. controls (1.0017246179, 
    4.5741162433) and (0.9872137831, 4.5840041194) .. (0.9871307304, 
    4.5840947392).. controls (0.9728314035, 4.5962487979) and (0.9652550333, 
    4.6091885104) .. (0.9643787598, 4.6230421202).. controls (0.963079391, 
    4.6416999548) and (0.9746670033, 4.6546773969) .. (0.9751429094, 
    4.6552061667).. controls (0.9937931504, 4.6727611121) and (1.0187207633, 
    4.659995125) .. (1.0197631687, 4.6594285992).. controls (1.0244238571, 
    4.6556819404) and (1.0263424302, 4.6519807371) .. (1.0258665506, 
    4.6480602735).. controls (1.0249902506, 4.6406953848) and (1.0160617538, 
    4.6361176698) .. (1.0149740252, 4.6356871133).. controls (1.0055620552, 
    4.6331188029) and (1.0005463225, 4.6296668106) .. (0.9995869698, 
    4.6250892015).. controls (0.9974115656, 4.6146725556) and (1.0144905785, 
    4.5991267502) .. (1.0212436183, 4.593891969) -- (1.0215534454, 4.5936803023) 
    -- (1.0217730496, 4.5932799877).. controls (1.0357550617, 4.574501265) and 
    (1.0404610467, 4.5575958188) .. (1.0357550617, 4.5430924981).. controls 
    (1.0290019954, 4.5223044502) and (1.0052528631, 4.5146298931) .. 
    (1.0042104312, 4.5142446333).. controls (1.0027298229, 4.5138517271) and 
    (1.0011587006, 4.5146298667) .. (1.000682821, 4.5161178833) -- cycle;
  
  
  
      \end{scope}
      \begin{scope}[shift={(0.8731702173, -3.4894512063)}]
        \path[fill=c005093,even odd rule] (0.997630085, 4.5521422535).. controls 
    (0.9970181037, 4.5535019473) and (0.99780471, 4.5552846304) .. (0.99928426, 
    4.5556774573).. controls (0.9995038642, 4.5557594781) and (1.0142709215, 
    4.560859216) .. (1.0185010269, 4.5725373952).. controls (1.0207671831, 
    4.5792451383) and (1.019369654, 4.586874584) .. (1.0140972225, 4.595501006).. 
    controls (1.0140152017, 4.5956835685) and (1.0088247117, 4.6042256415) .. 
    (0.9967537321, 4.6147178787).. controls (0.9958395967, 4.6157148288) and 
    (0.9744926958, 4.6364576065) .. (0.9788437158, 4.6575099996).. controls 
    (0.9814194346, 4.6692259083) and (0.9914357923, 4.6784642852) .. 
    (1.0086055838, 4.684749039).. controls (1.0100861921, 4.6853533473) and 
    (1.0116121235, 4.6845664765) .. (1.0123523483, 4.683087191).. controls 
    (1.0128357421, 4.6816065827) and (1.0120049504, 4.6800430275) .. 
    (1.0105695858, 4.6794689346).. controls (0.9952278535, 4.6738413529) and 
    (0.9864276267, 4.666136501) .. (0.9843805454, 4.6562864075).. controls 
    (0.9806339131, 4.638467011) and (1.0003718827, 4.6190386833) .. (1.0005455817,
     4.6188120942).. controls (1.0132284102, 4.6078364102) and (1.0185841854, 
    4.5988171613) .. (1.0188032075, 4.5984243344).. controls (1.0250803677, 
    4.5884457327) and (1.0267346221, 4.5790338685) .. (1.0236828915, 
    4.5705811454).. controls (1.0185840796, 4.5563800054) and (1.0018524381, 
    4.550669371) .. (1.00115751, 4.5504050258).. controls (0.9995864142, 
    4.5500121196) and (0.9981058852, 4.5506696092) .. (0.9976300056, 4.55214228) 
    -- cycle;
  
  
  
      \end{scope}
      \begin{scope}[shift={(0.8731702173, -3.5944262313)}]
        \path[fill=c005093,even odd rule] (1.11571659, 4.6613547394).. controls 
    (1.0900035085, 4.6789096848) and (1.035661505, 4.6849679552) .. (1.0351025463,
     4.685096331).. controls (1.0336145296, 4.685270956) and (1.0324436425, 
    4.686539104) .. (1.032617315, 4.6881026592).. controls (1.0328369192, 
    4.689673755) and (1.0342262463, 4.6908446421) .. (1.0356614785, 
    4.6906709696).. controls (1.0380183869, 4.6904593029) and (1.0919598377, 
    4.6843937035) .. (1.1189797492, 4.6659700517).. controls (1.120112695, 
    4.6649654288) and (1.1205056013, 4.6631827456) .. (1.1196369742, 
    4.6620043973).. controls (1.1187606742, 4.6607354556) and (1.1169781762, 
    4.6604785452) .. (1.1157166429, 4.6613548452) -- cycle;
  
  
  
      \end{scope}
      \begin{scope}[shift={(0.8731702173, -3.5757749583)}]
        \path[fill=c005093,even odd rule] (1.1411121185, 4.6390406277).. controls 
    (1.1405078102, 4.639425861) and (1.09422565, 4.6813494871) .. (1.0322242235, 
    4.6779049296).. controls (1.0306531277, 4.6777699921) and (1.0293915944, 
    4.6789473879) .. (1.0293009746, 4.6805562662).. controls (1.0291713288, 
    4.6820067121) and (1.0303963496, 4.6833888954) .. (1.0319598783, 
    4.6833888954).. controls (1.096235081, 4.6869240727) and (1.1443828713, 
    4.6435729137) .. (1.1448210213, 4.6431725198).. controls (1.1460370463, 
    4.6420848177) and (1.1460370463, 4.6403398112) .. (1.1450326879, 
    4.6391690035).. controls (1.1440354733, 4.6379982223) and (1.142290705, 
    4.6379074702) .. (1.1411123567, 4.6390393577) -- cycle;
  
  
  
      \end{scope}
      \begin{scope}[shift={(0.8731702173, -3.5699363979)}]
        \path[fill=c005093,even odd rule] (1.1593703529, 4.6386483035).. controls 
    (1.1588943675, 4.6390787806) and (1.1153166988, 4.6837821192) .. 
    (1.0242029565, 4.6866601508).. controls (1.022639269, 4.6866601508) and 
    (1.0213703538, 4.6879669279) .. (1.0214144598, 4.6894928594).. controls 
    (1.0214144598, 4.6911094635) and (1.0227225598, 4.6922801654) .. 
    (1.0242937085, 4.6922801654).. controls (1.1178473325, 4.6893266217) and 
    (1.1615987002, 4.644404261) .. (1.163472056, 4.6425687406).. controls 
    (1.1644692706, 4.6413450427) and (1.1644692706, 4.6395623331) .. 
    (1.1633820977, 4.6386483035).. controls (1.1620753206, 4.6374246056) and 
    (1.1604209075, 4.6375153577) .. (1.159370935, 4.6386483035) -- cycle;
  
  
  
      \end{scope}
      \begin{scope}[shift={(0.8731702173, -3.5575097125)}]
        \path[fill=c005093,even odd rule] (1.1054279231, 4.5935895767).. controls 
    (1.1039473148, 4.5947151142) and (1.07784138, 4.613758235) .. (1.0862941296, 
    4.6383306712).. controls (1.0867775233, 4.6398112796) and (1.0884772067, 
    4.6405968275) .. (1.0897839838, 4.6399848462).. controls (1.0913098358, 
    4.63959194) and (1.09205014, 4.637937765) .. (1.0915667463, 4.6365025327).. 
    controls (1.0845114994, 4.6158052104) and (1.108041636, 4.5986353131) .. 
    (1.1083438167, 4.5984163175).. controls (1.1466943771, 4.5798113731) and 
    (1.2063465417, 4.6241294783) .. (1.2069130675, 4.6245222787).. controls 
    (1.2080838488, 4.625473985) and (1.2098740725, 4.6250887517) .. (1.2108787219,
     4.6239028892).. controls (1.2117550219, 4.6227773517) and (1.2114377865, 
    4.6209871015) .. (1.2102670052, 4.6199899927).. controls (1.2076910219, 
    4.6180712344) and (1.147298659, 4.5733151908) .. (1.105428029, 4.5935894973) 
    -- cycle;
  
  
  
      \end{scope}
      \begin{scope}[shift={(0.8731702173, -3.5852245521)}]
        \path[fill=c005093,even odd rule] (1.1258163973, 4.6299676683).. controls 
    (1.1016820846, 4.6327096512) and (1.0997181883, 4.6584227063) .. 
    (1.0997181883, 4.6586417283).. controls (1.0996361675, 4.6602959033) and 
    (1.1007228113, 4.661565189) .. (1.1022486633, 4.6616406746).. controls 
    (1.1038575946, 4.6618232371) and (1.1052551767, 4.6605603808) .. 
    (1.1053004998, 4.6591101996).. controls (1.1054301456, 4.6582415725) and 
    (1.1070829977, 4.6376800344) .. (1.126390649, 4.6355423333).. controls 
    (1.1266023156, 4.6354603125) and (1.1632908429, 4.633502925) .. (1.203212605, 
    4.6662712269).. controls (1.2043833862, 4.6672681769) and (1.20617361, 
    4.666920779) .. (1.2071329363, 4.6658709123).. controls (1.2080846425, 
    4.664609379) and (1.2079184842, 4.6629476633) .. (1.2067326217, 4.661943014)..
     controls (1.1650357171, 4.6277017237) and (1.1274782452, 4.6297941017) .. 
    (1.1258164238, 4.6299678006) -- cycle;
  
  
  
      \end{scope}
      \begin{scope}[shift={(0.8731702173, -3.5770814708)}]
        \path[fill=c005093,even odd rule] (1.1133619835, 4.6447881979).. controls 
    (1.1121835294, 4.6457928208) and (1.1118813752, 4.647447234) .. (1.1127500023,
     4.6487162285).. controls (1.1137546252, 4.649970089) and (1.1155373877, 
    4.650279916) .. (1.1167157625, 4.6493657806).. controls (1.1449517783, 
    4.6284040337) and (1.1937944967, 4.6676231948) .. (1.1942779433, 
    4.6680084546).. controls (1.1954487246, 4.6690130775) and (1.1972767838, 
    4.6687940025) .. (1.1981908133, 4.6676232213).. controls (1.1992861883, 
    4.6664447671) and (1.1990745217, 4.6647073542) .. (1.1978886592, 
    4.6637026783).. controls (1.1956301758, 4.6619577513) and (1.145034884, 
    4.6213032785) .. (1.1133618777, 4.6447879862) -- cycle;
  
  
  
      \end{scope}
      \begin{scope}[shift={(0.8731702173, -3.5545940042)}]
        \path[fill=c005093,even odd rule] (1.1849536033, 4.6045272135).. controls 
    (1.1847710408, 4.6061361448) and (1.1859505533, 4.6074430013) .. 
    (1.1875142408, 4.6074430013).. controls (1.2029843754, 4.6084476242) and 
    (1.2041628296, 4.6167264896) .. (1.2041628296, 4.6169833735) -- (1.2042527879,
     4.6174667673) -- (1.2044723921, 4.6178596735).. controls (1.2061720754, 
    4.6216894908) and (1.2061720754, 4.6254361231) .. (1.2044723921, 
    4.6292734017).. controls (1.1984142275, 4.6434745152) and (1.1727917129, 
    4.6531962594) .. (1.1726179875, 4.6534152815).. controls (1.1611589627, 
    4.656421689) and (1.1441628702, 4.6659997908) .. (1.1434679421, 
    4.6664001848).. controls (1.1421990004, 4.667140489) and (1.141723015, 
    4.6688780077) .. (1.1425085629, 4.670237675).. controls (1.143158115, 
    4.6716201229) and (1.1449485504, 4.6720201729) .. (1.1463006242, 
    4.6712422979).. controls (1.14643027, 4.6712422979) and (1.1632889379, 
    4.6616037388) .. (1.1743628352, 4.6586426544).. controls (1.1757526915, 
    4.6581214252) and (1.2025913898, 4.6478785048) .. (1.2095635575, 
    4.6314943406).. controls (1.2118297138, 4.6262294233) and (1.2118297138, 
    4.6208207579) .. (1.2095635575, 4.615721946) -- (1.20974612, 4.6164622502).. 
    controls (1.2097011408, 4.615941021) and (1.2083033471, 4.6029560119) .. 
    (1.1878250204, 4.6019966592).. controls (1.1863444121, 4.6018220342) and 
    (1.1849546352, 4.6029560383) .. (1.1849546352, 4.6045273988) -- cycle;
  
  
  
      \end{scope}
      \begin{scope}[shift={(0.8731702173, -3.5413907667)}]
        \path[fill=c005093,even odd rule] (1.1786699873, 4.5836940542).. controls 
    (1.1771062998, 4.5839586375) and (1.176192429, 4.5855673042) .. (1.1765398269,
     4.5870025629).. controls (1.1770232206, 4.5884905796) and (1.1784585852, 
    4.5893594712) .. (1.1800674371, 4.5890044004).. controls (1.1947595908, 
    4.5852199062) and (1.20556147, 4.5912327212) .. (1.20556147, 4.5913234733).. 
    controls (1.2229049869, 4.6040817875) and (1.21440702, 4.6202544967) .. 
    (1.2140142196, 4.6209419371) -- (1.2137946154, 4.6212517642) -- (1.2136649696,
     4.6215539183).. controls (1.2093063825, 4.6384215292) and (1.1724893471, 
    4.652713289) .. (1.1586734669, 4.6567696954).. controls (1.1572381023, 
    4.6571549287) and (1.1563242315, 4.6588167767) .. (1.1568454606, 
    4.6601613627).. controls (1.1572383669, 4.6616948877) and (1.158764219, 
    4.6625635148) .. (1.1602448008, 4.6621253648).. controls (1.162420205, 
    4.6614758127) and (1.2131439256, 4.6462624298) .. (1.219111444, 4.6230874698) 
    -- (1.2188918398, 4.6236917781).. controls (1.2188918398, 4.623472174) and 
    (1.230267706, 4.6026091167) .. (1.2086941102, 4.5866177794).. controls 
    (1.2079159706, 4.5862174648) and (1.19575453, 4.5791622444) .. (1.1786678971, 
    4.5836945304) -- cycle;
  
  
  
      \end{scope}
      \begin{scope}[shift={(0.8731702173, -3.5532866979)}]
        \path[fill=c005093,even odd rule] (1.2266522542, 4.6080546385).. controls 
    (1.2252168896, 4.608530624) and (1.2243030188, 4.6099733969) .. (1.2247787396,
     4.6114463058).. controls (1.2281326773, 4.6234265598) and (1.2271356479, 
    4.6343342194) .. (1.2216363892, 4.6437387281).. controls (1.2088780485, 
    4.6654407262) and (1.1763212017, 4.6732361979) .. (1.176154964, 
    4.6732361979).. controls (1.1745383598, 4.6736365125) and (1.1736623244, 
    4.6752454438) .. (1.1739644785, 4.6765976763).. controls (1.1743647931, 
    4.6781613638) and (1.1758377285, 4.6791659867) .. (1.177325851, 
    4.6787656721).. controls (1.178760951, 4.6783804388) and (1.2123981181, 
    4.6705395381) .. (1.2264783171, 4.6465713573).. controls (1.2328385565, 
    4.6357618846) and (1.2340171165, 4.6234266656) .. (1.2300437892, 
    4.6099280208).. controls (1.2296963912, 4.6084400042) and (1.2281327037, 
    4.6076618646) .. (1.2266522277, 4.6080545063) -- cycle;
  
  
  
      \end{scope}
      \begin{scope}[shift={(0.8731702173, -3.6071005667)}]
        \path[fill=c005093,even odd rule] (1.1654271946, 4.643179981).. controls 
    (1.1641656613, 4.6439126123) and (1.1635991883, 4.6455667873) .. 
    (1.1643394925, 4.6470474485).. controls (1.1649512092, 4.6483994694) and 
    (1.1666053842, 4.6488831277) .. (1.1679953992, 4.6482259027).. controls 
    (1.1681700242, 4.6480962569) and (1.189002575, 4.6379829823) .. (1.2076906515,
     4.648618809).. controls (1.2090881806, 4.6494043569) and (1.2108783515, 
    4.6488833923) .. (1.21152793, 4.647614186).. controls (1.2123134779, 
    4.6462243298) and (1.2118300842, 4.6444792969) .. (1.2104780633, 
    4.6437843952).. controls (1.1890933006, 4.6317133892) and (1.1664319763, 
    4.6427345285) .. (1.1654273269, 4.6431800869) -- cycle;
  
  
  
      \end{scope}
      \begin{scope}[shift={(0.8731702173, -3.5261018188)}]
        \path[fill=c005093,even odd rule] (1.1188535958, 4.6250972713).. controls 
    (1.1174560667, 4.6258828192) and (1.1171089333, 4.6275372588) .. 
    (1.1178489729, 4.6287986863).. controls (1.1182418792, 4.6296297425) and 
    (1.1284771531, 4.6464066277) .. (1.1690484938, 4.6501079633).. controls 
    (1.1690484938, 4.6501079633) and (1.2176343281, 4.6660993006) .. 
    (1.2220382383, 4.6988676025).. controls (1.2222049258, 4.7000460567) and 
    (1.2232997717, 4.7090651204) .. (1.2136233506, 4.7113690592).. controls 
    (1.2122334944, 4.7119884488) and (1.2031235462, 4.7152441731) .. (1.197631934,
     4.702085465).. controls (1.1972467006, 4.7010808421) and (1.1867394085, 
    4.6719837398) .. (1.1546660083, 4.6669755742).. controls (1.1540617, 
    4.6668856158) and (1.1025977546, 4.6601695119) .. (1.0962752712, 
    4.6350758729).. controls (1.0958900379, 4.6335047771) and (1.0943189421, 
    4.6327115563) .. (1.0928382808, 4.6331118708).. controls (1.0914859954, 
    4.6335047771) and (1.0905268808, 4.6349851208) .. (1.0907911996, 
    4.6364657821).. controls (1.0980654685, 4.6651021919) and (1.1516143571, 
    4.6722405179) .. (1.1538351373, 4.6725955888).. controls (1.182645246, 
    4.6769012856) and (1.192057216, 4.7029541715) .. (1.1924500429, 
    4.7041401133).. controls (1.2002908377, 4.7228735129) and (1.2151870646, 
    4.7168153748) .. (1.2153683306, 4.716694566).. controls (1.2269558635, 
    4.7139374225) and (1.2285269594, 4.7032638398) .. (1.22752231, 4.6980366256)..
     controls (1.2227332458, 4.6617784431) and (1.1724023256, 4.6452734967) .. 
    (1.1702193015, 4.6447447004).. controls (1.13218595, 4.6410887144) and 
    (1.12277398, 4.6261774327) .. (1.1226833337, 4.626056624).. controls 
    (1.1218977858, 4.6247043385) and (1.1201150233, 4.6241831094) .. 
    (1.1188535165, 4.6250972448) -- cycle;
  
  
  
      \end{scope}
      \begin{scope}[shift={(0.8731702173, -3.5366089521)}]
        \path[fill=c005093,even odd rule] (1.1059937346, 4.640136426).. controls 
    (1.1046417138, 4.6406576552) and (1.1040297325, 4.6424401531) .. 
    (1.1047322013, 4.6438453288).. controls (1.1052534304, 4.6448425433) and 
    (1.11631984, 4.6672846075) .. (1.1545797806, 4.6694979265).. controls 
    (1.1547544056, 4.6694979265) and (1.2090877044, 4.6805265006) .. 
    (1.2115276125, 4.7248823354).. controls (1.2116175708, 4.7264081875) and 
    (1.2130082208, 4.7275789952) .. (1.2144963698, 4.7274583188).. controls 
    (1.2160600573, 4.7273233813) and (1.2172383527, 4.7260155458) .. 
    (1.2171025156, 4.7245350698).. controls (1.2144058558, 4.6758207538) and 
    (1.155841314, 4.66400642) .. (1.1553201113, 4.663832721).. controls 
    (1.1201119013, 4.6618763919) and (1.1101334054, 4.6422667981) .. 
    (1.1097405785, 4.6413905246).. controls (1.1090457827, 4.6401366642) and 
    (1.1073836702, 4.63939636) .. (1.1059938404, 4.6401366642) -- cycle;
  
  
  
      \end{scope}
      \begin{scope}[shift={(0.8731702173, -3.5647468604)}]
        \path[fill=c005093,even odd rule] (1.1137518735, 4.6961166238).. controls 
    (1.1125734194, 4.6967209321) and (1.0871701915, 4.7108841838) .. 
    (1.0840353817, 4.729443805).. controls (1.0828569275, 4.736514186) and 
    (1.0850323317, 4.7432597646) .. (1.0905240233, 4.749446305) -- (1.0906986483, 
    4.7497561321) -- (1.0910084754, 4.7498381529).. controls (1.0913937088, 
    4.7500127779) and (1.099151345, 4.7552466067) .. (1.1075662327, 
    4.7525801358).. controls (1.1133599463, 4.750631215) and (1.1177638565, 
    4.7456985879) .. (1.12080802, 4.7375556125).. controls (1.1207259992, 
    4.7377672792) and (1.1338383523, 4.7151511456) .. (1.1465211808, 
    4.7094480783).. controls (1.1478657933, 4.7086625304) and (1.1485682621, 
    4.7070459263) .. (1.1478657933, 4.705618261).. controls (1.1472993204, 
    4.7042585673) and (1.1456451454, 4.703654259) .. (1.1441644842, 
    4.7043493194).. controls (1.129918021, 4.7108380404) and (1.1164044538, 
    4.7342395896) .. (1.1157170133, 4.7352442654).. controls (1.11314103, 
    4.7419973052) and (1.10996095, 4.7458345837) .. (1.1057761677, 4.7470961171)..
     controls (1.1007226525, 4.7488788796) and (1.0956237612, 4.7461367379) .. 
    (1.0944453865, 4.7455250212).. controls (1.091008396, 4.741521505) and 
    (1.0892183046, 4.7375103158) .. (1.0892183046, 4.7331895906).. controls 
    (1.0892183046, 4.732192376) and (1.0893479504, 4.7313163406) .. (1.0895204588,
     4.7304022846).. controls (1.092179389, 4.7145393496) and (1.1161928665, 
    4.7010785402) .. (1.1164044273, 4.7009123025).. controls (1.1173638065, 
    4.7004212358) and (1.117764121, 4.6994242858) .. (1.117764121, 4.6985479858)..
     controls (1.117764121, 4.6980344296) and (1.1176741627, 4.6975510358) .. 
    (1.1174545585, 4.6970750504).. controls (1.1167142544, 4.6958061088) and 
    (1.115014571, 4.6952018004) .. (1.1137531435, 4.6961156713) -- cycle;
  
  
  
      \end{scope}
      \begin{scope}[shift={(0.8731702173, -3.5727994542)}]
        \path[fill=c005093,even odd rule] (1.1263913898, 4.7094809396).. controls 
    (1.1028537919, 4.7255176) and (1.09283754, 4.7550603665) .. (1.0924447131, 
    4.7563672229).. controls (1.0920518069, 4.757893075) and (1.092792111, 
    4.7593736304) .. (1.0943179631, 4.7598571035).. controls (1.0957985715, 
    4.7604688202) and (1.0973697996, 4.7595549494) .. (1.0978910023, 
    4.7580669327).. controls (1.0979730231, 4.7578923077) and (1.1076957992, 
    4.7290452631) .. (1.129526279, 4.7140510346).. controls (1.130833056, 
    4.7131824075) and (1.1312256977, 4.7113995921) .. (1.1302210748, 
    4.7102287844).. controls (1.1293524477, 4.7088767635) and (1.1276527644, 
    4.708476449) .. (1.1263912575, 4.7094810719) -- cycle;
  
  
  
      \end{scope}
      \begin{scope}[shift={(0.8731702173, -3.4406379625)}]
        \path[fill=c005093,even odd rule] (0.8263403696, 3.7278425094).. controls 
    (0.8239458904, 3.7327675694) and (0.8256831446, 3.7386594698) .. 
    (0.8305251519, 3.7410918108).. controls (0.8336146385, 3.7427536587) and 
    (1.149483694, 3.9054168463) .. (1.483217469, 3.7425724192).. controls 
    (1.4880593969, 3.7404422587) and (1.490098911, 3.7343387975) .. (1.4877496492,
     3.7295421398).. controls (1.4853927408, 3.7247152667) and (1.4793800846, 
    3.7227966406) .. (1.4746287765, 3.7251005265).. controls (1.1499595735, 
    3.8832692104) and (0.8426339404, 3.7255386765) .. (0.8395822098, 
    3.7238389931).. controls (0.8347930927, 3.7212630098) and (0.8289086535, 
    3.7231817681) .. (0.8263403696, 3.7278425094) -- cycle;
  
  
  
      \end{scope}
      \begin{scope}[shift={(0.8731702173, -2.4204917829)}]
        \path[fill=c005093,even odd rule] (0.8237187456, 2.7036102106) -- 
    (0.8187106329, 2.7059517731) -- (0.7441547538, 3.7368319133) -- (0.7597533171,
     3.7288626369).. controls (0.9696503181, 3.6203147579) and (1.1734966952, 
    3.6244088146) .. (1.3075688754, 3.6471532975).. controls (1.453107674, 
    3.6716878983) and (1.5504003631, 3.7227137467) .. (1.5511934781, 3.7231065735)
     -- (1.5669204702, 3.7315592967) -- (1.4952879458, 2.7120628279) -- 
    (1.4908010888, 2.7093963306).. controls (1.181896634, 2.5325320296) and 
    (0.8272539494, 2.7017368283) .. (0.8237187456, 2.7036101313) -- 
    cycle(1.4764564388, 2.7240885108).. controls (1.4772496596, 2.7348073373) and 
    (1.542650294, 3.6657128923) .. (1.5449163973, 3.6981866865).. controls 
    (1.5117552685, 3.6828902773) and (1.4277873135, 3.6475536915) .. 
    (1.3116176883, 3.6279893944).. controls (1.1777193129, 3.6051921271) and 
    (0.9756631596, 3.6010147002) .. (0.7662496052, 3.7040634792).. controls 
    (0.7685157615, 3.6718238412) and (0.8365374906, 2.7301013523) .. 
    (0.8374515202, 2.718891565).. controls (0.879450526, 2.6998635783) and 
    (1.1969286185, 2.5672568458) .. (1.4764564388, 2.7240885108) -- cycle;
  
  
  
      \end{scope}
      \begin{scope}[shift={(0.8731702173, -3.2878783938)}]
        \path[fill=c005093,even odd rule] (1.1595451631, 3.2515922977).. controls 
    (1.033094835, 3.2499379904) and (1.0249895098, 3.4670262606) .. (1.0245967094,
     3.4762645317).. controls (1.0245967094, 3.5139125704) and (1.0038992813, 
    3.5700371513) .. (1.0037709054, 3.5705054902).. controls (1.0018521471, 
    3.5756118692) and (1.0043828867, 3.5813225565) .. (1.0094816456, 
    3.5831505627).. controls (1.0145804575, 3.5850238127) and (1.020155149, 
    3.582448094) .. (1.0221569071, 3.5774021194).. controls (1.0230332071, 
    3.5749094798) and (1.0442062767, 3.5174401542) .. (1.0442062767, 
    3.4765212833).. controls (1.0442882975, 3.4748594354) and (1.0520471773, 
    3.2691095135) .. (1.159854805, 3.2711565683) -- (1.1609425071, 3.2712473204) 
    -- (1.1620302092, 3.2710282454).. controls (1.1628157571, 3.2708544142) and 
    (1.1839512294, 3.2673268833) .. (1.2075265833, 3.285886531).. controls 
    (1.2324465233, 3.3057605494) and (1.2634247338, 3.354255764) .. (1.2693166342,
     3.4725178465).. controls (1.2693166342, 3.4805701756) and (1.2711898842, 
    3.5549145469) .. (1.2957170502, 3.5805747646).. controls (1.2995468675, 
    3.584631171) and (1.3056048998, 3.5847595469) .. (1.3095329304, 
    3.5809674063).. controls (1.3134458004, 3.5772660442) and (1.3136724954, 
    3.5710795038) .. (1.3099181638, 3.5672497394).. controls (1.2935415931, 
    3.5499893546) and (1.2891376829, 3.49264097) .. (1.2888355023, 3.4719058652)..
     controls (1.2838197431, 3.3701186196) and (1.2604256817, 3.3022705629) .. 
    (1.2188949883, 3.2698950085).. controls (1.190832645, 3.2481099048) and 
    (1.164734436, 3.2507235119) .. (1.1595450044, 3.2515922183) -- cycle;
  
  
  
      \end{scope}
      \begin{scope}[shift={(0.8731702173, -3.3454377042)}]
        \path[fill=c005093,even odd rule] (1.1320923669, 3.380006575).. controls 
    (1.0875628067, 3.4253217096) and (1.0929637198, 3.6170064873) .. 
    (1.0936661885, 3.6386632152).. controls (1.0938408135, 3.6440641548) and 
    (1.0983721735, 3.6483773658) .. (1.1038562452, 3.6480752381).. controls 
    (1.1092193758, 3.6478926756) and (1.11344202, 3.6433692531) .. (1.113268321, 
    3.6379683135).. controls (1.1112664835, 3.5812696133) and (1.1133582794, 
    3.4270665837) .. (1.1460820254, 3.3937770792).. controls (1.1506143115, 
    3.38919947) and (1.1541871654, 3.3886329442) .. (1.1564911042, 3.3887160763)..
     controls (1.1568914188, 3.3887160763) and (1.1639919623, 3.389418545) .. 
    (1.1727014635, 3.3996161688).. controls (1.1882622179, 3.4182664627) and 
    (1.2118301106, 3.4733864471) .. (1.1934064323, 3.6337305881).. controls 
    (1.1927947156, 3.6391390948) and (1.19676037, 3.6440641548) .. (1.2020253138, 
    3.6445476279).. controls (1.2074337675, 3.6452426883) and (1.2123136631, 
    3.6414958708) .. (1.2129255121, 3.6360043115).. controls (1.2314473508, 
    3.4736885748) and (1.2084838458, 3.4098968981) .. (1.1859508973, 
    3.3850147406).. controls (1.1718328365, 3.3694161773) and (1.1581076025, 
    3.3692423725) .. (1.1566270471, 3.3692423725).. controls (1.1476909567, 
    3.3690232975) and (1.1394572556, 3.3727700621) .. (1.1320923669, 3.380006575) 
    -- cycle;
  
  
  
      \end{scope}
      \begin{scope}[shift={(0.8731702173, -3.3390324063)}]
        \path[fill=c005093,even odd rule] (0.5508589546, 3.9156977344) -- 
    (0.528764156, 3.910901209).. controls (0.528764156, 3.910901209) and 
    (0.5244585121, 3.8946151788) .. (0.5341726627, 3.8914728019).. controls 
    (0.5341726627, 3.8914728019) and (0.5826226071, 3.8957029073) .. 
    (0.5982287375, 3.9156977344).. controls (0.5982287375, 3.9156977344) and 
    (0.602066069, 3.9226774956) .. (0.602066069, 3.9344387275) -- (0.603063019, 
    4.0980612412).. controls (0.603063019, 4.0980612412) and (0.603063019, 
    4.1723979925) .. (0.5885068875, 4.1739691942).. controls (0.5885068875, 
    4.1739691942) and (0.5750837813, 4.1772324592) .. (0.577704929, 4.1276947598) 
    -- (0.5847602287, 3.9307298248).. controls (0.5847602287, 3.9307298248) and 
    (0.5879403617, 3.9188855402) .. (0.5682023127, 3.9167931358) -- (0.5583974894,
     4.1326122792).. controls (0.5583974894, 4.1326122792) and (0.5545601579, 
    4.1665513888) .. (0.5401399429, 4.1659471333).. controls (0.5401399429, 
    4.1659471333) and (0.5293758727, 4.1654259042) .. (0.533167881, 4.1309126223) 
    -- (0.5508588223, 3.9156977344) -- cycle;
  
  
  
      \end{scope}
      \begin{scope}[shift={(0.8731702173, -3.3293105563)}]
        \path[fill=c005093,even odd rule] (0.6483359525, 3.9091636902).. controls 
    (0.6483359525, 3.9091636902) and (0.630516556, 3.907547086) .. (0.630516556, 
    3.9259406813) -- (0.6374433742, 4.1340549992).. controls (0.6374433742, 
    4.1340549992) and (0.6380553554, 4.1614979) .. (0.6250702404, 4.1631974775).. 
    controls (0.6250702404, 4.1631974775) and (0.6116471606, 4.1658563813) .. 
    (0.6137848617, 4.1330956465) -- (0.6133467117, 3.9225414469).. controls 
    (0.6133467117, 3.9225414469) and (0.6105594056, 3.9027581542) .. 
    (0.6316495548, 3.8952194342).. controls (0.6316495548, 3.8952194342) and 
    (0.6789286385, 3.870374874) .. (0.6869885613, 3.8898183358).. controls 
    (0.6869885613, 3.8898183358) and (0.6886427363, 3.9037174804) .. 
    (0.6661174342, 3.9055077042) -- (0.6821088773, 4.1324837181).. controls 
    (0.6821088773, 4.1324837181) and (0.6848131042, 4.1551376606) .. 
    (0.6714731035, 4.156829671).. controls (0.6714731035, 4.156829671) and 
    (0.6606258483, 4.1605385473) .. (0.6568791631, 4.1265993583) -- (0.6483358996,
     3.9091636902) -- cycle;
  
  
  
      \end{scope}
      \begin{scope}[shift={(0.8731702173, -3.5502434604)}]
        \path[fill=c005093,even odd rule] (1.6809966856, 4.120466105).. controls 
    (1.6809966856, 4.120466105) and (1.699125724, 4.1176335023) .. (1.7018678127, 
    4.1040819673).. controls (1.7018678127, 4.1040819673) and (1.7148073135, 
    4.0832561633) .. (1.7261833915, 4.1040819673).. controls (1.7261833915, 
    4.1040819673) and (1.7223007898, 4.1170291146) .. (1.748754096, 
    4.1215085104).. controls (1.748754096, 4.1215085104) and (1.8007393765, 
    4.1306562675) .. (1.8002181738, 4.1430822121).. controls (1.8002181738, 
    4.1430822121) and (1.8025219008, 4.1578045548) .. (1.7691416442, 
    4.150356481).. controls (1.7691416442, 4.150356481) and (1.7261833915, 
    4.1401588573) .. (1.7137271521, 4.1300972029) -- (1.7092327279, 
    4.1306562675).. controls (1.7092327279, 4.1306562675) and (1.68265094, 
    4.1458695181) .. (1.6566432715, 4.1498806015).. controls (1.6566432715, 
    4.1498806015) and (1.629449661, 4.1571548704) .. (1.624479225, 4.1442605869)..
     controls (1.624479225, 4.1442605869) and (1.62282505, 4.1340629631) .. 
    (1.6447913933, 4.1288735315) -- (1.6713656142, 4.1232989194) -- (1.6809966856,
     4.1204662108) -- cycle;
  
  
  
      \end{scope}
      \begin{scope}[shift={(0.8731702173, -3.5592546396)}]
        \path[fill=c005093,even odd rule] (1.7463594052, 4.1752227371) -- 
    (1.67712136, 4.1752227371).. controls (1.6595588475, 4.1752227371) and 
    (1.6453124108, 4.1895672812) .. (1.6453124108, 4.2072583019) -- (1.6453124108,
     4.208784154).. controls (1.6453124108, 4.2264751746) and (1.6595588475, 
    4.2407593673) .. (1.67712136, 4.2407593673) -- (1.7463594052, 4.2407593673).. 
    controls (1.7639596738, 4.2407593673) and (1.778168275, 4.2264751746) .. 
    (1.778168275, 4.208784154) -- (1.778168275, 4.2072583019).. controls 
    (1.778168275, 4.1895672813) and (1.7639596738, 4.1752227371) .. (1.7463594052,
     4.1752227371) -- cycle(1.7395155075, 4.2245187131) -- (1.6829529354, 
    4.2245187131).. controls (1.6748930127, 4.2245187131) and (1.6683136719, 
    4.2181509067) .. (1.6683136719, 4.2101740631).. controls (1.6683136719, 
    4.20215949) and (1.6748930127, 4.1956632019) .. (1.6829529354, 4.1956632019) 
    -- (1.7395155075, 4.1956632019).. controls (1.7475754302, 4.1956632019) and 
    (1.7541094479, 4.20215949) .. (1.7541094479, 4.2101740631).. controls 
    (1.7541094479, 4.2181509067) and (1.7475754302, 4.2245187131) .. 
    (1.7395155075, 4.2245187131) -- cycle;
  
  
  
      \end{scope}
      \begin{scope}[shift={(0.8731702173, -3.4581255979)}]
        \path[fill=c005093,even odd rule] (1.699125724, 4.1345845627) -- 
    (1.7000848385, 4.2799419894).. controls (1.7000848385, 4.2799419894) and 
    (1.699125724, 4.2961977777) .. (1.7114986725, 4.2961977777).. controls 
    (1.7114986725, 4.2961977777) and (1.7193018171, 4.3001182148) .. 
    (1.7210844473, 4.2776759125) -- (1.7239625848, 4.1334515375) -- (1.6991255917,
     4.1345844833) -- cycle;
  
  
  
      \end{scope}
      \begin{scope}[shift={(0.8731702173, -3.461470725)}]
        \path[fill=c005093,even odd rule] (1.6633884531, 4.2882576583).. controls 
    (1.6633884531, 4.2882576583) and (1.6628219802, 4.3052537508) .. 
    (1.6531983965, 4.3030631596).. controls (1.6531983965, 4.3030631596) and 
    (1.6431293867, 4.3046419283) .. (1.6453124108, 4.2787550685).. controls 
    (1.6453124108, 4.2787550685) and (1.6527151615, 4.1429757438) .. 
    (1.7143313281, 4.1514284669).. controls (1.7143313281, 4.1514284669) and 
    (1.7708485771, 4.1419332325) .. (1.779860259, 4.2832872487).. controls 
    (1.779860259, 4.2832872487) and (1.7815221069, 4.2962345019) .. (1.7718911942,
     4.2977981365).. controls (1.7718911942, 4.2977981365) and (1.7601221042, 
    4.3001550448) .. (1.7577805152, 4.2797143948).. controls (1.7577805152, 
    4.2797143948) and (1.7493657333, 4.1656370681) .. (1.711929176, 
    4.1645946627).. controls (1.711929176, 4.1645946627) and (1.6810341506, 
    4.1497816737) .. (1.6646122304, 4.281497025) -- (1.6633885325, 4.2882576583) 
    -- cycle;
  
  
  
      \end{scope}
    \end{scope}
  
  \end{tikzpicture}

% \end{document}
};
        % \node[opacity=0.2] at (current page.west) {% \documentclass{standalone}
% \usepackage{graphicx} 
% \usepackage{pgfplots}
% \pgfplotsset{compat=1.18}
% \begin{document}
  \definecolor{c005093}{RGB}{0,80,147}
  \def \globalscale {1.2}
  \begin{tikzpicture}[y=1cm, x=1cm, yscale=\globalscale,xscale=\globalscale, every node/.append style={scale=\globalscale}, inner sep=0pt, outer sep=0pt]
    \begin{scope}[cm={ 1.3333333,-0.0,-0.0,-1.3333333,(-1.143, 2.5124065869)}]
      \begin{scope}[scale=154.0]
        % Image image16 not included. Base64 still not supported;
  
  
  
      \end{scope}
      \begin{scope}[shift={(0.3439583333, -0.7672916667)}]
        \path[fill=white,nonzero rule] (0.7408333333, -0.2232202271) -- 
    (2.6458333333, -0.2232202271).. controls (3.1397220458, -0.2232202271) and 
    (3.3866666667, 0.0237243937) .. (3.3866666667, 0.5176131062) -- (3.3866666667,
     2.4226131063).. controls (3.3866666667, 2.9165019246) and (3.1397220458, 
    3.1634464396) .. (2.6458333333, 3.1634464396) -- (0.7408333333, 
    3.1634464396).. controls (0.2469444621, 3.1634464396) and (0.0, 2.9165019246) 
    .. (0.0, 2.4226131063) -- (0.0, 0.5176131062).. controls (0.0, 0.0237243937) 
    and (0.2469444621, -0.2232202271) .. (0.7408333333, -0.2232202271) -- cycle;
  
  
  
      \end{scope}
      \begin{scope}[shift={(0.8731702173, -1.9486754588)}]
        \path[fill=c005093,even odd rule] (0.5631957875, 2.0542641054).. controls 
    (0.4048911342, 2.2125234885) and (0.317697829, 2.4230776881) .. (0.317697829, 
    2.646873675).. controls (0.317697829, 3.1090670769) and (0.6937856806, 
    3.4849283394) .. (1.1558431396, 3.4849283394).. controls (1.6179005721, 
    3.4849283394) and (1.9939882385, 3.1090670769) .. (1.9939882385, 
    2.646873675).. controls (1.9939882385, 2.1848162425) and (1.6179005721, 
    1.8088114435) .. (1.1558431396, 1.8088114435).. controls (0.9319187504, 
    1.8088114435) and (0.7214097946, 1.895959505) .. (0.5631957875, 2.0542641054) 
    -- cycle(0.3570908213, 2.646873675).. controls (0.3570908213, 2.4335321694) 
    and (0.4401824763, 2.2329640327) .. (0.5910391088, 2.0821074002).. controls 
    (0.741805095, 1.9312582554) and (0.9423732317, 1.8480760335) .. (1.1558431396,
     1.8480760335).. controls (1.5962438177, 1.8480760335) and (1.9545122729, 
    2.2064729969) .. (1.9545122729, 2.646873675).. controls (1.9545122729, 
    3.0872743796) and (1.5962438177, 3.4455805908) .. (1.1558431396, 
    3.4455805908).. controls (0.7154423821, 3.4455805908) and (0.3570908213, 
    3.0872743796) .. (0.3570908213, 2.646873675) -- cycle;
  
  
  
      \end{scope}
      \begin{scope}[shift={(0.8731702173, -1.2964583333)}]
        \path[fill=c005093,even odd rule] (0.0, 1.9993628256).. controls (0.0, 
    2.6412532635) and (0.5221856619, 3.1634464396) .. (1.1641666667, 
    3.1634464396).. controls (1.806049511, 3.1634464396) and (2.3282425019, 
    2.6412532635) .. (2.3282425019, 1.9993628256).. controls (2.3282425019, 
    1.3573894408) and (1.806049511, 0.8351131062) .. (1.1641666667, 
    0.8351131062).. controls (0.5221856619, 0.8351131062) and (0.0, 1.3573894408) 
    .. (0.0, 1.9993628256) -- cycle(0.0590403685, 1.9993628256).. controls 
    (0.0590403685, 1.3899839644) and (0.5547802119, 0.8942366069) .. 
    (1.1641666667, 0.8942366069).. controls (1.7734624752, 0.8942366069) and 
    (2.2691190012, 1.3899839644) .. (2.2691190012, 1.9993628256).. controls 
    (2.2691190012, 2.6086587665) and (1.7734624752, 3.1043078048) .. 
    (1.1641666667, 3.1043078048).. controls (0.5547802119, 3.1043078048) and 
    (0.0590403685, 2.6086587665) .. (0.0590403685, 1.9993628256) -- cycle;
  
  
  
      \end{scope}
      \begin{scope}[shift={(0.8731702173, -3.4438328063)}]
        \path[fill=c005093,even odd rule] (0.316519851, 4.787532335) -- 
    (0.2691576087, 4.7155069044).. controls (0.2612261412, 4.7034812215) and 
    (0.2581744106, 4.6923772675) .. (0.2599193377, 4.682005839).. controls 
    (0.2616189152, 4.6716722723) and (0.2668914525, 4.6636577785) .. 
    (0.2756085208, 4.657901821).. controls (0.2831018912, 4.6530673279) and 
    (0.2908596333, 4.6510656756) .. (0.2990479319, 4.6522818065).. controls 
    (0.3045848408, 4.6530673544) and (0.3103332842, 4.6556357177) .. 
    (0.3162629935, 4.6602133269).. controls (0.3124709852, 4.649796681) and 
    (0.3119422419, 4.6407774321) .. (0.3143897435, 4.6331556594).. controls 
    (0.3167842227, 4.6255263196) and (0.3218831404, 4.6190374927) .. 
    (0.3299883333, 4.613810305).. controls (0.3363486256, 4.6096179556) and 
    (0.3431017183, 4.6073594458) .. (0.350247585, 4.606966619).. controls 
    (0.3574388012, 4.6065737127) and (0.364191894, 4.6079715065) .. (0.3706881556,
     4.6110684279).. controls (0.3746085662, 4.6131533446) and (0.3790955027, 
    4.6171642427) .. (0.3842849608, 4.623131761).. controls (0.3910758096, 
    4.6309725294) and (0.3954796404, 4.6363735219) .. (0.3973529962, 4.639130639) 
    -- (0.4409685269, 4.7056566785) -- (0.3165199569, 4.787532335) -- 
    cycle(0.3397930508, 4.7165115538) -- (0.3288551758, 4.6998707438).. controls 
    (0.324844119, 4.6938956583) and (0.3210143281, 4.690322725) .. (0.317524474, 
    4.6893180492).. controls (0.3138608944, 4.6884039138) and (0.3101595588, 
    4.6891434242) .. (0.3062391217, 4.6917199367).. controls (0.3025755421, 
    4.6940239283) and (0.3005284608, 4.6970833848) .. (0.3000525548, 
    4.7009130962).. controls (0.2995313256, 4.7044406271) and (0.3011855006, 
    4.7092826608) .. (0.3051513667, 4.7150688073) -- (0.3162629935, 4.7321026558) 
    -- (0.3397930508, 4.7165115008) -- cycle(0.3887189013, 4.6843476396) -- 
    (0.3758699144, 4.6648285862).. controls (0.3715567033, 4.6581585727) and 
    (0.3672434923, 4.6542458085) .. (0.3631493298, 4.6530674338).. controls 
    (0.3590098706, 4.6518889796) and (0.3549534906, 4.6526745275) .. 
    (0.3509424338, 4.6552505108).. controls (0.3473317179, 4.6576904983) and 
    (0.3451486673, 4.6609158485) .. (0.3447558669, 4.6651383869).. controls 
    (0.3442798815, 4.6693609252) and (0.3462364752, 4.6747318346) .. 
    (0.3506402531, 4.6815226304) -- (0.3635345896, 4.7009132021) -- (0.3887189013,
     4.6843476925) -- cycle;
  
  
  
      \end{scope}
      \begin{scope}[shift={(0.8731702173, -3.4523682646)}]
        \path[fill=c005093,even odd rule] (0.2075404631, 4.6148984304) -- 
    (0.1611375735, 4.5005566527) -- (0.1905897731, 4.4885689375) -- (0.2196114163,
     4.5602089231) -- (0.2415702454, 4.5512349973) -- (0.2146032506, 4.4848675225)
     -- (0.2427486467, 4.4734537944) -- (0.2697685317, 4.539775946) -- 
    (0.2969168454, 4.5287852073) -- (0.2669811992, 4.4551433842) -- (0.298261405, 
    4.4424756898) -- (0.3455405417, 4.5588115527) -- cycle;
  
  
  
      \end{scope}
      \begin{scope}[shift={(0.8731702173, -3.5448724188)}]
        \path[fill=c005093,even odd rule] (0.1418312717, 4.5237774385) -- 
    (0.131029366, 4.4789005069) -- (0.2758278604, 4.4438584287) -- (0.2866751156, 
    4.4887807894) -- cycle;
  
  
  
      \end{scope}
      \begin{scope}[shift={(0.8731702173, -3.5054950104)}]
        \path[fill=c005093,even odd rule] (0.1105510923, 4.3034338938) -- 
    (0.1063209869, 4.257416211) -- (0.18706338, 4.2501042125).. controls 
    (0.2039234237, 4.2484045292) and (0.2169537031, 4.2487897625) .. 
    (0.2260635719, 4.2509728396).. controls (0.2352112231, 4.2531559167) and 
    (0.2432333369, 4.2585567504) .. (0.2501979375, 4.2671001198).. controls 
    (0.2571700788, 4.2755075463) and (0.2612717554, 4.2868834125) .. 
    (0.2626616381, 4.3009939062).. controls (0.2640590879, 4.3159050821) and 
    (0.2630620056, 4.3277947956) .. (0.2598289825, 4.336338165).. controls 
    (0.256565744, 4.34495702) and (0.251202534, 4.3519291083) .. (0.2438451594, 
    4.3571261335).. controls (0.2363895715, 4.362300431) and (0.2270606542, 
    4.3658280412) .. (0.2156922227, 4.3677014235) -- (0.2055852717, 
    4.3244333173).. controls (0.2121646125, 4.3237385215) and (0.2169537031, 
    4.3227791423) .. (0.2200054338, 4.3213891537).. controls (0.22314027, 
    4.3200368683) and (0.225497046, 4.317899379) .. (0.227234406, 4.3152404488).. 
    controls (0.2283297281, 4.3133671987) and (0.2287224756, 4.310798809) .. 
    (0.2285487502, 4.3074901944).. controls (0.2279369013, 4.3023838154) and 
    (0.2257991738, 4.2988637723) .. (0.2216597146, 4.296809124).. controls 
    (0.2176486577, 4.294716799) and (0.2110768575, 4.2941577873) .. (0.2020500944,
     4.294935874) -- (0.1105510923, 4.3034338408) -- cycle;
  
  
  
      \end{scope}
      \begin{scope}[shift={(0.8731702173, -3.5708576771)}]
        \path[fill=c005093,even odd rule] (0.1071066142, 4.2481850044) -- 
    (0.1095464958, 4.2021749152) -- (0.2582200777, 4.21001571) -- (0.2559085983, 
    4.2561089842) -- cycle;
  
  
  
      \end{scope}
      \begin{scope}[shift={(0.8731702173, -3.4535472479)}]
        \path[fill=c005093,even odd rule] (0.1162538156, 4.0058145796) -- 
    (0.1254467635, 3.9638080067) -- (0.2180486019, 3.9268246277) -- (0.1376083365,
     3.90899777) -- (0.1469372802, 3.8667343131) -- (0.2923778654, 3.8988907925) 
    -- (0.2831018119, 3.9411466558) -- (0.1911571192, 3.9780619046) -- 
    (0.2710308587, 3.9955716327) -- (0.2616641325, 4.037978626) -- cycle;
  
  
  
      \end{scope}
      \begin{scope}[shift={(0.8731702173, -3.4633521771)}]
        \path[fill=c005093,even odd rule] (0.2913343488, 3.7691239917) -- 
    (0.2631436031, 3.7560483363) -- (0.292905524, 3.6915163819) -- (0.350639089, 
    3.7181434929).. controls (0.353388639, 3.7343085556) and (0.3540760529, 
    3.747814741) .. (0.352731414, 3.7588432623).. controls (0.3513338848, 
    3.7697358406) and (0.3476779252, 3.7818067673) .. (0.3417104333, 
    3.7947011038).. controls (0.3343454917, 3.8108284369) and (0.3255377771, 
    3.8225896688) .. (0.3153477204, 3.8302567646).. controls (0.3052332023, 
    3.8378861044) and (0.2934719704, 3.841942405) .. (0.2801319698, 
    3.8425014167).. controls (0.2667164835, 3.8430678896) and (0.2529912496, 
    3.8400616938) .. (0.2389109977, 3.833655946).. controls (0.2242264375, 
    3.8269029062) and (0.2127295773, 3.8178836308) .. (0.2045790348, 
    3.8069004327).. controls (0.1965644881, 3.7957057267) and (0.1923796794, 
    3.782985195) .. (0.1920317523, 3.768640545).. controls (0.1919867731, 
    3.757445839) and (0.1953860869, 3.7441058383) .. (0.2026603823, 
    3.7286054883).. controls (0.2095041213, 3.7135281277) and (0.2159927894, 
    3.7031115083) .. (0.2221415738, 3.6968418358).. controls (0.2282374944, 
    3.6907308604) and (0.2353002819, 3.6864251635) .. (0.2431411296, 
    3.684279869).. controls (0.2511179202, 3.6819306335) and (0.2600011469, 
    3.6816284794) .. (0.2699419396, 3.6831014148) -- (0.2585659146, 
    3.7268605083).. controls (0.2526437458, 3.7261580396) and (0.2474543406, 
    3.7271552542) .. (0.2429673777, 3.7297762696).. controls (0.2384351181, 
    3.7325183054) and (0.234733756, 3.7370581056) .. (0.2319011004, 
    3.7431087031).. controls (0.2276709685, 3.7523469742) and (0.2275878365, 
    3.7611017985) .. (0.2313798713, 3.7693807169).. controls (0.2352096092, 
    3.7778787102) and (0.2440173502, 3.7851983552) .. (0.2578709865, 
    3.7916115113).. controls (0.2724649269, 3.7983570633) and (0.2841430267, 
    3.8004041975) .. (0.2929885238, 3.7977527815).. controls (0.3017962119, 
    3.7951392273) and (0.3085040079, 3.7888166646) .. (0.3129909444, 
    3.7790571644).. controls (0.3151287777, 3.7744418256) and (0.3165185546, 
    3.7696451415) .. (0.3171757267, 3.7648107277).. controls (0.3176969558, 
    3.7598403446) and (0.31756096, 3.7540087956) .. (0.3169111433, 3.7469384675) 
    -- (0.3041905852, 3.741061569) -- (0.2913340577, 3.7691239387) -- cycle;
  
  
  
      \end{scope}
      \begin{scope}[shift={(0.8731702173, -3.4192000979)}]
        \path[fill=c005093,even odd rule] (0.3381757323, 3.4664061302) -- 
    (0.3673257512, 3.4348539856) -- (0.4658950815, 3.4494554667) -- (0.4052835112,
     3.3936783102) -- (0.4346904406, 3.3617786619) -- (0.5442806985, 3.4626141219)
     -- (0.5148738485, 3.4945968494) -- (0.4168710442, 3.4801313642) -- 
    (0.4770444644, 3.5355534498) -- (0.4476829375, 3.567370019) -- cycle;
  
  
  
      \end{scope}
      \begin{scope}[shift={(0.8731702173, -3.4606478708)}]
        \path[fill=c005093,even odd rule] (0.53574442, 3.4124272144).. controls 
    (0.521800111, 3.3925985456) and (0.5165276002, 3.3731248419) .. (0.5198361354,
     3.3542554994).. controls (0.5230616179, 3.3353408338) and (0.5347397442, 
    3.3190473423) .. (0.5545608988, 3.3052314621).. controls (0.5748730142, 
    3.2910303485) and (0.594301289, 3.2854934131) .. (0.6129138269, 
    3.2887641923).. controls (0.6315641208, 3.2918990021) and (0.6476838869, 
    3.3033581063) .. (0.6614997671, 3.3229677265).. controls (0.6714707488, 
    3.3373048094) and (0.6772720823, 3.3506372429) .. (0.6788432575, 
    3.3631009965).. controls (0.6804447804, 3.3755647235) and (0.6784429429, 
    3.3874997073) .. (0.6729513042, 3.398921029).. controls (0.6674672325, 
    3.4104632123) and (0.6578362404, 3.4208798581) .. (0.6443300815, 
    3.430291881).. controls (0.6305594979, 3.4398776558) and (0.6176576208, 
    3.4456713694) .. (0.6055866677, 3.4476353715).. controls (0.593432609, 
    3.449637209) and (0.5812332535, 3.4477173923) .. (0.5690716806, 
    3.4422721879).. controls (0.5568269756, 3.4366521733) and (0.5457606983, 
    3.4267642973) .. (0.53574442, 3.4124272144) -- cycle(0.5736039667, 
    3.3862382269).. controls (0.58210196, 3.3985206879) and (0.5906453294, 
    3.4057572008) .. (0.5990073533, 3.4080233571).. controls (0.6073316213, 
    3.4101610581) and (0.6154821108, 3.4085067773) .. (0.6233229321, 
    3.403015165).. controls (0.6314734746, 3.3973045306) and (0.6358773848, 
    3.3902870398) .. (0.6365723394, 3.3819250158).. controls (0.6373578873, 
    3.3736007479) and (0.6329539771, 3.362489121) .. (0.6234135783, 
    3.3487639135).. controls (0.6153914646, 3.3374332117) and (0.6071578165, 
    3.3306348223) .. (0.5988788717, 3.3285424179).. controls (0.5904715246, 
    3.3264953367) and (0.582321035, 3.3282855075) .. (0.5744801873, 
    3.3337771727).. controls (0.5668508475, 3.3390874925) and (0.5626283092, 
    3.3460672273) .. (0.5617596027, 3.3547314054).. controls (0.560928811, 
    3.363320045) and (0.5648491423, 3.3737744998) .. (0.5736039667, 3.3862382269) 
    -- cycle;
  
  
  
      \end{scope}
      \begin{scope}[shift={(0.8731702173, -3.4475269188)}]
        \path[fill=c005093,even odd rule] (0.7466394558, 3.350198405) -- 
    (0.6884677408, 3.2129839273) -- (0.7591031829, 3.183048281).. controls 
    (0.772216515, 3.1774282665) and (0.7827238335, 3.1743311862) .. (0.7905646813,
     3.1736362581).. controls (0.798405529, 3.1729337365) and (0.8057704706, 
    3.1747693098) .. (0.8126142096, 3.1790371977).. controls (0.8193672758, 
    3.1833050856) and (0.8246775956, 3.1898844529) .. (0.8284242808, 
    3.1985184421).. controls (0.8315666577, 3.2061855115) and (0.832745059, 
    3.2135429125) .. (0.8318310294, 3.220386625).. controls (0.8309547294, 
    3.2272756871) and (0.8284242808, 3.2335453067) .. (0.8240733137, 
    3.2394297194).. controls (0.8214520867, 3.2431310815) and (0.8173201946, 
    3.2470514921) .. (0.8116548833, 3.2511078721).. controls (0.8173201946, 
    3.2507150717) and (0.8217618608, 3.2508887971) .. (0.8248135915, 
    3.2515837517).. controls (0.8267320852, 3.2519765521) and (0.8300861023, 
    3.2535477537) .. (0.8346183883, 3.2559423123).. controls (0.8391053512, 
    3.2582537652) and (0.842376104, 3.26030082) .. (0.8442494333, 3.2619173183) --
     (0.881504751, 3.2929785021) -- (0.8336137125, 3.3132830769) -- (0.7932236379,
     3.280915116).. controls (0.7879888831, 3.2766849842) and (0.7839853669, 
    3.2743810719) .. (0.781062065, 3.2735954975).. controls (0.7770132521, 
    3.2727645735) and (0.7730928415, 3.2732026971) .. (0.7693008331, 3.2748569779)
     -- (0.7655541479, 3.2764281531) -- (0.7890766117, 3.3321146898) -- 
    cycle(0.7545708969, 3.2504960231) -- (0.7724356429, 3.242874224).. controls 
    (0.7743919721, 3.2421717287) and (0.7777988, 3.2399509221) .. (0.7828069656, 
    3.2362948831).. controls (0.7854658694, 3.2345046329) and (0.7871201767, 
    3.2322007206) .. (0.7879510742, 3.2290130206).. controls (0.7887744575, 
    3.2260972858) and (0.7885175471, 3.222954909) .. (0.7872030971, 
    3.2199938246).. controls (0.7853750908, 3.2155068881) and (0.7825047056, 
    3.2127648523) .. (0.7787126708, 3.2115864775).. controls (0.7748753658, 
    3.2104080762) and (0.7690816258, 3.2113674025) .. (0.7613239102, 3.2147212608)
     -- (0.7427566954, 3.2225696756) -- (0.7545707646, 3.2504960231) -- cycle;
  
  
  
      \end{scope}
      \begin{scope}[shift={(0.8731702173, -3.4479727417)}]
        \path[fill=c005093,even odd rule] (0.888477051, 3.1364795506) -- 
    (0.9480839983, 3.1254963525) -- (0.987431641, 3.2104235544) -- (0.9939202827, 
    3.1170814383) -- (1.0532250494, 3.1061057808) -- (1.0803281194, 3.2526038527) 
    -- (1.0432464477, 3.259402295) -- (1.0226397452, 3.1476742567) -- 
    (1.014579849, 3.2646295092) -- (0.9811544808, 3.2708613462) -- (0.9319188563, 
    3.164360575) -- (0.9525255323, 3.2761716925) -- (0.9153986433, 3.2829247588) 
    -- cycle;
  
  
  
      \end{scope}
      \begin{scope}[shift={(0.8731702173, -3.471668825)}]
        \path[fill=c005093,even odd rule] (1.223298634, 3.2537669081) -- 
    (1.1709659556, 3.2516291806) -- (1.1625962852, 3.2758995156) -- (1.1157175954,
     3.2740714829) -- (1.1774999998, 3.1274449823) -- (1.2277477615, 3.1294089579)
     -- (1.277723505, 3.2805677446) -- (1.2296967352, 3.2785584194) -- 
    cycle(1.2149742867, 3.2212176548) -- (1.2006825269, 3.166973944) -- 
    (1.1821229056, 3.2198277721) -- cycle;
  
  
  
      \end{scope}
      \begin{scope}[shift={(0.8731702173, -3.4529805104)}]
        \path[fill=c005093,even odd rule] (1.347089926, 3.1242718079) -- 
    (1.391929234, 3.1346431571) -- (1.3663067988, 3.2440974985) -- (1.4362850688, 
    3.2603985835) -- (1.4279229654, 3.2960826467) -- (1.3130601967, 3.269198731) 
    -- cycle;
  
  
  
      \end{scope}
      \begin{scope}[shift={(0.8731702173, -3.4429115271)}]
        \path[fill=c005093,even odd rule] (1.7101021223, 3.2534644894) -- 
    (1.7499256975, 3.2765186671) -- (1.7053961373, 3.3532953).. controls 
    (1.7009998206, 3.3609623958) and (1.6956819867, 3.3674511169) .. 
    (1.6894123935, 3.3729427292).. controls (1.6831352069, 3.3783059127) and 
    (1.6759440171, 3.3821432177) .. (1.6679218504, 3.3842733781).. controls 
    (1.6599072508, 3.3863657825) and (1.6525423621, 3.3868492292) .. 
    (1.6457515398, 3.3856708544).. controls (1.6363320027, 3.3840996527) and 
    (1.6262627813, 3.3800885958) .. (1.615332394, 3.3738189498).. controls 
    (1.6091911767, 3.3701553702) and (1.6026119152, 3.3658421592) .. 
    (1.5955491013, 3.360652701).. controls (1.5886675533, 3.3554254869) and 
    (1.5834781217, 3.3501529231) .. (1.5800335642, 3.3447142275).. controls 
    (1.576513521, 3.3391772921) and (1.5741566921, 3.3329001054) .. (1.5727591894,
     3.3255351373).. controls (1.5714145769, 3.3183892706) and (1.5714148415, 
    3.3118099298) .. (1.5728888352, 3.3058424379).. controls (1.5751623998, 
    3.2962113665) and (1.5780782933, 3.2881892527) .. (1.5818249256, 3.2819196067)
     -- (1.62635438, 3.2049692219) -- (1.6660495265, 3.2279704829) -- 
    (1.6204322642, 3.3067111179).. controls (1.616466504, 3.3137739054) and 
    (1.6152428325, 3.3203532462) .. (1.6168517638, 3.3266228923).. controls 
    (1.6185136117, 3.3329001054) and (1.6226983675, 3.3379158117) .. 
    (1.6296705352, 3.3419193279).. controls (1.6365520831, 3.3458397385) and 
    (1.6430029688, 3.3470181398) .. (1.64910643, 3.3454922612).. controls 
    (1.6552400802, 3.3437095781) and (1.660474835, 3.3393888) .. (1.6646142942, 
    3.3321145046) -- (1.7101031806, 3.2534644894) -- cycle;
  
  
  
      \end{scope}
      \begin{scope}[shift={(0.8731702173, -3.4188151292)}]
        \path[fill=c005093,even odd rule] (1.817125234, 3.3015672475) -- 
    (1.8488888865, 3.3304604619) -- (1.8349447362, 3.4291582208) -- (1.8904119069,
     3.3682369558) -- (1.922439984, 3.3973870012) -- (1.8223146921, 3.5075438115) 
    -- (1.7901505398, 3.478393766) -- (1.8039665258, 3.3803909617) -- (1.7489297, 
    3.4409648287) -- (1.7169922692, 3.4117694073) -- cycle;
  
  
  
      \end{scope}
      \begin{scope}[shift={(0.8731702173, -3.496332225)}]
        \path[fill=c005093,even odd rule] (1.972942804, 3.5346469344) -- 
    (2.0015261383, 3.5709428729) -- (1.8846994469, 3.6631065612) -- (1.8560327423,
     3.6270221306) -- cycle;
  
  
  
      \end{scope}
      \begin{scope}[shift={(0.8731702173, -3.4592881771)}]
        \path[fill=c005093,even odd rule] (2.0401338738, 3.5833087835) -- 
    (2.0636185815, 3.6254135963) -- (1.9864945506, 3.7070171552) -- (2.0961678877,
     3.6834493154) -- (2.1190484194, 3.7241415971) -- (1.9620051671, 3.7488122467)
     -- (1.9376895883, 3.7052344721) -- cycle;
  
  
  
      \end{scope}
      \begin{scope}[shift={(0.8731702173, -3.4589862875)}]
        \path[fill=c005093,even odd rule] (2.1430555469, 3.7844508015) -- 
    (2.181481831, 3.9016175354) -- (2.1512892744, 3.9115582752) -- (2.1271549617, 
    3.8382261998) -- (2.1047049865, 3.8455835215) -- (2.1270641302, 3.9135600333) 
    -- (2.0982163713, 3.9230627554) -- (2.0757288517, 3.8550862435) -- 
    (2.0478930181, 3.8642338683) -- (2.0728204194, 3.9396205921) -- (2.040747125, 
    3.9502563658) -- (2.0015279375, 3.8309518779) -- cycle;
  
  
  
      \end{scope}
      \begin{scope}[shift={(0.8731702173, -3.4613802375)}]
        \path[fill=c005093,even odd rule] (2.0506790806, 3.998139414) -- 
    (2.198567141, 3.9800934283) -- (2.2077601154, 4.0564016694).. controls 
    (2.2094672071, 4.070429084) and (2.2095502863, 4.0813217152) .. (2.2081151863,
     4.0891700771).. controls (2.2064985821, 4.0967918498) and (2.2027973523, 
    4.1034920258) .. (2.1968299398, 4.1087796708).. controls (2.1908170983, 
    4.1141730963) and (2.1830517098, 4.1173986317) .. (2.1736397398, 
    4.118614736).. controls (2.1654060123, 4.1195288715) and (2.1581769606, 
    4.1186597152) .. (2.1518544773, 4.1158273506).. controls (2.1453658621, 
    4.1130853677) and (2.1400026256, 4.1089912052) .. (2.1355988213, 
    4.1032805708).. controls (2.1328567325, 4.0995792352) and (2.1301146438, 
    4.0944804233) .. (2.1278485669, 4.0879010825).. controls (2.1266701127, 
    4.0936117169) and (2.1252799919, 4.0977889321) .. (2.1238373777, 
    4.1005384821).. controls (2.1228401631, 4.1022834092) and (2.1204836517, 
    4.1050707681) .. (2.1169938769, 4.1087796708).. controls (2.1134284048, 
    4.1124356833) and (2.1105052617, 4.1148681038) .. (2.1084959365, 4.1162277446)
     -- (2.0683247781, 4.1435044077) -- (2.062009756, 4.09191206) -- 
    (2.1041902396, 4.0621048425).. controls (2.1096742054, 4.0581844054) and 
    (2.1132017098, 4.0549967054) .. (2.1146898323, 4.0523906919).. controls 
    (2.1165630823, 4.0487724354) and (2.1172584073, 4.044935051) .. (2.1167824219,
     4.0407125127) -- (2.1161704406, 4.0367015352) -- (2.056253825, 4.0440664504) 
    -- (2.0506790277, 3.998139414) -- cycle(2.1440967617, 4.0333400569) -- 
    (2.1465819929, 4.0526927667).. controls (2.1467645554, 4.054611525) and 
    (2.1479337492, 4.0585243421) .. (2.1500263388, 4.0644541044).. controls 
    (2.1509854533, 4.0674227823) and (2.1527306715, 4.069643536) .. (2.1555179775,
     4.0713506806).. controls (2.1582600663, 4.0728691244) and (2.1612738556, 
    4.073480841) .. (2.1644087448, 4.0731786869).. controls (2.1692431056, 
    4.0724762181) and (2.1726876631, 4.0706479473) .. (2.1748253906, 
    4.0672866542).. controls (2.1772198698, 4.063978066) and (2.1776580992, 
    4.0581844583) .. (2.1766989052, 4.0497317087) -- (2.174213674, 4.0297293146) 
    -- (2.1440968146, 4.033340004) -- cycle;
  
  
  
      \end{scope}
      \begin{scope}[shift={(0.8731702173, -3.4895414292)}]
        \path[fill=c005093,even odd rule] (2.1167843533, 4.2033464373) -- 
    (2.1161726367, 4.2471735023).. controls (2.10902677, 4.2475664085) and 
    (2.1033990825, 4.2491375044) .. (2.0996222612, 4.2517964081).. controls 
    (2.093254349, 4.2560189465) and (2.089590981, 4.2624774785) .. (2.0890243494, 
    4.2713153556).. controls (2.0884123681, 4.2778041031) and (2.0895910869, 
    4.2831220165) .. (2.0924689069, 4.2868309192).. controls (2.0953016154, 
    4.29066071) and (2.0987007704, 4.2927152525) .. (2.1027571504, 4.2929721365)..
     controls (2.1065794006, 4.2933272073) and (2.1101597952, 4.2917860094) .. 
    (2.1134305744, 4.2887041956).. controls (2.1167844856, 4.2856524385) and 
    (2.1203122017, 4.2779778021) .. (2.1238397325, 4.2658237433).. controls 
    (2.1298977913, 4.2460403448) and (2.136960711, 4.2320507656) .. (2.1452396294,
     4.2241646477).. controls (2.1532541231, 4.2160594283) and (2.1632779421, 
    4.2124940621) .. (2.1748276131, 4.2133627685).. controls (2.1824872213, 
    4.2139292415) and (2.189557814, 4.2166712773) .. (2.1960464292, 
    4.2216794429).. controls (2.2024973148, 4.2266120171) and (2.2072864583, 
    4.2336220467) .. (2.2104288352, 4.2428603177).. controls (2.2136921267, 
    4.252015536) and (2.2146515852, 4.2642602146) .. (2.2134731046, 
    4.2797227821).. controls (2.2120755754, 4.29876585) and (2.2073696433, 
    4.3128914779) .. (2.199619389, 4.3223034479).. controls (2.1918162444, 
    4.3318363325) and (2.180145844, 4.3367992015) .. (2.1646227398, 4.3376376396) 
    -- (2.1654082877, 4.2941958079).. controls (2.172078169, 4.2934480954) and 
    (2.1772224098, 4.2915369042) .. (2.1805308921, 4.2880092675).. controls 
    (2.1837564275, 4.284647895) and (2.1857203502, 4.2795489773) .. (2.1862340917,
     4.2730980123).. controls (2.1866344063, 4.2677876396) and (2.1858112875, 
    4.2636861217) .. (2.1836734542, 4.2608156571).. controls (2.1817094521, 
    4.2579829485) and (2.1789218285, 4.2564117733) .. (2.1757417221, 
    4.2561927512).. controls (2.1734755658, 4.2560181262) and (2.1712020542, 
    4.2568875471) .. (2.1691172433, 4.2589876244).. controls (2.16697941, 
    4.2608987098) and (2.1647962006, 4.2659975481) .. (2.1624923677, 
    4.2739743917).. controls (2.1565701725, 4.293667144) and (2.1510787454, 
    4.3077849931) .. (2.1461083094, 4.3160639379).. controls (2.1410094975, 
    4.3244787992) and (2.135042085, 4.3303556977) .. (2.12815297, 4.333890875).. 
    controls (2.1213167452, 4.3372448127) and (2.1138157813, 4.3385969658) .. 
    (2.105801314, 4.33803044).. controls (2.0962532688, 4.3373356442) and 
    (2.0875439792, 4.3341024358) .. (2.0800052856, 4.3282255637).. controls 
    (2.0723381898, 4.3223033421) and (2.0668540388, 4.3144928421) .. 
    (2.0634094813, 4.3046577769).. controls (2.0598742775, 4.2948454129) and 
    (2.0587034963, 4.2829104821) .. (2.0598742775, 4.2683921067).. controls 
    (2.0618382796, 4.2430718785) and (2.0679341738, 4.226000274) .. (2.0783962485,
     4.2169734579).. controls (2.0888960529, 4.2080524223) and (2.1016619342, 
    4.2034748131) .. (2.116784565, 4.2033464373) -- cycle;
  
  
  
      \end{scope}
      \begin{scope}[shift={(0.8731702173, -3.5462329062)}]
        \path[fill=c005093,even odd rule] (2.1922961456, 4.4696630296) -- 
    (2.1819247171, 4.5144569085) -- (2.0367335281, 4.4808199002) -- (2.047104745, 
    4.4358976454) -- cycle;
  
  
  
      \end{scope}
      \begin{scope}[shift={(0.8731702173, -3.4812086417)}]
        \path[fill=c005093,even odd rule] (2.1655796319, 4.51379844) -- 
    (2.1138135588, 4.6437387281) -- (2.0796552419, 4.6302249492) -- (2.096953515, 
    4.5866471746) -- (1.9927718167, 4.5450786458) -- (2.0097602363, 4.5022939863) 
    -- (2.1139950365, 4.5438096248) -- (2.1313837442, 4.5001562852) -- cycle;
  
  
  
      \end{scope}
      \begin{scope}[shift={(0.8731702173, -3.4566746229)}]
        \path[fill=c005093,even odd rule] (2.0892866308, 4.6693544958) -- 
    (2.0616096796, 4.7123128544) -- (2.003052599, 4.7103036085) -- (2.0290678346, 
    4.7628175529) -- (2.0015267204, 4.8055719442) -- (1.9600489173, 4.7099259158) 
    -- (1.9076256192, 4.6761982877) -- (1.9327720425, 4.6374549798) -- 
    (1.9851424769, 4.6711373112) -- (2.0892866308, 4.6693545488) -- cycle;
  
  
  
      \end{scope}
      \begin{scope}[shift={(0.8731702173, -3.4655733542)}]
        \path[fill=c005093,even odd rule] (0.9512183848, 5.0405175521).. controls 
    (0.9437249879, 5.091770069) and (0.9362317233, 5.1430149129) .. (0.9287383265,
     5.1942068931).. controls (0.9130037673, 5.1918954931) and (0.8972315579, 
    5.1893725323) .. (0.8815498625, 5.1866229823).. controls (0.887426761, 
    5.1532050754) and (0.8933110944, 5.1197869833) .. (0.8991879929, 
    5.0863161862).. controls (0.8911356373, 5.0903271637) and (0.8834610008, 
    5.0932050896) .. (0.8760129271, 5.0952144148).. controls (0.8686857679, 
    5.0970876648) and (0.8596287629, 5.098515674) .. (0.8489930156, 5.099308789)..
     controls (0.851342251, 5.0879781135) and (0.8536612446, 5.0766547142) .. 
    (0.8560105065, 5.0653240652).. controls (0.8715636937, 5.063707461) and 
    (0.8839822035, 5.0603461679) .. (0.8933110944, 5.0553834048).. controls 
    (0.9025040687, 5.050549044) and (0.9101711381, 5.043659929) .. (0.9162746258, 
    5.0349881573).. controls (0.9279074554, 5.0369974031) and (0.9395855287, 
    5.0388255417) .. (0.9512183848, 5.0405175521) -- cycle;
  
  
  
      \end{scope}
      \begin{scope}[shift={(0.8731702173, -3.4649153354)}]
        \path[fill=c005093,even odd rule] (1.0089134, 5.1688024746).. controls 
    (1.0242928883, 5.1685458288) and (1.0396723767, 5.1680544975) .. 
    (1.0549688123, 5.1674049454).. controls (1.0558374394, 5.1734252481) and 
    (1.0575371227, 5.177693189) .. (1.0602867521, 5.1803446315).. controls 
    (1.0628550625, 5.1828751065) and (1.0663373231, 5.1844464404) .. 
    (1.0705674285, 5.1847032185).. controls (1.0779323173, 5.1851413685) and 
    (1.0839451587, 5.1819159125) .. (1.088386825, 5.1752987098).. controls 
    (1.0916575777, 5.1703283002) and (1.0944902863, 5.1597378496) .. 
    (1.0966279873, 5.1434444375).. controls (1.0910079727, 5.1482787983) and 
    (1.0852520152, 5.151934705) .. (1.07932982, 5.1539440302).. controls 
    (1.0734907833, 5.1560289469) and (1.0668207962, 5.156905141) .. (1.0593273994,
     5.1563840177).. controls (1.0445974896, 5.1553793948) and (1.0325266158, 
    5.1494949027) .. (1.0235073669, 5.138866696).. controls (1.0142690694, 
    5.1283216744) and (1.010696136, 5.1155557931) .. (1.0121766915, 
    5.1007353429).. controls (1.0132265581, 5.0906737944) and (1.0166258983, 
    5.0817828154) .. (1.022283669, 5.073896724).. controls (1.0279112506, 
    5.0660483621) and (1.0351855196, 5.0603453742) .. (1.0439856671, 
    5.0569462456).. controls (1.0527857881, 5.0533807735) and (1.0635046146, 
    5.0520210269) .. (1.0763687092, 5.0528066806).. controls (1.0915669315, 
    5.053720816) and (1.10381161, 5.0568176581) .. (1.1128837756, 5.0624376727).. 
    controls (1.1220238333, 5.0680576873) and (1.1293509925, 5.0765104104) .. 
    (1.134751985, 5.0881507277).. controls (1.1401529775, 5.099647535) and 
    (1.1428570985, 5.1148535096) .. (1.1425927533, 5.1335490738).. controls 
    (1.1421622762, 5.1610901879) and (1.1355375065, 5.1812585287) .. 
    (1.1224618775, 5.1937299288).. controls (1.1093937892, 5.2061936294) and 
    (1.0911741046, 5.2118968025) .. (1.0683012458, 5.2105447552).. controls 
    (1.0549234892, 5.2097213719) and (1.0443784675, 5.2075382154) .. 
    (1.0368850971, 5.2041843042).. controls (1.0292633244, 5.2008229052) and 
    (1.0231598631, 5.1960339733) .. (1.018582254, 5.1902023979).. controls 
    (1.014057535, 5.1844011173) and (1.0107414592, 5.1772552242) .. (1.0089134, 
    5.1688024746) -- cycle(1.0966279873, 5.106929424).. controls (1.0970132206, 
    5.0986957229) and (1.0951021352, 5.0920258417) .. (1.0911364544, 
    5.0871914544).. controls (1.0870800744, 5.0821758275) and (1.0820264798, 
    5.0796527608) .. (1.0759229921, 5.07929769).. controls (1.0701292785, 
    5.0789502921) and (1.0652948119, 5.0807404629) .. (1.0612006758, 
    5.0846607148).. controls (1.0570612167, 5.08870966) and (1.0547043877, 
    5.094767904) .. (1.0540547827, 5.1032658708).. controls (1.0534428015, 
    5.1118092137) and (1.0549236744, 5.1185168775) .. (1.0586702273, 
    5.123306021).. controls (1.06237913, 5.1281026523) and (1.0673872956, 
    5.1306711215) .. (1.0735361065, 5.1310714096).. controls (1.0797150533, 
    5.1314188075) and (1.0851688567, 5.1295000492) .. (1.0895198502, 
    5.1253605635).. controls (1.0938406019, 5.121085135) and (1.0962729165, 
    5.1149893467) .. (1.0966279873, 5.106929424) -- cycle;
  
  
  
      \end{scope}
      \begin{scope}[shift={(0.8731702173, -3.4656937396)}]
        \path[fill=c005093,even odd rule] (1.1839502504, 5.1345466588).. controls 
    (1.1836936046, 5.105479666) and (1.189003845, 5.085227955) .. (1.1993298446, 
    5.0734969915).. controls (1.2097918135, 5.061516764) and (1.2255188056, 
    5.0554207375) .. (1.2469564585, 5.0542877123).. controls (1.2571918119, 
    5.0537135665) and (1.2657352871, 5.0543326915) .. (1.2726242963, 
    5.0565009519).. controls (1.2794150921, 5.0585935415) and (1.285080324, 
    5.0613884146) .. (1.28961261, 5.06518037).. controls (1.2940618433, 
    5.0687458421) and (1.2978539046, 5.0727568196) .. (1.3006411842, 
    5.0770322481).. controls (1.3034284902, 5.0813454062) and (1.3059062338, 
    5.0863988156) .. (1.3080439348, 5.092147206).. controls (1.3119265365, 
    5.1031756744) and (1.3145401435, 5.1149823352) .. (1.3156731952, 
    5.127272231).. controls (1.318332099, 5.1548056723) and (1.3153633681, 
    5.1754578831) .. (1.3061327171, 5.1891529546).. controls (1.2969397692, 
    5.2028251925) and (1.2794604152, 5.210673766) .. (1.2540117054, 
    5.2120183521).. controls (1.2397275127, 5.2128115729) and (1.2281324127, 
    5.2111947042) .. (1.2192869421, 5.2074028017).. controls (1.2103583394, 
    5.2035202) and (1.2030765033, 5.1976809781) .. (1.1973205194, 5.1898855065).. 
    controls (1.1932717065, 5.1844015406) and (1.1900009538, 5.1766817663) .. 
    (1.1876063688, 5.1669221867).. controls (1.1853402125, 5.1573440848) and 
    (1.1840786527, 5.1464817483) .. (1.1839502504, 5.1345467117) -- 
    cycle(1.2285251338, 5.1336325233).. controls (1.2292654379, 5.1530156337) and 
    (1.2316600229, 5.1662649881) .. (1.2356635392, 5.173199294).. controls 
    (1.2397275127, 5.1802091383) and (1.245256881, 5.1836082669) .. (1.2524027742,
     5.18321544).. controls (1.2569803569, 5.1829508567) and (1.2609838731, 
    5.1812136025) .. (1.2642168698, 5.1780334696).. controls (1.2674347588, 
    5.1746795583) and (1.2697008356, 5.1698377098) .. (1.270962369, 
    5.1630921312).. controls (1.2720122356, 5.1565051969) and (1.272178394, 
    5.146315246) .. (1.2713097669, 5.132582445).. controls (1.269957746, 
    5.112330734) and (1.2673064623, 5.0988623575) .. (1.2633784581, 
    5.0921470473).. controls (1.2592842956, 5.085393981) and (1.2538908702, 
    5.0821763302) .. (1.2468280563, 5.0825614577).. controls (1.2398105654, 
    5.082954364) and (1.2347496154, 5.0867008375) .. (1.2317960717, 
    5.0938467042).. controls (1.2288727169, 5.1012570219) and (1.2277773154, 
    5.1144609738) .. (1.2285251338, 5.1336325233) -- cycle;
  
  
  
      \end{scope}
      \begin{scope}[shift={(0.8731702173, -3.4586388896)}]
        \path[fill=c005093,even odd rule] (1.502477019, 5.1703672998).. controls 
    (1.456111859, 5.1799530746) and (1.4094369513, 5.1875369854) .. (1.3626034523,
     5.1931192704).. controls (1.3627780773, 5.180481765) and (1.3661764915, 
    5.1680634669) .. (1.3728011819, 5.15607554).. controls (1.379335226, 
    5.1441330421) and (1.3924939342, 5.1291916773) .. (1.4115823252, 
    5.1108964012).. controls (1.423124535, 5.0997772867) and (1.43014192, 
    5.0914152098) .. (1.4328461733, 5.0861425931).. controls (1.4355881563, 
    5.0807416006) and (1.4365021594, 5.0759449694) .. (1.4357165321, 
    5.0717147846).. controls (1.4349309842, 5.0672655513) and (1.4325363992, 
    5.0635189983) .. (1.4284875863, 5.0608223385).. controls (1.4244765029, 
    5.0583447802) and (1.4198084592, 5.0572946225) .. (1.4144981129, 
    5.0581632231).. controls (1.4090517973, 5.0590773585) and (1.4048669092, 
    5.061570051) .. (1.4019057983, 5.0653922748).. controls (1.398944714, 
    5.0694561425) and (1.3975020204, 5.0755973598) .. (1.3975473171, 
    5.084306861).. controls (1.3827721106, 5.0853567277) and (1.3680497679, 
    5.0862256194) .. (1.3532745615, 5.0868376006).. controls (1.353584124, 
    5.0748496738) and (1.3555407177, 5.0651357083) .. (1.3592496469, 
    5.0579444127).. controls (1.3627772571, 5.05076079) and (1.3684427006, 
    5.0449140804) .. (1.3760644733, 5.0400797196).. controls (1.3837315427, 
    5.0354189519) and (1.3946239623, 5.0321104431) .. (1.4088779925, 5.02975372)..
     controls (1.423690955, 5.0273592408) and (1.4353692135, 5.0269662023) .. 
    (1.444305304, 5.0286657533).. controls (1.45318861, 5.0304183533) and 
    (1.460681901, 5.0343311704) .. (1.4669213317, 5.0402986094).. controls 
    (1.473062549, 5.0462736948) and (1.4770282033, 5.0534573175) .. (1.4788109394,
     5.0620459042).. controls (1.480684454, 5.0711028033) and (1.4796795665, 
    5.0805602023) .. (1.4757667758, 5.089889199).. controls (1.4717934485, 
    5.0993465979) and (1.4631292704, 5.1104505783) .. (1.4497892698, 
    5.1232239208).. controls (1.4418578552, 5.1306342385) and (1.436328381, 
    5.1357253775) .. (1.433714774, 5.1385581125).. controls (1.4310105471, 
    5.1413453921) and (1.4278304142, 5.1450921567) .. (1.424302804, 
    5.1495413635).. controls (1.4480896923, 5.1457039792) and (1.4717935279, 
    5.1413456037) .. (1.4954595281, 5.1364280579).. controls (1.4977709281, 
    5.1477512721) and (1.50012019, 5.1590365185) .. (1.502477019, 5.1703671675) --
     cycle;
  
  
  
      \end{scope}
      \begin{scope}[shift={(0.8731702173, -3.5322052271)}]
        \path[fill=c005093,even odd rule] (1.3181497217, 4.0392482292) -- 
    (1.2864768213, 4.0449511906).. controls (1.2781148238, 4.046212724) and 
    (1.2764528965, 4.0513191029) .. (1.2768079673, 4.0539780331) -- (1.2766730298,
     4.0705283027).. controls (1.2765830715, 4.0755440354) and (1.2776776527, 
    4.0787695973) .. (1.2801174815, 4.0803407725).. controls (1.2818624085, 
    4.0815946329) and (1.2842117233, 4.0813453954) .. (1.2854278277, 4.080809085) 
    -- (1.3121003412, 4.0750606946).. controls (1.3183775279, 4.0727114592) and 
    (1.319722114, 4.0686097296) .. (1.3195105531, 4.0664795692) -- (1.3193359281, 
    4.0574149971).. controls (1.3193359281, 4.0544085896) and (1.3180669865, 
    4.0527090121) .. (1.3169792844, 4.0518780617).. controls (1.3142373015, 
    4.0501405429) and (1.3107021242, 4.0511001867) .. (1.3103092973, 4.0513189971)
     -- (1.2908809696, 4.0550203327).. controls (1.2896120279, 4.0552399369) and 
    (1.2887431362, 4.0563194369) .. (1.2887431362, 4.0578151794) -- (1.2887431362,
     4.0666909183).. controls (1.2887431362, 4.0676050538) and (1.2891360425, 
    4.0683527663) .. (1.2897855946, 4.0689644829).. controls (1.2905258987, 
    4.0694780392) and (1.2912738758, 4.0696595433) .. (1.2921877467, 4.0694783037)
     -- (1.3050367071, 4.0662225794).. controls (1.3064720717, 4.0658222648) and 
    (1.30735578, 4.0642585773) .. (1.3069554654, 4.0627780219).. controls 
    (1.3065625592, 4.0612900052) and (1.3050367071, 4.060330626) .. (1.3036392837,
     4.0607309406).. controls (1.3036392837, 4.0607309406) and (1.2979437571, 
    4.0621663052) .. (1.2943179335, 4.063080176).. controls (1.2943179335, 
    4.0616448115) and (1.2943179335, 4.0612521698) .. (1.2943179335, 
    4.0599452869).. controls (1.2978984604, 4.0593333056) and (1.3117972875, 
    4.0567198308) .. (1.3117972875, 4.0567198308).. controls (1.3126659146, 
    4.0564102683) and (1.3134440542, 4.0564102683) .. (1.3136708021, 
    4.0565848933).. controls (1.3136708021, 4.0567198308) and (1.3137528229, 
    4.0574611933) .. (1.3137528229, 4.0574611933) -- (1.3140549771, 4.066827946)..
     controls (1.3140549771, 4.066827946) and (1.3138433104, 4.0683537981) .. 
    (1.3106104196, 4.0695246058) -- (1.2836433983, 4.0755449085).. controls 
    (1.2831674129, 4.075727471) and (1.2822458692, 4.0748424398) .. (1.2822458692,
     4.0705291758) -- (1.2822458692, 4.0536692379).. controls (1.2822458692, 
    4.0527098588) and (1.2833412442, 4.0511009275) .. (1.287344681, 4.0504513754) 
    -- (1.319508701, 4.044649989).. controls (1.3199016073, 4.0445679681) and 
    (1.3212083844, 4.0442570827) .. (1.3220770115, 4.0448695931).. controls 
    (1.3229456385, 4.045436066) and (1.3234668677, 4.0474000681) .. (1.3234668677,
     4.0499230819) -- (1.3234668677, 4.0728035342).. controls (1.3234668677, 
    4.0728934925) and (1.3234668677, 4.0728934925) .. (1.3234668677, 
    4.0728934925).. controls (1.3234668677, 4.0738451987) and (1.3230816344, 
    4.0762398631) .. (1.3181490338, 4.0777657152).. controls (1.3181490338, 
    4.0777657152) and (1.2850786306, 4.0852591121) .. (1.284421485, 4.085432811)..
     controls (1.2833411913, 4.085432811) and (1.2809014154, 4.0855227694) .. 
    (1.2789374133, 4.0874346485).. controls (1.2777213883, 4.0885752673) and 
    (1.2762784038, 4.090576999) .. (1.2764521821, 4.0944596006) -- (1.2764521821, 
    4.1221215235).. controls (1.2763701613, 4.122256461) and (1.2761955363, 
    4.1258228592) .. (1.2789374133, 4.1282250113).. controls (1.2809014154, 
    4.1299699383) and (1.2836432925, 4.1305742467) .. (1.2871331731, 4.1297508633)
     -- (1.3206796675, 4.1221215235).. controls (1.3231648988, 4.1220315652) and 
    (1.3287469456, 4.1199461194) .. (1.329132179, 4.1128833583) -- (1.329132179, 
    4.0948449663).. controls (1.329132179, 4.0944144892) and (1.3289125748, 
    4.0916269979) .. (1.3263903019, 4.0895799167).. controls (1.3242072248, 
    4.0877897458) and (1.3211177117, 4.0873060875) .. (1.3170688723, 4.0879633125)
     -- (1.2935388415, 4.0944142775).. controls (1.2912726852, 4.0944592567) and 
    (1.2864761333, 4.0965822733) .. (1.2864761333, 4.1025195233) -- (1.2864761333,
     4.1110552992).. controls (1.2865660917, 4.1117048512) and (1.2869518542, 
    4.1144923954) .. (1.2895277581, 4.116237296).. controls (1.2914388435, 
    4.1176348252) and (1.2938409163, 4.1178160648) .. (1.296756704, 4.1173249981) 
    -- (1.3145383444, 4.1127095271).. controls (1.3157998777, 4.1124978604) and 
    (1.3166759131, 4.1111458396) .. (1.3165854256, 4.1098768185) -- (1.3157998777,
     4.101333555).. controls (1.3157998777, 4.1005554154) and (1.315316484, 
    4.0997698675) .. (1.3147121756, 4.0991581508).. controls (1.3139266277, 
    4.0987652446) and (1.3130580006, 4.0986821654) .. (1.3122724527, 4.0987652446)
     -- (1.3042503125, 4.1011218883).. controls (1.3028604562, 4.1015147946) and 
    (1.3018558333, 4.1030858904) .. (1.3023770625, 4.1045664458).. controls 
    (1.3027244604, 4.1060393812) and (1.3043410646, 4.1069080083) .. 
    (1.3057837581, 4.1065304479).. controls (1.3057837581, 4.1065304479) and 
    (1.3082237456, 4.1057372271) .. (1.310489849, 4.1051707542).. controls 
    (1.3107015156, 4.1065304479) and (1.3107015156, 4.1068326021) .. 
    (1.3107015156, 4.1080034627).. controls (1.3073098483, 4.1088720898) and 
    (1.2954957262, 4.1118408471) .. (1.2954957262, 4.1118408471).. controls 
    (1.292836796, 4.1124528283) and (1.2921871117, 4.1113196179) .. (1.2920511688,
     4.11053407) -- (1.2920511688, 4.1025194969).. controls (1.2920511688, 
    4.1005554948) and (1.2936222646, 4.0999890219) .. (1.2944985646, 4.0997699469)
     -- (1.3183760198, 4.0935002479).. controls (1.3211180027, 4.0930168542) and 
    (1.3222965627, 4.0935003273) .. (1.3228555215, 4.0938854813).. controls 
    (1.3234672381, 4.0944141188) and (1.3234672381, 4.0948448604) .. 
    (1.3234672381, 4.0949279396) -- (1.3234672381, 4.1128454435).. controls 
    (1.3234672381, 4.11602555) and (1.3210272506, 4.1164109156) .. (1.3199018454, 
    4.1165468585) -- (1.2860836769, 4.1244782731).. controls (1.284421829, 
    4.1246899398) and (1.2832509419, 4.124560294) .. (1.2826391194, 
    4.1241308752).. controls (1.2819895673, 4.1235191585) and (1.282034811, 
    4.1222576252) .. (1.282034811, 4.1222576252) -- (1.282034811, 4.0942859281).. 
    controls (1.282034811, 4.0934096281) and (1.2820349169, 4.0920954427) .. 
    (1.2826391194, 4.0915363781).. controls (1.2833415881, 4.0909622323) and 
    (1.2844218819, 4.0911434719) .. (1.2844218819, 4.0911434719).. controls 
    (1.2846938735, 4.0911434719) and (1.2849957631, 4.0911433925) .. 
    (1.2854265048, 4.0910085344) -- (1.319720791, 4.0832583858).. controls 
    (1.3304396175, 4.0795494831) and (1.3291326552, 4.0725396652) .. 
    (1.3291326552, 4.0721921615) -- (1.3291326552, 4.0499235581).. controls 
    (1.3291326552, 4.045345949) and (1.3278713865, 4.0420827633) .. (1.3254314519,
     4.0403378892).. controls (1.3221682663, 4.0379886538) and (1.318549904, 
    4.0391671079) .. (1.31814951, 4.0392501871) -- cycle;
  
  
  
      \end{scope}
      \begin{scope}[shift={(0.8731702173, -3.5324161)}]
        \path[fill=c005093,even odd rule] (1.3797749635, 4.0252576446) -- 
    (1.3480868231, 4.0309229558).. controls (1.3397248256, 4.0322146517) and 
    (1.3379345225, 4.0374116769) .. (1.3383273494, 4.0399799608) -- (1.3381606619,
     4.0564093423).. controls (1.3381606619, 4.061462831) and (1.3392409556, 
    4.0647714192) .. (1.341726134, 4.0663426208).. controls (1.343471061, 
    4.0676041542) and (1.3456993554, 4.0674379958) .. (1.3470062119, 4.0668109333)
     -- (1.3737616988, 4.0609416283).. controls (1.3800312919, 4.0586680637) and 
    (1.3813003923, 4.0546190127) .. (1.3812550692, 4.0525342019) -- (1.3808999983,
     4.0433790894).. controls (1.3808999983, 4.0403349258) and (1.3795929567, 
    4.0387561835) .. (1.3784676838, 4.0379705298).. controls (1.3758011865, 
    4.0361425235) and (1.3724020579, 4.0371019027) .. (1.372009231, 4.0373209777) 
    -- (1.3525807975, 4.040893911).. controls (1.3513192642, 4.0411135152) and 
    (1.3503601496, 4.0422911756) .. (1.3503601496, 4.0435452477) -- (1.3503601496,
     4.0527003602).. controls (1.3503601496, 4.0535766602) and (1.3507530558, 
    4.054490531) .. (1.3513195287, 4.054876029).. controls (1.3520598329, 
    4.0554501748) and (1.3529737038, 4.0556163331) .. (1.3536684996, 4.0554501748)
     -- (1.3666083179, 4.0520887758).. controls (1.3681794137, 4.0517035425) and 
    (1.3689649617, 4.0502607696) .. (1.36857232, 4.0487877548).. controls 
    (1.3681794137, 4.0472997381) and (1.3666083179, 4.046340359) .. (1.3652637319,
     4.0467330006).. controls (1.3652637319, 4.0467330006) and (1.359553071, 
    4.0481681006) .. (1.3558517354, 4.0490367277).. controls (1.3558517354, 
    4.047609036) and (1.3558517354, 4.0472087215) .. (1.3558517354, 
    4.0460379931).. controls (1.3594171281, 4.0453355244) and (1.3734973271, 
    4.0426764883) .. (1.3734973271, 4.0426764883).. controls (1.3742754667, 
    4.0424198425) and (1.3750684229, 4.0424198425) .. (1.3752800896, 
    4.0425018633).. controls (1.3752800896, 4.0426764883) and (1.3754547146, 
    4.0434612425) .. (1.3754547146, 4.0434612425) -- (1.3756372771, 
    4.0527825663).. controls (1.3756372771, 4.0527825663) and (1.3754547146, 
    4.0544896579) .. (1.3722380427, 4.0556151954) -- (1.3452182371, 4.061461905)..
     controls (1.3448253308, 4.0615915508) and (1.3437828725, 4.0607216008) .. 
    (1.3438736246, 4.056529225) -- (1.3438736246, 4.0396314252).. controls 
    (1.3438736246, 4.038672046) and (1.3449086746, 4.0371009502) .. (1.3490480279,
     4.0364059956) -- (1.3811213223, 4.0306121763).. controls (1.3815142285, 
    4.0306121763) and (1.3829114931, 4.0301287825) .. (1.3836896327, 
    4.0307418221).. controls (1.3846942556, 4.0314821263) and (1.385170241, 
    4.0332648887) .. (1.385170241, 4.0359236867) -- (1.385170241, 4.058668196).. 
    controls (1.385170241, 4.0588031335) and (1.385170241, 4.0588957377) .. 
    (1.385170241, 4.0589777585).. controls (1.385170241, 4.0599371377) and 
    (1.3846945202, 4.0621200296) .. (1.379776895, 4.0639027656).. controls 
    (1.379776895, 4.0639027656) and (1.3468273271, 4.0712601931) .. (1.3460492404,
     4.0714414856).. controls (1.3449086217, 4.0713515273) and (1.3425140631, 
    4.0715235065) .. (1.340512305, 4.0733978148).. controls (1.3392433633, 
    4.0744402731) and (1.3379363217, 4.0766233502) .. (1.3379363217, 4.0803699825)
     -- (1.3379363217, 4.1081225517).. controls (1.3379363217, 4.10821251) and 
    (1.33789928, 4.1117483752) .. (1.340512305, 4.1140976371).. controls 
    (1.3423858196, 4.1160616392) and (1.3452183958, 4.1164468725) .. 
    (1.3488819754, 4.1158878079) -- (1.3822998558, 4.1080016106).. controls 
    (1.3847398433, 4.1080016106) and (1.3904503983, 4.1058640419) .. 
    (1.3907526054, 4.0989749004) -- (1.3907526054, 4.0807627829).. controls 
    (1.3906705846, 4.0803698767) and (1.3904504513, 4.0776279731) .. 
    (1.3880029496, 4.075663971).. controls (1.3858275454, 4.0738812085) and 
    (1.3826851156, 4.0733147356) .. (1.3785909796, 4.0740550398) -- (1.3551138125,
     4.080369956).. controls (1.3529759792, 4.0804149352) and (1.3480887546, 
    4.0825077894) .. (1.3480887546, 4.0886035777) -- (1.3480887546, 
    4.0970637621).. controls (1.3481787129, 4.0977133142) and (1.3485721483, 
    4.1003724825) .. (1.3511479465, 4.1021173567).. controls (1.3529759527, 
    4.1035072129) and (1.3554612369, 4.1039075275) .. (1.3583769981, 4.1032955463)
     -- (1.3760225898, 4.0988010163).. controls (1.377420119, 4.0984157829) and 
    (1.378288746, 4.0972373288) .. (1.3782056669, 4.0958777673) -- (1.377420119, 
    4.0872966419).. controls (1.377420119, 4.0864732585) and (1.3768081377, 
    4.0856424669) .. (1.376196421, 4.0852572335).. controls (1.3756373565, 
    4.0846830877) and (1.3746779773, 4.0846830877) .. (1.373885021, 4.084856919) 
    -- (1.3658325596, 4.087039996).. controls (1.3643519512, 4.0875159815) and 
    (1.3634759158, 4.0890039981) .. (1.3638685575, 4.0905675004).. controls 
    (1.3643519512, 4.0919197858) and (1.3658325596, 4.0929244088) .. 
    (1.3673960883, 4.0924862588).. controls (1.3673960883, 4.0924862588) and 
    (1.3697453238, 4.0917459546) .. (1.3721021792, 4.0911718087).. controls 
    (1.3722768042, 4.0924862588) and (1.3722768042, 4.0928336567) .. 
    (1.3724043333, 4.0938837085).. controls (1.3689673429, 4.0948809231) and 
    (1.357198544, 4.0977964727) .. (1.357198544, 4.0977964727).. controls 
    (1.354456561, 4.0984158623) and (1.3536708015, 4.0972374081) .. (1.3535877488,
     4.0965426123) -- (1.3536697696, 4.0886036306).. controls (1.3536697696, 
    4.0864734702) and (1.3552408654, 4.0859521881) .. (1.3561097571, 4.0858615419)
     -- (1.3798589158, 4.0793729267).. controls (1.3826840573, 4.0789800204) and 
    (1.3839530519, 4.0795925308) .. (1.3844818217, 4.079977235).. controls 
    (1.3850482946, 4.0803701413) and (1.3851692092, 4.0809366142) .. 
    (1.3851692092, 4.0809366142) -- (1.3851692092, 4.0988012279).. controls 
    (1.3850475008, 4.1019436048) and (1.382517846, 4.1025026429) .. (1.3815131967,
     4.1025026429) -- (1.3475731875, 4.110396275).. controls (1.3460925792, 
    4.1106529208) and (1.3448235317, 4.1104862333) .. (1.3442570058, 
    4.1099959604).. controls (1.34356221, 4.1093916521) and (1.34356221, 
    4.1082131979) .. (1.34356221, 4.1082131979) -- (1.34356221, 4.0802869298).. 
    controls (1.34356221, 4.0793727944) and (1.3436442308, 4.0781038527) .. 
    (1.3442570058, 4.0774995179).. controls (1.3448234788, 4.0768499658) and 
    (1.3460471767, 4.0770161242) .. (1.3460471767, 4.0770161242).. controls 
    (1.3463038225, 4.0770161242) and (1.3466136496, 4.0770160183) .. 
    (1.3470065558, 4.0769261658) -- (1.3812554131, 4.0692666635).. controls 
    (1.3920950483, 4.0656484071) and (1.3907504623, 4.0584949469) .. 
    (1.3906673831, 4.0581927663) -- (1.390749404, 4.0359242423).. controls 
    (1.390749404, 4.0313013365) and (1.3894878706, 4.028000289) .. (1.3869195867, 
    4.0262554148).. controls (1.3837772098, 4.0240723377) and (1.3801136302, 
    4.025122469) .. (1.3797736935, 4.0252584648) -- cycle;
  
  
  
      \end{scope}
      \begin{scope}[shift={(0.8731702173, -3.5268413292)}]
        \path[fill=c005093,even odd rule] (1.4409891487, 4.0000727508) -- 
    (1.4092707135, 4.0102703481).. controls (1.399156169, 4.0131861358) and 
    (1.3998587435, 4.0200601167) .. (1.3998587435, 4.0202867058) -- (1.3998587435,
     4.0367161667).. controls (1.399676181, 4.0417243587) and (1.4007273706, 
    4.0448667092) .. (1.403295734, 4.0463396975).. controls (1.405259736, 
    4.0475181517) and (1.4075258129, 4.0471252454) .. (1.4087873463, 4.0462576767)
     -- (1.435671209, 4.0365433937).. controls (1.4429077219, 4.0328874077) and 
    (1.4428169963, 4.0280907765) .. (1.4427339171, 4.0275242506) -- (1.4425143129,
     4.0182784125).. controls (1.4425143129, 4.015445704) and (1.4411244567, 
    4.0138367462) .. (1.4399838379, 4.0130965215).. controls (1.4373325013, 
    4.0113968381) and (1.4338879438, 4.0129218965) .. (1.4333138508, 4.0130965215)
     -- (1.4137949033, 4.01968335).. controls (1.4126241221, 4.0200609104) and 
    (1.4119213887, 4.0210657979) .. (1.4119213887, 4.022334819) -- (1.4119213887, 
    4.0312255862).. controls (1.4119213887, 4.0322225362) and (1.412314295, 
    4.0330988362) .. (1.4131072513, 4.0336200654).. controls (1.4138853908, 
    4.03417913) and (1.41484477, 4.03417913) .. (1.4157133971, 4.0338846487) -- 
    (1.4283961198, 4.0288688896).. controls (1.4299672156, 4.0283024167) and 
    (1.430571524, 4.0266405687) .. (1.4300502948, 4.0251674746).. controls 
    (1.429566901, 4.0237776183) and (1.4277844031, 4.0229843975) .. (1.4264321442,
     4.0235963787).. controls (1.4264321442, 4.0235963787) and (1.4211595275, 
    4.0257265392) .. (1.4174128952, 4.027131556).. controls (1.4174128952, 
    4.0257266185) and (1.4174128952, 4.025167554) .. (1.4174128952, 
    4.0242536302).. controls (1.4205552721, 4.0232490073) and (1.4353682346, 
    4.0182785448) .. (1.4353682346, 4.0182785448).. controls (1.4356703887, 
    4.0182785448) and (1.4371507325, 4.0178933115) .. (1.4371507325, 
    4.0178933115).. controls (1.4371507325, 4.0178933115) and (1.4369390658, 
    4.0185883719) .. (1.4369390658, 4.0184976198) -- (1.4371957117, 
    4.0277811081).. controls (1.4371957117, 4.0277811081) and (1.4370210867, 
    4.0296998665) .. (1.4334037562, 4.0314372) -- (1.4064368144, 4.0412494581).. 
    controls (1.4061272519, 4.0413394165) and (1.4059155852, 4.0414161456) .. 
    (1.4059155852, 4.0414161456).. controls (1.4059155852, 4.0412494581) and 
    (1.4054321915, 4.0403813602) .. (1.4054321915, 4.0375032756).. controls 
    (1.4054321915, 4.0373207131) and (1.4054321915, 4.0198501169) .. 
    (1.4054321915, 4.0198501169).. controls (1.4054321915, 4.0196834294) and 
    (1.4056517956, 4.0171080546) .. (1.4109615598, 4.0155368529) -- (1.4431255004,
     4.0052108533).. controls (1.4439110483, 4.00482562) and (1.4449535067, 
    4.0045613012) .. (1.4455654879, 4.00494627).. controls (1.4460414733, 
    4.0053391762) and (1.4467439421, 4.0067367054) .. (1.4467439421, 4.0095238792)
     -- (1.4467439421, 4.0323060123).. controls (1.4467439421, 4.0323959706) and 
    (1.4467439421, 4.032533554) .. (1.4467439421, 4.0326232477).. controls 
    (1.4467439421, 4.0336201977) and (1.4462153046, 4.0361432908) .. 
    (1.4410786044, 4.0387191419).. controls (1.4410786044, 4.0387191419) and 
    (1.4072150069, 4.0510467937) .. (1.4070413079, 4.0510467937).. controls 
    (1.4059157704, 4.0512214187) and (1.4038612015, 4.051658775) .. (1.401942496, 
    4.0535775333).. controls (1.4008624669, 4.0547483146) and (1.3995480169, 
    4.0567576662) .. (1.3995480169, 4.0602022237).. controls (1.3995480169, 
    4.0602921821) and (1.3995480169, 4.0604588696) .. (1.3995480169, 4.0606782092)
     -- (1.3995480169, 4.0883099696).. controls (1.3995480169, 4.0885216362) and 
    (1.3995480169, 4.0885216362) .. (1.3995480169, 4.0886952029).. controls 
    (1.3995480169, 4.08956383) and (1.399857844, 4.0925324815) .. (1.4022070529, 
    4.0944058637).. controls (1.4042541342, 4.0959695512) and (1.4072151127, 
    4.0961507908) .. (1.4108333692, 4.0950101721) -- (1.4441681969, 
    4.0824255569).. controls (1.4466534281, 4.082070486) and (1.4520545, 
    4.07928318) .. (1.4523188452, 4.0725301402) -- (1.4523188452, 4.0543633458).. 
    controls (1.4522288869, 4.0539704396) and (1.4521442202, 4.0511378104) .. 
    (1.4494861367, 4.0493929098).. controls (1.4472199804, 4.0477387348) and 
    (1.4439944979, 4.0476101473) .. (1.4398550387, 4.0490455119) -- (1.4163249021,
     4.0585858842).. controls (1.4142778208, 4.0590692779) and (1.4097002381, 
    4.0617283404) .. (1.4097002381, 4.0674390012) -- (1.4097002381, 
    4.0759823706).. controls (1.4097452173, 4.0767605102) and (1.4101383881, 
    4.0792833917) .. (1.4127067515, 4.08089989).. controls (1.4146707535, 
    4.0820706712) and (1.4173295779, 4.0822066671) .. (1.4203813085, 4.0810665775)
     -- (1.4380195183, 4.074109544).. controls (1.4392053808, 4.0736184773) and 
    (1.4398549329, 4.0724021877) .. (1.4397718537, 4.0712768354) -- (1.4389786329,
     4.0628166246).. controls (1.4389786329, 4.0619479975) and (1.4384198329, 
    4.0610716975) .. (1.4376340204, 4.0606787912).. controls (1.4368029642, 
    4.0602028058) and (1.4359343371, 4.0600744829) .. (1.43506571, 4.0604591871) 
    -- (1.4270436492, 4.0639037446).. controls (1.4256537929, 4.0645609696) and 
    (1.42508732, 4.0661699008) .. (1.4255630408, 4.0675295681).. controls 
    (1.4262655096, 4.0688665077) and (1.4277836887, 4.069697564) .. (1.4293096731,
     4.0690932556).. controls (1.4293096731, 4.0690932556) and (1.4315758294, 
    4.0680886327) .. (1.433796636, 4.0671747619).. controls (1.4338865944, 
    4.068481539) and (1.4338865944, 4.068700614) .. (1.4341061985, 4.0696978285)..
     controls (1.4309713887, 4.0708762827) and (1.4183717452, 4.0758465071) .. 
    (1.4183717452, 4.0758465071).. controls (1.4166268181, 4.0764508154) and 
    (1.4158486785, 4.0761486613) .. (1.415539116, 4.0761486613) -- (1.4151917181, 
    4.0753631133) -- (1.4152816765, 4.06743916).. controls (1.4152816765, 
    4.064908685) and (1.4176309119, 4.0641155171) .. (1.418114385, 4.0639039562) 
    -- (1.441727495, 4.0542804254).. controls (1.4447792256, 4.0533210462) and 
    (1.446040759, 4.0538348671) .. (1.4462144579, 4.0539708629).. controls 
    (1.4466526079, 4.0542804254) and (1.4467430954, 4.0545825796) .. 
    (1.4467430954, 4.0545825796) -- (1.4467430954, 4.0724020554).. controls 
    (1.4465605329, 4.0761487142) and (1.4432985379, 4.0769343415) .. 
    (1.4427320121, 4.0771533635) -- (1.4088759817, 4.0897001433).. controls 
    (1.4070858108, 4.0903044517) and (1.4059148708, 4.0903044517) .. 
    (1.4056505521, 4.0899647267).. controls (1.4052576458, 4.0897001433) and 
    (1.4051671583, 4.0889225329) .. (1.4051671583, 4.0885219537) -- (1.4051671583,
     4.060678659).. controls (1.4050375125, 4.0572416685) and (1.4069120854, 
    4.0566222525) .. (1.4079998933, 4.0565467669).. controls (1.4082194975, 
    4.0565467669) and (1.4085663662, 4.056546661) .. (1.4088761933, 4.0564118294) 
    -- (1.4432081298, 4.0438197265).. controls (1.4535341294, 4.0388946135) and 
    (1.4523179985, 4.0321415473) .. (1.4522273787, 4.0318393667) -- (1.4523173371,
     4.009525546).. controls (1.4523173371, 4.0049479369) and (1.4509650517, 
    4.0018131535) .. (1.4484798733, 4.0002419519).. controls (1.4450806125, 
    3.9982025435) and (1.4414701348, 3.9998567185) .. (1.4409866881, 4.0000752644)
     -- cycle;
  
  
  
      \end{scope}
      \begin{scope}[shift={(0.8731702173, -3.5323182042)}]
        \path[fill=c005093,even odd rule] (1.3297068804, 4.2375798429) -- 
    (1.2981546829, 4.2432905037).. controls (1.2897851183, 4.2446351163) and 
    (1.2881309698, 4.2498245215) .. (1.2883499654, 4.2523928054) -- (1.2883499654,
     4.2689959917).. controls (1.2881303613, 4.2738682144) and (1.28930908, 
    4.2772295869) .. (1.2916660413, 4.2788007885).. controls (1.2936223704, 
    4.2800623219) and (1.2957526367, 4.2798884906) .. (1.2970594667, 4.2791936948)
     -- (1.3238602767, 4.2733999812).. controls (1.3301374369, 4.2711264167) and 
    (1.3313988644, 4.2669867723) .. (1.3313158117, 4.2649472581) -- (1.3310962075,
     4.2558373629).. controls (1.3310962075, 4.2528687115) and (1.3297063513, 
    4.2513430446) .. (1.3285202242, 4.2505573908).. controls (1.3259140783, 
    4.2484650658) and (1.3224697325, 4.2495601763) .. (1.3220769056, 4.2496432554)
     -- (1.3026484985, 4.2533899142).. controls (1.3012888048, 4.2536619058) and 
    (1.3004201777, 4.2548779308) .. (1.3004201777, 4.2561394642) -- (1.3004201777,
     4.265067961).. controls (1.3004201777, 4.2660273402) and (1.300684761, 
    4.266820561) .. (1.301462636, 4.2674248694).. controls (1.3021574319, 
    4.267825184) and (1.3030337319, 4.2681196652) .. (1.3038647881, 4.267825184) 
    -- (1.3166760983, 4.2646828071).. controls (1.3181490337, 4.2642899008) and 
    (1.3190253337, 4.2626735613) .. (1.3185493483, 4.2613213288).. controls 
    (1.3183747233, 4.2596671537) and (1.3166760983, 4.2587908537) .. 
    (1.3151954371, 4.2591911683).. controls (1.3151954371, 4.2591911683) and 
    (1.3096964429, 4.26061886) .. (1.3059951073, 4.2615025683).. controls 
    (1.3059951073, 4.2601505475) and (1.3059951073, 4.2595840746) .. 
    (1.3059951073, 4.2584509435).. controls (1.309530311, 4.2577863102) and 
    (1.3234668942, 4.2551348677) .. (1.3234668942, 4.2551348677).. controls 
    (1.3243431942, 4.2548782219) and (1.3251287421, 4.2548782219) .. 
    (1.3252118212, 4.254923201).. controls (1.3252118212, 4.2551348677) and 
    (1.3254314254, 4.2558373365) .. (1.3254314254, 4.2558373365) -- (1.3256060504,
     4.2651587925).. controls (1.3256060504, 4.2652487508) and (1.3254314254, 
    4.2669412904) .. (1.3222975417, 4.2679913952) -- (1.2953229533, 
    4.2738682937).. controls (1.2948395596, 4.2741778562) and (1.293842345, 
    4.2731809063) .. (1.293842345, 4.268996071) -- (1.293842345, 4.2520905983).. 
    controls (1.293842345, 4.2511312192) and (1.2950207992, 4.2495601233) .. 
    (1.2991602848, 4.2489105977) -- (1.3313166583, 4.243063994).. controls 
    (1.3315812417, 4.2428443898) and (1.3329708333, 4.2426710877) .. (1.333839725,
     4.2432915356).. controls (1.3346252729, 4.2438503356) and (1.3352372542, 
    4.2458521731) .. (1.3352372542, 4.2483826481) -- (1.3352372542, 
    4.2711271575).. controls (1.3352372542, 4.271262095) and (1.3352372542, 
    4.27130972) .. (1.3352372542, 4.2715274721).. controls (1.3352372542, 
    4.2723056117) and (1.3348443479, 4.2745717415) .. (1.32970778, 4.2762335629)..
     controls (1.32970778, 4.2762335629) and (1.2967581856, 4.2837646894) .. 
    (1.2961463367, 4.2839006323).. controls (1.2950207992, 4.2839006323) and 
    (1.2926640496, 4.2839826531) .. (1.2905716452, 4.2858495531).. controls 
    (1.289400864, 4.2868089323) and (1.287958091, 4.2890826556) .. (1.2881316577, 
    4.2927839648) -- (1.2881316577, 4.3205818571).. controls (1.2880496369, 
    4.3207564821) and (1.2879570327, 4.3242001135) .. (1.2905716452, 
    4.3266928854).. controls (1.292580891, 4.3284304042) and (1.2954135996, 
    4.3288682896) .. (1.2989336427, 4.3284301396) -- (1.3323591167, 
    4.3205817777).. controls (1.3348443479, 4.3204521319) and (1.340509765, 
    4.3184063735) .. (1.3407287871, 4.3113435067) -- (1.3407287871, 
    4.2933125758).. controls (1.3407287871, 4.29273843) and (1.3405938496, 
    4.2899512827) .. (1.3381604767, 4.2880779004).. controls (1.3358869121, 
    4.2862498942) and (1.332706594, 4.2857739087) .. (1.3287484008, 4.2864311337) 
    -- (1.3052108823, 4.2927385887).. controls (1.3030354781, 4.2929211512) and 
    (1.2980196925, 4.2950499887) .. (1.2980196925, 4.3009268608) -- (1.2980196925,
     4.3093796104).. controls (1.29815463, 4.3103011542) and (1.2987676696, 
    4.3129146819) .. (1.3010789902, 4.3146595825).. controls (1.3030353194, 
    4.3158756075) and (1.3056036033, 4.3163592658) .. (1.3084363383, 4.3157472846)
     -- (1.3262179787, 4.3111697812).. controls (1.3274795121, 4.3108297917) and 
    (1.3283933829, 4.30956085) .. (1.3281743079, 4.3082085646) -- (1.3275625912, 
    4.2997560531).. controls (1.3273959037, 4.2989705052) and (1.3270035267, 
    4.298230201) .. (1.3262632225, 4.2976258927).. controls (1.3256059975, 
    4.2972255781) and (1.3248204496, 4.2970591552) .. (1.3239518225, 4.2973160656)
     -- (1.315974979, 4.2994988781).. controls (1.3144491269, 4.2999748635) and 
    (1.3136709873, 4.3015382865) .. (1.3140562206, 4.3029358421).. controls 
    (1.3144491269, 4.3044616942) and (1.3159747144, 4.305254915) .. (1.3174629162,
     4.3048546004).. controls (1.3174629162, 4.3048546004) and (1.3199029037, 
    4.3042955358) .. (1.3222973829, 4.3035930671).. controls (1.3222973829, 
    4.3048546004) and (1.3222974887, 4.3052096713) .. (1.3224720079, 
    4.3064635317).. controls (1.3189897208, 4.3073398317) and (1.3071754929, 
    4.3103838629) .. (1.3071754929, 4.3103838629).. controls (1.3045241298, 
    4.3108296858) and (1.3038668783, 4.3096964754) .. (1.3036402892, 4.3090016796)
     -- (1.3037302475, 4.3009266227).. controls (1.3037302475, 4.2989702935) and 
    (1.3053013433, 4.298403556) .. (1.306170235, 4.2983207415) -- (1.3299192879, 
    4.2919150731).. controls (1.332706594, 4.2915147585) and (1.3340134504, 
    4.2919150731) .. (1.3346252729, 4.2923984669).. controls (1.3351465021, 
    4.2927384565) and (1.3352372542, 4.2933123377) .. (1.3352372542, 4.2935240044)
     -- (1.3352372542, 4.3111697019).. controls (1.3351472958, 4.3144857777) and 
    (1.3327065146, 4.3148785781) .. (1.331490516, 4.3148785781) -- (1.2977628615, 
    4.3228554217).. controls (1.2961462573, 4.3231120675) and (1.2948848298, 
    4.3230300467) .. (1.294318304, 4.3224095988).. controls (1.2936232435, 
    4.3219412863) and (1.2936232435, 4.3207551592) .. (1.2936232435, 4.3205815925)
     -- (1.2936232435, 4.2927384035).. controls (1.2936232435, 4.2918242681) and 
    (1.2938428477, 4.2905553265) .. (1.294318304, 4.2899509917).. controls 
    (1.2950207727, 4.2893845188) and (1.2961084748, 4.2895657583) .. 
    (1.2961084748, 4.2895657583).. controls (1.2963730581, 4.2895657583) and 
    (1.2967580269, 4.2895656525) .. (1.2969847748, 4.2894758) -- (1.3313167112, 
    4.2816350052).. controls (1.3421185904, 4.2780167487) and (1.3407287871, 
    4.2709087175) .. (1.3407287871, 4.2705687808) -- (1.3407287871, 
    4.2483831244).. controls (1.3407287871, 4.2436771394) and (1.3395503329, 
    4.2404139538) .. (1.3370198844, 4.238585921).. controls (1.3338397515, 
    4.236395171) and (1.3302289827, 4.2375812981) .. (1.32970778, 4.2375812981) --
     cycle;
  
  
  
      \end{scope}
      \begin{scope}[shift={(0.8731702173, -3.5321898813)}]
        \path[fill=c005093,even odd rule] (1.3914463694, 4.2241108579) -- 
    (1.3597732573, 4.2297309783).. controls (1.3514036927, 4.23099992) and 
    (1.3497494648, 4.2361819433) .. (1.349915676, 4.2387577944) -- (1.349915676, 
    4.2553155517).. controls (1.3498389469, 4.2603237173) and (1.3509200344, 
    4.2635491733) .. (1.3532769692, 4.2651203485).. controls (1.3552409712, 
    4.2664271256) and (1.3573786988, 4.2660723194) .. (1.3587686079, 4.265596334) 
    -- (1.3853503694, 4.2597118948).. controls (1.3917106087, 4.2574457385) and 
    (1.3929797092, 4.2534874925) .. (1.3928890894, 4.2512138485) -- (1.3925795269,
     4.242202246).. controls (1.3925795269, 4.2392335946) and (1.3913179935, 
    4.2374509644) .. (1.3901472123, 4.2367560363).. controls (1.3876167373, 
    4.23492803) and (1.3840890477, 4.2358797363) .. (1.3836886538, 4.2360911383) 
    -- (1.3642604319, 4.2398001204).. controls (1.3628629027, 4.2400197246) and 
    (1.3619035235, 4.2411976496) .. (1.3619035235, 4.2425421033) -- (1.3619035235,
     4.2514781938).. controls (1.3619035235, 4.2523468208) and (1.3622961652, 
    4.2531323688) .. (1.362998634, 4.25374435).. controls (1.3637389381, 
    4.2541372563) and (1.3645623215, 4.2543563313) .. (1.3654309485, 4.2541372563)
     -- (1.3782043175, 4.2510024465).. controls (1.3797754133, 4.2506095402) and 
    (1.380644305, 4.2490384444) .. (1.3802513988, 4.247557889).. controls 
    (1.3798584925, 4.2461225244) and (1.3783781487, 4.2451179015) .. 
    (1.3768068413, 4.2455031348).. controls (1.3768068413, 4.2455031348) and 
    (1.3712321498, 4.2468628285) .. (1.3675685702, 4.2478674515).. controls 
    (1.3675685702, 4.2463794348) and (1.3675685702, 4.2459034494) .. 
    (1.3675685702, 4.2448081538).. controls (1.3710962069, 4.2441207663) and 
    (1.38504803, 4.2414995392) .. (1.38504803, 4.2414995392).. controls 
    (1.3859621654, 4.241197385) and (1.3867400404, 4.2411976496) .. (1.3868307925,
     4.241279935).. controls (1.3868307925, 4.241279935) and (1.3868307925, 
    4.2412798292) .. (1.3868307925, 4.2414148725).. controls (1.3868307925, 
    4.2414968933) and (1.3870054175, 4.2422004204) .. (1.3870054175, 4.2422004204)
     -- (1.3872250217, 4.2516124169).. controls (1.3872250217, 4.2516124169) and 
    (1.3870054175, 4.253130596) .. (1.3839089458, 4.2543543998) -- (1.3569872477, 
    4.2603219975).. controls (1.3564207748, 4.2604119558) and (1.3554613956, 
    4.259536185) .. (1.355506904, 4.255313726) -- (1.355506904, 4.2384536823).. 
    controls (1.355506904, 4.2374490594) and (1.3566324415, 4.235877699) .. 
    (1.3607340917, 4.2353113054) -- (1.3928904652, 4.2294269721).. controls 
    (1.3931926194, 4.2294269721) and (1.3945901485, 4.2289509867) .. 
    (1.3953756965, 4.2296015971).. controls (1.396372911, 4.2302209867) and 
    (1.3968033881, 4.2321699075) .. (1.3968033881, 4.2347005148) -- (1.3968033881,
     4.2574448919).. controls (1.3968033881, 4.2576565585) and (1.3968033881, 
    4.2576565585) .. (1.3968033881, 4.2577922898).. controls (1.3968033881, 
    4.2586609169) and (1.396372911, 4.2610253129) .. (1.3914477981, 
    4.2625435979).. controls (1.3914477981, 4.2625435979) and (1.3585056915, 
    4.2701277204) .. (1.3577654667, 4.2702183667).. controls (1.3566325208, 
    4.2702183667) and (1.3542377506, 4.2702183667) .. (1.3520622935, 
    4.2721746958).. controls (1.3509216748, 4.2732700708) and (1.3495318185, 
    4.2753095056) .. (1.3497508935, 4.279063705) -- (1.3497508935, 4.3069069204)..
     controls (1.3497508935, 4.3071185871) and (1.3494942477, 4.3105252563) .. 
    (1.3520622935, 4.3129575973).. controls (1.354064131, 4.3147477681) and 
    (1.3569874065, 4.3153597494) .. (1.3605603398, 4.3145286931) -- (1.393970759, 
    4.3068163006).. controls (1.3964181548, 4.3067342798) and (1.4021287892, 
    4.3047239756) .. (1.4023856467, 4.2977065113) -- (1.4023856467, 
    4.2795395317).. controls (1.4023856467, 4.2790635463) and (1.4022110217, 
    4.2763972606) .. (1.3996814198, 4.2743501794).. controls (1.3974226719, 
    4.2725674169) and (1.3943711529, 4.272000944) .. (1.3903600696, 4.2727790835) 
    -- (1.3667832869, 4.2790647104).. controls (1.364562639, 4.2793213562) and 
    (1.3597732573, 4.2814139458) .. (1.3597732573, 4.2872983321) -- (1.3597732573,
     4.2958341079).. controls (1.3597732573, 4.2964913329) and (1.3603323219, 
    4.2991501838) .. (1.3628174208, 4.3008875967).. controls (1.3646531, 
    4.3022927987) and (1.3670929552, 4.3026777675) .. (1.3700540396, 4.3020660508)
     -- (1.3877902775, 4.2974961146).. controls (1.3890518108, 4.2973135521) and 
    (1.3899204379, 4.2959247542) .. (1.3897542796, 4.2945273573) -- (1.3890518108,
     4.2861199308).. controls (1.3890068317, 4.2852513037) and (1.3886136608, 
    4.2844657558) .. (1.3878810296, 4.2839823621).. controls (1.3872238046, 
    4.2835971287) and (1.3863551775, 4.2833780537) .. (1.3855619567, 4.2836349642)
     -- (1.377509601, 4.2859011204).. controls (1.3760215844, 4.286301435) and 
    (1.3751529573, 4.2879029579) .. (1.375636351, 4.2894363242).. controls 
    (1.3760215844, 4.2908261804) and (1.377509601, 4.2916948075) .. (1.3790732356,
     4.2912188221).. controls (1.3790732356, 4.2912188221) and (1.3815132231, 
    4.2906068408) .. (1.3837793265, 4.2898212929).. controls (1.3839539515, 
    4.2912188221) and (1.3839539515, 4.2916114637) .. (1.384088889, 
    4.2926540015).. controls (1.3805083885, 4.2935303015) and (1.3687849921, 
    4.2965820056) .. (1.3687849921, 4.2965820056).. controls (1.3660430092, 
    4.297186314) and (1.3654310808, 4.2960155327) .. (1.3652648431, 4.2953127994) 
    -- (1.3652648431, 4.2872982262).. controls (1.3652648431, 4.285251145) and 
    (1.3670020973, 4.2847299158) .. (1.3679161798, 4.2845108144) -- (1.3915368569,
     4.2782713837).. controls (1.3942788398, 4.2777578275) and (1.3956309929, 
    4.2783613421) .. (1.3961597892, 4.2786716983).. controls (1.3967262621, 
    4.2790646046) and (1.3968019329, 4.2797594004) .. (1.3968019329, 4.2797594004)
     -- (1.3968019329, 4.297586364).. controls (1.3966352454, 4.3007514421) and 
    (1.3942788662, 4.3012424823) .. (1.3931912435, 4.301287779) -- (1.3593729427, 
    4.3091663298).. controls (1.3577640115, 4.3094309131) and (1.3564194254, 
    4.3092483506) .. (1.3558528996, 4.3088718485).. controls (1.3552409183, 
    4.308252459) and (1.3552409183, 4.3071192485) .. (1.3552409183, 4.3071192485) 
    -- (1.3552409183, 4.2790644988).. controls (1.3552409183, 4.2782712779) and 
    (1.3552409183, 4.2768814217) .. (1.3560264662, 4.2762695463).. controls 
    (1.3566307746, 4.2757030733) and (1.3575901537, 4.2758390692) .. 
    (1.3575901537, 4.2758390692).. controls (1.35798306, 4.2758840483) and 
    (1.3583757017, 4.2758390692) .. (1.3585495329, 4.2758390692) -- (1.3928889306,
     4.26807376).. controls (1.4036983769, 4.264417774) and (1.4023839533, 
    4.2572793421) .. (1.4023462237, 4.2569697002) -- (1.4023832654, 
    4.2347012027).. controls (1.4023832654, 4.2300858375) and (1.4011217321, 
    4.226852735) .. (1.3985987712, 4.2250246758).. controls (1.3954563944, 
    4.2227585196) and (1.3918382437, 4.2238991383) .. (1.3914454169, 4.2241105404)
     -- cycle;
  
  
  
      \end{scope}
      \begin{scope}[shift={(0.8731702173, -3.5267058625)}]
        \path[fill=c005093,even odd rule] (1.4527151381, 4.1995921056) -- 
    (1.4209515121, 4.2097443004).. controls (1.410836941, 4.2127507344) and 
    (1.4115395156, 4.219541636) .. (1.4115395156, 4.2198059813) -- (1.4115395156, 
    4.2363186271).. controls (1.4113569531, 4.2412890367) and (1.4125367302, 
    4.2444237406) .. (1.4148480244, 4.24586654).. controls (1.416932941, 
    4.2470902379) and (1.4191612883, 4.2464782567) .. (1.420558791, 4.2457316025) 
    -- (1.4472236052, 4.2361081775).. controls (1.4547170021, 4.2325276771) and 
    (1.4544978742, 4.2275723752) .. (1.4544978742, 4.2270360383) -- (1.45419572, 
    4.2178883871).. controls (1.45419572, 4.2149273027) and (1.4529265138, 
    4.2133636946) .. (1.4517557325, 4.2126611994).. controls (1.4490137496, 
    4.2109236806) and (1.4456975679, 4.2123590452) .. (1.4449573431, 4.2126611994)
     -- (1.425476231, 4.2191498146).. controls (1.4243054498, 4.2195427208) and 
    (1.423602981, 4.2205925875) .. (1.423602981, 4.2218163119) -- (1.423602981, 
    4.2307827765).. controls (1.423602981, 4.2316590765) and (1.4240863748, 
    4.232527439) .. (1.4246980915, 4.2331470931).. controls (1.425476231, 
    4.2336606494) and (1.4265260977, 4.2337514015) .. (1.4273116456, 4.2333587598)
     -- (1.4402134962, 4.2283505942).. controls (1.441512865, 4.2277841213) and 
    (1.4422529046, 4.2260770296) .. (1.44173194, 4.2247399048).. controls 
    (1.4411276317, 4.2233347027) and (1.4394657838, 4.2225570923) .. 
    (1.4380230373, 4.2231158923).. controls (1.4380230373, 4.2231158923) and 
    (1.4328864958, 4.2252082173) .. (1.4291020017, 4.2266887462).. controls 
    (1.4291020017, 4.2252081379) and (1.4291020017, 4.2249059838) .. 
    (1.4291020017, 4.2238560377).. controls (1.4323198642, 4.2227760085) and 
    (1.447049774, 4.2178885988) .. (1.447049774, 4.2178885988).. controls 
    (1.447224399, 4.2177589529) and (1.4488325365, 4.2175033654) .. (1.4488325365,
     4.2175033654) -- (1.4486208698, 4.2180245946) -- (1.4488325365, 
    4.2273912415).. controls (1.4488325365, 4.2273912415) and (1.4487108281, 
    4.2291361685) .. (1.4450858777, 4.2310096037) -- (1.4181187506, 
    4.2408142948).. controls (1.4178165965, 4.2408963156) and (1.4176353569, 
    4.2408963156) .. (1.4176353569, 4.2408963156).. controls (1.4176353569, 
    4.2408142948) and (1.4170233756, 4.2398538573) .. (1.4170233756, 
    4.2370589577).. controls (1.4170233756, 4.236847291) and (1.4170233756, 
    4.2195417419) .. (1.4170233756, 4.219413366).. controls (1.4170233756, 
    4.2190281327) and (1.4173329381, 4.216663816) .. (1.4225979612, 4.2151002079) 
    -- (1.4547166846, 4.2046834562).. controls (1.4554569887, 4.2045088313) and 
    (1.4566354429, 4.2041243917) .. (1.4571640804, 4.2045088313).. controls 
    (1.457723145, 4.2049091458) and (1.4582896179, 4.2062915937) .. (1.4582896179,
     4.2091243023) -- (1.4582896179, 4.2319669985).. controls (1.4582896179, 
    4.2319669985) and (1.4582896179, 4.2319668927) .. (1.4582896179, 
    4.2321866027).. controls (1.4582896179, 4.2331459819) and (1.4580329721, 
    4.2357142129) .. (1.4527148206, 4.238146554).. controls (1.4527148206, 
    4.238146554) and (1.4187757375, 4.2506103604) .. (1.4187304144, 
    4.2507009802).. controls (1.4176350394, 4.2507909385) and (1.4155427144, 
    4.2510938865) .. (1.4135862794, 4.2530954594).. controls (1.4125364127, 
    4.2542284052) and (1.411237044, 4.2562755923) .. (1.411237044, 4.2596369648)..
     controls (1.411237044, 4.2598036523) and (1.411237044, 4.2600674419) .. 
    (1.411237044, 4.2602034377) -- (1.411237044, 4.2880014358).. controls 
    (1.411237044, 4.2880014358) and (1.411237044, 4.288046415) .. (1.411237044, 
    4.2883112629).. controls (1.411237044, 4.2892706421) and (1.4114037315, 
    4.2920125985) .. (1.413797946, 4.2937500115).. controls (1.4158527002, 
    4.2955401823) and (1.4187758433, 4.2957592573) .. (1.4225224756, 4.2945355594)
     -- (1.4558498421, 4.2819886737).. controls (1.4582443212, 4.2815957675) and 
    (1.4637435271, 4.2788992929) .. (1.4639097383, 4.2721839827) -- (1.4639097383,
     4.25397168).. controls (1.4639097383, 4.2535787738) and (1.4637430508, 
    4.2507915471) .. (1.461167676, 4.2490466729).. controls (1.4589015198, 
    4.2473924979) and (1.4556760373, 4.2472565021) .. (1.4515441452, 4.2484725271)
     -- (1.4280065473, 4.2581563829).. controls (1.425959466, 4.2586776121) and 
    (1.4213365337, 4.2613290546) .. (1.4213365337, 4.2669943658) -- (1.4213365337,
     4.2755376558).. controls (1.4213365337, 4.2762401246) and (1.4219559233, 
    4.27889924) .. (1.4243958579, 4.280417525).. controls (1.4264352662, 
    4.2815959792) and (1.4290111967, 4.2815957146) .. (1.4320176042, 4.2806371292)
     -- (1.4496633017, 4.2735819881).. controls (1.4507510037, 4.2731890819) and 
    (1.4515442246, 4.2720032194) .. (1.4514534725, 4.2707493854) -- (1.4506602517,
     4.262379715).. controls (1.4505385433, 4.2613750921) and (1.4501842662, 
    4.2605895442) .. (1.4493987183, 4.2602040462).. controls (1.4486205787, 
    4.2596375733) and (1.4476162204, 4.2596375733) .. (1.4468304079, 4.2600691087)
     -- (1.4388158083, 4.2633851581).. controls (1.4374182792, 4.263951631) and 
    (1.4366327312, 4.2656513144) .. (1.4373352, 4.2670864937).. controls 
    (1.4378111854, 4.2685671021) and (1.4395939479, 4.269133575) .. (1.4409913183,
     4.2685671021).. controls (1.4409913183, 4.2685671021) and (1.4432122308, 
    4.2675701521) .. (1.4454782812, 4.2666938521).. controls (1.4455682396, 
    4.2678723063) and (1.4456978854, 4.268257275) .. (1.4456978854, 4.269133575)..
     controls (1.4426461548, 4.2704025167) and (1.4300540519, 4.2753276825) .. 
    (1.4300540519, 4.2753276825).. controls (1.4281808019, 4.2760224783) and 
    (1.427395254, 4.2757129158) .. (1.4272290163, 4.2756224283) -- (1.4268739454, 
    4.2749273679) -- (1.4268739454, 4.2669959269).. controls (1.4268739454, 
    4.2643370231) and (1.4293136683, 4.2635514752) .. (1.4297971944, 4.2633852375)
     -- (1.453538786, 4.253754166).. controls (1.4564621408, 4.2527569515) and 
    (1.4576329221, 4.2533612598) .. (1.458033316, 4.253579541).. controls 
    (1.4582899619, 4.253754166) and (1.4582899619, 4.2540555265) .. (1.4582899619,
     4.2541386056) -- (1.4582899619, 4.2720033252).. controls (1.4582449827, 
    4.2757122015) and (1.4549812679, 4.2764977494) .. (1.4543316894, 4.2765808285)
     -- (1.4205588175, 4.2892712242).. controls (1.4187308112, 4.2898302887) and 
    (1.4176355685, 4.2896490492) .. (1.4173333615, 4.2894379117).. controls 
    (1.417023799, 4.2892712242) and (1.416857376, 4.2884406971) .. (1.4169327823, 
    4.2880480554) -- (1.4169327823, 4.2601142202).. controls (1.4168560531, 
    4.2566696627) and (1.4185038781, 4.2561936773) .. (1.4195540094, 
    4.2561030575).. controls (1.4199014073, 4.2561030575) and (1.4202490698, 
    4.2561031633) .. (1.4205588969, 4.2558913908) -- (1.4548908069, 4.243344611)..
     controls (1.4651714833, 4.23845736) and (1.4639100558, 4.2316134623) .. 
    (1.4639100558, 4.2313944402) -- (1.4639100558, 4.2091259427).. controls 
    (1.4639100558, 4.2045105775) and (1.4627392746, 4.2013757677) .. 
    (1.4601633177, 4.1998045925).. controls (1.4568547296, 4.1976744321) and 
    (1.4531080708, 4.1994193592) .. (1.452715244, 4.1995929258) -- cycle;
  
  
  
      \end{scope}
      \begin{scope}[shift={(0.8731702173, -3.5322052271)}]
        \path[fill=c005093,even odd rule] (0.9944876023, 4.0403358519).. controls 
    (0.9919192919, 4.042080779) and (0.99074097, 4.0453439117) .. (0.99074097, 
    4.0499215208) -- (0.99074097, 4.0727113269).. controls (0.99074097, 
    4.0725367019) and (0.9892603617, 4.0796834152) .. (0.9998960825, 4.0832185925)
     -- (1.0344924435, 4.0910064971).. controls (1.0347945977, 4.0911414346) and 
    (1.0351041602, 4.0911414346) .. (1.0353610706, 4.0911414346).. controls 
    (1.0353610706, 4.0911414346) and (1.0364487727, 4.0910064971) .. (1.037060754,
     4.0915343408).. controls (1.0377555498, 4.0920934054) and (1.0377558144, 
    4.0934983429) .. (1.0377558144, 4.094147895) -- (1.0377558144, 4.1221195921)..
     controls (1.0377558144, 4.1223471337) and (1.037755735, 4.1235171212) .. 
    (1.037151506, 4.1239928421).. controls (1.0365395248, 4.124559315) and 
    (1.035187504, 4.1246876379) .. (1.0334879265, 4.1243931567) -- (1.0003721208, 
    4.1166355469).. controls (0.998801025, 4.1165455885) and (0.9963610375, 
    4.1160235656) .. (0.9962326617, 4.1127075427) -- (0.9962326617, 4.094842929)..
     controls (0.9962326617, 4.0949249498) and (0.9963623075, 4.0944124519) .. 
    (0.9969277221, 4.0938835498).. controls (0.9974111158, 4.0935813956) and 
    (0.9986726492, 4.0930146581) .. (1.0017243269, 4.0934983165) -- (1.02481634, 
    4.0997680154).. controls (1.0261683608, 4.0999876196) and (1.0278679648, 
    4.1005535633) .. (1.0278679648, 4.1025175654) -- (1.0278679648, 
    4.1110533677).. controls (1.027738319, 4.1112359302) and (1.026954094, 
    4.1124508969) .. (1.0240760094, 4.1118389156).. controls (1.0240760094, 
    4.1118389156) and (1.0123978302, 4.1088701583) .. (1.0089985958, 
    4.1080014519).. controls (1.0090885542, 4.1068306706) and (1.0090885542, 
    4.1065285165) .. (1.0093081583, 4.1051688227).. controls (1.0114835625, 
    4.1058712915) and (1.0140142492, 4.1065285165) .. (1.0140142492, 
    4.1065285165).. controls (1.0154493492, 4.1069060769) and (1.0170584127, 
    4.1060374498) .. (1.0174512131, 4.1044737623).. controls (1.0179346069, 
    4.103083906) and (1.0170583069, 4.1015126515) .. (1.0154493756, 4.101119851) 
    -- (1.0075178287, 4.0987632073).. controls (1.0067775246, 4.0985885823) and 
    (1.0058633892, 4.0987632073) .. (1.0051685933, 4.099246601).. controls 
    (1.0045113683, 4.0997678302) and (1.0039901392, 4.1005533781) .. 
    (1.0038996517, 4.1013315177) -- (1.0032048558, 4.1098747813).. controls 
    (1.0029852517, 4.1111437229) and (1.0038996517, 4.1124960083) .. 
    (1.0053879329, 4.1127075162) -- (1.0228146083, 4.117322855).. controls 
    (1.0259493917, 4.1180253238) and (1.0283061412, 4.1176326821) .. 
    (1.0303078994, 4.1162351529).. controls (1.0328383744, 4.1144902258) and 
    (1.0333143069, 4.1117028669) .. (1.0334426827, 4.1115366292) -- (1.0334426827,
     4.102517486).. controls (1.0334426827, 4.096580236) and (1.0283061412, 
    4.0944575633) .. (1.0258662331, 4.0942838644) -- (1.002895161, 4.0879612752)..
     controls (0.9986726227, 4.0874326377) and (0.995575569, 4.0877866502) .. 
    (0.9933094921, 4.0895778794).. controls (0.9908771775, 4.0916249606) and 
    (0.9907411817, 4.0944121608) .. (0.9907411817, 4.0946690183) -- (0.9907411817,
     4.1127075427).. controls (0.9909158067, 4.1199440556) and (0.9966254356, 
    4.1220288665) .. (0.99958652, 4.1221195127) -- (1.032445574, 4.1298773342).. 
    controls (1.0360638304, 4.13057213) and (1.038934295, 4.1299672925) .. 
    (1.040852974, 4.1282228946).. controls (1.0435119042, 4.1257829071) and 
    (1.0433382052, 4.1223461019) .. (1.0433382052, 4.1219457344) -- (1.0433382052,
     4.094283785).. controls (1.0435128302, 4.0905749087) and (1.0420314281, 
    4.0885731506) .. (1.040852974, 4.0874325319).. controls (1.0388058927, 
    4.0855214465) and (1.036366011, 4.0854306944) .. (1.0352782825, 
    4.0854306944).. controls (1.0346210575, 4.0852560694) and (1.0013315, 
    4.0777635985) .. (1.0013315, 4.0777635985).. controls (0.9964517367, 
    4.0761017506) and (0.9962327146, 4.0737601087) .. (0.9962327146, 4.0732690685)
     -- (0.9962327146, 4.0499203037).. controls (0.9962327146, 4.0473972371) and 
    (0.9967539437, 4.0455163935) .. (0.9977133229, 4.044866815).. controls 
    (0.9984988708, 4.0442548338) and (0.9998964, 4.0445646608) .. (0.9998964, 
    4.0445646608) -- (1.0322265519, 4.0504491).. controls (1.036366011, 
    4.0510986521) and (1.0375369246, 4.0527075833) .. (1.0375369246, 4.0535762369)
     -- (1.0375369246, 4.0705269004).. controls (1.0377565287, 4.074840085) and 
    (1.0366682975, 4.0756332794) .. (1.0366682975, 4.0757239256) -- (1.008915649, 
    4.0695223569).. controls (1.005946971, 4.0683515756) and (1.0057731662, 
    4.0668256971) .. (1.0057731662, 4.0668256971) -- (1.0060377496, 
    4.0575873202).. controls (1.0060377496, 4.0574126952) and (1.0061673954, 
    4.0564996181) .. (1.0061673954, 4.0564996181).. controls (1.0063869996, 
    4.0564096598) and (1.0071267746, 4.0564096598) .. (1.0076480037, 
    4.056581639).. controls (1.0076480037, 4.056581639) and (1.0219397635, 
    4.059331189) .. (1.0254749408, 4.0600260906).. controls (1.0254749408, 
    4.0611137927) and (1.0254749408, 4.061725774) .. (1.0254749408, 
    4.0630778213).. controls (1.0217736317, 4.0621184421) and (1.0162367756, 
    4.0608569087) .. (1.0162367756, 4.0608569087).. controls (1.0147940027, 
    4.0603282713) and (1.013139616, 4.0612873858) .. (1.0128374354, 
    4.0627756671).. controls (1.0123992854, 4.0642562754) and (1.0133134208, 
    4.0657292108) .. (1.0148845167, 4.0662202246) -- (1.0276503979, 4.069475949)..
     controls (1.0283907021, 4.0695209281) and (1.029395325, 4.0694758431) .. 
    (1.0300070417, 4.0689621281).. controls (1.0306565937, 4.0683501469) and 
    (1.0310495, 4.0676024344) .. (1.0310495, 4.0666885635) -- (1.0310495, 
    4.057812851).. controls (1.0310495, 4.0563171615) and (1.0300901208, 
    4.0552371323) .. (1.0287833437, 4.0549348988) -- (1.0089999452, 
    4.0512260225).. controls (1.0089999452, 4.0510963767) and (1.0055629548, 
    4.0501380558) .. (1.002896484, 4.0518755746).. controls (1.0017257027, 
    4.0528801975) and (1.0003734173, 4.0544060496) .. (1.0003734173, 4.0574125365)
     -- (1.0001538131, 4.0665980231).. controls (1.0000717923, 4.068607269) and 
    (1.0013322673, 4.0727089456) .. (1.0072996798, 4.0748390796) -- (1.0350145194,
     4.0809424615).. controls (1.0356640715, 4.0814184469) and (1.0378019048, 
    4.0815922781) .. (1.03981123, 4.0803381531).. controls (1.0421604654, 
    4.0787670573) and (1.0432481146, 4.0754507963) .. (1.0430290925, 4.0705256833)
     -- (1.0430290925, 4.0536657454).. controls (1.043338655, 4.0514903413) and 
    (1.0416392362, 4.0462101046) .. (1.0332318627, 4.0449485712) -- (1.0012415417,
     4.0391624248).. controls (1.0012415417, 4.0391624248) and (0.9976308523, 
    4.0379839706) .. (0.9944884754, 4.040333206) -- cycle;
  
  
  
      \end{scope}
      \begin{scope}[shift={(0.8731702173, -3.5300147417)}]
        \path[fill=c005093,even odd rule] (0.932484456, 4.0237249927).. controls 
    (0.9300370602, 4.0253791677) and (0.9286471775, 4.028521624) .. (0.9286471775,
     4.0332729321) -- (0.9286471775, 4.0560022808).. controls (0.9286471775, 
    4.0557906142) and (0.9273856442, 4.0628459669) .. (0.9377948287, 4.0668646173)
     -- (0.9723080312, 4.0765787944).. controls (0.9726175938, 4.0765787944) and 
    (0.9728745042, 4.0765787944) .. (0.9732674104, 4.0765787944).. controls 
    (0.9732674104, 4.0765787944) and (0.9743551125, 4.0765789002) .. (0.975140925,
     4.0772812631).. controls (0.9757529063, 4.0779686506) and (0.9758357208, 
    4.0791545131) .. (0.9758357208, 4.0800231402) -- (0.9758357208, 
    4.1079948373).. controls (0.9758357208, 4.1083422352) and (0.9757537, 
    4.1093471227) .. (0.9752692479, 4.109732356).. controls (0.9746649396, 
    4.1103065019) and (0.9733959979, 4.1103062373) .. (0.9717416377, 4.1098672935)
     -- (0.938580509, 4.1003268419).. controls (0.9371906527, 4.1000993002) and 
    (0.9344410233, 4.099541294) .. (0.9342220277, 4.0960136837) -- (0.9342220277, 
    4.07814907).. controls (0.9342220277, 4.0782707783) and (0.9343569652, 
    4.0777487554) .. (0.9348337444, 4.0773558492).. controls (0.9352266506, 
    4.0771891617) and (0.9365334277, 4.0765779742) .. (0.9396681317, 4.07727912) 
    -- (0.9626769598, 4.0847270879).. controls (0.9637646619, 4.0849837337) and 
    (0.9658646333, 4.0855052275) .. (0.9658646333, 4.0879450562).. controls 
    (0.9658646333, 4.0879450562) and (0.9658646333, 4.095921794) .. (0.9658646333,
     4.0961861392).. controls (0.9658646333, 4.0961861392) and (0.9656370917, 
    4.0965713725) .. (0.9651243292, 4.0967981204).. controls (0.9645956917, 
    4.0971833538) and (0.9635908042, 4.0971833538) .. (0.9622009744, 
    4.0967981204).. controls (0.9622009744, 4.0967981204) and (0.9503491227, 
    4.0930891383) .. (0.9470330469, 4.0920920296).. controls (0.947078026, 
    4.0909966546) and (0.947078026, 4.0906870921) .. (0.9472076719, 
    4.0893424796).. controls (0.9495645802, 4.0900827837) and (0.9519136569, 
    4.0909061671) .. (0.9519136569, 4.0909061671).. controls (0.9534016735, 
    4.0913895608) and (0.9549653875, 4.0906865629) .. (0.9554488606, 
    4.0891686483).. controls (0.955924846, 4.0877711192) and (0.9551842773, 
    4.0861999969) .. (0.953794421, 4.085595715) -- (0.9456439844, 4.0828536263).. 
    controls (0.9448584365, 4.0826710638) and (0.9439895448, 4.0828086471) .. 
    (0.943287076, 4.0832388596).. controls (0.9425015281, 4.0837222533) and 
    (0.9419877073, 4.0845078013) .. (0.9418972198, 4.0853764283) -- (0.941202424, 
    4.0938743687).. controls (0.9411124656, 4.0950073146) and (0.9418972198, 
    4.0963974354) .. (0.9431587531, 4.0967978294) -- (0.9607591275, 
    4.1021534194).. controls (0.9638561812, 4.1029768027) and (0.9663413067, 
    4.1026746485) .. (0.9682977417, 4.1015038673).. controls (0.9708282167, 
    4.0997136965) and (0.9712210965, 4.0971831421) .. (0.9713494723, 4.0967978823)
     -- (0.9713494723, 4.087944871).. controls (0.9713494723, 4.0819773527) and 
    (0.9666434873, 4.0797565725) .. (0.9640752033, 4.0793184225) -- (0.9411116983,
     4.0718703752).. controls (0.9367985137, 4.070910996) and (0.9337013806, 
    4.0712586585) .. (0.9314352773, 4.072950669).. controls (0.9289952898, 
    4.0748315919) and (0.9287763735, 4.0777473002) .. (0.9287763735, 4.0779663223)
     -- (0.9286467277, 4.0961406838).. controls (0.929039634, 4.1030675548) and 
    (0.9345764106, 4.1054167108) .. (0.9374921983, 4.1058095113) -- (0.9703435529,
     4.1153122333).. controls (0.9740071325, 4.1162263688) and (0.9767944385, 
    4.1158334625) .. (0.9788415198, 4.1140432917).. controls (0.9811529198, 
    4.1120869625) and (0.9814175031, 4.1093447679) .. (0.9814175031, 
    4.1081663931).. controls (0.9814175031, 4.1079917681) and (0.9814175031, 
    4.1077737515) .. (0.9814175031, 4.1077737515) -- (0.9814175031, 
    4.0801041556).. controls (0.9814175031, 4.0800221348) and (0.9814175031, 
    4.0797567577) .. (0.9814175031, 4.0797567577).. controls (0.9814175031, 
    4.0763197673) and (0.9800654823, 4.0743104421) .. (0.9789777802, 
    4.0732604696).. controls (0.9770137781, 4.0712586321) and (0.9747476748, 
    4.0710773925) .. (0.9736146231, 4.0710773925).. controls (0.9732217169, 
    4.0709953717) and (0.9394564121, 4.061499079) .. (0.9394564121, 4.061499079)..
     controls (0.9345766223, 4.0595803206) and (0.9342216573, 4.0570499515) .. 
    (0.9342216573, 4.0565740719).. controls (0.9342216573, 4.056492051) and 
    (0.9342216573, 4.0332706302) .. (0.9342216573, 4.0332706302).. controls 
    (0.9342216573, 4.0296901033) and (0.9352262802, 4.0285193221) .. 
    (0.9356189219, 4.028338056).. controls (0.9363213906, 4.0279377415) and 
    (0.9374093573, 4.028163431) .. (0.9377945906, 4.0282084102) -- (0.9701700656, 
    4.0359661258).. controls (0.9743548479, 4.0367895092) and (0.9755332227, 
    4.0384891925) .. (0.9755332227, 4.0394031163) -- (0.9755332227, 
    4.0563991029).. controls (0.9756628685, 4.0604932654) and (0.9747476748, 
    4.0614979148) .. (0.9747476748, 4.0614979148) -- (0.9470328617, 
    4.0536119027).. controls (0.9439886981, 4.052342961) and (0.9437695702, 
    4.0506884685) .. (0.9437695702, 4.0505147431) -- (0.9440341535, 
    4.0413672242).. controls (0.9440341535, 4.0412772658) and (0.9441637994, 
    4.0403623367) .. (0.9441637994, 4.0403623367).. controls (0.944375466, 
    4.0403623367) and (0.9450324265, 4.0404919825) .. (0.9455536556, 
    4.0406718992).. controls (0.9455536556, 4.0406718992) and (0.9601549779, 
    4.0441996152) .. (0.9634635925, 4.0451136448).. controls (0.9634635925, 
    4.0462465906) and (0.9634635925, 4.046631824) .. (0.9634635925, 
    4.0481200788).. controls (0.9598906327, 4.0469871329) and (0.9543613702, 
    4.0452873438) .. (0.9543613702, 4.0452873438).. controls (0.9529185973, 
    4.0448944375) and (0.9513096131, 4.04568025) .. (0.9509168127, 4.0471608583)..
     controls (0.9504786627, 4.0486414667) and (0.9512189669, 4.05016716) .. 
    (0.9526995752, 4.0506883892) -- (0.9655107796, 4.0546163933).. controls 
    (0.9663418358, 4.0548730392) and (0.9672104629, 4.05482806) .. (0.9679960108, 
    4.0542160788).. controls (0.9686984796, 4.0536117704) and (0.9691742004, 
    4.0528262225) .. (0.9691742004, 4.0518668433) -- (0.9691742004, 
    4.0429307529).. controls (0.9691742004, 4.0416768925) and (0.968298165, 
    4.0404907654) .. (0.9669083088, 4.0402717433) -- (0.9472079894, 
    4.0353919535).. controls (0.9469058352, 4.0351723494) and (0.943635056, 
    4.0341833369) .. (0.941029016, 4.0359660994).. controls (0.939714566, 
    4.0367894827) and (0.9384530327, 4.0382248473) .. (0.9384530327, 4.0412764456)
     -- (0.9382784077, 4.0505147167).. controls (0.9381963869, 4.0509076229) and 
    (0.9380588035, 4.0559986825) .. (0.9451599556, 4.0589220373) -- (0.9728294721,
     4.0665514829).. controls (0.97361502, 4.0669517975) and (0.9758358796, 
    4.0673293579) .. (0.9777998817, 4.0660678246).. controls (0.9801491171, 
    4.06463246) and (0.9813275183, 4.0614978883) .. (0.9811538194, 4.0561800544) 
    -- (0.9811538194, 4.0394937096).. controls (0.9811538194, 4.0394937096) and 
    (0.9817202923, 4.0326424565) .. (0.9713489167, 4.0303914079) -- (0.9393208396,
     4.0229357671).. controls (0.9393208396, 4.0228537463) and (0.9357478798, 
    4.0214477504) .. (0.9324846942, 4.023721315) -- cycle;
  
  
  
      \end{scope}
      \begin{scope}[shift={(0.8731702173, -3.5267130063)}]
        \path[fill=c005093,even odd rule] (0.8713078417, 4.0001105862).. controls 
    (0.8688226104, 4.0016816821) and (0.8674704573, 4.0048165713) .. 
    (0.8674704573, 4.0095225563) -- (0.8674704573, 4.0322669598).. controls 
    (0.8674704573, 4.032010314) and (0.8661636802, 4.0387633538) .. (0.8764066535,
     4.0435523915) -- (0.910912289, 4.0562804373).. controls (0.9112218515, 
    4.0564153748) and (0.9114787619, 4.056492104) .. (0.911780916, 4.056492104).. 
    controls (0.912876291, 4.0566270415) and (0.9147495675, 4.0571945727) .. 
    (0.9146589477, 4.0603293825) -- (0.9146589477, 4.0883918052).. controls 
    (0.9146589477, 4.0887921198) and (0.914447281, 4.0896607469) .. (0.9142737144,
     4.0898345781).. controls (0.9137450769, 4.0901745677) and (0.9126571102, 
    4.0901743031) .. (0.9107385371, 4.0895699948) -- (0.8775774083, 
    4.0771062677).. controls (0.8764897062, 4.0768041135) and (0.8732189006, 
    4.0760185656) .. (0.8730526894, 4.0722719069) -- (0.8730526894, 
    4.0543618906).. controls (0.8737474852, 4.0537046656) and (0.8750090185, 
    4.0531911094) .. (0.878234554, 4.054150224) -- (0.9013264612, 4.0636528402).. 
    controls (0.9020667654, 4.0639851569) and (0.9045141348, 4.0647783777) .. 
    (0.9045141348, 4.0673995783) -- (0.9045141348, 4.0758523015).. controls 
    (0.903939989, 4.076018989) and (0.903154441, 4.076320614) .. (0.9013264612, 
    4.075717364).. controls (0.9013264612, 4.075717364) and (0.8888627342, 
    4.07083768) .. (0.8856826277, 4.0695686854).. controls (0.8858175652, 
    4.0685714708) and (0.8859022319, 4.0682693167) .. (0.8859921902, 
    4.0670456187).. controls (0.8882583465, 4.0679597542) and (0.8905244762, 
    4.0689641125) .. (0.8905244762, 4.0689641125).. controls (0.8920050846, 
    4.0695684208) and (0.8935686398, 4.0687365708) .. (0.8942258119, 
    4.0674760958).. controls (0.8948377931, 4.0660409958) and (0.894143791, 
    4.0644319058) .. (0.892654716, 4.0638654065) -- (0.8846402487, 4.0603302027)..
     controls (0.8838091925, 4.0600735569) and (0.8828500779, 4.0601635152) .. 
    (0.8821097737, 4.0605498069).. controls (0.8813694696, 4.0609879569) and 
    (0.8808934842, 4.0619018277) .. (0.8808029967, 4.0626876402) -- (0.8800174487,
     4.071147851).. controls (0.8799354279, 4.0722733885) and (0.8805915946, 
    4.0734894135) .. (0.8817623758, 4.0739805596) -- (0.8992342421, 
    4.0809375931).. controls (0.9024596717, 4.0820782119) and (0.9050280615, 
    4.0820329681) .. (0.9070751427, 4.0807709056).. controls (0.9096056177, 
    4.0792450535) and (0.9100437942, 4.0766313406) .. (0.9101344404, 4.0763216988)
     -- (0.9101344404, 4.0674006367).. controls (0.9101344404, 4.061599356) and 
    (0.9055115081, 4.0589404523) .. (0.9029809802, 4.0583663594) -- (0.8801081215,
     4.0489165275).. controls (0.8757949369, 4.0476097504) and (0.8725694015, 
    4.0477835817) .. (0.8703033246, 4.0493546775).. controls (0.8676443944, 
    4.0510921962) and (0.8674705896, 4.0538416404) .. (0.8674705896, 4.0541512823)
     -- (0.8674705896, 4.0722728594).. controls (0.8677727437, 4.0792450271) and 
    (0.8731359273, 4.0820322273) .. (0.8760970381, 4.0824703773) -- (0.9088200169,
     4.0948812935).. controls (0.9126572954, 4.0961502352) and (0.9155353271, 
    4.0958406727) .. (0.9175824083, 4.0942769852).. controls (0.9199393167, 
    4.0924942227) and (0.9202413121, 4.08966162) .. (0.9202413121, 4.0885663244)..
     controls (0.9202413121, 4.0883916994) and (0.9202413121, 4.088090339) .. 
    (0.9202413121, 4.088090339) -- (0.9202413121, 4.0604209283).. controls 
    (0.9202413121, 4.06033097) and (0.9202413121, 4.0601642825) .. (0.9202413121, 
    4.0600735304).. controls (0.9202413121, 4.056674296) and (0.9189797787, 
    4.0546196477) .. (0.91783916, 4.0534488665).. controls (0.9159280746, 
    4.0515301081) and (0.9138809669, 4.0510919581) .. (0.9127856712, 
    4.051008879).. controls (0.9126110462, 4.0509189206) and (0.878536814, 
    4.0383790465) .. (0.878536814, 4.0383790465).. controls (0.8733549494, 
    4.0360147298) and (0.8730527687, 4.0333632079) .. (0.8730527687, 
    4.032758926).. controls (0.8730527687, 4.0326689677) and (0.8730527687, 
    4.0095235088) .. (0.8730527687, 4.0095235088).. controls (0.8730527687, 
    4.006607721) and (0.8737478292, 4.0052103506) .. (0.8743521375, 
    4.0048175237).. controls (0.8747902875, 4.0045608779) and (0.8758856625, 
    4.0048175237) .. (0.8764068917, 4.0050371279) -- (0.9087823931, 
    4.0154084506).. controls (0.9143570846, 4.017100461) and (0.9143570846, 
    4.0195554769) .. (0.9143570846, 4.0197217146) -- (0.9143570846, 
    4.0365893254).. controls (0.9144470429, 4.0406381648) and (0.9136620242, 
    4.0415068713) .. (0.9136620242, 4.0415068713) -- (0.8859924283, 
    4.0313092475).. controls (0.8827669987, 4.0295717287) and (0.8825479767, 
    4.0277815315) .. (0.8825479767, 4.0276531292) -- (0.8828501308, 
    4.0183696408).. controls (0.8828501308, 4.0183696408) and (0.8829850683, 
    4.0177653325) .. (0.8829850683, 4.0176745804).. controls (0.8829850683, 
    4.0179841429) and (0.8841558496, 4.0181505658) .. (0.8841558496, 
    4.0181505658).. controls (0.8841558496, 4.0181505658) and (0.8991501046, 
    4.0231209754) .. (0.9022848879, 4.0241256512).. controls (0.9022848879, 
    4.0252210262) and (0.9022848879, 4.0255985867) .. (0.9022848879, 
    4.0270035771).. controls (0.8985382556, 4.0255986396) and (0.893356391, 
    4.0235590196) .. (0.893356391, 4.0235590196).. controls (0.8920041056, 
    4.0229925467) and (0.890304634, 4.0236489779) .. (0.8897381346, 
    4.0250396279).. controls (0.8892621492, 4.0266033154) and (0.8898201554, 
    4.0281744113) .. (0.8913092304, 4.0286578844) -- (0.9041128677, 
    4.0338851779).. controls (0.9048984156, 4.0341418237) and (0.9059030385, 
    4.0340518654) .. (0.906688851, 4.0334922717).. controls (0.9073836469, 
    4.0330541217) and (0.907904876, 4.0320947425) .. (0.907904876, 4.0310977925) 
    -- (0.907904876, 4.0222070252).. controls (0.907904876, 4.0209380835) and 
    (0.9070740844, 4.019933196) .. (0.9059863823, 4.0195555827) -- (0.8862105508, 
    4.0129687542).. controls (0.8858176446, 4.0128391083) and (0.8825015687, 
    4.0112690708) .. (0.8797142627, 4.0129687542).. controls (0.8785358085, 
    4.0138450542) and (0.8772745398, 4.0152272375) .. (0.8772745398, 4.0183697467)
     -- (0.8769723856, 4.0275626152).. controls (0.8768824273, 4.0279629298) and 
    (0.8768824273, 4.0328804227) .. (0.8838085046, 4.0363249802) -- (0.9113949419,
     4.046386661).. controls (0.912263569, 4.0469986423) and (0.9145297517, 
    4.047391284) .. (0.9164937537, 4.046212036).. controls (0.918978985, 
    4.0448221798) and (0.920021364, 4.04167975) .. (0.9199383112, 4.0365885052) --
     (0.9199383112, 4.0199324552).. controls (0.9199383112, 4.0199324552) and 
    (0.9206331071, 4.0130584744) .. (0.9104355892, 4.0101426867) -- (0.878452941, 
    3.9999450629).. controls (0.8783232952, 3.9997254588) and (0.8746156625, 
    3.9980718129) .. (0.8713070479, 4.0001117504) -- cycle;
  
  
  
      \end{scope}
      \begin{scope}[shift={(0.8731702173, -3.5348560875)}]
        \path[fill=c005093,even odd rule] (0.9789792354, 4.2408124427).. controls 
    (0.9765392479, 4.2425499615) and (0.9753608731, 4.245828096) .. (0.9753608731,
     4.2502244127) -- (0.9753608731, 4.2724023962).. controls (0.9753608731, 
    4.2721002421) and (0.9738805294, 4.2791554625) .. (0.9844708477, 4.2825999935)
     -- (1.0189764831, 4.2901386077).. controls (1.0193693894, 4.2903502744) and 
    (1.0197244602, 4.2903502744) .. (1.019981106, 4.2903502744).. controls 
    (1.0202832602, 4.2901386077) and (1.0212426394, 4.2903502744) .. 
    (1.0217712769, 4.290750589).. controls (1.0223303415, 4.2913548973) and 
    (1.0223300769, 4.2922687681) .. (1.0223300769, 4.2931829035) -- (1.0223300769,
     4.320240571).. controls (1.0223300769, 4.3204601752) and (1.0223301563, 
    4.321555021) .. (1.0217712769, 4.322113821).. controls (1.0212423748, 
    4.3226048877) and (1.019890354, 4.3227332106) .. (1.0182360731, 4.3223784044) 
    -- (0.9849465685, 4.3148850075).. controls (0.983511204, 4.31475007) and 
    (0.9809807025, 4.3142806992) .. (0.9808523002, 4.3111382694) -- (0.9808523002,
     4.2937571288).. controls (0.9808523002, 4.2938867746) and (0.980981946, 
    4.2934020579) .. (0.9815021169, 4.2930924954).. controls (0.982023346, 
    4.292714935) and (0.9833753669, 4.2921029537) .. (0.9863440448, 4.292714935) 
    -- (1.0094813546, 4.298757939).. controls (1.0105690567, 4.298757939) and 
    (1.0123517133, 4.2991582535) .. (1.0123517133, 4.3013339223) -- (1.0123517133,
     4.3095750052).. controls (1.0123517133, 4.3097866719) and (1.01152833, 
    4.3108365385) .. (1.0086957273, 4.3102700656).. controls (1.0086957273, 
    4.3102700656) and (0.9971460298, 4.3073921398) .. (0.9935730965, 
    4.3065234069).. controls (0.9937027423, 4.3054735402) and (0.9937027423, 
    4.305171386) .. (0.9937927006, 4.3039098527).. controls (0.996141936, 
    4.3045594048) and (0.9986346285, 4.3052544652) .. (0.9986346285, 
    4.3052544652).. controls (1.0001530723, 4.305609536) and (1.0016334954, 
    4.304733236) .. (1.002071619, 4.303418786).. controls (1.002509769, 
    4.3018175277) and (1.0015506544, 4.3002538137) .. (1.0001531252, 4.2998534198)
     -- (0.9921385256, 4.2977232594).. controls (0.9913074694, 4.2974211052) and 
    (0.9904388423, 4.2975883219) .. (0.9897816173, 4.2981087573).. controls 
    (0.9890868215, 4.298592151) and (0.9886939152, 4.299158624) .. (0.9886031631, 
    4.3001635115) -- (0.9878250235, 4.3083971067).. controls (0.9877350652, 
    4.3097869629) and (0.9886028985, 4.3109201733) .. (0.9900004277, 4.3113128944)
     -- (1.0075177229, 4.3156260525).. controls (1.0105694535, 4.3164116004) and 
    (1.0131377375, 4.3158456567) .. (1.0148826381, 4.3145835942).. controls 
    (1.0174509485, 4.3129217462) and (1.0179343687, 4.3103608706) .. 
    (1.0179343687, 4.3100512023) -- (1.0179343687, 4.3013342398).. controls 
    (1.0179343687, 4.2954573677) and (1.0131376317, 4.2934026135) .. 
    (1.0104787279, 4.2931835915) -- (0.9875153023, 4.2871784229).. controls 
    (0.9832927904, 4.2864381187) and (0.980067255, 4.2870487771) .. (0.9779295275,
     4.2887873542).. controls (0.9754442962, 4.2906153604) and (0.9753612171, 
    4.2934027194) .. (0.9753612171, 4.2936746317) -- (0.9753612171, 
    4.3113128415).. controls (0.9755358421, 4.3182850092) and (0.9812456562, 
    4.3202413648) .. (0.984161444, 4.3203319846) -- (1.0170581481, 4.3278252756)..
     controls (1.0207670244, 4.3284825006) and (1.0235997594, 4.3280078381) .. 
    (1.025601491, 4.3262615881).. controls (1.0282150452, 4.3239046798) and 
    (1.027912891, 4.3204603075) .. (1.027912891, 4.3200750477) -- (1.027912891, 
    4.2932742642).. controls (1.0281324952, 4.2897466275) and (1.0267344369, 
    4.2877826254) .. (1.025601491, 4.2866495737).. controls (1.023509166, 
    4.2846477362) and (1.0209861523, 4.2845646571) .. (1.0198908567, 
    4.2846477362).. controls (1.0191960608, 4.2844731112) and (0.9859517471, 
    4.2772375244) .. (0.9859517471, 4.2772375244).. controls (0.9810266342, 
    4.2756285931) and (0.9807623154, 4.2732717642) .. (0.9807623154, 4.2728865044)
     -- (0.9808522738, 4.2502250742).. controls (0.9808522738, 4.2477850867) and 
    (0.9814187467, 4.2461309381) .. (0.9822498029, 4.2453831462).. controls 
    (0.98311843, 4.2447335942) and (0.9844707154, 4.2452085212) .. (0.9844707154, 
    4.2452085212) -- (1.0169292696, 4.2508359971).. controls (1.020985676, 
    4.2514932221) and (1.0221564573, 4.2529661575) .. (1.0221564573, 4.2537442971)
     -- (1.0221564573, 4.2701360283).. controls (1.0222464156, 4.2742301908) and 
    (1.0213709094, 4.2751441146) .. (1.0213709094, 4.2751441146) -- (0.9933992123,
     4.2691388402).. controls (0.9905665037, 4.2682172965) and (0.9904381015, 
    4.2666083652) .. (0.9904381015, 4.2666083652) -- (0.9907402556, 
    4.2576722748).. controls (0.9907402556, 4.2574606081) and (0.9908751931, 
    4.2566221435) .. (0.9908751931, 4.2566221435).. controls (0.9908751931, 
    4.2566221435) and (0.9917438202, 4.2566222494) .. (0.9922198056, 
    4.256804706).. controls (0.9922198056, 4.256804706) and (1.0065568885, 
    4.259335181) .. (1.010175145, 4.2601132148).. controls (1.010175145, 
    4.2612461606) and (1.010175145, 4.2615935585) .. (1.010175145, 4.2630817604)..
     controls (1.0064285127, 4.2621223812) and (1.0008462542, 4.2607251167) .. 
    (1.0008462542, 4.2607251167).. controls (0.9992825667, 4.2603248021) and 
    (0.9978020906, 4.2612460812) .. (0.9974092637, 4.2626891187).. controls 
    (0.9970163575, 4.2642528062) and (0.9980135721, 4.2658240079) .. 
    (0.9995015887, 4.2660355687) -- (1.0122222527, 4.2692684596).. controls 
    (1.0131363881, 4.2694801262) and (1.013921936, 4.2693584179) .. (1.0146619756,
     4.2687398221).. controls (1.0154022798, 4.2682638367) and (1.0156665985, 
    4.2673952096) .. (1.0156665985, 4.2666096617) -- (1.0156665985, 
    4.2578094083).. controls (1.0156665985, 4.2565026312) and (1.0147902985, 
    4.2552410979) .. (1.0134004423, 4.2550145617) -- (0.9935718265, 4.251449169)..
     controls (0.9933072431, 4.2513142315) and (0.9900441104, 4.2503085502) .. 
    (0.9874305033, 4.2521818002).. controls (0.9863428012, 4.2530959356) and 
    (0.9850812679, 4.2546289315) .. (0.9850812679, 4.2574617723) -- (0.9848537262,
     4.2664357246).. controls (0.9846791012, 4.2683544829) and (0.9860321804, 
    4.2724032429) .. (0.9919165402, 4.2744502977) -- (1.019586136, 4.2805537854)..
     controls (1.020371684, 4.2809390187) and (1.022547141, 4.28111285) .. 
    (1.0244205233, 4.2798134812).. controls (1.0267319233, 4.2781516333) and 
    (1.027864975, 4.2750621731) .. (1.0277365992, 4.2701370602) -- (1.0277365992, 
    4.2538811925).. controls (1.0277365992, 4.2538811925) and (1.0283030721, 
    4.2469997504) .. (1.017841156, 4.245337929) -- (0.985813079, 4.2397633433).. 
    controls (0.985813079, 4.2397633433) and (0.9822477127, 4.2385470537) .. 
    (0.97897696, 4.24081321) -- cycle;
  
  
  
      \end{scope}
      \begin{scope}[shift={(0.8731702173, -3.5301041708)}]
        \path[fill=c005093,even odd rule] (0.9189787469, 4.2224260442).. controls 
    (0.9164935156, 4.2240802192) and (0.9152321146, 4.2273434577) .. 
    (0.9152321146, 4.2319588229) -- (0.9152321146, 4.2545823913).. controls 
    (0.9152321146, 4.2543998287) and (0.9138800938, 4.2614639392) .. 
    (0.9241607173, 4.2653843763) -- (0.958756999, 4.2751891731).. controls 
    (0.9591499052, 4.275318819) and (0.9594973031, 4.275318819) .. (0.9597616219, 
    4.275318819).. controls (0.9597616219, 4.275318819) and (0.9609324031, 
    4.2753187396) .. (0.9615443844, 4.2759757794).. controls (0.9622016094, 
    4.2765877606) and (0.9623299323, 4.2777659502) .. (0.9623299323, 4.2787253294)
     -- (0.9623299323, 4.3066970265).. controls (0.9623299323, 4.3070822598) and 
    (0.9622002865, 4.3080871473) .. (0.961717951, 4.308351466).. controls 
    (0.9612419656, 4.3089557744) and (0.9598899448, 4.3089557744) .. 
    (0.9581903408, 4.3085710702) -- (0.9250292121, 4.2989400781).. controls 
    (0.9236845996, 4.2987125365) and (0.92089732, 4.2981545302) .. (0.9208066738, 
    4.2946268142) -- (0.9208066738, 4.2768074442).. controls (0.9208066738, 
    4.2768074442) and (0.9208066738, 4.2763616213) .. (0.9214109821, 
    4.2759763879).. controls (0.9216755654, 4.2758017629) and (0.9229445071, 
    4.27519084) .. (0.9261169671, 4.2759763879) -- (0.9491710919, 4.2834244617).. 
    controls (0.950258794, 4.2836361283) and (0.9521775258, 4.2842176825) .. 
    (0.9521775258, 4.286604674).. controls (0.9521775258, 4.286604674) and 
    (0.9521775258, 4.2946267348) .. (0.9521775258, 4.2948836188).. controls 
    (0.9521775258, 4.2948836188) and (0.9521325467, 4.2952310167) .. 
    (0.9516110529, 4.2954122562).. controls (0.9511276592, 4.2958051625) and 
    (0.9500399571, 4.2958051625) .. (0.9487405883, 4.2954122562).. controls 
    (0.9487405883, 4.2954122562) and (0.9367980904, 4.2917939998) .. 
    (0.9335726344, 4.2907062713).. controls (0.9335726344, 4.2896108963) and 
    (0.9337022802, 4.2893010692) .. (0.9337022802, 4.2879641825).. controls 
    (0.9360136802, 4.2887423221) and (0.938408371, 4.2896109492) .. (0.938408371, 
    4.2896109492).. controls (0.9398437356, 4.2901395867) and (0.9414146727, 
    4.2893013867) .. (0.9418981458, 4.2877376992).. controls (0.9423741312, 
    4.2863930867) and (0.9417155833, 4.2848218056) .. (0.9401532187, 4.2842175502)
     -- (0.9321763752, 4.2816870752).. controls (0.9313077481, 4.281294169) and 
    (0.9304388565, 4.2815124502) .. (0.9297363877, 4.2818987419).. controls 
    (0.9290413273, 4.2823821356) and (0.9285656065, 4.283258171) .. (0.9284748544,
     4.2840363106) -- (0.9276893065, 4.2924437106).. controls (0.9276072856, 
    4.2938412398) and (0.9283917752, 4.295012021) .. (0.9296533085, 4.2954123885) 
    -- (0.947208254, 4.3008585983).. controls (0.9503506308, 4.3015914942) and 
    (0.9529190206, 4.3013798275) .. (0.9547923765, 4.3001182942).. controls 
    (0.9573228515, 4.2985471983) and (0.9577533815, 4.2958052419) .. 
    (0.9577533815, 4.2956239494) -- (0.957888319, 4.2866047004).. controls 
    (0.957888319, 4.2806825052) and (0.9530010679, 4.2783256763) .. (0.9504781071,
     4.2779328494) -- (0.9276052219, 4.2705680665).. controls (0.9232844967, 
    4.2694803644) and (0.9201496075, 4.2699182498) .. (0.9179741504, 
    4.2716557685).. controls (0.9153981671, 4.2736197706) and (0.9152321675, 
    4.2763617535) .. (0.9152321675, 4.2765883427) -- (0.9152321675, 
    4.2948457569).. controls (0.915398855, 4.3018556806) and (0.9211089337, 
    4.3042502656) .. (0.9240700446, 4.3044239646) -- (0.9568383465, 
    4.3139266867).. controls (0.9604112798, 4.3148408221) and (0.9632893379, 
    4.3144479158) .. (0.9653363927, 4.3126651533).. controls (0.967602549, 
    4.3107011513) and (0.967912376, 4.307959089) .. (0.967912376, 4.3067882813).. 
    controls (0.967912376, 4.3066983229) and (0.967912376, 4.3064784542) .. 
    (0.967912376, 4.3063879667) -- (0.967912376, 4.2788545254).. controls 
    (0.967912376, 4.2786349213) and (0.967912376, 4.2785521067) .. (0.967912376, 
    4.2783256233).. controls (0.967912376, 4.2750170087) and (0.9666431698, 
    4.2729247367) .. (0.9654649802, 4.2718823313).. controls (0.9634631427, 
    4.2700921604) and (0.9613330881, 4.2696992542) .. (0.9601547398, 
    4.2696992542).. controls (0.959671346, 4.2696092958) and (0.9259889352, 
    4.2600682621) .. (0.9259889352, 4.2600682621).. controls (0.9210260927, 
    4.2582402558) and (0.9208070706, 4.2557549981) .. (0.9208070706, 4.2551960394)
     -- (0.9208070706, 4.2319605163).. controls (0.9208070706, 4.2283045038) and 
    (0.9216756977, 4.227345151) .. (0.9221969269, 4.2269145946).. controls 
    (0.9227710727, 4.22651428) and (0.923941854, 4.2267346779) .. (0.9242440081, 
    4.2268696154) -- (0.9567101294, 4.2346272252).. controls (0.9609326677, 
    4.2353675294) and (0.9620280427, 4.2371502919) .. (0.9620280427, 4.2381095387)
     -- (0.9620280427, 4.2549696883).. controls (0.962118001, 4.2591091475) and 
    (0.9613329823, 4.2600230713) .. (0.9613329823, 4.2600230713) -- (0.933572846, 
    4.2523183252).. controls (0.9304380627, 4.2510038752) and (0.930302014, 
    4.2492211921) .. (0.930302014, 4.2492211921) -- (0.9305216181, 4.240118864).. 
    controls (0.9305216181, 4.2398468723) and (0.9306962431, 4.239114241) .. 
    (0.9306962431, 4.239114241).. controls (0.930825889, 4.239114241) and 
    (0.9315272994, 4.239114241) .. (0.9319577765, 4.2393338452).. controls 
    (0.9319577765, 4.2393338452) and (0.9466045277, 4.242906805) .. (0.9500415181,
     4.2438208081).. controls (0.9500415181, 4.2447801873) and (0.9500415181, 
    4.2453844956) .. (0.9500415181, 4.2467440571).. controls (0.9462948594, 
    4.2455581946) and (0.94093919, 4.2440398302) .. (0.94093919, 4.2440398302).. 
    controls (0.9393225858, 4.2434204406) and (0.9377968131, 4.2444250635) .. 
    (0.9374115533, 4.2457847573).. controls (0.9368829158, 4.2473484448) and 
    (0.9376681992, 4.2489195406) .. (0.9392850679, 4.2492670444) -- (0.9520056525,
     4.2532327252).. controls (0.9529197879, 4.2535348794) and (0.9537958233, 
    4.2534443919) .. (0.9544453754, 4.2528324106).. controls (0.9551856796, 
    4.2523188544) and (0.9556238296, 4.2514499627) .. (0.9556238296, 4.2506570065)
     -- (0.9556238296, 4.2416378369).. controls (0.9556238296, 4.2402479806) and 
    (0.9547930379, 4.2391147702) .. (0.9534859962, 4.2389335042) -- (0.9337025977,
     4.2340084177).. controls (0.9333551998, 4.2338787719) and (0.9300390181, 
    4.2328450448) .. (0.9273876815, 4.2346278073).. controls (0.9263378148, 
    4.2353681115) and (0.924947694, 4.236931799) .. (0.924947694, 4.2398475337) --
     (0.9247280898, 4.2491311279).. controls (0.924646069, 4.2496145217) and 
    (0.9245534648, 4.2544943644) .. (0.931700284, 4.2575385544) -- (0.959369774, 
    4.2652131908).. controls (0.9601553219, 4.2656135054) and (0.9623308848, 
    4.2659910658) .. (0.9642042406, 4.2647750408).. controls (0.9667347156, 
    4.263203945) and (0.9676941212, 4.2600238385) .. (0.9676035015, 4.2549703498) 
    -- (0.9676035015, 4.2381102002).. controls (0.9676035015, 4.2381102002) and 
    (0.9682982973, 4.231221191) .. (0.9578892185, 4.2290986242) -- (0.9259443, 
    4.2215599042).. controls (0.9259443, 4.2215599042) and (0.922288314, 
    4.2201549667) .. (0.9189796994, 4.2224285312) -- cycle;
  
  
  
      \end{scope}
      \begin{scope}[shift={(0.8731702173, -3.5267058625)}]
        \path[fill=c005093,even odd rule] (0.8597123448, 4.1984585512).. controls 
    (0.8570534146, 4.1999767304) and (0.8557919077, 4.2031646421) .. 
    (0.8557919077, 4.2077422513) -- (0.8557919077, 4.2304790877).. controls 
    (0.8557919077, 4.2302224419) and (0.8545756181, 4.2370583227) .. 
    (0.8647279717, 4.2418928158) -- (0.8992337129, 4.2545378883).. controls 
    (0.8994083379, 4.254749555) and (0.8998078588, 4.254749555) .. (0.9001478483, 
    4.254749555).. controls (0.901197715, 4.2548395133) and (0.9030710973, 
    4.2553237008) .. (0.9028520752, 4.2586775592) -- (0.9028520752, 4.286520854)..
     controls (0.902981721, 4.2871251623) and (0.902852181, 4.2879107102) .. 
    (0.9025954294, 4.288046706).. controls (0.9020667919, 4.288394104) and 
    (0.9010243335, 4.2884396123) .. (0.8991508719, 4.2879117685) -- (0.8658991233,
     4.2753271269).. controls (0.8648114213, 4.2751525019) and (0.8616764792, 
    4.2742772602) .. (0.8614574571, 4.2705305221) -- (0.8614574571, 
    4.2525751827).. controls (0.8620694383, 4.2519632015) and (0.86333838, 
    4.2514874806) .. (0.8666015921, 4.2525751827) -- (0.8896028531, 
    4.2619872585).. controls (0.8905243969, 4.2622148002) and (0.8928358498, 
    4.2629918815) .. (0.8928358498, 4.2656508381) -- (0.8928358498, 
    4.2741487785).. controls (0.892261704, 4.2743683827) and (0.8915213998, 
    4.2746623348) .. (0.8896028531, 4.2740588202).. controls (0.8896028531, 
    4.2740588202) and (0.8771844492, 4.26908849) .. (0.8741402856, 4.2677892271)..
     controls (0.8742223065, 4.2669129271) and (0.8742223065, 4.2666107729) .. 
    (0.8743969315, 4.2652965875).. controls (0.8765800085, 4.2662107229) and 
    (0.8787177625, 4.2671698375) .. (0.8787177625, 4.2671698375).. controls 
    (0.8803266938, 4.2677892271) and (0.8818904077, 4.2670878167) .. 
    (0.8825022567, 4.2656969021).. controls (0.8830234858, 4.2643372083) and 
    (0.8823726108, 4.2625922813) .. (0.8809764046, 4.2619880258) -- (0.8730525306,
     4.2586794113).. controls (0.8721762306, 4.2583320133) and (0.8711716077, 
    4.2584598071) .. (0.8704313035, 4.2588460988).. controls (0.8696909994, 
    4.2592464133) and (0.8691245265, 4.2601150404) .. (0.8691245265, 4.2609913404)
     -- (0.8683389785, 4.2694893867).. controls (0.8683389785, 4.2706223325) and 
    (0.8689131244, 4.2717933783) .. (0.8700839056, 4.2722767721) -- (0.8876917413,
     4.2792489663).. controls (0.8907812279, 4.2803289954) and (0.8933571583, 
    4.2802385079) .. (0.8954419956, 4.2790293621).. controls (0.8979724706, 
    4.2774958371) and (0.8984559702, 4.2749276325) .. (0.8985390494, 4.2745348321)
     -- (0.8985390494, 4.2656516319).. controls (0.8985390494, 4.259940971) and 
    (0.8938330644, 4.2572819615) .. (0.8913101035, 4.2566247894) -- (0.868512704, 
    4.2472581425).. controls (0.8641163873, 4.2458230425) and (0.8608003115, 
    4.2458682863) .. (0.8585416958, 4.2474698092).. controls (0.8560943, 
    4.2493128967) and (0.8557921458, 4.2521757942) .. (0.8557921458, 4.2524930031)
     -- (0.8557921458, 4.2706221473).. controls (0.8560943, 4.2774961281) and 
    (0.8615481033, 4.280238111) .. (0.8644185679, 4.2807291248) -- (0.8972776219, 
    4.2931852848).. controls (0.9009336079, 4.294363739) and (0.9039400419, 
    4.2941443994) .. (0.9059870967, 4.2926185473).. controls (0.9082984967, 
    4.2907450327) and (0.9085628154, 4.2879125623) .. (0.9085628154, 
    4.2867870777).. controls (0.9085628154, 4.2865224944) and (0.9085628154, 
    4.2863941715) .. (0.9085628154, 4.286303684) -- (0.9085628154, 4.2587625698)..
     controls (0.9085628154, 4.2585429656) and (0.9085628154, 4.2584604156) .. 
    (0.9085628154, 4.2582865844).. controls (0.9085628154, 4.2549326731) and 
    (0.9073012821, 4.2528403746) .. (0.9061230925, 4.2517979692).. controls 
    (0.9042495779, 4.2498339671) and (0.9022026554, 4.2493127379) .. 
    (0.9011073333, 4.2493127379).. controls (0.9009327083, 4.2492677588) and 
    (0.866858476, 4.236720635) .. (0.866858476, 4.236720635).. controls 
    (0.8616766115, 4.2341446517) and (0.8613744308, 4.2316143619) .. 
    (0.8613744308, 4.2310931592) -- (0.8614564517, 4.2077443944).. controls 
    (0.8614564517, 4.2049116858) and (0.8619398454, 4.2035519392) .. 
    (0.8626727413, 4.2031666794).. controls (0.863156135, 4.2027663648) and 
    (0.8643345892, 4.2030370335) .. (0.8647272308, 4.2033413044) -- (0.8971028381,
     4.2137127065).. controls (0.9025944504, 4.2153293106) and (0.9026775296, 
    4.2179353506) .. (0.9026775296, 4.2179353506) -- (0.9026775296, 
    4.2349313373).. controls (0.9027674879, 4.2389424206) and (0.9019827338, 
    4.2398563444) .. (0.9018919817, 4.2398563444) -- (0.8743960848, 
    4.2295227777).. controls (0.8710874702, 4.2278230944) and (0.8709592002, 
    4.225995141) .. (0.8709592002, 4.2259120883) -- (0.8711708669, 4.2166662502)..
     controls (0.8711708669, 4.2165842294) and (0.8713058044, 4.2161071856) .. 
    (0.8713058044, 4.2159788627).. controls (0.871342846, 4.2161085085) and 
    (0.8725673377, 4.216364096) .. (0.8725673377, 4.216364096).. controls 
    (0.8725673377, 4.216364096) and (0.8874710523, 4.2213723675) .. (0.8906889148,
     4.2224676631).. controls (0.8906889148, 4.223472286) and (0.8906889148, 
    4.2239480069) .. (0.8906889148, 4.2253002658).. controls (0.8868590975, 
    4.2238574929) and (0.8816772329, 4.2218178729) .. (0.8816772329, 
    4.2218178729).. controls (0.8803249475, 4.2212058917) and (0.878534856, 
    4.2218179788) .. (0.8781420292, 4.2233363167).. controls (0.8774472333, 
    4.2247412542) and (0.8782319875, 4.2265240167) .. (0.8796300458, 4.2269545731)
     -- (0.8923959535, 4.2319703058).. controls (0.8932645806, 4.2323555392) and 
    (0.8942239598, 4.2323555392) .. (0.8950095077, 4.2317877433).. controls 
    (0.895832891, 4.2312286788) and (0.8963162848, 4.2303978871) .. (0.8963162848,
     4.2295215871) -- (0.8963162848, 4.2204117977).. controls (0.8963162848, 
    4.2191957727) and (0.8954854931, 4.2181607227) .. (0.8943522827, 4.217760461) 
    -- (0.8745311281, 4.2111810938).. controls (0.8742215656, 4.2110461563) and 
    (0.8709127658, 4.2095266542) .. (0.8681707829, 4.2114006979).. controls 
    (0.8669923288, 4.2119597625) and (0.8655950642, 4.2135308583) .. 
    (0.8655950642, 4.2165826683) -- (0.8653384183, 4.2259115856).. controls 
    (0.8653384183, 4.2261682315) and (0.8651188142, 4.2310934502) .. 
    (0.8721745373, 4.2346210869) -- (0.8997156515, 4.2446072556).. controls 
    (0.9005842785, 4.2452189723) and (0.9028504613, 4.2457325285) .. 
    (0.9048522194, 4.2445622765).. controls (0.9072996152, 4.2431647473) and 
    (0.9084327198, 4.2398941533) .. (0.9082967769, 4.2348029085) -- (0.9082967769,
     4.2180258381).. controls (0.9082967769, 4.2180258381) and (0.908818006, 
    4.2114012535) .. (0.8987563252, 4.2083495229) -- (0.8667736506, 
    4.1980612529).. controls (0.8666387131, 4.1980612529) and (0.8629362663, 
    4.1962784904) .. (0.8597108367, 4.1984615675) -- cycle;
  
  
  
      \end{scope}
      \begin{scope}[shift={(0.8731702173, -3.5863580271)}]
        \path[fill=c005093,even odd rule] (0.9074754044, 3.9753959365).. controls 
    (0.8960163531, 3.9734470156) and (0.8855089817, 3.9795806129) .. 
    (0.8838094042, 3.9892192779).. controls (0.8822835521, 3.9988428088) and 
    (0.8903055865, 4.0080433237) .. (0.901681585, 4.009916706).. controls 
    (0.9131406362, 4.0118807081) and (0.9237689223, 4.0058149765) .. 
    (0.9253325304, 3.9961838785).. controls (0.9269036262, 3.9865603742) and 
    (0.9189797523, 3.9773599385) .. (0.9074754044, 3.9753959365) -- 
    cycle(0.9067729356, 3.9987597296).. controls (0.9030716265, 3.99814034) and 
    (0.9004126169, 3.9949223452) .. (0.9009338196, 3.9916515925).. controls 
    (0.9014550487, 3.9884790531) and (0.9050279556, 3.9864244048) .. 
    (0.9087821815, 3.9869456075).. controls (0.9124834906, 3.9875575887) and 
    (0.9152331465, 3.9907375629) .. (0.9146590535, 3.9939252629).. controls 
    (0.9140547452, 3.9972716071) and (0.9106102142, 3.9993640115) .. 
    (0.9067729356, 3.9987597296) -- cycle;
  
  
  
      \end{scope}
      \begin{scope}[shift={(0.8731702173, -3.5881934417)}]
        \path[fill=c005093,even odd rule] (0.9614155058, 3.9987211798).. controls 
    (0.9498657025, 3.9979356319) and (0.9399325298, 4.005081525) .. (0.9392829513,
     4.0147050294).. controls (0.9385804825, 4.0243361273) and (0.9474258208, 
    4.0328793908) .. (0.9589755977, 4.0337480973).. controls (0.9705706977, 
    4.0345336452) and (0.9804586531, 4.0273877785) .. (0.9811535813, 
    4.0176283048).. controls (0.9818483771, 4.0080500971) and (0.9730030388, 
    3.9995898862) .. (0.9614155058, 3.9987211798) -- cycle(0.9629866017, 
    4.0219415687).. controls (0.9590585975, 4.0217669438) and (0.9560974867, 
    4.018844515) .. (0.9564525575, 4.0155887906).. controls (0.9566271825, 
    4.0122651477) and (0.9599347652, 4.0098781298) .. (0.9637645825, 
    4.0102180665).. controls (0.967601861, 4.0103926915) and (0.9705706183, 
    4.013315226) .. (0.9702155475, 4.0165784117).. controls (0.9700409225, 
    4.0198038412) and (0.9668238802, 4.0223343692) .. (0.9629866017, 4.0219415687)
     -- cycle;
  
  
  
      \end{scope}
      \begin{scope}[shift={(0.8731702173, -3.5862749479)}]
        \path[fill=c005093,even odd rule] (1.018891499, 4.0131792831).. controls 
    (1.0075155269, 4.0111322019) and (0.9969250763, 4.0173263094) .. 
    (0.9952254987, 4.0269045171).. controls (0.9935713238, 4.0363240542) and 
    (1.0015480879, 4.0458116421) .. (1.0130071392, 4.0476849979).. controls 
    (1.0244661904, 4.049649) and (1.0351019642, 4.0435153498) .. (1.0366655723, 
    4.033967331).. controls (1.0383197473, 4.024464609) and (1.0304337087, 
    4.0151432852) .. (1.018891499, 4.0131792831) -- cycle(1.018105951, 
    4.0363996456).. controls (1.0142686725, 4.0358408456) and (1.0117456058, 
    4.0326151515) .. (1.0123498877, 4.029397289).. controls (1.0129163606, 
    4.0262096154) and (1.0164516967, 4.0239888352) .. (1.0201530323, 
    4.0246459808).. controls (1.0238996646, 4.0253410412) and (1.0265585948, 
    4.0285210154) .. (1.0259467458, 4.03182963).. controls (1.0253802729, 
    4.0349719804) and (1.0219357683, 4.0371021408) .. (1.018105951, 4.0363996456) 
    -- cycle;
  
  
  
      \end{scope}
      \begin{scope}[shift={(0.8731702173, -3.5863580271)}]
        \path[fill=c005093,even odd rule] (1.4121516821, 3.9753952221).. controls 
    (1.4006397406, 3.9773592242) and (1.3925799238, 3.9865596598) .. 
    (1.3942795013, 3.9961831906).. controls (1.3958431888, 4.0058142621) and 
    (1.4064410213, 4.0118798879) .. (1.4178169933, 4.0100368004).. controls 
    (1.4294045263, 4.0080426358) and (1.4373360467, 3.9988420944) .. 
    (1.4357648715, 3.9891279437).. controls (1.4341106965, 3.9795798985) and 
    (1.423482331, 3.9734463542) .. (1.4121516821, 3.9753952221) -- 
    cycle(1.4128012342, 3.9987590152).. controls (1.4089639556, 3.9993633235) and 
    (1.4055723942, 3.9971802465) .. (1.4050436244, 3.9939245485).. controls 
    (1.4044771515, 3.990736875) and (1.4070454619, 3.9875567421) .. (1.4108373379,
     3.9869449196).. controls (1.4145839702, 3.9864236904) and (1.418119174, 
    3.9884784446) .. (1.4187309965, 3.9916509046).. controls (1.4191614735, 
    3.9949216573) and (1.4166309985, 3.9980565465) .. (1.4128012342, 3.9987590152)
     -- cycle;
  
  
  
      \end{scope}
      \begin{scope}[shift={(0.8731702173, -3.5880116729)}]
        \path[fill=c005093,even odd rule] (1.36073261, 4.0025351485).. controls 
    (1.34913751, 4.0033283694) and (1.340337389, 4.0118716065) .. (1.3409038883, 
    4.0215857835).. controls (1.3417272717, 4.0311715583) and (1.351713361, 
    4.0383627481) .. (1.3632555708, 4.0374864746).. controls (1.3748431038, 
    4.036708335) and (1.3835676333, 4.0281651244) .. (1.3829104877, 
    4.0184509738).. controls (1.3822985065, 4.0088198758) and (1.3722748198, 
    4.0015833629) .. (1.36073261, 4.0025351485) -- cycle(1.3592520017, 
    4.0258989417).. controls (1.3554600463, 4.0260735667) and (1.3520608119, 
    4.0236327854) .. (1.3518795458, 4.0204074087).. controls (1.3517128583, 
    4.0170534975) and (1.3546216346, 4.0141377098) .. (1.358375834, 
    4.0139111206).. controls (1.362296271, 4.0136089665) and (1.3655218065, 
    4.0160941977) .. (1.3657408285, 4.0193195744).. controls (1.3660429827, 
    4.0226735121) and (1.363081819, 4.0255061148) .. (1.3592520017, 4.0258989417) 
    -- cycle;
  
  
  
      \end{scope}
      \begin{scope}[shift={(0.8731702173, -3.5862297042)}]
        \path[fill=c005093,even odd rule] (1.3026511179, 4.015098359).. controls 
    (1.2911844996, 4.0171454402) and (1.2833438106, 4.0264742517) .. (1.284998065,
     4.0359693802).. controls (1.2865691608, 4.0456004781) and (1.2971520179, 
    4.0516963723) .. (1.3085733396, 4.0497399373).. controls (1.3201155494, 
    4.0476399394) and (1.3279563177, 4.0384092883) .. (1.3263927096, 
    4.0288612431).. controls (1.3247308617, 4.0193207915) and (1.3140573848, 
    4.0131343569) .. (1.3026511179, 4.015098359) -- cycle(1.3034745012, 
    4.0384092883).. controls (1.299727869, 4.0392326717) and (1.2961472627, 
    4.0370572675) .. (1.29562606, 4.033831785).. controls (1.2951125038, 
    4.0305231704) and (1.2977638933, 4.0273429581) .. (1.3014652025, 
    4.0266933796).. controls (1.3052118348, 4.02607399) and (1.308739445, 
    4.0281287442) .. (1.309351294, 4.031399444).. controls (1.3099177669, 
    4.0347004915) and (1.3073118856, 4.0377596833) .. (1.3034745012, 4.0384092883)
     -- cycle;
  
  
  
      \end{scope}
      \begin{scope}[shift={(0.8731702173, -3.5863201917)}]
        \path[fill=c005093,even odd rule] (1.4257463177, 4.1785917031).. controls 
    (1.4142796994, 4.1804649531) and (1.4063028558, 4.1897863033) .. 
    (1.4079193542, 4.1992890254).. controls (1.4095735292, 4.208784154) and 
    (1.4200808742, 4.2151066373) .. (1.4315777608, 4.2131049056).. controls 
    (1.4430443792, 4.2111861473) and (1.4509683325, 4.2019554963) .. 
    (1.4493971573, 4.1923621544).. controls (1.4478260615, 4.1829501844) and 
    (1.4371978813, 4.1766352681) .. (1.4257463177, 4.1785917031) -- 
    cycle(1.4265241927, 4.2018649558).. controls (1.4226869142, 4.2025597517) and 
    (1.419159304, 4.2003692663) .. (1.418728721, 4.1971588915).. controls 
    (1.4181170044, 4.1938805452) and (1.4207681294, 4.1907079) .. (1.4246055138, 
    4.190088484).. controls (1.4283068229, 4.1893936881) and (1.4318420267, 
    4.191622009) .. (1.4323179063, 4.194885221).. controls (1.4328843792, 
    4.1981182177) and (1.430270825, 4.2011624871) .. (1.4265241927, 4.2018649558) 
    -- cycle;
  
  
  
      \end{scope}
      \begin{scope}[shift={(0.8731702173, -3.5882312771)}]
        \path[fill=c005093,even odd rule] (1.3707042531, 4.1991538762).. controls 
    (1.3591620433, 4.2000301763) and (1.3503618958, 4.2085582792) .. 
    (1.3509208546, 4.218151621).. controls (1.3517140754, 4.227820475) and 
    (1.361601925, 4.2349285063) .. (1.3732801042, 4.2341807144).. controls 
    (1.3848676371, 4.233266579) and (1.3937131871, 4.2248593642) .. (1.3929804235,
     4.2152282662).. controls (1.3921872027, 4.2055141156) and (1.3822993531, 
    4.1983682225) .. (1.3707042531, 4.1991538762) -- cycle(1.3692688885, 
    4.2225101023).. controls (1.36543161, 4.2226847273) and (1.3620852394, 
    4.2203270252) .. (1.3619039733, 4.2170184635).. controls (1.3616843692, 
    4.2137477108) and (1.3645629829, 4.2108243825) .. (1.3684002615, 
    4.2105221754).. controls (1.3722754019, 4.2103105088) and (1.3756368008, 
    4.2127885963) .. (1.3757198535, 4.2160214871).. controls (1.3760220077, 
    4.2193225081) and (1.3731062994, 4.2222004604) .. (1.3692688885, 4.2225101023)
     -- cycle;
  
  
  
      \end{scope}
      \begin{scope}[shift={(0.8731702173, -3.5863201917)}]
        \path[fill=c005093,even odd rule] (1.3107037646, 4.2117458733).. controls 
    (1.2993353596, 4.2137098754) and (1.2914414894, 4.2231218454) .. 
    (1.2930579877, 4.2325338154).. controls (1.2946290835, 4.2421724804) and 
    (1.3053027721, 4.2482608075) .. (1.3166787442, 4.2463874517).. controls 
    (1.3281756573, 4.2443403704) and (1.3361977181, 4.2349737235) .. 
    (1.3345358967, 4.2254785685).. controls (1.33288913, 4.21598344) and 
    (1.3222988646, 4.2097894383) .. (1.3107037646, 4.2117458733) -- 
    cycle(1.3116176354, 4.2350567763).. controls (1.30782568, 4.2358044888) and 
    (1.3043434723, 4.2336290846) .. (1.3037316498, 4.2303961144).. controls 
    (1.3032556644, 4.227223575) and (1.3057861394, 4.223990446) .. (1.3095328246, 
    4.2233408675).. controls (1.3133701031, 4.2227288862) and (1.316844929, 
    4.22468548) .. (1.317464345, 4.2280544196).. controls (1.3179855742, 
    4.2312648208) and (1.3154095908, 4.2344071977) .. (1.3116176354, 4.2350567763)
     -- cycle;
  
  
  
      \end{scope}
      \begin{scope}[shift={(0.8731702173, -3.5863201917)}]
        \path[fill=c005093,even odd rule] (1.4363295717, 4.3745820454).. controls 
    (1.4248629533, 4.3763269725) and (1.4169314329, 4.3857465625) .. 
    (1.4185026346, 4.3953247702).. controls (1.4200814033, 4.4048728154) and 
    (1.4307548802, 4.4110516829) .. (1.4420402325, 4.4090952215).. controls 
    (1.4535371192, 4.4072219715) and (1.46155918, 4.3979384831) .. (1.4599955719, 
    4.388398005).. controls (1.4582884802, 4.3789028765) and (1.4477508669, 
    4.3725802873) .. (1.4363295717, 4.3745820454) -- cycle(1.4370243675, 
    4.3979383773).. controls (1.4333154648, 4.3985125231) and (1.4296594523, 
    4.396329446) .. (1.4293119485, 4.3932398535).. controls (1.4286999673, 
    4.3898331579) and (1.4313667027, 4.3867361042) .. (1.435067906, 
    4.3861318223).. controls (1.4388145648, 4.3854822702) and (1.4423422808, 
    4.3875290869) .. (1.4429088067, 4.3908378073).. controls (1.4435207879, 
    4.3940632369) and (1.4408540525, 4.3973340954) .. (1.4370243675, 4.3979383773)
     -- cycle;
  
  
  
      \end{scope}
      \begin{scope}[shift={(0.8731702173, -3.5881556063)}]
        \path[fill=c005093,even odd rule] (1.3830843719, 4.3995547698).. controls 
    (1.3714892719, 4.4005593927) and (1.3628175267, 4.4089667398) .. 
    (1.3634746988, 4.418688484).. controls (1.3640790071, 4.4284102017) and 
    (1.3740651229, 4.4354125583) .. (1.3856073327, 4.4346344981).. controls 
    (1.3972477558, 4.433894194) and (1.4060477975, 4.4252226075) .. (1.4053906254,
     4.4156745887).. controls (1.4046579942, 4.4060133283) and (1.394762551, 
    4.3987767096) .. (1.3830843719, 4.3995547698) -- cycle(1.381732351, 
    4.4229185894).. controls (1.377811914, 4.4233038227) and (1.374450515, 
    4.4207355123) .. (1.374276816, 4.4175931883).. controls (1.3740651494, 
    4.4141561979) and (1.377026366, 4.4113234894) .. (1.3807352688, 4.411013821)..
     controls (1.3846934619, 4.4107651127) and (1.3878810296, 4.4132873856) .. 
    (1.3881832367, 4.416512921).. controls (1.3883949033, 4.4197761067) and 
    (1.3854412538, 4.4226995673) .. (1.381732351, 4.4229185894) -- cycle;
  
  
  
      \end{scope}
      \begin{scope}[shift={(0.8731702173, -3.5863654354)}]
        \path[fill=c005093,even odd rule] (1.3231649517, 4.4122763598).. controls 
    (1.3116606038, 4.4142403619) and (1.3038650263, 4.4235691733) .. 
    (1.3055117135, 4.4330643283).. controls (1.3071735615, 4.4426953998) and 
    (1.3177639592, 4.4487459973) .. (1.3291323642, 4.4467895623).. controls 
    (1.3405914154, 4.444870804) and (1.348568259, 4.4354210515) .. (1.3469140046, 
    4.4260090815).. controls (1.3452143213, 4.4164233067) and (1.3346240029, 
    4.4103199248) .. (1.3231649517, 4.4122763598) -- cycle(1.3240412517, 
    4.43550421).. controls (1.3202492962, 4.4362897579) and (1.3166762571, 
    4.43406911) .. (1.3161928104, 4.4309266008).. controls (1.3155808292, 
    4.4276709029) and (1.3182398917, 4.4244378798) .. (1.3220772496, 
    4.4238336244).. controls (1.3258239083, 4.423297314) and (1.3293061954, 
    4.4252160723) .. (1.3299180444, 4.4285396094).. controls (1.3304392735, 
    4.4317953073) and (1.327787884, 4.4350283304) .. (1.3240412517, 4.43550421) --
     cycle;
  
  
  
      \end{scope}
      \begin{scope}[shift={(0.8731702173, -3.5863127833)}]
        \path[fill=c005093,even odd rule] (0.8906905552, 4.177625286).. controls 
    (0.8792389652, 4.1756234485) and (0.8687316996, 4.18193855) .. (0.8671680915, 
    4.1914412721).. controls (0.8655062435, 4.2009817238) and (0.8734376846, 
    4.2103936938) .. (0.884942059, 4.2122292142).. controls (0.8964010837, 
    4.2141027288) and (0.9069915344, 4.2079009219) .. (0.9085627096, 
    4.198413334).. controls (0.9101338054, 4.1889106383) and (0.9022024702, 
    4.1795439915) .. (0.8906905552, 4.177625286) -- cycle(0.8899124156, 
    4.2009817238).. controls (0.8862111065, 4.2003623342) and (0.8836351496, 
    4.1971518271) .. (0.8841563523, 4.1938734808).. controls (0.8847228252, 
    4.1907009415) and (0.8881296796, 4.1885632404) .. (0.89200482, 4.1891674958)..
     controls (0.8957061292, 4.1898699646) and (0.8983650594, 4.1928688315) .. 
    (0.8978891533, 4.1962830942).. controls (0.8972774367, 4.1994934954) and 
    (0.8938328527, 4.2016313023) .. (0.8899124156, 4.2009817238) -- cycle;
  
  
  
      \end{scope}
      \begin{scope}[shift={(0.8731702173, -3.5873174063)}]
        \path[fill=c005093,even odd rule] (0.9472987415, 4.197453796).. controls 
    (0.9357489646, 4.1960562669) and (0.9255136113, 4.20277163) .. (0.9243352365, 
    4.2124404575).. controls (0.9232853698, 4.2220715554) and (0.931700231, 
    4.2310076458) .. (0.943287764, 4.2322616121).. controls (0.9547921383, 
    4.2335305538) and (0.9650804877, 4.2269060221) .. (0.9661228931, 
    4.2172371946).. controls (0.9672105952, 4.2076438527) and (0.9587577927, 
    4.1987153294) .. (0.9472987415, 4.197453796) -- cycle(0.9478199706, 
    4.2208100221).. controls (0.9438617775, 4.22028112) and (0.9411121217, 
    4.2172371946) .. (0.9415049485, 4.21392858).. controls (0.9418978548, 
    4.2107862296) and (0.9452516073, 4.2084370471) .. (0.9490813981, 
    4.2088223069).. controls (0.9529186767, 4.2093057006) and (0.9557060885, 
    4.2123574048) .. (0.9553585848, 4.215711316).. controls (0.9549656785, 
    4.2188083698) and (0.9515288733, 4.2212859281) .. (0.9478199706, 4.2208100221)
     -- cycle;
  
  
  
      \end{scope}
      \begin{scope}[shift={(0.8731702173, -3.5863201917)}]
        \path[fill=c005093,even odd rule] (1.008604446, 4.2108768494).. controls 
    (0.997236041, 4.2087845244) and (0.9866002938, 4.2149784731) .. (0.9849460394,
     4.2245189248).. controls (0.9832918644, 4.2340971325) and (0.9912609556, 
    4.2433430765) .. (1.002727574, 4.2453825906).. controls (1.0141866252, 
    4.2472558406) and (1.0248149113, 4.2411675135) .. (1.0263407898, 
    4.2316648179).. controls (1.0279949648, 4.2220714496) and (1.0201088204, 
    4.2128408515) .. (1.008604446, 4.2108768494) -- cycle(1.0078188981, 
    4.2340971325).. controls (1.0041175625, 4.2335380679) and (1.0014662258, 
    4.230267421) .. (1.0020704812, 4.2270495585).. controls (1.0025086313, 
    4.2238165619) and (1.0060361621, 4.2217316188) .. (1.0098734406, 
    4.2223434677).. controls (1.0136200729, 4.2230308552) and (1.01623368, 
    4.2262639048) .. (1.0156671542, 4.2293911475).. controls (1.0150551729, 
    4.2327601135) and (1.0115277744, 4.2347996012) .. (1.0078188981, 4.2340971325)
     -- cycle;
  
  
  
      \end{scope}
      \begin{scope}[shift={(0.8731702173, -3.5863201917)}]
        \path[fill=c005093,even odd rule] (0.8762698904, 4.3708354925).. controls 
    (0.8647276806, 4.3688791633) and (0.854220415, 4.3751561383) .. (0.8525284046,
     4.3846512669).. controls (0.8510025525, 4.3941917185) and (0.8590171256, 
    4.4035206358) .. (0.8704308538, 4.4054770708).. controls (0.881889905, 
    4.4073505854) and (0.8924350325, 4.4011714798) .. (0.8940515308, 
    4.391623461).. controls (0.8957057058, 4.382120739) and (0.8876911856, 
    4.3727994946) .. (0.8762698904, 4.3708354925) -- cycle(0.8755295863, 
    4.3941917185).. controls (0.871699769, 4.3936329185) and (0.8691239444, 
    4.3904450863) .. (0.8696451471, 4.3871364717).. controls (0.8700832971, 
    4.383865719) and (0.873746956, 4.3817354792) .. (0.8774482917, 4.3824304073)..
     controls (0.8812855702, 4.3829892073) and (0.883808531, 4.3861318223) .. 
    (0.8833779481, 4.3894932213).. controls (0.8827659669, 4.3927110838) and 
    (0.879238489, 4.3947658115) .. (0.8755295863, 4.3941917185) -- cycle;
  
  
  
      \end{scope}
      \begin{scope}[shift={(0.8731702173, -3.5882312771)}]
        \path[fill=c005093,even odd rule] (0.93457776, 4.3917449842).. controls 
    (0.9228996073, 4.3910120883) and (0.9130041642, 4.398150626) .. (0.9122639394,
     4.407781724).. controls (0.9116143873, 4.417458145) and (0.920414376, 
    4.4260391381) .. (0.9319566123, 4.4268247919).. controls (0.943634765, 
    4.4275725044) and (0.9535756106, 4.4204191235) .. (0.9542705123, 
    4.4108409158).. controls (0.9549655727, 4.4011191981) and (0.9461199698, 
    4.3925306115) .. (0.93457776, 4.3917449842) -- cycle(0.9359223725, 
    4.415018131).. controls (0.9321757137, 4.4149281727) and (0.929259926, 
    4.4119665062) .. (0.9294336515, 4.4086578917).. controls (0.9297358056, 
    4.4053418158) and (0.9329613675, 4.4029472573) .. (0.936798646, 
    4.4031662794).. controls (0.9406284369, 4.4035591856) and (0.9435894419, 
    4.4063011685) .. (0.9432872612, 4.4096625675).. controls (0.9431126362, 
    4.4129635885) and (0.9397143279, 4.4153277994) .. (0.9359223725, 4.415018131) 
    -- cycle;
  
  
  
      \end{scope}
      \begin{scope}[shift={(0.8731702173, -3.58628315)}]
        \path[fill=c005093,even odd rule] (0.9948786035, 4.4052208219).. controls 
    (0.9832910706, 4.403271901) and (0.9727007523, 4.4094509008) .. (0.9710464979,
     4.418998946).. controls (0.9694298937, 4.4284562392) and (0.9774068167, 
    4.4379136381) .. (0.9888658679, 4.4397416708).. controls (1.000279596, 
    4.4417810792) and (1.0109607723, 4.4356400735) .. (1.0125244069, 
    4.4260014085).. controls (1.0141862548, 4.4165138206) and (1.0061639558, 
    4.4073131998) .. (0.9948786035, 4.4052208219) -- cycle(0.9939647327, 
    4.428456345).. controls (0.9901727773, 4.4278747908) and (0.9875136619, 
    4.4247549035) .. (0.9881255108, 4.4214312606).. controls (0.9886919838, 
    4.4182964508) and (0.9922197527, 4.4161210202) .. (0.9960570313, 
    4.4168159219).. controls (0.9998036635, 4.4175107177) and (1.0024172706, 
    4.4206532004) .. (1.0018507448, 4.4238710629).. controls (1.0012387635, 
    4.4270964925) and (0.997711365, 4.4292797283) .. (0.9939647327, 4.428456345) 
    -- cycle;
  
  
  
      \end{scope}
      \begin{scope}[shift={(0.8731702173, -3.614556525)}]
        \path[fill=c005093,even odd rule] (0.8862944238, 3.9662101323).. controls 
    (0.8862944238, 3.9689974913) and (0.8840282675, 3.971263621) .. (0.8811956119,
     3.971263621).. controls (0.878453629, 3.971263621) and (0.8761874463, 
    3.9689974648) .. (0.8761874463, 3.9662101323).. controls (0.8761874463, 
    3.963294371) and (0.8784536025, 3.9610282677) .. (0.8811956119, 
    3.9610282677).. controls (0.8840283204, 3.9610282677) and (0.8862944238, 
    3.963294424) .. (0.8862944238, 3.9662101323) -- cycle;
  
  
  
      \end{scope}
      \begin{scope}[shift={(0.8731702173, -3.6145943604)}]
        \path[fill=c005093,even odd rule] (0.8695159246, 3.9825414063).. controls 
    (0.8695159246, 3.985419385) and (0.8672497683, 3.9875646267) .. (0.864507759, 
    3.9875646267).. controls (0.861675024, 3.9875646267) and (0.8594089471, 
    3.985419385) .. (0.8594089471, 3.9825414063).. controls (0.8594089471, 
    3.9797087506) and (0.8616751033, 3.9773670029) .. (0.864507759, 
    3.9773670029).. controls (0.8672497419, 3.9773670029) and (0.8695159246, 
    3.97970883) .. (0.8695159246, 3.9825414063) -- cycle;
  
  
  
      \end{scope}
      \begin{scope}[shift={(0.8731702173, -3.6145943604)}]
        \path[fill=c005093,even odd rule] (0.8667739417, 4.0150073158).. controls 
    (0.8667739417, 4.0177492988) and (0.8645077854, 4.0201061277) .. 
    (0.8616751298, 4.0201061277).. controls (0.8589255798, 4.0201061277) and 
    (0.8565763179, 4.0177492194) .. (0.8565763179, 4.0150073158).. controls 
    (0.8565763179, 4.0121746073) and (0.8589255533, 4.0099085304) .. 
    (0.8616751298, 4.0099085304).. controls (0.8645078648, 4.0099085304) and 
    (0.8667739417, 4.0121746867) .. (0.8667739417, 4.0150073158) -- cycle;
  
  
  
      \end{scope}
      \begin{scope}[shift={(0.8731702173, -3.6145943604)}]
        \path[fill=c005093,even odd rule] (0.8649002683, 4.046424179).. controls 
    (0.8649002683, 4.0491284058) and (0.8626719475, 4.051394615) .. (0.8598845356,
     4.051394615).. controls (0.8570518006, 4.051394615) and (0.8547026446, 
    4.0491284588) .. (0.8547026446, 4.046424179).. controls (0.8547026446, 
    4.0435537144) and (0.85705188, 4.0411969912) .. (0.8598845356, 4.0411969912)..
     controls (0.8626718152, 4.0411969912) and (0.8649002683, 4.043553635) .. 
    (0.8649002683, 4.046424179) -- cycle;
  
  
  
      \end{scope}
      \begin{scope}[shift={(0.8731702173, -3.6146316667)}]
        \path[fill=c005093,even odd rule] (1.061768604, 4.0379254183).. controls 
    (1.061768604, 4.0407581269) and (1.0595024477, 4.0429411775) .. (1.0567528713,
     4.0429411775).. controls (1.0539201363, 4.0429411775) and (1.0516616265, 
    4.0407581004) .. (1.0516616265, 4.0379254183).. controls (1.0516616265, 
    4.0350473867) and (1.0539201098, 4.0327812833) .. (1.0567528713, 
    4.0327812833).. controls (1.0595024213, 4.0327812833) and (1.061768604, 
    4.0350474396) .. (1.061768604, 4.0379254183) -- cycle;
  
  
  
      \end{scope}
      \begin{scope}[shift={(0.8731702173, -3.6146396042)}]
        \path[fill=c005093,even odd rule] (1.0588905723, 4.0686019508).. controls 
    (1.0588905723, 4.0713439337) and (1.0567530035, 4.0736630067) .. 
    (1.0539202421, 4.0736630067).. controls (1.0510951006, 4.0736630067) and 
    (1.0488289973, 4.0713439338) .. (1.0488289973, 4.0686019508).. controls 
    (1.0488289973, 4.0657692423) and (1.0510951535, 4.063510706) .. (1.0539202421,
     4.063510706).. controls (1.0567529506, 4.063510706) and (1.0588905723, 
    4.0657691894) .. (1.0588905723, 4.0686019508) -- cycle;
  
  
  
      \end{scope}
      \begin{scope}[shift={(0.8731702173, -3.6147679271)}]
        \path[fill=c005093,even odd rule] (1.056191981, 4.1001084283).. controls 
    (1.056191981, 4.1028957079) and (1.0539184165, 4.1050334354) .. (1.0510931956,
     4.1050334354).. controls (1.0483436456, 4.1050334354) and (1.0460774365, 
    4.1028956021) .. (1.0460774365, 4.1001084283).. controls (1.0460774365, 
    4.0971926406) and (1.0483435927, 4.0950096165) .. (1.0510931956, 
    4.0950096165).. controls (1.0539183371, 4.0950096165) and (1.056191981, 
    4.0971926935) .. (1.056191981, 4.1001084283) -- cycle;
  
  
  
      \end{scope}
      \begin{scope}[shift={(0.8731702173, -3.6146316667)}]
        \path[fill=c005093,even odd rule] (1.2598960388, 4.0379254183).. controls 
    (1.2598960388, 4.0407581269) and (1.262162195, 4.0429411775) .. (1.2649948506,
     4.0429411775).. controls (1.2678275592, 4.0429411775) and (1.2700936625, 
    4.0407581004) .. (1.2700936625, 4.0379254183).. controls (1.2700936625, 
    4.0350473867) and (1.2678275063, 4.0327812833) .. (1.2649948506, 
    4.0327812833).. controls (1.2621621421, 4.0327812833) and (1.2598960388, 
    4.0350474396) .. (1.2598960388, 4.0379254183) -- cycle;
  
  
  
      \end{scope}
      \begin{scope}[shift={(0.8731702173, -3.6146396042)}]
        \path[fill=c005093,even odd rule] (1.2628579433, 4.0686019508).. controls 
    (1.2628579433, 4.0713439337) and (1.2649955121, 4.0736630067) .. 
    (1.2678282735, 4.0736630067).. controls (1.2705778235, 4.0736630067) and 
    (1.2728364392, 4.0713439338) .. (1.2728364392, 4.0686019508).. controls 
    (1.2728364392, 4.0657692423) and (1.2705776912, 4.063510706) .. (1.2678282735,
     4.063510706).. controls (1.264995565, 4.063510706) and (1.2628579433, 
    4.0657691894) .. (1.2628579433, 4.0686019508) -- cycle;
  
  
  
      \end{scope}
      \begin{scope}[shift={(0.8731702173, -3.6147679271)}]
        \path[fill=c005093,even odd rule] (1.2673501715, 4.1001084283).. controls 
    (1.2673501715, 4.1028957079) and (1.2696994069, 4.1050334354) .. 
    (1.2724489833, 4.1050334354).. controls (1.2751909663, 4.1050334354) and 
    (1.2775403075, 4.1028956021) .. (1.2775403075, 4.1001084283).. controls 
    (1.2775403075, 4.0971926406) and (1.2751910721, 4.0950096165) .. 
    (1.2724489833, 4.0950096165).. controls (1.2696994333, 4.0950096165) and 
    (1.2673501715, 4.0971926935) .. (1.2673501715, 4.1001084283) -- cycle;
  
  
  
      \end{scope}
      \begin{scope}[shift={(0.8731702173, -3.6145943604)}]
        \path[fill=c005093,even odd rule] (1.4478260615, 3.9825414063).. controls 
    (1.4478260615, 3.985419385) and (1.4501827052, 3.9875646267) .. (1.4529247675,
     3.9875646267).. controls (1.4557575025, 3.9875646267) and (1.4580312523, 
    3.985419385) .. (1.4580312523, 3.9825414063).. controls (1.4580312523, 
    3.9797087506) and (1.4557574231, 3.9773670029) .. (1.4529247675, 
    3.9773670029).. controls (1.450182811, 3.9773670029) and (1.4478260615, 
    3.97970883) .. (1.4478260615, 3.9825414063) -- cycle;
  
  
  
      \end{scope}
      \begin{scope}[shift={(0.8731702173, -3.6146848479)}]
        \path[fill=c005093,even odd rule] (1.4506585848, 4.0142291498).. controls 
    (1.4506585848, 4.0169711327) and (1.4529244765, 4.0192373154) .. 
    (1.4557572908, 4.0192373154).. controls (1.4586353225, 4.0192373154) and 
    (1.46077305, 4.0169711592) .. (1.46077305, 4.0142291498).. controls 
    (1.46077305, 4.0113964412) and (1.4586352167, 4.0091303379) .. (1.4557572908, 
    4.0091303379).. controls (1.4529245823, 4.0091303379) and (1.4506585848, 
    4.0113964942) .. (1.4506585848, 4.0142291498) -- cycle;
  
  
  
      \end{scope}
      \begin{scope}[shift={(0.8731702173, -3.6145943604)}]
        \path[fill=c005093,even odd rule] (1.4525322581, 4.046424179).. controls 
    (1.4525322581, 4.0491284058) and (1.4548889019, 4.051394615) .. (1.4576309642,
     4.051394615).. controls (1.4603805142, 4.051394615) and (1.4626014002, 
    4.0491284588) .. (1.4626014002, 4.046424179).. controls (1.4626014002, 
    4.0435537144) and (1.4603804877, 4.0411969912) .. (1.4576309642, 
    4.0411969912).. controls (1.4548889813, 4.0411969912) and (1.4525322581, 
    4.043553635) .. (1.4525322581, 4.046424179) -- cycle;
  
  
  
      \end{scope}
      \begin{scope}[shift={(0.8731702173, -3.6146017688)}]
        \path[fill=c005093,even odd rule] (0.9186248931, 3.9791119829).. controls 
    (0.9186248931, 3.981906909) and (0.9163209015, 3.9842032277) .. (0.9135260813,
     3.9842032277).. controls (0.9107387752, 3.9842032277) and (0.9084348365, 
    3.981906909) .. (0.9084348365, 3.9791119829).. controls (0.9084348365, 
    3.9762793273) and (0.9107388281, 3.974013171) .. (0.9135260813, 3.974013171)..
     controls (0.9163209544, 3.974013171) and (0.9186248931, 3.9762793273) .. 
    (0.9186248931, 3.9791119829) -- cycle;
  
  
  
      \end{scope}
      \begin{scope}[shift={(0.8731702173, -3.6146769104)}]
        \path[fill=c005093,even odd rule] (0.954271015, 3.9883342202).. controls 
    (0.954271015, 3.9912575221) and (0.9520048587, 3.9934330321) .. (0.9491722296,
     3.9934330321).. controls (0.9464226796, 3.9934330321) and (0.9441640375, 
    3.9912576279) .. (0.9441640375, 3.9883342202).. controls (0.9441640375, 
    3.9855921844) and (0.9464227854, 3.983318461) .. (0.9491722296, 3.983318461)..
     controls (0.9520049381, 3.983318461) and (0.954271015, 3.9855922902) .. 
    (0.954271015, 3.9883342202) -- cycle;
  
  
  
      \end{scope}
      \begin{scope}[shift={(0.8731702173, -3.6146017688)}]
        \path[fill=c005093,even odd rule] (0.9896049815, 3.9969765967).. controls 
    (0.9896049815, 3.9997261996) and (0.9872105023, 4.0020679208) .. (0.98446844, 
    4.0020679208).. controls (0.98171889, 4.0020679208) and (0.9793696281, 
    3.9997260938) .. (0.9793696281, 3.9969765967).. controls (0.9793696281, 
    3.9942345608) and (0.9817188635, 3.9918777848) .. (0.98446844, 3.9918777848)..
     controls (0.9872104229, 3.9918777848) and (0.9896049815, 3.9942346931) .. 
    (0.9896049815, 3.9969765967) -- cycle;
  
  
  
      \end{scope}
      \begin{scope}[shift={(0.8731702173, -3.6146840542)}]
        \path[fill=c005093,even odd rule] (1.0244665873, 4.0057764267).. controls 
    (1.0244665873, 4.0085184096) and (1.0222456748, 4.0107845923) .. 
    (1.0194961777, 4.0107845923).. controls (1.0166634692, 4.0107845923) and 
    (1.01426899, 4.008518436) .. (1.01426899, 4.0057764267).. controls 
    (1.01426899, 4.0029437181) and (1.0166634692, 4.0006776148) .. (1.0194961777, 
    4.0006776148).. controls (1.0222457277, 4.0006776148) and (1.0244665873, 
    4.002943771) .. (1.0244665873, 4.0057764267) -- cycle;
  
  
  
      \end{scope}
      \begin{scope}[shift={(0.8731702173, -3.6147679271)}]
        \path[fill=c005093,even odd rule] (1.0641164635, 4.0127417417).. controls 
    (1.0641164635, 4.0155744502) and (1.0617672281, 4.0177574744) .. 
    (1.0590252188, 4.0177574744).. controls (1.0561924838, 4.0177574744) and 
    (1.0539189456, 4.0155743973) .. (1.0539189456, 4.0127417417).. controls 
    (1.0539189456, 4.0099090331) and (1.0561925102, 4.007733576) .. (1.0590252188,
     4.007733576).. controls (1.0617672017, 4.007733576) and (1.0641164635, 
    4.0099089802) .. (1.0641164635, 4.0127417417) -- cycle;
  
  
  
      \end{scope}
      \begin{scope}[shift={(0.8731702173, -3.6156061271)}]
        \path[fill=c005093,even odd rule] (0.8550059629, 4.190534201).. controls 
    (0.8550059629, 4.1931855377) and (0.8529136379, 4.1951041373) .. 
    (0.8503905977, 4.1951041373).. controls (0.8478222873, 4.1951041373) and 
    (0.8456846127, 4.1931856435) .. (0.8456846127, 4.190534201).. controls 
    (0.8456846127, 4.1879658906) and (0.8478221815, 4.1859187565) .. 
    (0.8503905977, 4.1859187565).. controls (0.8529136644, 4.1859187565) and 
    (0.8550059629, 4.1879658377) .. (0.8550059629, 4.190534201) -- cycle;
  
  
  
      \end{scope}
      \begin{scope}[shift={(0.8731702173, -3.6154632521)}]
        \path[fill=c005093,even odd rule] (0.8534422225, 4.2194963129).. controls 
    (0.8534422225, 4.2220193796) and (0.8513874683, 4.2241192452) .. 
    (0.8488268838, 4.2241192452).. controls (0.8462964088, 4.2241192452) and 
    (0.8441133317, 4.2220192473) .. (0.8441133317, 4.2194963129).. controls 
    (0.8441133317, 4.2169732463) and (0.8462964088, 4.2147903279) .. 
    (0.8488268838, 4.2147903279).. controls (0.8513875213, 4.2147903279) and 
    (0.8534422225, 4.216973405) .. (0.8534422225, 4.2194963129) -- cycle;
  
  
  
      \end{scope}
      \begin{scope}[shift={(0.8731702173, -3.6153719708)}]
        \path[fill=c005093,even odd rule] (0.8524358533, 4.2460778892).. controls 
    (0.8524358533, 4.2486461996) and (0.8504718513, 4.2507837683) .. 
    (0.8478205146, 4.2507837683).. controls (0.8452900396, 4.2507837683) and 
    (0.8432429054, 4.2486461996) .. (0.8432429054, 4.2460778892).. controls 
    (0.8432429054, 4.2435474142) and (0.8452899867, 4.2413642313) .. 
    (0.8478205146, 4.2413642313).. controls (0.8504718513, 4.2413642313) and 
    (0.8524358533, 4.2435473083) .. (0.8524358533, 4.2460778892) -- cycle;
  
  
  
      \end{scope}
      \begin{scope}[shift={(0.8731702173, -3.6153793792)}]
        \path[fill=c005093,even odd rule] (1.4624283098, 4.1923544815).. controls 
    (1.4624283098, 4.1948775481) and (1.4646113869, 4.1970605723) .. 
    (1.4671342948, 4.1970605723).. controls (1.4697102781, 4.1970605723) and 
    (1.4717496335, 4.1948774952) .. (1.4717496335, 4.1923544815).. controls 
    (1.4717496335, 4.189786171) and (1.4697102252, 4.1876484965) .. (1.4671342948,
     4.1876484965).. controls (1.4646112281, 4.1876484965) and (1.4624283098, 
    4.1897860652) .. (1.4624283098, 4.1923544815) -- cycle;
  
  
  
      \end{scope}
      \begin{scope}[shift={(0.8731702173, -3.6154246229)}]
        \path[fill=c005093,even odd rule] (1.4641360365, 4.2214665856).. controls 
    (1.4641360365, 4.2240801398) and (1.4661754448, 4.2261726765) .. 
    (1.4687060785, 4.2261726765).. controls (1.471274389, 4.2261726765) and 
    (1.4734120635, 4.2240800869) .. (1.4734120635, 4.2214665856).. controls 
    (1.4734120635, 4.2189813544) and (1.4712742302, 4.2168060296) .. 
    (1.4687060785, 4.2168060296).. controls (1.4661756035, 4.2168060296) and 
    (1.4641360365, 4.2189814337) .. (1.4641360365, 4.2214665856) -- cycle;
  
  
  
      \end{scope}
      \begin{scope}[shift={(0.8731702173, -3.6154703958)}]
        \path[fill=c005093,even odd rule] (1.4650876104, 4.2483513479).. controls 
    (1.4650876104, 4.2508818229) and (1.4670442042, 4.2529667131) .. 
    (1.4695746792, 4.2529667131).. controls (1.4721506625, 4.2529667131) and 
    (1.4742806642, 4.2508817965) .. (1.4742806642, 4.2483513479).. controls 
    (1.4742806642, 4.245647121) and (1.4721505038, 4.2436453629) .. (1.4695746792,
     4.2436453629).. controls (1.4670442042, 4.2436453629) and (1.4650876104, 
    4.2456472004) .. (1.4650876104, 4.2483513479) -- cycle;
  
  
  
      \end{scope}
      \begin{scope}[shift={(0.8731702173, -3.6154254167)}]
        \path[fill=c005093,even odd rule] (1.2683130431, 4.2409413213).. controls 
    (1.2683130431, 4.2434717963) and (1.2703524515, 4.2456021154) .. 
    (1.2729736256, 4.2456021154).. controls (1.2755871798, 4.2456021154) and 
    (1.2776797165, 4.243471955) .. (1.2776797165, 4.2409413213).. controls 
    (1.2776797165, 4.2382824175) and (1.2755871269, 4.2362353363) .. 
    (1.2729736256, 4.2362353363).. controls (1.2703523985, 4.2362353363) and 
    (1.2683130431, 4.2382824175) .. (1.2683130431, 4.2409413213) -- cycle;
  
  
  
      \end{scope}
      \begin{scope}[shift={(0.8731702173, -3.6155156396)}]
        \path[fill=c005093,even odd rule] (1.2699210748, 4.2699629908).. controls 
    (1.2699210748, 4.272621921) and (1.2721041519, 4.2746689758) .. (1.274498684, 
    4.2746689758).. controls (1.2770669944, 4.2746689758) and (1.2791140227, 
    4.2726218946) .. (1.2791140227, 4.2699629908).. controls (1.2791140227, 
    4.2674399242) and (1.2770669415, 4.2653929752) .. (1.274498684, 
    4.2653929752).. controls (1.2721042048, 4.2653929752) and (1.2699210748, 
    4.2674400565) .. (1.2699210748, 4.2699629908) -- cycle;
  
  
  
      \end{scope}
      \begin{scope}[shift={(0.8731702173, -3.6154703958)}]
        \path[fill=c005093,even odd rule] (1.2707951258, 4.2968090975).. controls 
    (1.2707951258, 4.2993774079) and (1.27297053, 4.3014245685) .. (1.2755011108, 
    4.3014245685).. controls (1.2780694213, 4.3014245685) and (1.280116476, 
    4.2993774873) .. (1.280116476, 4.2968090975).. controls (1.280116476, 
    4.2942333788) and (1.2780693948, 4.2921031125) .. (1.2755011108, 
    4.2921031125).. controls (1.2729706358, 4.2921031125) and (1.2707951258, 
    4.2942332729) .. (1.2707951258, 4.2968090975) -- cycle;
  
  
  
      \end{scope}
      \begin{scope}[shift={(0.8731702173, -3.6155156396)}]
        \path[fill=c005093,even odd rule] (1.0532263458, 4.2420368285).. controls 
    (1.0532263458, 4.2446881652) and (1.0511792646, 4.2466068706) .. 
    (1.0485204667, 4.2466068706).. controls (1.0459521563, 4.2466068706) and 
    (1.0439050221, 4.2446881123) .. (1.0439050221, 4.2420368285).. controls 
    (1.0439050221, 4.2393779248) and (1.0459521033, 4.2373308435) .. 
    (1.0485204667, 4.2373308435).. controls (1.0511793704, 4.2373308435) and 
    (1.0532263458, 4.2393779248) .. (1.0532263458, 4.2420368285) -- cycle;
  
  
  
      \end{scope}
      \begin{scope}[shift={(0.8731702173, -3.6154632521)}]
        \path[fill=c005093,even odd rule] (1.0526569096, 4.2701294402).. controls 
    (1.0526569096, 4.2726977506) and (1.0504812408, 4.274752346) .. (1.0479508188,
     4.274752346).. controls (1.0453825083, 4.274752346) and (1.0432448338, 
    4.2726975919) .. (1.0432448338, 4.2701294402).. controls (1.0432448338, 
    4.2675688027) and (1.0453824025, 4.2654234552) .. (1.0479508188, 
    4.2654234552).. controls (1.0504812938, 4.2654234552) and (1.0526569096, 
    4.2675686969) .. (1.0526569096, 4.2701294402) -- cycle;
  
  
  
      \end{scope}
      \begin{scope}[shift={(0.8731702173, -3.6154703958)}]
        \path[fill=c005093,even odd rule] (1.0516626319, 4.2977681592).. controls 
    (1.0516626319, 4.3003364696) and (1.0496155506, 4.3023836038) .. 
    (1.0470396202, 4.3023836038).. controls (1.044426066, 4.3023836038) and 
    (1.0423337146, 4.3003365225) .. (1.0423337146, 4.2977681592).. controls 
    (1.0423337146, 4.2952376842) and (1.0444260396, 4.2930621742) .. 
    (1.0470396202, 4.2930621742).. controls (1.0496156035, 4.2930621742) and 
    (1.0516626319, 4.2952375783) .. (1.0516626319, 4.2977681592) -- cycle;
  
  
  
      \end{scope}
      \begin{scope}[shift={(0.8731702173, -3.6155981896)}]
        \path[fill=c005093,even odd rule] (0.8368748873, 4.4170647096).. controls 
    (0.8368748873, 4.4197236133) and (0.8348730498, 4.4216801542) .. 
    (0.8322141725, 4.4216801542).. controls (0.8296911058, 4.4216801542) and 
    (0.8275082404, 4.4197235604) .. (0.8275082404, 4.4170647096).. controls 
    (0.8275082404, 4.4144511554) and (0.8296913175, 4.4124872062) .. 
    (0.8322141725, 4.4124872062).. controls (0.8348730763, 4.4124872062) and 
    (0.8368748873, 4.4144512083) .. (0.8368748873, 4.4170647096) -- cycle;
  
  
  
      \end{scope}
      \begin{scope}[shift={(0.8731702173, -3.6155981896)}]
        \path[fill=c005093,even odd rule] (0.8401106621, 4.395407426).. controls 
    (0.8401106621, 4.3980663563) and (0.837927585, 4.4001134375) .. (0.8354046771,
     4.4001134375).. controls (0.8327457733, 4.4001134375) and (0.8306986392, 
    4.3980663563) .. (0.8306986392, 4.395407426).. controls (0.8306986392, 
    4.3928391156) and (0.8327457204, 4.3909203573) .. (0.8354046771, 
    4.3909203573).. controls (0.8379277438, 4.3909203573) and (0.8401106621, 
    4.3928391156) .. (0.8401106621, 4.395407426) -- cycle;
  
  
  
      \end{scope}
      \begin{scope}[shift={(0.8731702173, -3.6154254167)}]
        \path[fill=c005093,even odd rule] (1.0354890231, 4.4794219213).. controls 
    (1.0354890231, 4.4820431483) and (1.0334419419, 4.4841280121) .. 
    (1.0307830381, 4.4841280121).. controls (1.0282147277, 4.4841280121) and 
    (1.0261676729, 4.4820430954) .. (1.0261676729, 4.4794219213).. controls 
    (1.0261676729, 4.4768459379) and (1.0282147542, 4.4747613387) .. 
    (1.0307830381, 4.4747613387).. controls (1.0334419419, 4.4747613387) and 
    (1.0354890231, 4.4768459908) .. (1.0354890231, 4.4794219213) -- cycle;
  
  
  
      \end{scope}
      \begin{scope}[shift={(0.8731702173, -3.6156368188)}]
        \path[fill=c005093,even odd rule] (1.0387148496, 4.4571530004).. controls 
    (1.0387148496, 4.4598043371) and (1.0364486933, 4.4617230425) .. 
    (1.0338803829, 4.4617230425).. controls (1.0313573163, 4.4617230425) and 
    (1.0293027738, 4.4598042842) .. (1.0293027738, 4.4571530004).. controls 
    (1.0293027738, 4.4544865296) and (1.0313575279, 4.45256793) .. (1.0338803829, 
    4.45256793).. controls (1.0364486933, 4.45256793) and (1.0387148496, 
    4.4544864237) .. (1.0387148496, 4.4571530004) -- cycle;
  
  
  
      \end{scope}
      \begin{scope}[shift={(0.8731702173, -3.6155084958)}]
        \path[fill=c005093,even odd rule] (1.2908830333, 4.4818169031).. controls 
    (1.2908830333, 4.4843852135) and (1.2888359521, 4.4865228881) .. 
    (1.2863054242, 4.4865228881).. controls (1.2837371137, 4.4865228881) and 
    (1.281554116, 4.4843850548) .. (1.281554116, 4.4818169031).. controls 
    (1.281554116, 4.4792938365) and (1.2837371931, 4.4772393998) .. (1.2863054242,
     4.4772393998).. controls (1.2888358992, 4.4772393998) and (1.2908830333, 
    4.4792938894) .. (1.2908830333, 4.4818169031) -- cycle;
  
  
  
      \end{scope}
      \begin{scope}[shift={(0.8731702173, -3.6155989833)}]
        \path[fill=c005093,even odd rule] (1.2885244846, 4.4585961173).. controls 
    (1.2885244846, 4.4611644277) and (1.2863414075, 4.463083186) .. (1.2838183938,
     4.463083186).. controls (1.2812879188, 4.463083186) and (1.2791124088, 
    4.4611644277) .. (1.2791124088, 4.4585961173).. controls (1.2791124088, 
    4.456020134) and (1.2812878129, 4.4538901323) .. (1.2838183938, 
    4.4538901323).. controls (1.2863414604, 4.4538901323) and (1.2885244846, 
    4.4560202927) .. (1.2885244846, 4.4585961173) -- cycle;
  
  
  
      \end{scope}
      \begin{scope}[shift={(0.8731702173, -3.6154703958)}]
        \path[fill=c005093,even odd rule] (1.4792065973, 4.4161969027).. controls 
    (1.4792065973, 4.4188558329) and (1.481163191, 4.4209028877) .. (1.483693666, 
    4.4209028877).. controls (1.4863525698, 4.4209028877) and (1.4883995452, 
    4.4188558065) .. (1.4883995452, 4.4161969027).. controls (1.4883995452, 
    4.4136362652) and (1.486352464, 4.4115814581) .. (1.483693666, 4.4115814581)..
     controls (1.481163191, 4.4115814581) and (1.4792065973, 4.4136362123) .. 
    (1.4792065973, 4.4161969027) -- cycle;
  
  
  
      \end{scope}
      \begin{scope}[shift={(0.8731702173, -3.6154703958)}]
        \path[fill=c005093,even odd rule] (1.4769382715, 4.3937163417).. controls 
    (1.4769382715, 4.3962394083) and (1.4791213485, 4.3983317069) .. 
    (1.4816442565, 4.3983317069).. controls (1.4842202398, 4.3983317069) and 
    (1.486259701, 4.3962393819) .. (1.486259701, 4.3937163417).. controls 
    (1.486259701, 4.3910574379) and (1.4842200281, 4.3890103567) .. (1.4816442565,
     4.3890103567).. controls (1.4791211898, 4.3890103567) and (1.4769382715, 
    4.3910574379) .. (1.4769382715, 4.3937163417) -- cycle;
  
  
  
      \end{scope}
      \begin{scope}[shift={(0.8731702173, -3.6154246229)}]
        \path[fill=c005093,even odd rule] (0.8346167744, 4.4403754273).. controls 
    (0.8346167744, 4.4429437377) and (0.8325696931, 4.4449983331) .. (0.829993895,
     4.4449983331).. controls (0.8273349913, 4.4449983331) and (0.8252878571, 
    4.442943579) .. (0.8252878571, 4.4403754273).. controls (0.8252878571, 
    4.4377692815) and (0.8273349383, 4.4356315804) .. (0.829993895, 
    4.4356315804).. controls (0.8325698783, 4.4356315804) and (0.8346167744, 
    4.4377694137) .. (0.8346167744, 4.4403754273) -- cycle;
  
  
  
      \end{scope}
      \begin{scope}[shift={(0.8731702173, -3.6154246229)}]
        \path[fill=c005093,even odd rule] (1.4904903092, 4.4403754273).. controls 
    (1.4904903092, 4.4429437377) and (1.4883979842, 4.4449983331) .. 
    (1.4858674033, 4.4449983331).. controls (1.4833067658, 4.4449983331) and 
    (1.4811614183, 4.442943579) .. (1.4811614183, 4.4403754273).. controls 
    (1.4811614183, 4.4377692815) and (1.48330666, 4.4356315804) .. (1.4858674033, 
    4.4356315804).. controls (1.4883978783, 4.4356315804) and (1.4904903092, 
    4.4377694137) .. (1.4904903092, 4.4403754273) -- cycle;
  
  
  
      \end{scope}
      \begin{scope}[shift={(0.8731702173, -3.6154703958)}]
        \path[fill=c005093,even odd rule] (0.8664646967, 4.4524472271).. controls 
    (0.8664646967, 4.4551136979) and (0.8644176154, 4.4571532121) .. 
    (0.8617586058, 4.4571532121).. controls (0.8592355392, 4.4571532121) and 
    (0.8571885638, 4.4551138037) .. (0.8571885638, 4.4524472271).. controls 
    (0.8571885638, 4.4500149125) and (0.859235645, 4.4478318619) .. (0.8617586058,
     4.4478318619).. controls (0.8644175096, 4.4478318619) and (0.8664646967, 
    4.450014939) .. (0.8664646967, 4.4524472271) -- cycle;
  
  
  
      \end{scope}
      \begin{scope}[shift={(0.8731702173, -3.6155002938)}]
        \path[fill=c005093,even odd rule] (0.8992331044, 4.4643741675).. controls 
    (0.8992331044, 4.4669046425) and (0.8971407794, 4.4689594496) .. 
    (0.8946177656, 4.4689594496).. controls (0.8920494552, 4.4689594496) and 
    (0.8899117542, 4.4669046954) .. (0.8899117542, 4.4643741675).. controls 
    (0.8899117542, 4.4617228308) and (0.8920495875, 4.4596681825) .. 
    (0.8946177656, 4.4596681825).. controls (0.8971408323, 4.4596681825) and 
    (0.8992331044, 4.4617229367) .. (0.8992331044, 4.4643741675) -- cycle;
  
  
  
      \end{scope}
      \begin{scope}[shift={(0.8731702173, -3.6155907813)}]
        \path[fill=c005093,even odd rule] (0.9320928463, 4.4768000327).. controls 
    (0.9320928463, 4.4794589629) and (0.9300380921, 4.4815059383) .. 
    (0.9274775075, 4.4815059383).. controls (0.9248185773, 4.4815059383) and 
    (0.9227715225, 4.4794588571) .. (0.9227715225, 4.4768000327).. controls 
    (0.9227715225, 4.4742695577) and (0.9248186038, 4.4723055027) .. 
    (0.9274775075, 4.4723055027).. controls (0.930038145, 4.4723055027) and 
    (0.9320928463, 4.4742695048) .. (0.9320928463, 4.4768000327) -- cycle;
  
  
  
      \end{scope}
      \begin{scope}[shift={(0.8731702173, -3.6154703958)}]
        \path[fill=c005093,even odd rule] (0.9658657181, 4.4874891994).. controls 
    (0.9658657181, 4.490012266) and (0.963682641, 4.4920138919) .. (0.9611597331, 
    4.4920138919).. controls (0.9585837498, 4.4920138919) and (0.9564537481, 
    4.4900120544) .. (0.9564537481, 4.4874891994).. controls (0.9564537481, 
    4.4848302692) and (0.9585839085, 4.4826924623) .. (0.9611597331, 
    4.4826924623).. controls (0.9636827998, 4.4826924623) and (0.9658657181, 
    4.4848302956) .. (0.9658657181, 4.4874891994) -- cycle;
  
  
  
      \end{scope}
      \begin{scope}[shift={(0.8731702173, -3.6154703958)}]
        \path[fill=c005093,even odd rule] (0.9986323267, 4.4969010635).. controls 
    (0.9986323267, 4.4995524002) and (0.9964492496, 4.5016070485) .. 
    (0.9939186688, 4.5016070485).. controls (0.9913051146, 4.5016070485) and 
    (0.9892580863, 4.4995522944) .. (0.9892580863, 4.4969010635).. controls 
    (0.9892580863, 4.494461076) and (0.9913051675, 4.4922856983) .. (0.9939186688,
     4.4922856983).. controls (0.9964491438, 4.4922856983) and (0.9986323267, 
    4.4944611025) .. (0.9986323267, 4.4969010635) -- cycle;
  
  
  
      \end{scope}
      \begin{scope}[shift={(0.8731702173, -3.6153793792)}]
        \path[fill=c005093,even odd rule] (1.0313551996, 4.5044395719).. controls 
    (1.0313551996, 4.5070078823) and (1.0293004454, 4.5091456627) .. 
    (1.0267322938, 4.5091456627).. controls (1.0242018187, 4.5091456627) and 
    (1.0221546846, 4.5070078294) .. (1.0221546846, 4.5044395719).. controls 
    (1.0221546846, 4.501901424) and (1.0242017658, 4.4997335869) .. (1.0267322938,
     4.4997335869).. controls (1.0293006042, 4.4997335869) and (1.0313551996, 
    4.5019015827) .. (1.0313551996, 4.5044395719) -- cycle;
  
  
  
      \end{scope}
      \begin{scope}[shift={(0.8731702173, -3.6153801729)}]
        \path[fill=c005093,even odd rule] (1.0641187654, 4.5099693371).. controls 
    (1.0641187654, 4.5125376475) and (1.0620716842, 4.5146754279) .. 
    (1.0595034267, 4.5146754279).. controls (1.0569729517, 4.5146754279) and 
    (1.0549258175, 4.5125375946) .. (1.0549258175, 4.5099693371).. controls 
    (1.0549258175, 4.5073933538) and (1.0569728988, 4.5052633521) .. 
    (1.0595034267, 4.5052633521).. controls (1.0620717371, 4.5052633521) and 
    (1.0641187654, 4.5073935125) .. (1.0641187654, 4.5099693371) -- cycle;
  
  
  
      \end{scope}
      \begin{scope}[shift={(0.8731702173, -3.6155537396)}]
        \path[fill=c005093,even odd rule] (1.096974274, 4.515407874).. controls 
    (1.096974274, 4.5178931052) and (1.0949271927, 4.5199400806) .. (1.0923513681,
     4.5199400806).. controls (1.089737814, 4.5199400806) and (1.0876453831, 
    4.5178929994) .. (1.0876453831, 4.515407874).. controls (1.0876453831, 
    4.5127112142) and (1.0897377081, 4.510701889) .. (1.0923513681, 4.510701889)..
     controls (1.0949273515, 4.510701889) and (1.096974274, 4.5127111348) .. 
    (1.096974274, 4.515407874) -- cycle;
  
  
  
      \end{scope}
      \begin{scope}[shift={(0.8731702173, -3.6154254167)}]
        \path[fill=c005093,even odd rule] (1.1315738894, 4.5198114931).. controls 
    (1.1315738894, 4.52251572) and (1.1295268081, 4.524517584) .. (1.1269962802, 
    4.524517584).. controls (1.12433735, 4.524517584) and (1.1222902952, 
    4.5225157465) .. (1.1222902952, 4.5198114931).. controls (1.1222902952, 
    4.5172884265) and (1.1243373765, 4.5151509106) .. (1.1269962802, 
    4.5151509106).. controls (1.1295267552, 4.5151509106) and (1.1315738894, 
    4.5172884794) .. (1.1315738894, 4.5198114931) -- cycle;
  
  
  
      \end{scope}
      \begin{scope}[shift={(0.8731702173, -3.6156061271)}]
        \path[fill=c005093,even odd rule] (1.1644676302, 4.5201735225).. controls 
    (1.1644676302, 4.5228248592) and (1.162420549, 4.5248720462) .. (1.1597616452,
     4.5248720462).. controls (1.1572385785, 4.5248720462) and (1.1551463065, 
    4.522824965) .. (1.1551463065, 4.5201735225).. controls (1.1551463065, 
    4.5176504558) and (1.1572386315, 4.5156866654) .. (1.1597616452, 
    4.5156866654).. controls (1.1624205754, 4.5156866654) and (1.1644676302, 
    4.5176506675) .. (1.1644676302, 4.5201735225) -- cycle;
  
  
  
      \end{scope}
      \begin{scope}[shift={(0.8731702173, -3.6155608833)}]
        \path[fill=c005093,even odd rule] (1.197273794, 4.51934326).. controls 
    (1.197273794, 4.5219115704) and (1.1951436335, 4.5238680848) .. (1.192575376, 
    4.5238680848).. controls (1.190044901, 4.5238680848) and (1.187869391, 
    4.521911491) .. (1.187869391, 4.51934326).. controls (1.187869391, 
    4.5168201933) and (1.1900447952, 4.5146373808) .. (1.192575376, 
    4.5146373808).. controls (1.1951436865, 4.5146373808) and (1.197273794, 
    4.5168201933) .. (1.197273794, 4.51934326) -- cycle;
  
  
  
      \end{scope}
      \begin{scope}[shift={(0.8731702173, -3.6155989833)}]
        \path[fill=c005093,even odd rule] (1.229091739, 4.5164583756).. controls 
    (1.229091739, 4.5189814423) and (1.2270446577, 4.5209453385) .. (1.224385754, 
    4.5209453385).. controls (1.2218100352, 4.5209453385) and (1.2197628481, 
    4.5189813365) .. (1.2197628481, 4.5164583756).. controls (1.2197628481, 
    4.513807039) and (1.2218099294, 4.5117523906) .. (1.224385754, 4.5117523906)..
     controls (1.2270446842, 4.5117523906) and (1.229091739, 4.5138071448) .. 
    (1.229091739, 4.5164583756) -- cycle;
  
  
  
      \end{scope}
      \begin{scope}[shift={(0.8731702173, -3.6155077021)}]
        \path[fill=c005093,even odd rule] (1.2609395554, 4.5115323631).. controls 
    (1.2609395554, 4.5141912669) and (1.2589001471, 4.5162383481) .. 
    (1.2563695398, 4.5162383481).. controls (1.2537106096, 4.5162383481) and 
    (1.2516635548, 4.5141912669) .. (1.2516635548, 4.5115323631).. controls 
    (1.2516635548, 4.5090092965) and (1.253710636, 4.506954754) .. (1.2563695398, 
    4.506954754).. controls (1.2589000148, 4.506954754) and (1.2609395554, 
    4.5090095081) .. (1.2609395554, 4.5115323631) -- cycle;
  
  
  
      \end{scope}
      \begin{scope}[shift={(0.8731702173, -3.6155537396)}]
        \path[fill=c005093,even odd rule] (1.2938431123, 4.5066989019).. controls 
    (1.2938431123, 4.509320129) and (1.2916677081, 4.5113141348) .. (1.2891371273,
     4.5113141348).. controls (1.2865688169, 4.5113141348) and (1.2844235752, 
    4.5093199702) .. (1.2844235752, 4.5066989019).. controls (1.2844235752, 
    4.5041229185) and (1.2865688169, 4.5020759696) .. (1.2891371273, 
    4.5020759696).. controls (1.2916676023, 4.5020759696) and (1.2938431123, 
    4.5041230508) .. (1.2938431123, 4.5066989019) -- cycle;
  
  
  
      \end{scope}
      \begin{scope}[shift={(0.8731702173, -3.6155537396)}]
        \path[fill=c005093,even odd rule] (1.3266022331, 4.5012897338).. controls 
    (1.3266022331, 4.5038202087) and (1.3244265644, 4.5058220992) .. 
    (1.3218961688, 4.5058220992).. controls (1.3193355313, 4.5058220992) and 
    (1.3172354804, 4.5038202617) .. (1.3172354804, 4.5012897338).. controls 
    (1.3172354804, 4.4986383706) and (1.3193354783, 4.4965837487) .. 
    (1.3218961688, 4.4965837487).. controls (1.3244266438, 4.4965837487) and 
    (1.3266022331, 4.4986385029) .. (1.3266022331, 4.5012897338) -- cycle;
  
  
  
      \end{scope}
      \begin{scope}[shift={(0.8731702173, -3.6155537396)}]
        \path[fill=c005093,even odd rule] (1.3593733925, 4.4916669967).. controls 
    (1.3593733925, 4.49424298) and (1.3572884758, 4.4961993885) .. (1.3546220844, 
    4.4961993885).. controls (1.3520614469, 4.4961993885) and (1.3500066133, 
    4.4942427948) .. (1.3500066133, 4.4916669967).. controls (1.3500066133, 
    4.4891365217) and (1.3520613675, 4.4869610117) .. (1.3546220844, 
    4.4869610117).. controls (1.3572885552, 4.4869610117) and (1.3593733925, 
    4.4891364158) .. (1.3593733925, 4.4916669967) -- cycle;
  
  
  
      \end{scope}
      \begin{scope}[shift={(0.8731702173, -3.6155156396)}]
        \path[fill=c005093,even odd rule] (1.3930987188, 4.4834324754).. controls 
    (1.3930987188, 4.4859629504) and (1.3909233146, 4.4880025175) .. 
    (1.3883927338, 4.4880025175).. controls (1.3858244233, 4.4880025175) and 
    (1.3837773685, 4.4859631092) .. (1.3837773685, 4.4834324754).. controls 
    (1.3837773685, 4.480864165) and (1.3858244498, 4.4787264904) .. (1.3883927338,
     4.4787264904).. controls (1.3909232088, 4.4787264904) and (1.3930987188, 
    4.4808643237) .. (1.3930987188, 4.4834324754) -- cycle;
  
  
  
      \end{scope}
      \begin{scope}[shift={(0.8731702173, -3.615591575)}]
        \path[fill=c005093,even odd rule] (1.4258678408, 4.472306905).. controls 
    (1.4258678408, 4.4748828883) and (1.4238284325, 4.4768014615) .. 
    (1.4211619352, 4.4768014615).. controls (1.4185936248, 4.4768014615) and 
    (1.4166296492, 4.4748827031) .. (1.4166296492, 4.472306905).. controls 
    (1.4166296492, 4.4697462675) and (1.4185936513, 4.46760092) .. (1.4211619352, 
    4.46760092).. controls (1.4238284325, 4.46760092) and (1.4258678408, 
    4.4697461617) .. (1.4258678408, 4.472306905) -- cycle;
  
  
  
      \end{scope}
      \begin{scope}[shift={(0.8731702173, -3.6153801729)}]
        \path[fill=c005093,even odd rule] (1.4595965008, 4.4576746).. controls 
    (1.4595965008, 4.460205075) and (1.4574210967, 4.4623806908) .. (1.4548905158,
     4.4623806908).. controls (1.4523222054, 4.4623806908) and (1.4502750448, 
    4.4602050221) .. (1.4502750448, 4.4576746).. controls (1.4502750448, 
    4.4550988813) and (1.452322126, 4.452968615) .. (1.4548905158, 4.452968615).. 
    controls (1.4574209908, 4.452968615) and (1.4595965008, 4.4550987754) .. 
    (1.4595965008, 4.4576746) -- cycle;
  
  
  
      \end{scope}
      \begin{scope}[shift={(0.8731702173, -3.614556525)}]
        \path[fill=c005093,even odd rule] (1.0973672067, 4.0181495869).. controls 
    (1.0973672067, 4.0208538402) and (1.0951010504, 4.0232483987) .. 
    (1.0922683948, 4.0232483987).. controls (1.0895188448, 4.0232483987) and 
    (1.0871695829, 4.0208539196) .. (1.0871695829, 4.0181495869).. controls 
    (1.0871695829, 4.0151431794) and (1.0895188183, 4.0130130454) .. 
    (1.0922683948, 4.0130130454).. controls (1.0951011033, 4.0130130454) and 
    (1.0973672067, 4.0151432058) .. (1.0973672067, 4.0181495869) -- cycle;
  
  
  
      \end{scope}
      \begin{scope}[shift={(0.8731702173, -3.6146769104)}]
        \path[fill=c005093,even odd rule] (1.1313092267, 4.0213675552).. controls 
    (1.1313092267, 4.0240264854) and (1.1291338225, 4.0262926681) .. (1.126301061,
     4.0262926681).. controls (1.1234683525, 4.0262926681) and (1.1212022492, 
    4.0240265119) .. (1.1212022492, 4.0213675552).. controls (1.1212022492, 
    4.018451794) and (1.1234684054, 4.0161781235) .. (1.126301061, 4.0161781235)..
     controls (1.1291337696, 4.0161781235) and (1.1313092267, 4.0184516881) .. 
    (1.1313092267, 4.0213675552) -- cycle;
  
  
  
      \end{scope}
      \begin{scope}[shift={(0.8731702173, -3.6145943604)}]
        \path[fill=c005093,even odd rule] (1.1664336696, 4.021889684).. controls 
    (1.1664336696, 4.024639234) and (1.1641675133, 4.0269960629) .. (1.1613348577,
     4.0269960629).. controls (1.1585853077, 4.0269960629) and (1.1563265862, 
    4.0246391546) .. (1.1563265862, 4.021889684).. controls (1.1563265862, 
    4.0190116523) and (1.1585853342, 4.0167984392) .. (1.1613348577, 
    4.0167984392).. controls (1.1641675662, 4.0167984392) and (1.1664336696, 
    4.0190116788) .. (1.1664336696, 4.021889684) -- cycle;
  
  
  
      \end{scope}
      \begin{scope}[shift={(0.8731702173, -3.6146769104)}]
        \path[fill=c005093,even odd rule] (1.201250164, 4.0211938563).. controls 
    (1.201250164, 4.0240265913) and (1.1989840077, 4.0262926681) .. (1.1961513521,
     4.0262926681).. controls (1.1934093692, 4.0262926681) and (1.1910525402, 
    4.0240265119) .. (1.1910525402, 4.0211938563).. controls (1.1910525402, 
    4.0184518733) and (1.1934094485, 4.0161781235) .. (1.1961513521, 
    4.0161781235).. controls (1.1989840606, 4.0161781235) and (1.201250164, 
    4.0184519527) .. (1.201250164, 4.0211938563) -- cycle;
  
  
  
      \end{scope}
      \begin{scope}[shift={(0.8731702173, -3.614556525)}]
        \path[fill=c005093,even odd rule] (1.2364573421, 4.0181495869).. controls 
    (1.2364573421, 4.0209823219) and (1.2342364296, 4.0232483987) .. 
    (1.2314491765, 4.0232483987).. controls (1.2286996265, 4.0232483987) and 
    (1.2263429033, 4.0209822425) .. (1.2263429033, 4.0181495869).. controls 
    (1.2263429033, 4.0153623073) and (1.2286995471, 4.0130130454) .. 
    (1.2314491765, 4.0130130454).. controls (1.2342364825, 4.0130130454) and 
    (1.2364573421, 4.0153622808) .. (1.2364573421, 4.0181495869) -- cycle;
  
  
  
      \end{scope}
      \begin{scope}[shift={(0.8731702173, -3.6145935667)}]
        \path[fill=c005093,even odd rule] (1.2681312744, 4.0143950965).. controls 
    (1.2681312744, 4.017182376) and (1.2658651181, 4.019539205) .. (1.2630324625, 
    4.019539205).. controls (1.2602904796, 4.019539205) and (1.2579412177, 
    4.0171822967) .. (1.2579412177, 4.0143950965).. controls (1.2579412177, 
    4.0116077904) and (1.2602904531, 4.0093415813) .. (1.2630324625, 
    4.0093415813).. controls (1.265865171, 4.0093415813) and (1.2681312744, 
    4.0116077375) .. (1.2681312744, 4.0143950965) -- cycle;
  
  
  
      \end{scope}
      \begin{scope}[shift={(0.8731702173, -3.6146848479)}]
        \path[fill=c005093,even odd rule] (1.3087387571, 4.0075140513).. controls 
    (1.3087387571, 4.0103920829) and (1.3064726008, 4.0126581598) .. 
    (1.3037306708, 4.0126581598).. controls (1.3008979623, 4.0126581598) and 
    (1.2987678283, 4.0103920035) .. (1.2987678283, 4.0075140513).. controls 
    (1.2987678283, 4.0047267452) and (1.3008979888, 4.0025511823) .. 
    (1.3037306708, 4.0025511823).. controls (1.3064726538, 4.0025511823) and 
    (1.3087387571, 4.004726851) .. (1.3087387571, 4.0075140513) -- cycle;
  
  
  
      \end{scope}
      \begin{scope}[shift={(0.8731702173, -3.614600975)}]
        \path[fill=c005093,even odd rule] (1.3464767806, 3.9989325819).. controls 
    (1.3464767806, 4.0015991056) and (1.3442558681, 4.0039406946) .. 
    (1.3413855358, 4.0039406946).. controls (1.3385528273, 4.0039406946) and 
    (1.3364076385, 4.0015991321) .. (1.3364076385, 3.9989325819).. controls 
    (1.3364076385, 3.9959790381) and (1.3385528802, 3.9937506644) .. 
    (1.3413855358, 3.9937506644).. controls (1.3442560004, 3.9937506644) and 
    (1.3464767806, 3.9959789852) .. (1.3464767806, 3.9989325819) -- cycle;
  
  
  
      \end{scope}
      \begin{scope}[shift={(0.8731702173, -3.6146396042)}]
        \path[fill=c005093,even odd rule] (1.3811207931, 3.9901706402).. controls 
    (1.3811207931, 3.9930033223) and (1.3787715577, 3.9952693992) .. 
    (1.3760219813, 3.9952693992).. controls (1.3731061935, 3.9952693992) and 
    (1.3709307365, 3.9930032429) .. (1.3709307365, 3.9901706402).. controls 
    (1.3709307365, 3.9872926615) and (1.3731061406, 3.9851171515) .. 
    (1.3760219813, 3.9851171515).. controls (1.3787715313, 3.9851171515) and 
    (1.3811207931, 3.9872925556) .. (1.3811207931, 3.9901706402) -- cycle;
  
  
  
      \end{scope}
      \begin{scope}[shift={(0.8731702173, -3.6145943604)}]
        \path[fill=c005093,even odd rule] (1.4136676388, 3.9801090388).. controls 
    (1.4136676388, 3.9828510746) and (1.4114014825, 3.98520031) .. (1.4087047963, 
    3.98520031).. controls (1.4058267646, 3.98520031) and (1.4036059844, 
    3.9828510746) .. (1.4036059844, 3.9801090388).. controls (1.4036059844, 
    3.9773670029) and (1.4058266323, 3.9750027127) .. (1.4087047963, 
    3.9750027127).. controls (1.411401456, 3.9750027127) and (1.4136676388, 
    3.9773670294) .. (1.4136676388, 3.9801090388) -- cycle;
  
  
  
      \end{scope}
      \begin{scope}[shift={(0.8731702173, -3.6146922563)}]
        \path[fill=c005093,even odd rule] (1.4439156256, 3.9671385023).. controls 
    (1.4439156256, 3.970008914) and (1.4415134735, 3.9722675296) .. (1.4388168137,
     3.9722675296).. controls (1.4359387821, 3.9722675296) and (1.4337180019, 
    3.9700090462) .. (1.4337180019, 3.9671385023).. controls (1.4337180019, 
    3.9644267083) and (1.4359386498, 3.9621680662) .. (1.4388168137, 
    3.9621680662).. controls (1.4415134735, 3.9621680662) and (1.4439156256, 
    3.9644268142) .. (1.4439156256, 3.9671385023) -- cycle;
  
  
  
      \end{scope}
      \begin{scope}[shift={(0.8731702173, -3.4588127208)}]
        \path[fill=c005093,even odd rule] (1.458244824, 3.9086884192).. controls 
    (1.3903589583, 3.9441533808) and (1.2701858169, 3.9573196825) .. 
    (1.2689242835, 3.9575312169).. controls (1.2665599669, 3.9576979044) and 
    (1.2649132002, 3.9598804523) .. (1.2649962794, 3.962101259).. controls 
    (1.2653891856, 3.9643674152) and (1.2673529231, 3.9661501777) .. 
    (1.2697022644, 3.9659839665).. controls (1.2750126106, 3.9654627373) and 
    (1.3858266723, 3.9530366869) .. (1.4563714681, 3.9188860429).. controls 
    (1.4570286931, 3.9294235769) and (1.4621726429, 4.0084210165) .. 
    (1.4625654698, 4.0136557712).. controls (1.4464381896, 4.0204918902) and 
    (1.3581571558, 4.0552243) .. (1.2713565981, 4.0628460727).. controls 
    (1.2690904419, 4.0632011435) and (1.2674360552, 4.0650291498) .. 
    (1.2674360552, 4.0669855319).. controls (1.2674360552, 4.0671680944) and 
    (1.2674360552, 4.067287686) .. (1.2674360552, 4.0673784381).. controls 
    (1.267518076, 4.0696520027) and (1.2694831365, 4.0716161106) .. (1.2718778008,
     4.0712989017).. controls (1.3710137098, 4.0601494923) and (1.467656794, 
    4.0206658802) .. (1.4686161467, 4.0202806204) -- (1.4713656967, 4.0191021663) 
    -- (1.4639101617, 3.9056367679) -- (1.458244824, 3.9086884985) -- cycle;
  
  
  
      \end{scope}
      \begin{scope}[shift={(0.8731702173, -3.4588127208)}]
        \path[fill=c005093,even odd rule] (0.8626718152, 3.9086884192).. controls 
    (0.9305199513, 3.9441533808) and (1.0515692604, 3.9582715475) .. 
    (1.0528307938, 3.9584905696).. controls (1.0551877021, 3.9587022363) and 
    (1.0568418771, 3.9607567258) .. (1.0567512573, 3.9629850202).. controls 
    (1.0564491031, 3.965372091) and (1.0544020219, 3.9671094246) .. (1.0521357863,
     3.9669432133).. controls (1.0467801169, 3.9663767404) and (0.9350068083, 
    3.9530366869) .. (0.8645073621, 3.9188860429).. controls (0.8637218142, 
    3.9294235769) and (0.8586230023, 4.0084210165) .. (0.8582301754, 
    4.0136557712).. controls (0.8745312604, 4.0204918902) and (0.9626764306, 
    4.0552243) .. (1.0494769883, 4.0628460727).. controls (1.0516600654, 
    4.0632011435) and (1.0534426427, 4.0650291498) .. (1.0534426427, 
    4.0669855319).. controls (1.0534426427, 4.0671680944) and (1.0534426427, 
    4.067287686) .. (1.0534426427, 4.0673784381).. controls (1.0533129969, 
    4.0696520027) and (1.051266974, 4.0716161106) .. (1.0490010029, 
    4.0712989017).. controls (0.9498651204, 4.0601494923) and (0.8532597394, 
    4.0206658802) .. (0.8522626571, 4.0202806204) -- (0.8495131071, 4.0191021663) 
    -- (0.8568402663, 3.9056367679) -- (0.8626718152, 3.9086884985) -- cycle;
  
  
  
      \end{scope}
      \begin{scope}[shift={(0.8731702173, -3.4590839188)}]
        \path[fill=c005093,even odd rule] (1.4721517473, 4.1071705279).. controls 
    (1.4041752354, 4.1429075077) and (1.2711001375, 4.1551521862) .. 
    (1.2697026612, 4.1551521862).. controls (1.267436505, 4.1554921758) and 
    (1.2657824094, 4.1574937488) .. (1.2660014315, 4.1597675515).. controls 
    (1.2661760565, 4.161988464) and (1.2681315919, 4.1637710677) .. (1.270488315, 
    4.16360483).. controls (1.2758967688, 4.1630761925) and (1.3996807054, 
    4.15157158) .. (1.470316121, 4.1173680458).. controls (1.470973346, 
    4.1280491163) and (1.4762004808, 4.2069106923) .. (1.4764573383, 
    4.2122360933).. controls (1.4602997633, 4.2191176412) and (1.3673958502, 
    4.25675806) .. (1.2815546187, 4.261464045).. controls (1.2793339708, 
    4.26163867) and (1.2775436413, 4.2633828033) .. (1.2775436413, 4.2656565002)..
     controls (1.2775436413, 4.2656565002) and (1.2775436413, 4.2658681669) .. 
    (1.2775436413, 4.2660038981).. controls (1.2777262038, 4.2683078898) and 
    (1.2797190454, 4.2699620913) .. (1.282037986, 4.269871445).. controls 
    (1.3774196692, 4.2646065277) and (1.4815561502, 4.2192837731) .. 
    (1.4825607996, 4.2188985133) -- (1.4853028883, 4.2176295717) -- (1.4778548146,
     4.1042547931) -- (1.4721517473, 4.1071705808) -- cycle;
  
  
  
      \end{scope}
      \begin{scope}[shift={(0.8731702173, -3.4590839188)}]
        \path[fill=c005093,even odd rule] (0.8489942327, 4.1071705279).. controls 
    (0.9169707446, 4.1429075077) and (1.0500004135, 4.1551521862) .. 
    (1.0513978898, 4.1551521862).. controls (1.0537092898, 4.1554921758) and 
    (1.0554920258, 4.1574937488) .. (1.0551898452, 4.1597675515).. controls 
    (1.054970241, 4.161988464) and (1.052923689, 4.1637710677) .. (1.050612236, 
    4.16360483).. controls (1.0452944285, 4.1630761925) and (0.9214122785, 
    4.15157158) .. (0.8508750763, 4.1173680458).. controls (0.8502178513, 
    4.1280491163) and (0.8451191452, 4.2069106923) .. (0.8446810217, 
    4.2122360933).. controls (0.8608007877, 4.2191176412) and (0.9538783998, 
    4.25675806) .. (1.0395836619, 4.261464045).. controls (1.0418572265, 
    4.26163867) and (1.0435116925, 4.2633828033) .. (1.0435116925, 4.2656565002)..
     controls (1.0435116925, 4.2656565002) and (1.0435116925, 4.2658681669) .. 
    (1.0435116925, 4.2660038981).. controls (1.0434217342, 4.2683078898) and 
    (1.0414572029, 4.2699620913) .. (1.0391985079, 4.269871445).. controls 
    (0.9436355852, 4.2646065277) and (0.8395822098, 4.2192837731) .. 
    (0.8386228571, 4.2188985133) -- (0.8358808212, 4.2176295717) -- (0.8433288421,
     4.1042547931) -- (0.8489942327, 4.1071705808) -- cycle;
  
  
  
      \end{scope}
      \begin{scope}[shift={(0.8731702173, -3.4857795833)}]
        \path[fill=c005093,even odd rule] (1.0924439194, 4.0174321956) -- 
    (1.1274783246, 4.0162614144).. controls (1.1436434402, 4.0129075031) and 
    (1.1403273908, 3.9929504056) .. (1.1403273908, 3.9929504056).. controls 
    (1.1403273908, 3.9775255942) and (1.1319652875, 3.9792251717) .. 
    (1.1245550756, 3.9762111706).. controls (1.1171070019, 3.9732953829) and 
    (1.0954050038, 3.9779182623) .. (1.0954050038, 3.9779182623).. controls 
    (1.0829487644, 3.9772157935) and (1.086295135, 3.9671542185) .. (1.086295135, 
    3.9671542185).. controls (1.0859099017, 3.9513365802) and (1.0738388956, 
    3.9508152717) .. (1.0738388956, 3.9508152717).. controls (1.0609823681, 
    3.9500220508) and (1.0604990008, 3.9745644304) .. (1.0604990008, 3.9745644304)
     -- (1.0609823946, 4.0627549238).. controls (1.0609823946, 4.0627549238) and 
    (1.0621605842, 4.0723934829) .. (1.0629463967, 4.0793958396).. controls 
    (1.0638150237, 4.0865039767) and (1.0729627279, 4.0882488508) .. (1.077933164,
     4.0885434644).. controls (1.0829488967, 4.089034531) and (1.0870807888, 
    4.0861111498) .. (1.0870807888, 4.0861111498) -- (1.1074759304, 
    4.0857182435).. controls (1.1282639783, 4.0835428394) and (1.131965314, 
    4.0764800783) .. (1.1345335979, 4.0698100912).. controls (1.1346235562, 
    4.0695153454) and (1.1347532021, 4.0690245433) .. (1.134843425, 
    4.0687223892).. controls (1.1367997542, 4.0614858762) and (1.1345335979, 
    4.0461970077) .. (1.1345335979, 4.0461970077).. controls (1.1307037806, 
    4.0270254581) and (1.0924439458, 4.0337333071) .. (1.0924439458, 
    4.0337333071).. controls (1.0812869694, 4.034518855) and (1.0824653442, 
    4.0261568575) .. (1.0824653442, 4.0261568575).. controls (1.0820271942, 
    4.0178325896) and (1.0924439458, 4.0174322221) .. (1.0924439458, 4.0174322221)
     -- cycle(1.0916583715, 4.0504270867) -- (1.1055120077, 4.0504270867).. 
    controls (1.1104824173, 4.0504270867) and (1.1146670938, 4.0543475237) .. 
    (1.1146670938, 4.0593103927).. controls (1.1146670938, 4.0641069975) and 
    (1.1104824173, 4.0681558369) .. (1.1055120077, 4.0681558369) -- (1.0916583715,
     4.0681558369).. controls (1.0866048827, 4.0681558369) and (1.0824653177, 
    4.0641069975) .. (1.0824653177, 4.0593103927).. controls (1.0824653177, 
    4.0543475237) and (1.0866048827, 4.0504270867) .. (1.0916583715, 4.0504270867)
     -- cycle(1.0837721742, 3.9988272512).. controls (1.0837721742, 3.9936982769) 
    and (1.0883421104, 3.9895511448) .. (1.093924369, 3.9895511448) -- 
    (1.1093492863, 3.9895511448).. controls (1.1149315448, 3.9895511448) and 
    (1.1195468835, 3.9936982769) .. (1.1195468835, 3.9988272512).. controls 
    (1.1195468835, 4.0037145023) and (1.1149315448, 4.0079370406) .. 
    (1.1093492863, 4.0079370406) -- (1.093924369, 4.0079370406).. controls 
    (1.0883421104, 4.0079370406) and (1.0837721742, 4.0037145023) .. 
    (1.0837721742, 3.9988272512) -- cycle;
  
  
  
      \end{scope}
      \begin{scope}[shift={(0.8731702173, -3.4974723146)}]
        \path[fill=c005093,even odd rule] (1.1730948725, 3.98050131) -- 
    (1.2511405265, 3.9812868579).. controls (1.2511405265, 3.9812868579) and 
    (1.2649110835, 3.982155485) .. (1.2649110835, 3.9925646431).. controls 
    (1.2649110835, 3.9925646431) and (1.2628563294, 4.0067732444) .. 
    (1.2511405265, 4.0071736383) -- (1.2304432042, 4.0083444196).. controls 
    (1.2304432042, 4.0083444196) and (1.2230404535, 4.0070828863) .. 
    (1.2230404535, 4.0165857142).. controls (1.2230404535, 4.0165857142) and 
    (1.2208499681, 4.0257785562) .. (1.233404315, 4.0270400896).. controls 
    (1.233404315, 4.0270400896) and (1.26577979, 4.027954225) .. (1.2686955777, 
    4.0416793531).. controls (1.2686955777, 4.0416793531) and (1.2799884971, 
    4.0578898183) .. (1.2721023792, 4.0795464669).. controls (1.2721023792, 
    4.0795464669) and (1.2624636083, 4.1039527977) .. (1.2537088104, 4.079848621) 
    -- (1.2508761813, 4.0640310092) -- (1.2487007771, 4.0578897919).. controls 
    (1.2487007771, 4.0578897919) and (1.2445613179, 4.0495201215) .. 
    (1.2346281452, 4.0495201215).. controls (1.2346281452, 4.0495201215) and 
    (1.2230406123, 4.047389961) .. (1.2232974698, 4.0636456435) -- (1.2213787115, 
    4.0902199438).. controls (1.2213787115, 4.0902199438) and (1.2163252227, 
    4.1061282283) .. (1.2080463044, 4.1061282283).. controls (1.2080463044, 
    4.1061282283) and (1.19379238, 4.1078202387) .. (1.19379238, 4.0857330867) -- 
    (1.1925760904, 4.0599367408).. controls (1.1925760904, 4.0599367408) and 
    (1.1950536488, 4.0499129483) .. (1.1800670667, 4.0495201215).. controls 
    (1.1800670667, 4.0495201215) and (1.165125596, 4.047389961) .. (1.1659112498, 
    4.068691671) -- (1.1646420435, 4.0807173275).. controls (1.1646420435, 
    4.0807173275) and (1.161809335, 4.09249372) .. (1.1539685402, 4.0919196271).. 
    controls (1.1539685402, 4.0919196271) and (1.1426831879, 4.0944953458) .. 
    (1.1430760148, 4.0749311017) -- (1.1426831085, 4.0491349675).. controls 
    (1.1426831085, 4.0491349675) and (1.1464298467, 4.031308004) .. (1.1629877098,
     4.0308321244) -- (1.1751493356, 4.0304392181).. controls (1.1751493356, 
    4.0304392181) and (1.1918355746, 4.0320933931) .. (1.1922208344, 
    4.0165856083).. controls (1.1922208344, 4.0165856083) and (1.1904383365, 
    4.0063049054) .. (1.1775816502, 4.0063049054) -- (1.1581080258, 
    4.0059045908).. controls (1.1581080258, 4.0059045908) and (1.1484768221, 
    4.0067732179) .. (1.148869649, 3.9925646167).. controls (1.148869649, 
    3.9925646167) and (1.1476914594, 3.9805012042) .. (1.1662963244, 3.9805012042)
     -- (1.1730946873, 3.9805012042) -- cycle;
  
  
  
      \end{scope}
      \begin{scope}[shift={(0.8731702173, -3.4325480625)}]
        \path[fill=c005093,even odd rule] (1.0891265206, 4.1074284967) -- 
    (1.1294335425, 4.1083047967).. controls (1.1294335425, 4.1083047967) and 
    (1.1427736225, 4.0825086625) .. (1.1577603104, 4.082078106).. controls 
    (1.1577603104, 4.082078106) and (1.1723163096, 4.0799479456) .. (1.1839491392,
     4.1083047967) -- (1.2417281596, 4.1066053779).. controls (1.2417281596, 
    4.1066053779) and (1.2671693023, 4.1062201446) .. (1.2693145704, 4.1253839948)
     -- (1.2775407308, 4.1573290192).. controls (1.2775407308, 4.1573290192) and 
    (1.2858649723, 4.209789994) .. (1.3008969569, 4.2384641598).. controls 
    (1.3008969569, 4.2384641598) and (1.3133682246, 4.2505275723) .. 
    (1.3034729402, 4.2596450346).. controls (1.3034729402, 4.2596450346) and 
    (1.2905257665, 4.2721918144) .. (1.2742322485, 4.2434270023).. controls 
    (1.2742322485, 4.2434270023) and (1.2588980833, 4.2251317262) .. 
    (1.2414639467, 4.1452580396).. controls (1.2414639467, 4.1452580396) and 
    (1.2426424008, 4.1282167298) .. (1.2285242342, 4.1290854362) -- (1.1927948479,
     4.1299617362).. controls (1.1927948479, 4.1299617362) and (1.2031209269, 
    4.1565434183) .. (1.2114376013, 4.1972885637).. controls (1.2114376013, 
    4.1972885637) and (1.2173295281, 4.219338145) .. (1.2230703515, 
    4.2314015575).. controls (1.2230703515, 4.2314015575) and (1.2301785944, 
    4.2484428937) .. (1.2168461873, 4.2554981406).. controls (1.2168461873, 
    4.2554981406) and (1.1981126554, 4.2580738594) .. (1.1927948479, 4.2364172108)
     -- (1.1744920048, 4.1672169217) -- (1.1661225196, 4.1395021879).. controls 
    (1.1661225196, 4.1395021879) and (1.158583826, 4.1215468485) .. (1.1523972856,
     4.136065224) -- (1.1453419329, 4.1649130887).. controls (1.1453419329, 
    4.1649130887) and (1.1328782323, 4.2288331942) .. (1.12028605, 4.2434271346)..
     controls (1.12028605, 4.2434271346) and (1.1057752681, 4.2626062512) .. 
    (1.09623479, 4.2554981406).. controls (1.09623479, 4.2554981406) and 
    (1.081678685, 4.2488734502) .. (1.0940969567, 4.2269221617).. controls 
    (1.0940969567, 4.2269221617) and (1.1078221906, 4.198157376) .. (1.1108286246,
     4.1748010442) -- (1.1191529719, 4.1299617362) -- (1.0900028206, 
    4.1290854362).. controls (1.0900028206, 4.1290854362) and (1.0829476531, 
    4.1233445071) .. (1.0783247473, 4.1429012106) -- (1.0616761585, 4.1959817867) 
    -- (1.0475580977, 4.2359790344) -- (1.0455110165, 4.2505275723).. controls 
    (1.0455110165, 4.2505275723) and (1.0428974623, 4.2534509271) .. 
    (1.0388032733, 4.256812326).. controls (1.0347091373, 4.2602039933) and 
    (1.0267323996, 4.2622057515) .. (1.023423785, 4.2604230154).. controls 
    (1.0201077092, 4.2588593279) and (1.0146613935, 4.2558982965) .. 
    (1.0142685667, 4.2534508212).. controls (1.0137925813, 4.2509203463) and 
    (1.0163608917, 4.2388569867) .. (1.0167082896, 4.2372027323).. controls 
    (1.0171011958, 4.2355408844) and (1.0234237056, 4.2196855165) .. 
    (1.0234237056, 4.2196855165) -- (1.0475580183, 4.1474410637) -- (1.0541827881,
     4.12116905).. controls (1.0541827881, 4.12116905) and (1.0616760792, 
    4.1095663829) .. (1.0728330292, 4.1083048496) -- (1.0841637046, 4.1066054308) 
    -- (1.0891264413, 4.1074288142) -- cycle;
  
  
  
      \end{scope}
      \begin{scope}[shift={(0.8731702173, -3.5647397167)}]
        \path[fill=c005093,even odd rule] (1.1699624969, 4.5179765548).. controls 
    (1.168957874, 4.518935934) and (1.169048626, 4.5208545071) .. (1.1701741635, 
    4.5218969125).. controls (1.2031237315, 4.5509185821) and (1.2453040298, 
    4.5697878452) .. (1.2923640915, 4.5764125092).. controls (1.2938446998, 
    4.5766321133) and (1.2953252817, 4.5755441467) .. (1.2954989806, 
    4.5740937008).. controls (1.2956286265, 4.5724469342) and (1.2946303535, 
    4.57113259) .. (1.293194989, 4.5708379765).. controls (1.2465655102, 
    4.5642586092) and (1.2064322931, 4.5464316721) .. (1.173754611, 
    4.5175838073).. controls (1.1727952319, 4.5165791844) and (1.1709672256, 
    4.5168059323) .. (1.1699625498, 4.5179767135) -- cycle;
  
  
  
      \end{scope}
      \begin{scope}[shift={(0.8731702173, -3.5940327958)}]
        \path[fill=c005093,even odd rule] (1.1335741129, 4.5021204725).. controls 
    (1.1321842567, 4.5029136933) and (1.1317008629, 4.5045752767) .. 
    (1.1323580879, 4.5058671312).. controls (1.1327054858, 4.5064412771) and 
    (1.1379779967, 4.5166387421) .. (1.1500036796, 4.530802126).. controls 
    (1.1511366254, 4.5319729073) and (1.152791065, 4.5321467385) .. (1.1540600596,
     4.5311873594).. controls (1.1552308408, 4.5300996573) and (1.1553215929, 
    4.5284452706) .. (1.1542796638, 4.5271837373).. controls (1.1426845638, 
    4.5135415825) and (1.1374120529, 4.5035253306) .. (1.1372836771, 
    4.5032987415).. controls (1.1366716958, 4.5020372081) and (1.1349270333, 
    4.501387656) .. (1.1335747744, 4.5021202873) -- cycle;
  
  
  
      \end{scope}
      \begin{scope}[shift={(0.8731702173, -3.5582656271)}]
        \path[fill=c005093,even odd rule] (1.1606390565, 4.5200009348).. controls 
    (1.1595436815, 4.5212246327) and (1.1597627565, 4.5229619398) .. 
    (1.1609412106, 4.5240572354).. controls (1.1891772265, 4.5487505862) and 
    (1.2353610146, 4.5795700202) .. (1.2944996494, 4.5850540919).. controls 
    (1.296146416, 4.5852657585) and (1.297370114, 4.5840571419) .. (1.2975438129, 
    4.5824781085).. controls (1.2977634171, 4.5809599294) and (1.2965844338, 
    4.5796908025) .. (1.2951114983, 4.5794792681).. controls (1.2373250167, 
    4.5740707879) and (1.1922214958, 4.544044469) .. (1.1645595729, 
    4.5197363515).. controls (1.1634718708, 4.5187394015) and (1.1617267585, 
    4.5188600515) .. (1.16063903, 4.5200009348) -- cycle;
  
  
  
      \end{scope}
      \begin{scope}[shift={(0.8731702173, -3.5859278146)}]
        \path[fill=c005093,even odd rule] (1.114666591, 4.4951938131).. controls 
    (1.1132690619, 4.4960248694) and (1.1128838285, 4.4977621235) .. 
    (1.1137524556, 4.4991216321).. controls (1.1141453619, 4.4996351883) and 
    (1.1232929337, 4.5137608956) .. (1.1405080483, 4.5321090619).. controls 
    (1.1417240733, 4.5332798431) and (1.1433784071, 4.5332798431) .. 
    (1.1445567554, 4.5321910827).. controls (1.1456897013, 4.5311033806) and 
    (1.1456899658, 4.5293584535) .. (1.1446467138, 4.5281876723).. controls 
    (1.1276507271, 4.5103229527) and (1.1184576469, 4.4961973248) .. 
    (1.1183670271, 4.4960236258).. controls (1.1175888875, 4.4948076008) and 
    (1.1157987167, 4.4944146946) .. (1.1146656915, 4.4951928342) -- cycle;
  
  
  
      \end{scope}
      \begin{scope}[shift={(0.8731702173, -3.5545940042)}]
        \path[fill=c005093,even odd rule] (1.1494809688, 4.525091571).. controls 
    (1.1483480229, 4.5263454315) and (1.1485668333, 4.5281358404) .. 
    (1.1496555937, 4.5291404898).. controls (1.1925761963, 4.5660029542) and 
    (1.2421214892, 4.5878787306) .. (1.2965843015, 4.5939293281).. controls 
    (1.298019666, 4.5941118906) and (1.2995454123, 4.592969949) .. (1.2997266783, 
    4.5914970135).. controls (1.299938345, 4.5899259177) and (1.2987672992, 
    4.5885284415) .. (1.2972866908, 4.5884907119).. controls (1.2438209608, 
    4.5823872242) and (1.1954464756, 4.5609874331) .. (1.1534021467, 
    4.5248727871).. controls (1.1521861217, 4.5239134079) and (1.1503958185, 
    4.5239134079) .. (1.149481789, 4.5250923913) -- cycle;
  
  
  
      \end{scope}
      \begin{scope}[shift={(0.8731702173, -3.5748166375)}]
        \path[fill=c005093,even odd rule] (1.0957970369, 4.4825266479).. controls 
    (1.0944902598, 4.4832140354) and (1.0940521098, 4.4849137188) .. 
    (1.0948376577, 4.486273386).. controls (1.0953136431, 4.4872250923) and 
    (1.1065611335, 4.5066609342) .. (1.1298268456, 4.5305837919).. controls 
    (1.1309145477, 4.5317545731) and (1.1327047715, 4.5317545731) .. 
    (1.1337471769, 4.5307055002).. controls (1.134925631, 4.529625471) and 
    (1.134925631, 4.5278726858) .. (1.1338371352, 4.5266565815).. controls 
    (1.111221134, 4.5033455727) and (1.0998073, 4.4836982229) .. (1.0997166538, 
    4.483524524).. controls (1.0989311058, 4.4821722385) and (1.0971483433, 
    4.4817795969) .. (1.0957962167, 4.4825273094) -- cycle;
  
  
  
      \end{scope}
      \begin{scope}[shift={(0.8731702173, -3.5564220104)}]
        \path[fill=c005093,even odd rule] (1.2011580096, 4.5647786477).. controls 
    (1.1984160267, 4.5806340421) and (1.2022081408, 4.5887089931) .. 
    (1.2024272158, 4.5890111738).. controls (1.2185017117, 4.6251334133) and 
    (1.2594203708, 4.6280567417) .. (1.2611577838, 4.6281322537).. controls 
    (1.2732287633, 4.6305796496) and (1.2787278633, 4.6220816562) .. 
    (1.2798080248, 4.6171490821).. controls (1.2799905873, 4.6111891308) and 
    (1.2785464915, 4.606739924) .. (1.2754117081, 4.6036504373).. controls 
    (1.2699578254, 4.5982041217) and (1.262162539, 4.5985063023) .. (1.2618603583,
     4.5985063023).. controls (1.2616857333, 4.5985063023) and (1.230179785, 
    4.596368469) .. (1.2296963119, 4.5767211721).. controls (1.2296963119, 
    4.57519532) and (1.2283064556, 4.573979295) .. (1.2268637092, 4.573979295).. 
    controls (1.2252547779, 4.573979295) and (1.2239932446, 4.5752783992) .. 
    (1.2240762973, 4.5768495479).. controls (1.2246882785, 4.6015579535) and 
    (1.2601983252, 4.6040356971) .. (1.2616864477, 4.6040356971).. controls 
    (1.262079354, 4.6042103221) and (1.2678278767, 4.6040356971) .. (1.2713555133,
     4.6075331448).. controls (1.2734102675, 4.6096180615) and (1.274324059, 
    4.6126697921) .. (1.2742334392, 4.6163711277).. controls (1.2741434808, 
    4.6167714423) and (1.2718767954, 4.6246122106) .. (1.2618603583, 
    4.6226481821).. controls (1.2611578896, 4.6226481821) and (1.2222029415, 
    4.6195964515) .. (1.2075259219, 4.5866544506).. controls (1.2074359635, 
    4.5864348465) and (1.2041645229, 4.5796747952) .. (1.206732701, 
    4.5657758623).. controls (1.2069152635, 4.564174604) and (1.2059545615, 
    4.5628147515) .. (1.2043834656, 4.5625201379).. controls (1.2028197781, 
    4.5622105754) and (1.2014676779, 4.5632149338) .. (1.201158036, 4.5647786213) 
    -- cycle;
  
  
  
      \end{scope}
      \begin{scope}[shift={(0.8731702173, -3.57524685)}]
        \path[fill=c005093,even odd rule] (1.209483574, 4.5887027225).. controls 
    (1.209483574, 4.5891787079) and (1.2090906677, 4.6332322827) .. (1.2670509012,
     4.6355965729).. controls (1.2685238367, 4.6355965729) and (1.2699214717, 
    4.6344107104) .. (1.2699214717, 4.6328016998).. controls (1.2699214717, 
    4.6312681748) and (1.2686977738, 4.6300143144) .. (1.2672624621, 
    4.6298859121).. controls (1.2149751863, 4.6279671538) and (1.2149751863, 
    4.5904476231) .. (1.2149751863, 4.5887933688).. controls (1.2150572071, 
    4.5873053521) and (1.2137967321, 4.5860513858) .. (1.2123161767, 
    4.5860513858).. controls (1.2108810767, 4.5860513858) and (1.209483574, 
    4.5872221671) .. (1.209483574, 4.5887027225) -- cycle;
  
  
  
      \end{scope}
      \begin{scope}[shift={(0.8731702173, -3.5804816313)}]
        \path[fill=c005093,even odd rule] (1.2544937498, 4.5358566202).. controls 
    (1.2267411806, 4.5412576127) and (1.2126231198, 4.5618643946) .. 
    (1.2120112708, 4.5627784242).. controls (1.2111349708, 4.5639190429) and 
    (1.211437125, 4.5657016731) .. (1.2127894104, 4.5665703796).. controls 
    (1.2140583521, 4.5674466796) and (1.215757956, 4.5670992817) .. (1.2167173088,
     4.5657925046).. controls (1.2167173088, 4.5655729004) and (1.230178224, 
    4.5462734513) .. (1.2556722569, 4.5413483383).. controls (1.276369685, 
    4.5374203342) and (1.2991592794, 4.5436975738) .. (1.3234673969, 
    4.5603008923).. controls (1.3240415427, 4.5607768777) and (1.342601111, 
    4.5764659285) .. (1.3568097917, 4.5747210279).. controls (1.3614251569, 
    4.5742450425) and (1.3649528729, 4.5721979612) .. (1.36756648, 4.5684060058)..
     controls (1.3686541821, 4.5674466267) and (1.3708371269, 4.5640929535) .. 
    (1.3702252779, 4.5596059642).. controls (1.3695680529, 4.5538122506) and 
    (1.3645599402, 4.5482374533) .. (1.3558504654, 4.5430404546).. controls 
    (1.3554575592, 4.5427383004) and (1.3411130944, 4.532238549) .. (1.3177643296,
     4.5489248937).. controls (1.3173714233, 4.5491894771) and (1.3173714233, 
    4.5495744458) .. (1.3171901837, 4.5499747604).. controls (1.2950197673, 
    4.5370275867) and (1.2739296181, 4.532155364) .. (1.2544937498, 4.5358566996) 
    -- cycle(1.3526703854, 4.5476556875).. controls (1.3619010894, 4.5531095702) 
    and (1.3642578125, 4.5574680515) .. (1.3646506394, 4.5602176015).. controls 
    (1.3650358727, 4.5631333892) and (1.3634721852, 4.5648330725) .. 
    (1.3632531102, 4.5649236923).. controls (1.361425104, 4.5675298381) and 
    (1.3591589742, 4.56884405) .. (1.3562356194, 4.5692292833).. controls 
    (1.3443384446, 4.5705889771) and (1.3271761144, 4.5559420935) .. 
    (1.3270024154, 4.5558967704).. controls (1.3252122446, 4.5546807454) and 
    (1.3236863396, 4.5538949329) .. (1.3220774083, 4.5528526069).. controls 
    (1.3411129092, 4.5397769515) and (1.35218678, 4.5474441531) .. (1.3526702531, 
    4.5476556875) -- cycle;
  
  
  
      \end{scope}
      \begin{scope}[shift={(0.8731702173, -3.5424784688)}]
        \path[fill=c005093,even odd rule] (1.3074302867, 4.5055275915).. controls 
    (1.3061687533, 4.5066076206) and (1.3061687533, 4.5084358121) .. 
    (1.3071736408, 4.5095311077).. controls (1.3082235075, 4.5106111369) and 
    (1.3100063494, 4.5108378848) .. (1.311101645, 4.5097427744).. controls 
    (1.319901766, 4.5019019796) and (1.3432580979, 4.4856990815) .. (1.3687823198,
     4.5074766181).. controls (1.369175226, 4.5078695244) and (1.3834744735, 
    4.5231733419) .. (1.3709275615, 4.5397690404).. controls (1.3705724906, 
    4.5404262654) and (1.3619008248, 4.55240644) .. (1.3375927337, 4.5368910088)..
     controls (1.3346240558, 4.5347156046) and (1.2604232475, 4.483123204) .. 
    (1.2246862942, 4.5297526827).. controls (1.2236816713, 4.531014216) and 
    (1.2239912338, 4.5326684704) .. (1.2252149317, 4.5338015221).. controls 
    (1.2263404692, 4.5345418262) and (1.2281307194, 4.5343227513) .. 
    (1.2290900456, 4.5331898054).. controls (1.2614655471, 4.4908734848) and 
    (1.3336192477, 4.5411288135) .. (1.3344123627, 4.541597126).. controls 
    (1.3637361335, 4.5602473671) and (1.3750668883, 4.5436063719) .. 
    (1.3756334142, 4.5430020635).. controls (1.3908316365, 4.5225992754) and 
    (1.3728007056, 4.5036543944) .. (1.3724910373, 4.5033447525).. controls 
    (1.3426460902, 4.4779337723) and (1.3151048702, 4.4986387675) .. 
    (1.3074302073, 4.5055278296) -- cycle;
  
  
  
      \end{scope}
      \begin{scope}[shift={(0.8731702173, -3.5050565958)}]
        \path[fill=c005093,even odd rule] (1.2935390267, 4.4637322619).. controls 
    (1.2923605725, 4.4647368848) and (1.2921867413, 4.4665272408) .. 
    (1.2932366079, 4.4677055892).. controls (1.2942335579, 4.4688840433) and 
    (1.295978776, 4.4688840433) .. (1.2971495838, 4.4679251933).. controls 
    (1.3446477954, 4.4259261875) and (1.3791609715, 4.4634003685) .. (1.380505584,
     4.4651378079).. controls (1.4019432104, 4.49027677) and (1.3771062438, 
    4.5124094304) .. (1.3760185152, 4.5134140798).. controls (1.3532740058, 
    4.5304553896) and (1.3275610567, 4.5081037335) .. (1.326390249, 
    4.5071443808).. controls (1.2707037652, 4.4653492365) and (1.2367570885, 
    4.496901434) .. (1.2364549079, 4.4972489377).. controls (1.2353595329, 
    4.4982988044) and (1.2352764538, 4.5000816727) .. (1.2363649496, 
    4.501252454).. controls (1.2373243288, 4.502347829) and (1.2391976581, 
    4.5023930727) .. (1.2402929538, 4.501252454).. controls (1.2415544871, 
    4.5000816727) and (1.2712633854, 4.4728123385) .. (1.3228557596, 
    4.5114499719).. controls (1.3229457179, 4.5114499719) and (1.3529727248, 
    4.5377221973) .. (1.3795922423, 4.5176441852).. controls (1.3800303923, 
    4.5173346227) and (1.4097469635, 4.4908358081) .. (1.3847363773, 
    4.4614364723).. controls (1.3842982273, 4.4609454056) and (1.34513193, 
    4.4179871) .. (1.2935395823, 4.463732791) -- cycle;
  
  
  
      \end{scope}
      \begin{scope}[shift={(0.8731702173, -3.5748682313)}]
        \path[fill=c005093,even odd rule] (1.315582549, 4.4796556542).. controls 
    (1.3118812133, 4.4804337938) and (1.3091390452, 4.4824881775) .. 
    (1.3074319006, 4.4854115323).. controls (1.3042517942, 4.4914092396) and 
    (1.3071752548, 4.4992651685) .. (1.3073551715, 4.4995597821).. controls 
    (1.3116155452, 4.5153396644) and (1.3065619506, 4.5230067337) .. 
    (1.3065619506, 4.5232182946).. controls (1.3056933235, 4.5244041571) and 
    (1.3059954777, 4.5261415435) .. (1.3073551715, 4.52701025).. controls 
    (1.3086167048, 4.5279243854) and (1.3103086888, 4.5275013167) .. 
    (1.3112680415, 4.526232375).. controls (1.3114955831, 4.5259302208) and 
    (1.3176737098, 4.5165182244) .. (1.3126655706, 4.4978225544).. controls 
    (1.3120160185, 4.4958133085) and (1.3107015685, 4.4911147319) .. 
    (1.3124830081, 4.4880629748).. controls (1.313276229, 4.486491879) and 
    (1.314711329, 4.4855324998) .. (1.3168414894, 4.48501135).. controls 
    (1.3169764269, 4.4849663708) and (1.3217741694, 4.48413505) .. (1.3264801544, 
    4.4876626867).. controls (1.332100169, 4.4918097394) and (1.3354919421, 
    4.5007382363) .. (1.3364059717, 4.5133833352).. controls (1.3365885342, 
    4.5149845935) and (1.3378939883, 4.5161252917) .. (1.33946519, 4.5159064019)..
     controls (1.34090029, 4.5159064019) and (1.3421165267, 4.5145996248) .. 
    (1.341995665, 4.5129831529).. controls (1.34090029, 4.4982985662) and 
    (1.3367986398, 4.4882822085) .. (1.3297056369, 4.4829568075).. controls 
    (1.3228543838, 4.4780015321) and (1.3159729417, 4.4794368702) .. 
    (1.3155801148, 4.4796558658) -- cycle;
  
  
  
      \end{scope}
      \begin{scope}[shift={(0.8731702173, -3.5857920833)}]
        \path[fill=c005093,even odd rule] (1.3138453213, 4.4937969454).. controls 
    (1.3126668671, 4.49479416) and (1.3126668671, 4.496455955) .. (1.3136706962, 
    4.4976267627).. controls (1.3187619146, 4.5036018481) and (1.3224708173, 
    4.5224711112) .. (1.3233395238, 4.5292241775).. controls (1.3236869217, 
    4.5307500296) and (1.3250013717, 4.5317546525) .. (1.3266028152, 
    4.5315355775).. controls (1.328173911, 4.531353015) and (1.3291332902, 
    4.5298814025) .. (1.3289142152, 4.5284385237).. controls (1.328521309, 
    4.5256512177) and (1.3249560221, 4.5022571298) .. (1.3177648323, 
    4.4940159675).. controls (1.3168054531, 4.4928830217) and (1.3150152823, 
    4.4927544342) .. (1.3138444746, 4.4937963633) -- cycle;
  
  
  
      \end{scope}
      \begin{scope}[shift={(0.8731702173, -3.5064009438)}]
        \path[fill=c005093,even odd rule] (1.4126274558, 4.4163625319).. controls 
    (1.3825180577, 4.4308808808) and (1.3890066729, 4.4790287769) .. 
    (1.3892635569, 4.481204234).. controls (1.3894910985, 4.4826922506) and 
    (1.3908346527, 4.4837725444) .. (1.3924059338, 4.4835611423).. controls 
    (1.3938939504, 4.4833865173) and (1.3949742442, 4.4818617235) .. 
    (1.3948533296, 4.480426359).. controls (1.3946257879, 4.4799429652) and 
    (1.3886138196, 4.4341064956) .. (1.4146743783, 4.4216049331).. controls 
    (1.4176354892, 4.4208572206) and (1.4281805902, 4.4191953727) .. (1.431927249,
     4.4233424519).. controls (1.4328866281, 4.424649229) and (1.4333171052, 
    4.4261750546) .. (1.4333171052, 4.4280484369).. controls (1.4333171052, 
    4.4307979869) and (1.4323201552, 4.4338497175) .. (1.431368449, 
    4.4365917004).. controls (1.4291853719, 4.439024015) and (1.4163287915, 
    4.4549322731) .. (1.4158528854, 4.486001024).. controls (1.4158528854, 
    4.4862206281) and (1.4153316562, 4.5102260623) .. (1.4014251827, 
    4.5196833554).. controls (1.3938940562, 4.5245631452) and (1.3835679508, 
    4.5245631452) .. (1.370575454, 4.5194637512).. controls (1.3691779248, 
    4.5188594429) and (1.3675237233, 4.5196833554) .. (1.3670025206, 
    4.5210348471).. controls (1.3663077248, 4.5223794596) and (1.367092479, 
    4.5241696304) .. (1.3684452935, 4.5246531035).. controls (1.3832960121, 
    4.5305374369) and (1.3954575321, 4.5303562767) .. (1.4046580471, 
    4.5241697098).. controls (1.4210346442, 4.5131940258) and (1.421427471, 
    4.487095711) .. (1.421427471, 4.4861363583).. controls (1.4219561085, 
    4.4551961156) and (1.4357646333, 4.4400280558) .. (1.4358930356, 4.4398090602)
     -- (1.4362859419, 4.4394235623) -- (1.4364605669, 4.4389401685).. controls 
    (1.4398598013, 4.4303967992) and (1.4397766427, 4.4239005904) .. 
    (1.4360676606, 4.4197309687).. controls (1.4294052406, 4.4121469521) and 
    (1.4144185527, 4.4158935844) .. (1.4126283554, 4.4163618969) -- cycle;
  
  
  
      \end{scope}
      \begin{scope}[shift={(0.8731702173, -3.5333752146)}]
        \path[fill=c005093,even odd rule] (1.4328850142, 4.4469626527).. controls 
    (1.42360142, 4.448352509) and (1.4162440983, 4.4526733135) .. (1.4111452865, 
    4.4597285604).. controls (1.4042562773, 4.4692765792) and (1.4027378865, 
    4.4819064117) .. (1.4027378865, 4.4915828327).. controls (1.4027378865, 
    4.4989477215) and (1.4036065135, 4.5044772221) .. (1.403863424, 
    4.5057462167).. controls (1.403296951, 4.5162006979) and (1.4005549417, 
    4.523693989) .. (1.3956298288, 4.528188519).. controls (1.390750039, 
    4.5324035696) and (1.3856060098, 4.532101389) .. (1.3854397721, 4.532101389)..
     controls (1.3839062471, 4.5320246598) and (1.3825164173, 4.533196764) .. 
    (1.3825164173, 4.534631864).. controls (1.3822968131, 4.5362407952) and 
    (1.3834757965, 4.5375098956) .. (1.3850471569, 4.5376005419).. controls 
    (1.385349311, 4.5378122085) and (1.3924951248, 4.5382048502) .. (1.3991575448,
     4.5324035167).. controls (1.4053894083, 4.5270025242) and (1.4087886427, 
    4.5179908423) .. (1.4093627356, 4.5057461638) -- (1.409484444, 4.5054363367) 
    -- (1.4093627356, 4.5051267742).. controls (1.4093627356, 4.5048780658) and 
    (1.4046567506, 4.4783335577) .. (1.4156701906, 4.4630824452).. controls 
    (1.4197643531, 4.4571980854) and (1.4259583548, 4.4536628023) .. 
    (1.4337160704, 4.4524466979).. controls (1.435151435, 4.4523699688) and 
    (1.4362843808, 4.4508756021) .. (1.4360653058, 4.4493572112).. controls 
    (1.4358906808, 4.44774828) and (1.4345016183, 4.4467889008) .. (1.4328851729, 
    4.4469627321) -- cycle;
  
  
  
      \end{scope}
      \begin{scope}[shift={(0.8731702173, -3.5034852354)}]
        \path[fill=c005093,even odd rule] (1.2921904983, 4.4876174694).. controls 
    (1.29070989, 4.4876174694) and (1.2894030865, 4.4888866756) .. (1.2894030865, 
    4.4904501779).. controls (1.2894030865, 4.4914095571) and (1.2895327323, 
    4.5177268675) .. (1.2983317156, 4.5476625667).. controls (1.2988151094, 
    4.5485311937) and (1.3193387856, 4.5867912137) .. (1.3012928, 4.6038249565).. 
    controls (1.3001220188, 4.6049203315) and (1.3001220188, 4.6066575592) .. 
    (1.3010811333, 4.6077981779).. controls (1.3021611625, 4.6089237154) and 
    (1.3039138419, 4.6089237154) .. (1.3050393265, 4.6079278237).. controls 
    (1.3262654185, 4.5878347833) and (1.3045255056, 4.5472710098) .. 
    (1.3036417973, 4.5453976275).. controls (1.2951135885, 4.5169424838) and 
    (1.2950229423, 4.49067047) .. (1.2950229423, 4.4904514479).. controls 
    (1.2950229423, 4.4888877604) and (1.2938444881, 4.4876187394) .. 
    (1.2921903396, 4.4876187394) -- cycle;
  
  
  
      \end{scope}
      \begin{scope}[shift={(0.8731702173, -3.4967777833)}]
        \path[fill=c005093,even odd rule] (1.3046463144, 4.4892262948).. controls 
    (1.3036416915, 4.4930635733) and (1.3032564581, 4.4971653029) .. 
    (1.3032564581, 4.5014709998).. controls (1.3032564581, 4.5245704477) and 
    (1.3163698696, 4.5508425673) .. (1.3177673458, 4.5535921173).. controls 
    (1.3198144271, 4.5622109988) and (1.3206831335, 4.5696666131) .. 
    (1.3206831335, 4.5760268525).. controls (1.3206831335, 4.600689935) and 
    (1.3081814652, 4.6089311238) .. (1.3074336733, 4.6093239506).. controls 
    (1.3065650463, 4.6098904235) and (1.30617214, 4.6107590506) .. (1.30617214, 
    4.6116276777).. controls (1.30617214, 4.6122393944) and (1.30617214, 
    4.6127230527) .. (1.3065650463, 4.6132442819).. controls (1.3073582671, 
    4.6145058152) and (1.3091410296, 4.6148987215) .. (1.3103116785, 
    4.614112909).. controls (1.3114070535, 4.6135011923) and (1.3344158023, 
    4.5989071196) .. (1.323168206, 4.5520207569) -- (1.3229486019, 4.5516278506) 
    -- (1.3228586435, 4.5513256965).. controls (1.3227766227, 4.551016134) and 
    (1.3036417708, 4.5139420294) .. (1.3100095773, 4.4907141792).. controls 
    (1.310492971, 4.4892261625) and (1.3095338565, 4.4876625544) .. (1.308181571, 
    4.4872621604).. controls (1.3065649669, 4.4868846) and (1.3050391942, 
    4.487662475) .. (1.3046463673, 4.4892261625) -- cycle;
  
  
  
      \end{scope}
      \begin{scope}[shift={(0.8731702173, -3.4974731083)}]
        \path[fill=c005093,even odd rule] (1.31658741, 4.4987297313).. controls 
    (1.3165053892, 4.4996438667) and (1.3146310808, 4.5206055077) .. (1.325606659,
     4.5481843779).. controls (1.3259162215, 4.5485846925) and (1.3463039813, 
    4.5950707406) .. (1.3151068546, 4.6180342456).. controls (1.3138453213, 
    4.6188197935) and (1.3136717546, 4.620610229) .. (1.3145403817, 
    4.6219547092).. controls (1.3154090088, 4.6231331633) and (1.3171917448, 
    4.6233519738) .. (1.3183776602, 4.6225664258).. controls (1.3532761225, 
    4.5968533708) and (1.3310982448, 4.5464016419) .. (1.3308338996, 
    4.545835116).. controls (1.3202510425, 4.5199482563) and (1.3221697215, 
    4.4995605758) .. (1.3221697215, 4.499250934).. controls (1.3222993673, 
    4.4978155694) and (1.3212103423, 4.496418199) .. (1.3196392465, 
    4.4961992033).. controls (1.3181512298, 4.4960325158) and (1.3166781621, 
    4.4972038262) .. (1.3165875158, 4.4987296783) -- cycle;
  
  
  
      \end{scope}
      \begin{scope}[shift={(0.8731702173, -3.6150700812)}]
        \path[fill=c005093,even odd rule] (1.3293065658, 4.6650027879).. controls 
    (1.3277883867, 4.6653880213) and (1.3267834992, 4.6668760379) .. 
    (1.3270025742, 4.6683490263).. controls (1.3273954804, 4.6699201221) and 
    (1.3287475013, 4.6708795012) .. (1.3304395646, 4.6704868596).. controls 
    (1.3329700396, 4.6700108742) and (1.3405088125, 4.6692179179) .. (1.343870185,
     4.6717105575).. controls (1.3451317183, 4.6725791846) and (1.3469143485, 
    4.6722770304) .. (1.3477377583, 4.6709702533).. controls (1.3487423813, 
    4.66970872) and (1.348349475, 4.6679185227) .. (1.3470882063, 4.6670498163).. 
    controls (1.3409922856, 4.6628574404) and (1.3305303431, 4.6648214954) .. 
    (1.3293066452, 4.665002735) -- cycle;
  
  
  
      \end{scope}
      \begin{scope}[shift={(0.8731702173, -3.6178265104)}]
        \path[fill=c005093,even odd rule] (1.3317032675, 4.683259911).. controls 
    (1.3314916008, 4.6848386798) and (1.3327078904, 4.6861832658) .. 
    (1.3341506633, 4.6864022615) -- (1.3419990252, 4.6874068844).. controls 
    (1.3434719606, 4.6876635302) and (1.3449073252, 4.686583501) .. (1.3451339144,
     4.6850124052).. controls (1.3453006019, 4.6835696323) and (1.3442576144, 
    4.6820891562) .. (1.3426486831, 4.6818776219) -- (1.3349288294, 
    4.6809182427).. controls (1.3333651419, 4.6806986385) and (1.3319677185, 
    4.6816963823) .. (1.3317033733, 4.6832598052) -- cycle;
  
  
  
      \end{scope}
      \begin{scope}[shift={(0.8731702173, -3.6184313479)}]
        \path[fill=c005093,even odd rule] (1.3327536633, 4.6992522538).. controls 
    (1.3327086842, 4.7007328621) and (1.3338413654, 4.7021302854) .. 
    (1.3354126729, 4.7021302854) -- (1.3432610348, 4.7026967583).. controls 
    (1.3448247223, 4.7027787792) and (1.3460936375, 4.7017373792) .. 
    (1.3462144727, 4.7001284479).. controls (1.346304431, 4.6985949229) and 
    (1.3451344435, 4.6972049873) .. (1.3435631096, 4.6970312883) -- (1.335805394, 
    4.696646055).. controls (1.3342342981, 4.6964185133) and (1.3328896063, 
    4.6976808404) .. (1.3327536633, 4.6992522008) -- cycle;
  
  
  
      \end{scope}
      \begin{scope}[shift={(0.8731702173, -3.350136175)}]
        \path[fill=c005093,even odd rule] (1.4723622763, 4.3133517206).. controls 
    (1.452533581, 4.3177103077) and (1.4382040652, 4.3429851333) .. (1.4376375394,
     4.3439897827) -- (1.437547581, 4.344375016) -- (1.4374179352, 4.3446395994)..
     controls (1.4305742492, 4.3702620081) and (1.3915363277, 4.3754438992) .. 
    (1.3911359338, 4.375534519).. controls (1.3501416038, 4.3787977046) and 
    (1.3411146819, 4.4081216606) .. (1.3407294221, 4.4094359783).. controls 
    (1.3341500548, 4.456971946) and (1.3775919658, 4.4722683552) .. (1.3779847927,
     4.4723590015).. controls (1.4152853806, 4.4903067738) and (1.4238286442, 
    4.5128774519) .. (1.4238664002, 4.5130511508).. controls (1.4449565494, 
    4.5585023077) and (1.4741067006, 4.574493645) .. (1.475285049, 4.5752792723) 
    -- (1.4810787625, 4.5783234358).. controls (1.4890102035, 4.5836866988) and 
    (1.4966772729, 4.5858619442) .. (1.5037401133, 4.5844721408).. controls 
    (1.5121475133, 4.5828179658) and (1.5164607773, 4.576623779) .. (1.5166797994,
     4.57645008).. controls (1.519905229, 4.5707469863) and (1.5206908827, 
    4.5656481744) .. (1.5188628765, 4.5611612115).. controls (1.5155014775, 
    4.5527084883) and (1.505001779, 4.5497399163) .. (1.5046089521, 4.54964927).. 
    controls (1.4727546798, 4.5392779737) and (1.4490960615, 4.4893927442) .. 
    (1.4488317163, 4.4889545942).. controls (1.4345022004, 4.4604994504) and 
    (1.4038640325, 4.4473407423) .. (1.4025118794, 4.4467742165) -- (1.4022097252,
     4.4466842581) -- (1.4019451419, 4.4465546123).. controls (1.382682946, 
    4.4433367498) and (1.3792458498, 4.4236441033) .. (1.3791627971, 
    4.4227678033).. controls (1.3775086221, 4.4180995744) and (1.3778107763, 
    4.4144887792) .. (1.3797747783, 4.4113539958).. controls (1.3850472892, 
    4.4029391081) and (1.4010765942, 4.40093735) .. (1.4012502931, 4.40093735).. 
    controls (1.4458780402, 4.3994114979) and (1.4641354544, 4.3716059327) .. 
    (1.4648682179, 4.370435125).. controls (1.4727543094, 4.3533107625) and 
    (1.4803836492, 4.3495640244) .. (1.4803836492, 4.3495640244).. controls 
    (1.4887910756, 4.3465122673) and (1.494328011, 4.3419799813) .. (1.4965487913,
     4.3357481442).. controls (1.4998119769, 4.3275069554) and (1.4955063329, 
    4.3198398596) .. (1.4952872579, 4.3194546263) -- (1.4950226746, 4.3189712325) 
    -- (1.4945014454, 4.3185785908).. controls (1.4875292777, 4.3135628581) and 
    (1.4799906635, 4.3116895817) .. (1.4723613238, 4.3133514031) -- 
    cycle(1.4735404658, 4.3188431477).. controls (1.4793795025, 4.3174911269) and 
    (1.4850901633, 4.3188432535) .. (1.4905817756, 4.3227636906).. controls 
    (1.491057761, 4.3235492385) and (1.492062384, 4.3259438235) .. (1.492062384, 
    4.329380814).. controls (1.492062384, 4.3308312598) and (1.4918427798, 
    4.3324324387) .. (1.4913599152, 4.3338752381).. controls (1.4896223965, 
    4.3384981704) and (1.4853017506, 4.3418520817) .. (1.4782465038, 
    4.3443750425).. controls (1.4777252746, 4.3445496675) and (1.4687060521, 
    4.3486958735) .. (1.4599965508, 4.3676936183).. controls (1.4599965508, 
    4.3676936183) and (1.442517091, 4.3939202296) .. (1.4007295402, 4.395363029)..
     controls (1.3998609131, 4.3954926748) and (1.3817317954, 4.3976217769) .. 
    (1.375069349, 4.4082951479).. controls (1.3721006975, 4.4130087265) and 
    (1.3717078706, 4.4182359935) .. (1.3736718198, 4.4242035119).. controls 
    (1.3736718198, 4.4240209494) and (1.377811279, 4.4481262902) .. (1.4009485094,
     4.45212978) -- (1.4003367927, 4.4519630925).. controls (1.4006389469, 
    4.4520398217) and (1.4303629, 4.4649858577) .. (1.4439144615, 4.4913863531).. 
    controls (1.4447830885, 4.4934411073) and (1.4688419685, 4.5437718158) .. 
    (1.5029548829, 4.5549665219).. controls (1.5031295079, 4.5550564802) and 
    (1.5114075002, 4.5574969969) .. (1.5137643292, 4.5630793348).. controls 
    (1.5148067875, 4.5660404192) and (1.5141948063, 4.5695680558) .. 
    (1.5117550833, 4.573269365).. controls (1.5115804583, 4.5737527588) and 
    (1.502645294, 4.5863450204) .. (1.4839118944, 4.5734969067) -- (1.477944376, 
    4.5701731579).. controls (1.4776422219, 4.5700435121) and (1.4492702102, 
    4.554355546) .. (1.4290109321, 4.5109137144).. controls (1.4287913279, 
    4.5101281665) and (1.4199011427, 4.4863867071) .. (1.3801153235, 
    4.4671623467).. controls (1.3782042381, 4.4666032821) and (1.3403369921, 
    4.4531348262) .. (1.3462138642, 4.4105692946).. controls (1.3464334683, 
    4.4098668258) and (1.354545805, 4.3840328298) .. (1.3917104235, 
    4.3810188552).. controls (1.393712261, 4.3808521677) and (1.4350690702, 
    4.3753535175) .. (1.442736166, 4.3461204194) -- (1.4425165619, 4.3468228881)..
     controls (1.4426462077, 4.3465133256) and (1.4562492835, 4.3225070977) .. 
    (1.4735398838, 4.3188435181) -- cycle;
  
  
  
      \end{scope}
      \begin{scope}[shift={(0.8731702173, -3.3587695292)}]
        \path[fill=c005093,even odd rule] (1.4894027923, 4.3254298704).. controls 
    (1.4882243381, 4.3261246663) and (1.459814491, 4.3418065733) .. (1.4474414367,
     4.3681239896).. controls (1.4474414367, 4.3679123229) and (1.4331043538, 
    4.3900978735) .. (1.3943685071, 4.3935801871).. controls (1.3941938821, 
    4.3935801871) and (1.3541899933, 4.3951060392) .. (1.3546205498, 
    4.4314019777).. controls (1.3548480915, 4.4318324548) and (1.3572870471, 
    4.4611638721) .. (1.3895341463, 4.4722302023).. controls (1.3897537504, 
    4.4722302023) and (1.4286099031, 4.4887427423) .. (1.4389812258, 4.520121135) 
    -- (1.4390711842, 4.5203407392) -- (1.4392907883, 4.5206352204).. controls 
    (1.4407713967, 4.5231733683) and (1.4773316275, 4.5811260083) .. 
    (1.5088385019, 4.5903567123).. controls (1.5102736019, 4.5907496185) and 
    (1.5117541837, 4.5898807269) .. (1.5122829535, 4.5884003831).. controls 
    (1.5123649744, 4.5880908206) and (1.5123649744, 4.5879169894) .. 
    (1.5123649744, 4.5876148352).. controls (1.5123649744, 4.586436381) and 
    (1.5115794265, 4.5852581915) .. (1.5102726494, 4.5848729581).. controls 
    (1.4809942017, 4.5764202085) and (1.4446906167, 4.5187620498) .. 
    (1.4440787677, 4.5176743213).. controls (1.4324912348, 4.4844301133) and 
    (1.3931888888, 4.4675701754) .. (1.3915346344, 4.4670036496).. controls 
    (1.3628151454, 4.4571988527) and (1.3603299406, 4.4321881606) .. 
    (1.3603299406, 4.4311835113).. controls (1.3603299406, 4.4309639071) and 
    (1.3603299406, 4.4306622821) .. (1.3603299406, 4.4304810425).. controls 
    (1.3603299406, 4.4004698577) and (1.3933249375, 4.3992007573) .. 
    (1.3946242269, 4.3991101375).. controls (1.4364571008, 4.3953634788) and 
    (1.4517535365, 4.3717428281) .. (1.4523200358, 4.3706626667).. controls 
    (1.4639075688, 4.345954261) and (1.4918034627, 4.3305217767) .. (1.4920603202,
     4.330438724).. controls (1.4930649431, 4.329872251) and (1.4935409285, 
    4.3288676281) .. (1.4935409285, 4.3278630052).. controls (1.4935409285, 
    4.3274777719) and (1.4934509702, 4.326903626) .. (1.4931027785, 
    4.3265260656).. controls (1.4925815494, 4.3251663719) and (1.4907990515, 
    4.3247282219) .. (1.4894015488, 4.3254306906) -- cycle;
  
  
  
      \end{scope}
      \begin{scope}[shift={(0.8731702173, -3.5274914104)}]
        \path[fill=c005093,even odd rule] (1.4394641698, 4.562520429).. controls 
    (1.4315780783, 4.5858767873) and (1.4313666233, 4.6029558267) .. 
    (1.4385955427, 4.6134933871).. controls (1.445348609, 4.6229053571) and 
    (1.4557575819, 4.6232604279) .. (1.4562485956, 4.6232604279).. controls 
    (1.4701852846, 4.624129055) and (1.4758960248, 4.6104114675) .. (1.4768553775,
     4.6031296315).. controls (1.4764550629, 4.5987257477) and (1.4756769233, 
    4.5935438567) .. (1.4717489985, 4.592410805).. controls (1.4678361285, 
    4.5913609383) and (1.4643010306, 4.5960668175) .. (1.4633417044, 
    4.5974642937).. controls (1.4625635648, 4.5988165792) and (1.4599953338, 
    4.600123224) .. (1.4585072113, 4.6005537804).. controls (1.4554554806, 
    4.6009996033) and (1.4530157048, 4.6003791554) .. (1.4517541715, 
    4.5988996054).. controls (1.4488308167, 4.5955910173) and (1.4500092444, 
    4.588218535) .. (1.4505381465, 4.5856502775).. controls (1.4524569048, 
    4.5802945817) and (1.4554555865, 4.5766385692) .. (1.4598141735, 
    4.5747198902).. controls (1.468614321, 4.5709279348) and (1.4799902931, 
    4.5764195735) .. (1.4802018275, 4.5764195735).. controls (1.505914909, 
    4.5952889425) and (1.4891378385, 4.6209869427) .. (1.4883976137, 
    4.6220369417).. controls (1.4721493925, 4.6504089004) and (1.4375455702, 
    4.635195544) .. (1.4360650148, 4.6347272315).. controls (1.4075268183, 
    4.6216062529) and (1.4150580506, 4.5758530477) .. (1.4151939935, 
    4.5754602208).. controls (1.415328931, 4.5739722042) and (1.414415854, 
    4.5725369719) .. (1.4127995144, 4.5722347912).. controls (1.4113567415, 
    4.5719326371) and (1.4098760538, 4.5730203392) .. (1.4097023548, 
    4.5745009475).. controls (1.4091811256, 4.5766387808) and (1.4013403573, 
    4.624914841) .. (1.4337158323, 4.6398108827).. controls (1.4341087385, 
    4.639900841) and (1.4742797117, 4.6576831429) .. (1.4931037575, 
    4.6249978937).. controls (1.4932334033, 4.6248232688) and (1.5135443281, 
    4.5938839521) .. (1.4829514569, 4.5717135621).. controls (1.4819545069, 
    4.5712754121) and (1.4689164752, 4.5648697438) .. (1.4576311229, 
    4.5695757287).. controls (1.4518448706, 4.57206096) and (1.447614871, 
    4.5768501035) .. (1.445303471, 4.5841319396).. controls (1.4449182377, 
    4.5853858) and (1.4420400737, 4.5964596973) .. (1.4474410398, 4.6025631585).. 
    controls (1.4501829962, 4.6057055354) and (1.454496366, 4.6067404796) .. 
    (1.4598142, 4.6058717731).. controls (1.4613022167, 4.6054788669) and 
    (1.4660384965, 4.6038321002) .. (1.4680478217, 4.6005085367).. controls 
    (1.4687050467, 4.5995944013) and (1.4695736737, 4.5987257742) .. 
    (1.4703592217, 4.5982500533).. controls (1.4705788258, 4.5987260387) and 
    (1.4709712029, 4.5996854179) .. (1.4711447696, 4.6029560383).. controls 
    (1.4711447696, 4.6029560383) and (1.4688407779, 4.6185999777) .. 
    (1.4564601829, 4.617640625).. controls (1.4559767892, 4.617640625) and 
    (1.4482265612, 4.6173307979) .. (1.4432108285, 4.6101849842).. controls 
    (1.4370242881, 4.601377296) and (1.4375455967, 4.5853859587) .. (1.444782189, 
    4.5644846427).. controls (1.4453034181, 4.5629135469) and (1.4445176056, 
    4.561387589) .. (1.4431277494, 4.5608588192).. controls (1.4415111452, 
    4.5603828337) and (1.4399853725, 4.561115465) .. (1.4394641698, 4.5625204025) 
    -- cycle;
  
  
  
      \end{scope}
      \begin{scope}[shift={(0.8731702173, -3.5517163958)}]
        \path[fill=c005093,even odd rule] (1.4238294379, 4.594284399).. controls 
    (1.4236945004, 4.5947603844) and (1.4109653433, 4.642258596) .. (1.4436883221,
     4.656157529).. controls (1.4441717158, 4.6563771331) and (1.4713652998, 
    4.6650406233) .. (1.4849168348, 4.6402944881).. controls (1.4853020681, 
    4.6391237069) and (1.49040088, 4.6234269831) .. (1.4830359119, 4.6138487754)..
     controls (1.4793345763, 4.6089236625) and (1.4733216554, 4.6065668335) .. 
    (1.4650880338, 4.6068690142).. controls (1.4586823919, 4.6070436392) and 
    (1.4466568148, 4.6107970183) .. (1.4462186648, 4.623124776).. controls 
    (1.4460440398, 4.6247337073) and (1.4472235523, 4.6260027283) .. 
    (1.4488324835, 4.626048025).. controls (1.4504035794, 4.6261829625) and 
    (1.4517557325, 4.6249150792) .. (1.4517557325, 4.6233513652).. controls 
    (1.4521486387, 4.6130631217) and (1.4646954185, 4.6124437056) .. 
    (1.4652619444, 4.6124437056).. controls (1.4715844277, 4.6122770181) and 
    (1.4760714171, 4.6138489077) .. (1.4786396746, 4.6171497965).. controls 
    (1.4840859902, 4.6243863094) and (1.4797726204, 4.6379378444) .. 
    (1.4797726204, 4.6381191369).. controls (1.4687062902, 4.6581592871) and 
    (1.4466568148, 4.6512777392) .. (1.4456974621, 4.6509226683).. controls 
    (1.4176351452, 4.6390330606) and (1.4291018694, 4.5962862365) .. 
    (1.4291849221, 4.595893436).. controls (1.4295701554, 4.5942845048) and 
    (1.4287920158, 4.592841679) .. (1.4272285929, 4.5924564456).. controls 
    (1.4256574971, 4.591965379) and (1.4242147242, 4.592841679) .. (1.4238294644, 
    4.5942844519) -- cycle;
  
  
  
      \end{scope}
      \begin{scope}[shift={(0.8731702173, -3.5439656917)}]
        \path[fill=c005093,even odd rule] (1.3630824804, 4.7047428606).. controls 
    (1.3628628763, 4.7059213148) and (1.3585503002, 4.7314153742) .. 
    (1.3684457169, 4.743788455).. controls (1.3718449513, 4.7481016131) and 
    (1.3764149933, 4.7502846373) .. (1.3819066321, 4.7505414948).. controls 
    (1.3820362779, 4.7505414948) and (1.3909257752, 4.751931351) .. (1.3960245871,
     4.7461377169).. controls (1.3995521973, 4.7417866969) and (1.3999450242, 
    4.7355548333) .. (1.397248285, 4.727019031).. controls (1.3970286808, 
    4.7262258102) and (1.3919303717, 4.7142002596) .. (1.4023848529, 
    4.7077039715).. controls (1.4029060821, 4.7073639819) and (1.4115399654, 
    4.7015249981) .. (1.4250915004, 4.7142002596).. controls (1.4262699546, 
    4.7153332054) and (1.4280072617, 4.7152424533) .. (1.4291025838, 
    4.7140706137).. controls (1.43005429, 4.7129829117) and (1.43005429, 
    4.711200255) .. (1.4290126254, 4.7102409023).. controls (1.4117521877, 
    4.6941589452) and (1.399681314, 4.7027853673) .. (1.3991601113, 
    4.7031706271).. controls (1.387882326, 4.7102410081) and (1.3895365804, 
    4.723180509) .. (1.3919311654, 4.729057381).. controls (1.3939705738, 
    4.7351608423) and (1.3938951675, 4.739995309) .. (1.3917115613, 4.742389894)..
     controls (1.3887957735, 4.7458420185) and (1.3825185075, 4.7449656127) .. 
    (1.3825185075, 4.7449656127).. controls (1.3782053229, 4.7446634585) and 
    (1.3751536981, 4.7431754419) .. (1.3728044363, 4.7403427069).. controls 
    (1.3645632475, 4.7299713842) and (1.3685743308, 4.7060938496) .. 
    (1.3685743308, 4.7059201506).. controls (1.3756371448, 4.6778578073) and 
    (1.403782435, 4.6837421671) .. (1.4048701635, 4.6840443477).. controls 
    (1.4064412594, 4.684437254) and (1.4080049733, 4.6835683623) .. (1.4084053673,
     4.681959431).. controls (1.4087075215, 4.6805166581) and (1.4076121465, 
    4.6790360763) .. (1.4062222902, 4.6786432494).. controls (1.405829384, 
    4.6786432494) and (1.3714899598, 4.6711045558) .. (1.3630825598, 4.7047415642)
     -- cycle;
  
  
  
      \end{scope}
      \begin{scope}[shift={(0.8731702173, -3.5622163854)}]
        \path[fill=c005093,even odd rule] (1.3817308165, 4.7185894325).. controls 
    (1.3817308165, 4.7186793908) and (1.3662229258, 4.7443931338) .. 
    (1.3891864308, 4.764002701).. controls (1.3903572121, 4.7649620802) and 
    (1.3920945456, 4.764788249) .. (1.3930992215, 4.7636023865).. controls 
    (1.3941038444, 4.7624316052) and (1.3939678485, 4.7606413021) .. 
    (1.3927970673, 4.7596819494).. controls (1.3734971419, 4.7431694094) and 
    (1.3860440275, 4.7222982823) .. (1.3866105533, 4.7214220088).. controls 
    (1.3938923894, 4.7087089912) and (1.4173316681, 4.7132790333) .. 
    (1.4176338752, 4.7133696796).. controls (1.4191218919, 4.7136718338) and 
    (1.4205571242, 4.7126219671) .. (1.4208593048, 4.7111413587).. controls 
    (1.421161459, 4.709622915) and (1.420156836, 4.7081726808) .. (1.4187291444, 
    4.7079158233).. controls (1.4175054465, 4.7075684254) and (1.390923685, 
    4.7023865344) .. (1.3817307371, 4.7185894325) -- cycle;
  
  
  
      \end{scope}
      \begin{scope}[shift={(0.8731702173, -3.4658374083)}]
        \path[fill=c005093,even odd rule] (1.3050379242, 4.5190333006).. controls 
    (1.3039123867, 4.5194262069) and (1.2801178783, 4.5270553615) .. 
    (1.2733648121, 4.5479189215).. controls (1.2686965831, 4.5623843802) and 
    (1.2733648121, 4.57929742) .. (1.2873468242, 4.598166789) -- (1.2875743658, 
    4.5983863931) -- (1.2879595992, 4.598643039).. controls (1.2946295863, 
    4.603787174) and (1.3116632231, 4.6193405729) .. (1.3095330892, 
    4.6297496252).. controls (1.30857371, 4.6344556102) and (1.3036412681, 
    4.637983141) .. (1.2946295863, 4.6402568115).. controls (1.2930584904, 
    4.6409063635) and (1.284212861, 4.6454765379) .. (1.2833441281, 
    4.6527206179).. controls (1.282822899, 4.6566863781) and (1.2846964135, 
    4.6604707665) .. (1.2888358727, 4.663922891).. controls (1.2904069685, 
    4.6648370265) and (1.3154099613, 4.6775120763) .. (1.3338412067, 
    4.6601308298).. controls (1.3344153525, 4.6593831173) and (1.3460481292, 
    4.6463602727) .. (1.3448244577, 4.6277024646).. controls (1.3438727515, 
    4.6139772306) and (1.3362886554, 4.6009997621) .. (1.3222990762, 
    4.5890572642).. controls (1.3218987617, 4.58875511) and (1.3074332235, 
    4.5788597463) .. (1.3059073715, 4.5657010381).. controls (1.305212311, 
    4.5595975769) and (1.3074332235, 4.55384916) .. (1.3124867123, 4.5481385256)..
     controls (1.3136725748, 4.5469677444) and (1.3134460915, 4.5452227379) .. 
    (1.3124867123, 4.5442180621).. controls (1.3112706873, 4.5430472808) and 
    (1.3094424429, 4.5432134392) .. (1.3083471473, 4.5444376663).. controls 
    (1.3021606069, 4.5511907325) and (1.2995470262, 4.5586388063) .. 
    (1.3004232998, 4.566404089).. controls (1.3021608185, 4.5819951381) and 
    (1.31815205, 4.5930614683) .. (1.3188545188, 4.593544915).. controls 
    (1.3314014044, 4.6041806888) and (1.3383282754, 4.6158963858) .. 
    (1.3391666871, 4.6279673919).. controls (1.3403374683, 4.6442684769) and 
    (1.329920849, 4.6561657575) .. (1.3297093146, 4.6563394565).. controls 
    (1.314450635, 4.6706689723) and (1.2927563098, 4.6597311238) .. (1.2918800363,
     4.6592628113).. controls (1.2900065217, 4.6576009633) and (1.2888358727, 
    4.6557729306) .. (1.2888358727, 4.6538995748).. controls (1.2888358727, 
    4.6537646373) and (1.2888358727, 4.6536720331) .. (1.2889178935, 
    4.6535066685).. controls (1.2893107998, 4.6497600098) and (1.2945001785, 
    4.6463986373) .. (1.2965849894, 4.6454770406).. controls (1.30717544, 
    4.6429164031) and (1.3136716223, 4.6379837496) .. (1.3150162083, 
    4.6308907467).. controls (1.3151987708, 4.6299766113) and (1.3154091146, 
    4.6289719883) .. (1.3154091146, 4.6279673125).. controls (1.3154091146, 
    4.6132449698) and (1.2952329156, 4.5972083094) .. (1.2916675229, 
    4.5945494056).. controls (1.2820364515, 4.5814813173) and (1.2772398202, 
    4.5697502744) .. (1.2772398202, 4.5593336285).. controls (1.2772398202, 
    4.5559872844) and (1.277723214, 4.5527619342) .. (1.2787278369, 
    4.5497177706).. controls (1.2845140892, 4.5316718115) and (1.306397221, 
    4.5245258125) .. (1.3065634588, 4.5243974367).. controls (1.3077794838, 
    4.5239971221) and (1.3085727046, 4.5228186679) .. (1.3085727046, 
    4.5216478867).. controls (1.3085727046, 4.5214732617) and (1.3085727046, 
    4.5211644929) .. (1.3084377671, 4.5208623387).. controls (1.3079617817, 
    4.5194272387) and (1.3063983588, 4.5185586117) .. (1.3050386121, 4.5190343325)
     -- cycle;
  
  
  
      \end{scope}
      \begin{scope}[shift={(0.8731702173, -3.4894512063)}]
        \path[fill=c005093,even odd rule] (1.3079627871, 4.5552851331).. controls 
    (1.307313235, 4.5555047373) and (1.2905286769, 4.5611619258) .. (1.285429865, 
    4.5754536856).. controls (1.2823857015, 4.583740171) and (1.2840400088, 
    4.5931522469) .. (1.290317116, 4.6032667915).. controls (1.2905287827, 
    4.6035234373) and (1.2958918075, 4.6125881417) .. (1.3086199592, 
    4.6237450917).. controls (1.3087416675, 4.6237450917) and (1.3285241665, 
    4.6431734194) .. (1.324732211, 4.6611287588).. controls (1.3226928027, 
    4.6708429358) and (1.3138471469, 4.678638381) .. (1.2985507113, 
    4.6843490419).. controls (1.2970626946, 4.6849610231) and (1.296284555, 
    4.686532119) .. (1.2969870238, 4.6878767579).. controls (1.2973722571, 
    4.6892742871) and (1.2990264321, 4.690059835) .. (1.3004693373, 
    4.6896669288).. controls (1.3177675046, 4.6832991223) and (1.3277913235, 
    4.6740683125) .. (1.3302312317, 4.6622617575).. controls (1.3346275483, 
    4.6413000106) and (1.3132805417, 4.6205120421) .. (1.3124874267, 
    4.6196433356).. controls (1.3002880713, 4.6090075883) and (1.2952344767, 
    4.6003888127) .. (1.2951136679, 4.6003434896).. controls (1.2897051877, 
    4.5917548235) and (1.2883530346, 4.5839970021) .. (1.2907098635, 
    4.5772439358).. controls (1.2948417556, 4.5656564029) and (1.3097001471, 
    4.5606859669) .. (1.30983609, 4.5604669448).. controls (1.31127119, 
    4.5599909594) and (1.3121022463, 4.5584201281) .. (1.3116185879, 
    4.5569394404).. controls (1.3111050317, 4.5555040758) and (1.3095336712, 
    4.5547185279) .. (1.3079625754, 4.5552850008) -- cycle;
  
  
  
      \end{scope}
      \begin{scope}[shift={(0.8731702173, -3.6006571688)}]
        \path[fill=c005093,even odd rule] (1.2282165767, 4.6784253385).. controls 
    (1.2273479496, 4.6797397885) and (1.2277482642, 4.6814770692) .. 
    (1.2289190454, 4.6823533427).. controls (1.256195735, 4.7009961225) and 
    (1.2744079319, 4.6996514306) .. (1.2751935856, 4.6996514306).. controls 
    (1.2768101898, 4.6994768056) and (1.2778448165, 4.6980803348) .. 
    (1.2777240606, 4.6964714035).. controls (1.2775414981, 4.6949455515) and 
    (1.2762812877, 4.6938578494) .. (1.2746724358, 4.6939409285).. controls 
    (1.2746724358, 4.6939409285) and (1.2574574271, 4.6951569535) .. 
    (1.2321446338, 4.6776474106).. controls (1.2308378567, 4.6768618627) and 
    (1.2290929031, 4.6772545044) .. (1.2282166296, 4.6784252856) -- cycle;
  
  
  
      \end{scope}
      \begin{scope}[shift={(0.8731702173, -3.5986934313)}]
        \path[fill=c005093,even odd rule] (1.2243851454, 4.6871889206).. controls 
    (1.2235995975, 4.6885335331) and (1.2242105204, 4.6902330842) .. 
    (1.2255635996, 4.6910187379).. controls (1.2267420538, 4.6915852108) and 
    (1.2559750725, 4.7074935219) .. (1.2765817485, 4.7111042112).. controls 
    (1.2781152735, 4.7114137737) and (1.2795051033, 4.7104546592) .. 
    (1.2798073104, 4.708838055).. controls (1.2801168729, 4.7073122029) and 
    (1.2791125146, 4.7059298608) .. (1.2775411542, 4.7056201925).. controls 
    (1.2578484019, 4.7022209581) and (1.2285245517, 4.6863202671) .. 
    (1.2282148833, 4.6860105988).. controls (1.2268628625, 4.6853158029) and 
    (1.2252161487, 4.6858280362) .. (1.2243851719, 4.6871887883) -- cycle;
  
  
  
      \end{scope}
      \begin{scope}[shift={(0.8731702173, -3.5951657417)}]
        \path[fill=c005093,even odd rule] (1.218245939, 4.6945449458).. controls 
    (1.2174529827, 4.6959046396) and (1.2178530327, 4.6974683006) .. 
    (1.2192053181, 4.6983370071).. controls (1.2205499306, 4.6992133071) and 
    (1.25401324, 4.7190418965) .. (1.2846060848, 4.7222220029).. controls 
    (1.2860866931, 4.7223886904) and (1.2875746569, 4.7214287821) .. 
    (1.2876577096, 4.7198198508).. controls (1.2878773137, 4.7181656758) and 
    (1.2867435742, 4.7168966019) .. (1.2850819908, 4.7167303642).. controls 
    (1.2558488662, 4.7138070094) and (1.2224233923, 4.6937668327) .. 
    (1.2220381325, 4.6935931338).. controls (1.2207691908, 4.6927620775) and 
    (1.2191148042, 4.6931549838) .. (1.2182460713, 4.69454484) -- cycle;
  
  
  
      \end{scope}
      \begin{scope}[shift={(0.8731702173, -3.5204278292)}]
        \path[fill=c005093,even odd rule] (1.1718345563, 4.4249510127).. controls 
    (1.1713133271, 4.425774396) and (1.1286496085, 4.5059124544) .. (1.0159697581,
     4.5216847431).. controls (1.0144891498, 4.5219043473) and (1.0134014477, 
    4.5233844265) .. (1.0137036019, 4.5248649819).. controls (1.0137935602, 
    4.526390834) and (1.0152746977, 4.5273954569) .. (1.0168459523, 
    4.5272216256).. controls (1.1323584319, 4.5111396685) and (1.1763214398, 
    4.4283501677) .. (1.1767142667, 4.4275192173).. controls (1.1774998146, 
    4.4261746048) and (1.1769338708, 4.4244297306) .. (1.1755358125, 
    4.4237272619).. controls (1.1741459563, 4.423032466) and (1.172491649, 
    4.4235605744) .. (1.1718344769, 4.4249509598) -- cycle;
  
  
  
      \end{scope}
      \begin{scope}[shift={(0.8731702173, -3.5063104563)}]
        \path[fill=c005093,even odd rule] (1.1906584963, 4.4116566527).. controls 
    (1.1900467796, 4.4126612756) and (1.1242078629, 4.5124921921) .. 
    (1.0140964023, 4.5226898158).. controls (1.0125253065, 4.5229014825) and 
    (1.0113544194, 4.5241627512) .. (1.0114828481, 4.5258623817).. controls 
    (1.011612494, 4.5273959067) and (1.0131370231, 4.5283931212) .. (1.0145799019,
     4.5281814546).. controls (1.1274787744, 4.5178101319) and (1.1946695531, 
    4.4157131385) .. (1.1954476133, 4.4146178429).. controls (1.1963617488, 
    4.4134018179) and (1.1958479279, 4.4116567585) .. (1.1945789863, 
    4.4108333487).. controls (1.1933174529, 4.409911805) and (1.1915348227, 
    4.4103121196) .. (1.1906585227, 4.4116567321) -- cycle;
  
  
  
      \end{scope}
      \begin{scope}[shift={(0.8731702173, -3.4922909792)}]
        \path[fill=c005093,even odd rule] (1.2094832565, 4.3962017581).. controls 
    (1.2088789481, 4.3972516248) and (1.1446491215, 4.5064417269) .. 
    (1.0119593892, 4.5212546894).. controls (1.0103504579, 4.5214372519) and 
    (1.0093081319, 4.5227882144) .. (1.0094818308, 4.5243970662).. controls 
    (1.0095638517, 4.5258398392) and (1.0110907621, 4.5269198683) .. 
    (1.0125259944, 4.5268367892).. controls (1.148259996, 4.5116309998) and 
    (1.213622795, 4.4001221423) .. (1.2142724, 4.3990344137).. controls 
    (1.2150579479, 4.3976368846) and (1.2145896354, 4.3958920633) .. 
    (1.2133130208, 4.3950686535).. controls (1.2118324125, 4.3944191015) and 
    (1.2102688573, 4.3948940285) .. (1.2094832035, 4.3962015994) -- cycle;
  
  
  
      \end{scope}
      \begin{scope}[shift={(0.8731702173, -3.5565585354)}]
        \path[fill=c005093,even odd rule] (1.104733921, 4.5578152642).. controls 
    (1.1032533127, 4.5580348683) and (1.1022489544, 4.5596056996) .. 
    (1.1025963523, 4.5612069315).. controls (1.1048170002, 4.5749397325) and 
    (1.1016822169, 4.5817758779) .. (1.1016822169, 4.5818664977).. controls 
    (1.0868692544, 4.6149445737) and (1.0480427879, 4.6179055787) .. 
    (1.0475593148, 4.6179055787).. controls (1.0372786383, 4.61995266) and 
    (1.0351031813, 4.6120665421) .. (1.0347935394, 4.6111978356).. controls 
    (1.0347935394, 4.6078817598) and (1.035752654, 4.604920649) .. (1.0377545179, 
    4.602964214).. controls (1.0412443985, 4.5993383904) and (1.0471665938, 
    4.5994668721) .. (1.0471665938, 4.5994668721).. controls (1.0489115208, 
    4.5993372263) and (1.0845125048, 4.5968985617) .. (1.0850337075, 
    4.5721522413).. controls (1.0850337075, 4.5706263892) and (1.0838174179, 
    4.5694026912) .. (1.0823746979, 4.5692364535).. controls (1.0807657667, 
    4.5692364535) and (1.0795042333, 4.570490314) .. (1.0794589102, 
    4.5719785423).. controls (1.0790207602, 4.591588136) and (1.0473478863, 
    4.5938543188) .. (1.0471665938, 4.5939373979).. controls (1.0470369479, 
    4.5938553771) and (1.0391066975, 4.5937627729) .. (1.0338341867, 
    4.5989531306).. controls (1.030571001, 4.6020879404) and (1.0291280958, 
    4.6064918242) .. (1.0293017948, 4.6118852496).. controls (1.0304347406, 
    4.6172862685) and (1.0359717819, 4.6259201519) .. (1.0483448627, 
    4.6234878373).. controls (1.0496971481, 4.6233528998) and (1.0906158867, 
    4.6204285131) .. (1.1067810288, 4.5842232208).. controls (1.1069106746, 
    4.5839586375) and (1.1107014658, 4.5761104344) .. (1.10795181, 4.5602097433)..
     controls (1.107777185, 4.5586460558) and (1.1062973704, 4.5576869413) .. 
    (1.1047337358, 4.5578152642) -- cycle;
  
  
  
      \end{scope}
      \begin{scope}[shift={(0.8731702173, -3.5741131104)}]
        \path[fill=c005093,even odd rule] (1.0968842627, 4.5800759035).. controls 
    (1.0953131669, 4.5800759035) and (1.0940968508, 4.5813826806) .. 
    (1.0940968508, 4.582863289).. controls (1.0940968508, 4.584389141) and 
    (1.0934396258, 4.6223922242) .. (1.0407671696, 4.6250435873).. controls 
    (1.0391960738, 4.6250435873) and (1.0380176196, 4.6264334435) .. 
    (1.0381459425, 4.6279668362).. controls (1.0381909217, 4.6294397717) and 
    (1.0395810425, 4.6307542217) .. (1.0410691915, 4.6306258458).. controls 
    (1.0992409065, 4.627702491) and (1.099716786, 4.5832561158) .. (1.099716786, 
    4.582863289).. controls (1.099716786, 4.5812921931) and (1.0983721735, 
    4.5800759035) .. (1.0968841833, 4.5800759035) -- cycle;
  
  
  
      \end{scope}
      \begin{scope}[shift={(0.8731702173, -3.5804363875)}]
        \path[fill=c005093,even odd rule] (0.9919196358, 4.5451326208).. controls 
    (0.9917450108, 4.5447397146) and (0.9917450108, 4.5442639938) .. 
    (0.9914362421, 4.5440449188).. controls (0.9679968575, 4.5274038971) and 
    (0.9536597746, 4.5379866748) .. (0.9530479256, 4.5383795017).. controls 
    (0.9445498794, 4.5433952344) and (0.9395418196, 4.54896982) .. (0.9388468915, 
    4.5548467185).. controls (0.9384087415, 4.5591674437) and (0.9404558227, 
    4.5627329158) .. (0.9418078965, 4.5638735346).. controls (0.9441648048, 
    4.567310525) and (0.9478207379, 4.5693576062) .. (0.952179325, 4.5699694287)..
     controls (0.9666447837, 4.5716236037) and (0.9850834904, 4.555896691) .. 
    (0.9858162275, 4.5552395454).. controls (1.0099581073, 4.5389913242) and 
    (1.0328386654, 4.532578221) .. (1.0534453415, 4.5365515483).. controls 
    (1.0790224535, 4.5414313381) and (1.0923547813, 4.5607311577) .. (1.092445401,
     4.5609501798).. controls (1.0933595365, 4.562347709) and (1.0951044371, 
    4.5626046194) .. (1.0962752183, 4.5617357277).. controls (1.0976727475, 
    4.5609501798) and (1.0979749017, 4.5591674173) .. (1.0971062746, 
    4.5578076971).. controls (1.0965019662, 4.5570295575) and (1.082467011, 
    4.5364153937) .. (1.0545783931, 4.531105021).. controls (1.0351878215, 
    4.5272677425) and (1.0140976723, 4.5321927231) .. (0.991919821, 4.5451324356) 
    -- cycle(0.9530025231, 4.5643872496).. controls (0.9500414388, 4.56404726) and
     (0.9478206585, 4.5627330746) .. (0.946121081, 4.5603460038).. controls 
    (0.9457358477, 4.5600364413) and (0.9440361644, 4.558208435) .. (0.9444213977,
     4.5553756735).. controls (0.9448595477, 4.5527167433) and (0.9472087831, 
    4.5484034) .. (0.9562355198, 4.5429949198).. controls (0.9569303156, 
    4.5425266073) and (0.9679967517, 4.5348066213) .. (0.9871304394, 
    4.548010679).. controls (0.9855593435, 4.5489700581) and (0.9839881683, 
    4.549755606) .. (0.9824245602, 4.5508432817).. controls (0.9820242456, 
    4.5510999275) and (0.9647713221, 4.5657468904) .. (0.9530024967, 4.5643872496)
     -- cycle;
  
  
  
      \end{scope}
      \begin{scope}[shift={(0.8731702173, -3.5425226542)}]
        \path[fill=c005093,even odd rule] (0.93670628, 4.4985012106).. controls 
    (0.9364041258, 4.4989015252) and (0.9183657073, 4.5178540263) .. 
    (0.9336998725, 4.53833992).. controls (0.9340927788, 4.53889872) and 
    (0.9453327285, 4.5554944715) .. (0.974490341, 4.5369347179).. controls 
    (0.9754420473, 4.5363756533) and (1.0475582035, 4.4861279181) .. 
    (1.0799790017, 4.5283913485).. controls (1.0809836246, 4.5297057985) and 
    (1.0827662813, 4.5298796298) .. (1.083944656, 4.5289202506).. controls 
    (1.0850776019, 4.5280061152) and (1.0854252644, 4.5261327594) .. 
    (1.0845113935, 4.5249543846).. controls (1.0487366577, 4.4784230927) and 
    (0.9745282029, 4.5300531965) .. (0.9714387163, 4.5322286535).. controls 
    (0.9472968365, 4.5475704123) and (0.9385797681, 4.5355826706) .. 
    (0.9382775875, 4.535106791).. controls (0.9257307019, 4.5183298265) and 
    (0.9399320271, 4.5031165494) .. (0.9405437438, 4.5024291354).. controls 
    (0.9658639983, 4.4810368056) and (0.9891750071, 4.4970281165) .. 
    (0.9980128842, 4.5049897729).. controls (0.9991913383, 4.5060396396) and 
    (1.000936239, 4.5059491521) .. (1.0019408883, 4.5047781062).. controls 
    (1.0028925946, 4.503607325) and (1.0028925946, 4.501816916) .. (1.0017212842, 
    4.5007291875).. controls (0.9939635421, 4.4937948552) and (0.9664224279, 
    4.4732259087) .. (0.9367058567, 4.4985008667) -- cycle;
  
  
  
      \end{scope}
      \begin{scope}[shift={(0.8731702173, -3.5050864937)}]
        \path[fill=c005093,even odd rule] (0.9244609929, 4.4565328965).. controls 
    (0.8993220308, 4.4860833358) and (0.9290386021, 4.5126574244) .. (0.929431429,
     4.5128386904).. controls (0.9560962431, 4.5328788671) and (0.9861679117, 
    4.5067806581) .. (0.9863416106, 4.506478451).. controls (1.0377980154, 
    4.4679541917) and (1.0675599098, 4.4952761779) .. (1.0688289044, 
    4.4965301442).. controls (1.0699166065, 4.4975347671) and (1.071699369, 
    4.4975347671) .. (1.0728324206, 4.4964084358).. controls (1.0738370435, 
    4.49527549) and (1.0738370435, 4.49353051) .. (1.0727503998, 4.4925711573).. 
    controls (1.0723574935, 4.4921330073) and (1.0385392985, 4.460460001) .. 
    (0.9828074652, 4.5020284504).. controls (0.981636684, 4.5033352275) and 
    (0.9558329827, 4.5255962373) .. (0.9328694512, 4.5084340129).. controls 
    (0.9319553158, 4.507648465) and (0.9072998269, 4.4853798617) .. (0.9286468335,
     4.4602861963).. controls (0.9300366898, 4.4587151004) and (0.964466866, 
    4.421157655) .. (1.0119575106, 4.4631189312).. controls (1.0131359648, 
    4.4641613896) and (1.0148808654, 4.4641158813) .. (1.0158402181, 
    4.4628913896).. controls (1.0169279202, 4.4617206083) and (1.016844841, 
    4.4600585488) .. (1.0156655931, 4.4590237104).. controls (0.9639901663, 
    4.4131722919) and (0.9248163019, 4.4561458904) .. (0.924461231, 4.4565310708) 
    -- cycle;
  
  
  
      \end{scope}
      \begin{scope}[shift={(0.8731702173, -3.5749441667)}]
        \path[fill=c005093,even odd rule] (0.9793711363, 4.4783260171).. controls 
    (0.9723083223, 4.483432396) and (0.9681688631, 4.493629914) .. (0.9671264577, 
    4.5081785312).. controls (0.9669968119, 4.5097496271) and (0.968168916, 
    4.511094319) .. (0.9696495244, 4.511275585).. controls (0.9712206202, 
    4.511275585) and (0.9726181758, 4.5102257183) .. (0.972701255, 4.5086620308)..
     controls (0.9737058779, 4.495888556) and (0.977014519, 4.4871337317) .. 
    (0.9826345335, 4.4829489494).. controls (0.9873405185, 4.4793306929) and 
    (0.9921371498, 4.4802900456) .. (0.9923108487, 4.480380639).. controls 
    (0.9944862529, 4.4809018681) and (0.9959668612, 4.4817781681) .. 
    (0.9966240069, 4.4832511035).. controls (0.9982781819, 4.4864010475) and 
    (0.997145236, 4.4913184875) .. (0.9965340485, 4.492844366).. controls 
    (0.9915258829, 4.511887434) and (0.997629159, 4.5211258108) .. (0.9978860694, 
    4.5216016904).. controls (0.9986716173, 4.5226893925) and (1.0004617881, 
    4.5230822988) .. (1.0017233479, 4.5222964863).. controls (1.0029848813, 
    4.5214201863) and (1.0033777875, 4.5197281758) .. (1.0025088958, 
    4.5183760492).. controls (1.0025088958, 4.5181643825) and (0.9976292119, 
    4.5106257948) .. (1.0019424229, 4.4944531915).. controls (1.0019424229, 
    4.4944531915) and (1.0047751315, 4.4867030429) .. (1.0016328604, 
    4.4806902015).. controls (1.0001522521, 4.4777215235) and (0.9973197023, 
    4.4756745481) .. (0.993489885, 4.4747982746).. controls (0.9931803225, 
    4.4747982746) and (0.9863440183, 4.4732345871) .. (0.9793718242, 4.4783259113)
     -- cycle;
  
  
  
      \end{scope}
      \begin{scope}[shift={(0.8731702173, -3.5856629667)}]
        \path[fill=c005093,even odd rule] (0.9913065698, 4.4890447377).. controls 
    (0.9841607031, 4.4972859) and (0.9805878492, 4.520861254) .. (0.9802403454, 
    4.523474861).. controls (0.9799757621, 4.5249554694) and (0.9810258933, 
    4.5263075696) .. (0.9825062371, 4.5267003965).. controls (0.9840773329, 
    4.5269200006) and (0.9854674273, 4.5257788527) .. (0.9857317725, 
    4.524260409).. controls (0.9867289871, 4.5175073427) and (0.9904378633, 
    4.4986832969) .. (0.9955744313, 4.4927535081).. controls (0.9965338104, 
    4.4916205623) and (0.9964507312, 4.4898377469) .. (0.9952270333, 
    4.488795315).. controls (0.9940940875, 4.4877906921) and (0.9923112456, 
    4.4878737713) .. (0.9913065698, 4.4890440233) -- cycle;
  
  
  
      \end{scope}
      \begin{scope}[shift={(0.8731702173, -3.506492225)}]
        \path[fill=c005093,even odd rule] (0.8730518163, 4.4149804015).. controls 
    (0.8692976169, 4.4191198606) and (0.8692145377, 4.4256539048) .. 
    (0.8727422538, 4.4341972742) -- (0.8729618579, 4.4348089908) -- (0.8732185038,
     4.4351942242).. controls (0.8734381079, 4.4351942242) and (0.8871703533, 
    4.4505359565) .. (0.887691556, 4.4812949596).. controls (0.887691556, 
    4.4823903346) and (0.8880844623, 4.5085791898) .. (0.9045138438, 
    4.5195549002).. controls (0.9136614685, 4.5256583615) and (0.9258155273, 
    4.5258320604) .. (0.9406284898, 4.5199478065).. controls (0.9419881835, 
    4.5193360898) and (0.9427663231, 4.5177647294) .. (0.9421995856, 
    4.5162917146).. controls (0.9415876044, 4.5148489417) and (0.9400694252, 
    4.5141993896) .. (0.9385813292, 4.5146298667).. controls (0.9256038871, 
    4.5199477006) and (0.9152326702, 4.51994778) .. (0.9076863831, 4.5148494708)..
     controls (0.8937496942, 4.5054375008) and (0.8932661417, 4.4815221308) .. 
    (0.8932661417, 4.4812955417).. controls (0.8927901563, 4.4503174371) and 
    (0.8799337081, 4.434197724) .. (0.8777960071, 4.4317578158).. controls 
    (0.8767990571, 4.4290989121) and (0.8757941696, 4.4260018583) .. 
    (0.8757941696, 4.4232598754).. controls (0.8757941696, 4.4214697046) and 
    (0.876270155, 4.4199059377) .. (0.8772747779, 4.4186822663).. controls 
    (0.8809761135, 4.4145805631) and (0.8915212146, 4.4161591996) .. 
    (0.8949204754, 4.4171564142).. controls (0.920497561, 4.4294464423) and 
    (0.91444699, 4.4751996475) .. (0.9143563438, 4.4755849073).. controls 
    (0.9141367396, 4.4771560031) and (0.9152326438, 4.4785535852) .. 
    (0.9167129875, 4.4788632271).. controls (0.9182311667, 4.4791198729) and 
    (0.9197572304, 4.4779946) .. (0.9198478767, 4.4763705875).. controls 
    (0.9202407829, 4.4744595021) and (0.9266009429, 4.426175716) .. (0.8968843717,
     4.4118386331).. controls (0.8948372904, 4.4112795685) and (0.8797145802, 
    4.4075329362) .. (0.8730521338, 4.41498101) -- cycle;
  
  
  
      \end{scope}
      \begin{scope}[shift={(0.8731702173, -3.5331564042)}]
        \path[fill=c005093,even odd rule] (0.8730527423, 4.444296076).. controls 
    (0.8728781173, 4.4457840927) and (0.8739213694, 4.4473553737) .. 
    (0.8755303006, 4.4475290727).. controls (0.8831596404, 4.4485260227) and 
    (0.8893462867, 4.452152005) .. (0.8934404227, 4.4579382308).. controls 
    (0.90451432, 4.4734083654) and (0.8998082292, 4.4998163485) .. (0.8997176094, 
    4.5001185292) -- (0.899627651, 4.5004206833) -- (0.8997176094, 4.5008209979)..
     controls (0.900329326, 4.5130657029) and (0.9037286663, 4.5219036858) .. 
    (0.910043609, 4.5273952981).. controls (0.9165776267, 4.5332796579) and 
    (0.9238594627, 4.532675376) .. (0.9241616698, 4.532675376).. controls 
    (0.9256044427, 4.5325854177) and (0.9268205735, 4.5312249302) .. 
    (0.9267299802, 4.5297067246).. controls (0.9266003344, 4.5281808725) and 
    (0.925332451, 4.5270101706) .. (0.9237688958, 4.5270931704).. controls 
    (0.9235492917, 4.5271751913) and (0.9183679033, 4.5273953246) .. (0.913525949,
     4.5230820871).. controls (0.9086461592, 4.518768929) and (0.9059041763, 
    4.5111923735) .. (0.9052923538, 4.500730431).. controls (0.9055119579, 
    4.4994162456) and (0.9063877288, 4.4940225821) .. (0.9063877288, 
    4.4866123438).. controls (0.9063877288, 4.4769812723) and (0.9048542038, 
    4.4642153646) .. (0.8979727352, 4.4548033946).. controls (0.8928361938, 
    4.4475291256) and (0.8855996808, 4.4432537235) .. (0.8762707635, 
    4.4418637879).. controls (0.8747523198, 4.4416441838) and (0.8732191123, 
    4.4428231671) .. (0.873052901, 4.4442961025) -- cycle;
  
  
  
      \end{scope}
      \begin{scope}[shift={(0.8731702173, -3.5032963229)}]
        \path[fill=c005093,even odd rule] (1.0140978046, 4.4853363377).. controls 
    (1.0140978046, 4.4855929835) and (1.0140157837, 4.5120465015) .. 
    (1.0053883033, 4.541068171).. controls (1.0046405908, 4.5422466252) and 
    (0.9828100317, 4.5828482606) .. (1.0039907742, 4.6029866242).. controls 
    (1.0051692283, 4.6040364908) and (1.0069518585, 4.6038629242) .. 
    (1.0080395871, 4.6028569783).. controls (1.0090894537, 4.6017240325) and 
    (1.0089989663, 4.5998429773) .. (1.0078199829, 4.598845895).. controls 
    (0.9898722371, 4.5816760771) and (1.0102671142, 4.5435445387) .. 
    (1.0104788073, 4.5432423581).. controls (1.0197246454, 4.5128307794) and 
    (1.0197625073, 4.4865134425) .. (1.0197625073, 4.4853350677).. controls 
    (1.0197625073, 4.4838092156) and (1.018500974, 4.4825401946) .. (1.0169297723,
     4.4825401946).. controls (1.0154039202, 4.4825401946) and (1.0140971696, 
    4.4838091363) .. (1.0140971696, 4.4853350677) -- cycle;
  
  
  
      \end{scope}
      \begin{scope}[shift={(0.8731702173, -3.4966872958)}]
        \path[fill=c005093,even odd rule] (1.0010666521, 4.4823000587).. controls 
    (0.9995029646, 4.482692965) and (0.9986721729, 4.4842188171) .. (0.9990648146,
     4.485835236).. controls (1.0055611292, 4.5089648729) and (0.986427415, 
    4.546137085) .. (0.9863443623, 4.5464392921) -- (0.9860345352, 4.5468396067) 
    -- (0.9859525144, 4.5470062942).. controls (0.9746671621, 4.5940210592) and 
    (0.9977136404, 4.6086301602) .. (0.9986729931, 4.6092344156).. controls 
    (1.0000705223, 4.6100125552) and (1.0017248296, 4.609619649) .. (1.002593536, 
    4.6083202802).. controls (1.0028956902, 4.6078368865) and (1.0028956902, 
    4.6072249052) .. (1.0028956902, 4.606831999).. controls (1.0028956902, 
    4.6058726198) and (1.0025104569, 4.604958749) .. (1.0016341569, 4.604392276)..
     controls (1.000939361, 4.6040901219) and (0.9883923696, 4.5958943356) .. 
    (0.9883923696, 4.5711556352).. controls (0.9883923696, 4.5647878288) and 
    (0.9892609967, 4.5572491352) .. (0.9914365331, 4.5486680098).. controls 
    (0.9927888185, 4.5459637829) and (1.005947315, 4.5195633933) .. (1.005947315, 
    4.4966300508).. controls (1.005947315, 4.492324354) and (1.0055620817, 
    4.4881320046) .. (1.0045119504, 4.4842191346).. controls (1.0041190442, 
    4.4827763617) and (1.0025101129, 4.4819077346) .. (1.0010673929, 4.4823003763)
     -- cycle;
  
  
  
      \end{scope}
      \begin{scope}[shift={(0.8731702173, -3.4974723146)}]
        \path[fill=c005093,even odd rule] (0.9894796219, 4.4914090279).. controls 
    (0.9879990135, 4.4914910488) and (0.9867300983, 4.4929724508) .. 
    (0.9869565552, 4.494415356).. controls (0.9869565552, 4.4946270227) and 
    (0.9888679052, 4.51524116) .. (0.9783226719, 4.5412563956).. controls 
    (0.977975274, 4.5416493019) and (0.9558350465, 4.5920178987) .. (0.9907410494,
     4.6176403075).. controls (0.9919195035, 4.6186449304) and (0.9937020544, 
    4.6182898596) .. (0.9945783279, 4.6170283262).. controls (0.9955755425, 
    4.615857545) and (0.9952281446, 4.6141201056) .. (0.9939666113, 
    4.6132363708).. controls (0.9628072406, 4.5902728394) and (0.9832099494, 
    4.5438775169) .. (0.98351213, 4.5433864767).. controls (0.9945784602, 
    4.5158529296) and (0.9926144581, 4.4948082094) .. (0.992531379, 4.493977259)..
     controls (0.9924414206, 4.4923683277) and (0.9909147748, 4.4913258958) .. 
    (0.9894795425, 4.4914089485) -- cycle;
  
  
  
      \end{scope}
      \begin{scope}[shift={(0.8731702173, -3.6149788)}]
        \path[fill=c005093,even odd rule] (0.962027831, 4.662207809).. controls 
    (0.9607588894, 4.663076436) and (0.9604112269, 4.6648213631) .. (0.9614161144,
     4.6660905165).. controls (0.9622016623, 4.6672991331) and (0.9639918331, 
    4.6676994477) .. (0.9652533929, 4.6667853123).. controls (0.9686979504, 
    4.6643379165) and (0.9761459977, 4.6651311373) .. (0.9788049279, 
    4.6656899373).. controls (0.9803760238, 4.6659095415) and (0.9817206098, 
    4.6651308727) .. (0.9820227904, 4.6633407019).. controls (0.982423105, 
    4.6619884165) and (0.9814184821, 4.6605080727) .. (0.9798473863, 
    4.6602058921).. controls (0.9787144404, 4.6600233296) and (0.9680861544, 
    4.6578868192) .. (0.9620279898, 4.6622077296) -- cycle;
  
  
  
      \end{scope}
      \begin{scope}[shift={(0.8731702173, -3.6177217354)}]
        \path[fill=c005093,even odd rule] (0.9742736473, 4.6758808671) -- 
    (0.9664253913, 4.6770214858).. controls (0.9648617038, 4.6771511317) and 
    (0.9637663817, 4.678539665) .. (0.9639930767, 4.6801109725).. controls 
    (0.9642047433, 4.6815915808) and (0.9656398433, 4.682777364) .. (0.967211045, 
    4.6825959392) -- (0.9750518133, 4.6815082371).. controls (0.97653983, 
    4.6812806954) and (0.9775822883, 4.6798540621) .. (0.9774087217, 
    4.6782828075).. controls (0.9771065675, 4.6768852783) and (0.9757542821, 
    4.6757070888) .. (0.9742738325, 4.6758806554) -- cycle;
  
  
  
      \end{scope}
      \begin{scope}[shift={(0.8731702173, -3.61843955)}]
        \path[fill=c005093,even odd rule] (0.9733964477, 4.6917216035) -- 
    (0.9655103563, 4.6923259119).. controls (0.9639921771, 4.6925084744) and 
    (0.9628968021, 4.6937760931) .. (0.9629872896, 4.6952491608).. controls 
    (0.9630772479, 4.6968202567) and (0.9643771458, 4.6979912496) .. (0.965948374,
     4.6979912496) -- (0.9737891688, 4.697432185).. controls (0.9753602646, 
    4.6973025392) and (0.9764027229, 4.6959515767) .. (0.9763196438, 
    4.6944710742).. controls (0.9762746646, 4.6929073867) and (0.9749676229, 
    4.6916384715) .. (0.9733963948, 4.6917215242) -- cycle;
  
  
  
      \end{scope}
      \begin{scope}[shift={(0.8731702173, -3.3501888271)}]
        \path[fill=c005093,even odd rule] (0.8146597298, 4.313796935) -- 
    (0.814176336, 4.3141972496) -- (0.8137834298, 4.3147108058).. controls 
    (0.8137384506, 4.3149753892) and (0.8094324627, 4.3228160252) .. 
    (0.8125218965, 4.3310949435).. controls (0.8147880527, 4.3371531081) and 
    (0.8203627177, 4.3419044163) .. (0.8289060342, 4.344820204).. controls 
    (0.8287763883, 4.344820204) and (0.8363616221, 4.3485215131) .. (0.8440287179,
     4.3652985042).. controls (0.8449880971, 4.3668696) and (0.863336131, 
    4.3946298156) .. (0.9081678454, 4.3961934237).. controls (0.9080778871, 
    4.3961934237) and (0.9241592885, 4.3981574258) .. (0.9294317994, 
    4.4066100696).. controls (0.9314336369, 4.4097448794) and (0.9316979556, 
    4.4133632417) .. (0.9301268598, 4.4175857535).. controls (0.9299072556, 
    4.4189002035) and (0.9263800952, 4.4385929294) .. (0.9072992977, 4.441901544) 
    -- (0.9069894706, 4.441901544) -- (0.9065965644, 4.4421211481).. controls 
    (0.9052897873, 4.4426045419) and (0.8747498327, 4.4558464879) .. 
    (0.8601937542, 4.4840825037).. controls (0.8600191292, 4.484694485) and 
    (0.8363615162, 4.5345040702) .. (0.804507191, 4.544913149).. controls 
    (0.8041142848, 4.545087774) and (0.7935692896, 4.5480102027) .. (0.7903891831,
     4.556281554).. controls (0.7885159331, 4.5608591631) and (0.789210729, 
    4.5659201131) .. (0.7925193435, 4.5714948046).. controls (0.7926489894, 
    4.5718877108) and (0.796968524, 4.5780289546) .. (0.805375871, 4.5796378858)..
     controls (0.8124386321, 4.5810883317) and (0.8201887806, 4.5789428254) .. 
    (0.8283393231, 4.5733228637) -- (0.8338309619, 4.570407076).. controls 
    (0.835009416, 4.5698406031) and (0.8641141648, 4.5537207312) .. (0.8851136677,
     4.5085264583).. controls (0.8852962302, 4.5080129021) and (0.8939591119, 
    4.4854799006) .. (0.9314334252, 4.467615181).. controls (0.9316080502, 
    4.4675252227) and (0.9750490352, 4.4523187719) .. (0.9685603142, 
    4.4050396617).. controls (0.9680843287, 4.4034685658) and (0.9590576979, 
    4.3740615571) .. (0.9178366465, 4.3707001846).. controls (0.9175800006, 
    4.3707001846) and (0.8785343004, 4.3653898119) .. (0.8717812606, 4.339895779) 
    -- (0.8716992398, 4.3395031373) -- (0.8714796356, 4.3393364498).. controls 
    (0.8709584065, 4.3381127519) and (0.8564928683, 4.3128378469) .. 
    (0.8368757075, 4.30853215).. controls (0.829125506, 4.3069533812) and 
    (0.8216321621, 4.3087438167) .. (0.8146600473, 4.3137971996) -- 
    cycle(0.8356592327, 4.3140165392).. controls (0.8529120504, 4.3178538177) and 
    (0.8664635854, 4.3416330331) .. (0.8665995283, 4.3420334271) -- (0.8662521304,
     4.3413764667).. controls (0.874002279, 4.3704887825) and (0.9155329723, 
    4.3761086912) .. (0.9173157083, 4.3763202521).. controls (0.9545785667, 
    4.3792436069) and (0.9626761131, 4.4051303608) .. (0.9630765071, 
    4.4062105223).. controls (0.9687796008, 4.4483909) and (0.9309124606, 
    4.4619499227) .. (0.9293866085, 4.4623427231).. controls (0.8892608527, 
    4.4815595958) and (0.8803246829, 4.5054748865) .. (0.8800150146, 
    4.5064342392).. controls (0.8598842181, 4.5496191869) and (0.8314366415, 
    4.5653915021) .. (0.8311722698, 4.5655198779) -- (0.8253785563, 
    4.5686162438).. controls (0.806418541, 4.5815558504) and (0.7973992921, 
    4.56900915) .. (0.7971424875, 4.5684416188).. controls (0.7949670833, 
    4.56469496) and (0.7944004517, 4.561257996) .. (0.7954806396, 4.558425261).. 
    controls (0.7978372833, 4.5527599233) and (0.8060257406, 4.5504032002) .. 
    (0.80616171, 4.5501916658).. controls (0.8403652442, 4.5389817992) and 
    (0.864333425, 4.4886962021) .. (0.8653380744, 4.4866566879).. controls 
    (0.878715831, 4.4601655727) and (0.908560884, 4.4473922037) .. (0.908779906, 
    4.4472185048) -- (0.9081679248, 4.4473931298).. controls (0.9313051288, 
    4.4434726927) and (0.93553534, 4.4191571933) .. (0.93553534, 4.418945659).. 
    controls (0.9363208879, 4.4168078256) and (0.9367061213, 4.414632395) .. 
    (0.9367061213, 4.4124116148).. controls (0.9367061213, 4.4093523171) and 
    (0.9359205733, 4.4064366352) .. (0.9341378108, 4.4036417621).. controls 
    (0.9273847446, 4.3928851796) and (0.909346326, 4.3907020496) .. (0.9085606988,
     4.3906189969).. controls (0.8665994225, 4.3891383885) and (0.8491275563, 
    4.3628663219) .. (0.8489538573, 4.3626472998).. controls (0.8403652177, 
    4.3439516562) and (0.831391239, 4.3398047094) .. (0.830953089, 4.3395931485)..
     controls (0.8239431652, 4.337153161) and (0.8194940377, 4.3336256567) .. 
    (0.8178397304, 4.3290933706).. controls (0.8172280138, 4.3275598456) and 
    (0.8169258596, 4.3260870425) .. (0.8169258596, 4.32459892).. controls 
    (0.8169258596, 4.3211543625) and (0.8180966408, 4.3187220479) .. 
    (0.8184895471, 4.3181555221).. controls (0.8240264825, 4.314197329) and 
    (0.8296086617, 4.3127470419) .. (0.8356592327, 4.3140160629) -- cycle;
  
  
  
      \end{scope}
      \begin{scope}[shift={(0.8731702173, -3.358784875)}]
        \path[fill=c005093,even odd rule] (0.815967724, 4.3217439071).. controls 
    (0.8156578969, 4.3222122196) and (0.8155748177, 4.3226125342) .. 
    (0.8155748177, 4.3230885196).. controls (0.8155748177, 4.3241762217) and 
    (0.8160508031, 4.3250448488) .. (0.817055426, 4.3255811592).. controls 
    (0.8173200094, 4.3257478467) and (0.8451178223, 4.3410814827) .. (0.8565769, 
    4.3655784333).. controls (0.8574455271, 4.3669230458) and (0.8726587777, 
    4.3905362352) .. (0.9143557088, 4.3943281906).. controls (0.9159268046, 
    4.394418149) and (0.9487406148, 4.3957783719) .. (0.9487406148, 
    4.4259709021).. controls (0.9487406148, 4.4259709021) and (0.9487406148, 
    4.4262733208) .. (0.9487406148, 4.4265828833).. controls (0.9487406148, 
    4.4274970187) and (0.9464215419, 4.4523714504) .. (0.9178002663, 
    4.4620931681).. controls (0.9159267517, 4.4628334723) and (0.8765792148, 
    4.4795197377) .. (0.8651201635, 4.5129074556).. controls (0.8645084469, 
    4.5138969973) and (0.8281217563, 4.5716381294) .. (0.798797959, 
    4.5801814988).. controls (0.7975290173, 4.5806045675) and (0.7967508777, 
    4.5815261112) .. (0.7967508777, 4.5827498092).. controls (0.7967508777, 
    4.5831427154) and (0.796840836, 4.5833996258) .. (0.796840836, 4.5835279487)..
     controls (0.7972337423, 4.5850990446) and (0.7988879173, 4.5859676717) .. 
    (0.8003609321, 4.5854464425).. controls (0.831830024, 4.5764347606) and 
    (0.8683374969, 4.5183915537) .. (0.8699086985, 4.5158157027) -- (0.870043636, 
    4.5156490152) -- (0.8700806777, 4.5152487006).. controls (0.8805879433, 
    4.4839609277) and (0.9194521658, 4.467493605) .. (0.9197543465, 
    4.4674029852).. controls (0.9517824235, 4.4564725979) and (0.9543582746, 
    4.4270506402) .. (0.9543582746, 4.4267484596).. controls (0.9549625829, 
    4.3902788221) and (0.9150483615, 4.3887076469) .. (0.9146555346, 
    4.3887076469).. controls (0.8760934133, 4.3853159796) and (0.8616732513, 
    4.3632212604) .. (0.8615448754, 4.3629190798).. controls (0.8494285727, 
    4.3371153785) and (0.8208904027, 4.321342984) .. (0.8197120015, 
    4.3206480558).. controls (0.8183144723, 4.3199606683) and (0.8166603237, 
    4.3203382287) .. (0.8159653427, 4.3217434308) -- cycle;
  
  
  
      \end{scope}
      \begin{scope}[shift={(0.8731702173, -3.5274914104)}]
        \path[fill=c005093,even odd rule] (0.8661235429, 4.5560700725).. controls 
    (0.8645069387, 4.5565460579) and (0.8638121429, 4.5582076413) .. 
    (0.8643333721, 4.5595976033).. controls (0.8715623179, 4.580687726) and 
    (0.8720910877, 4.5964601471) .. (0.8659875471, 4.6053962375).. controls 
    (0.8608887352, 4.6125798602) and (0.8531310196, 4.6128442054) .. 
    (0.8527457598, 4.6128442054).. controls (0.8403273294, 4.6137128325) and 
    (0.8379252831, 4.5982503444) .. (0.837842151, 4.5974647171).. controls 
    (0.8381443052, 4.594677411) and (0.8385369469, 4.5938012433) .. (0.8387107781,
     4.5934159835).. controls (0.8395418344, 4.5938012169) and (0.8404557052, 
    4.5946775169) .. (0.8411960094, 4.595757546).. controls (0.8430240156, 
    4.5990358658) and (0.8478207263, 4.6007732788) .. (0.8488254021, 
    4.6009998679).. controls (0.8547022742, 4.6019592471) and (0.8589247067, 
    4.6009099096) .. (0.8616742567, 4.5977215481).. controls (0.8670373873, 
    4.5915425748) and (0.8642047317, 4.5805593238) .. (0.8639404129, 
    4.5794262721).. controls (0.8615459338, 4.5720613833) and (0.8572326698, 
    4.5672269167) .. (0.8514767123, 4.5647870085).. controls (0.84019136, 
    4.5600810235) and (0.8271157046, 4.5664414481) .. (0.826639825, 4.566569771)..
     controls (0.7955710742, 4.5891027196) and (0.8158756225, 4.6199900721) .. 
    (0.8160493744, 4.620382899).. controls (0.8347903146, 4.652841559) and 
    (0.8750067431, 4.6351053475) .. (0.8755279458, 4.6348937867).. controls 
    (0.9079034473, 4.6199901779) and (0.8998058744, 4.5717140119) .. 
    (0.8994054804, 4.5697953065).. controls (0.8992308554, 4.568231619) and 
    (0.8977510408, 4.5672267315) .. (0.8963159144, 4.5674383981).. controls 
    (0.894835306, 4.5676130231) and (0.893747604, 4.5691909981) .. (0.8940497581, 
    4.5706639335).. controls (0.8940497581, 4.571139919) and (0.9015884517, 
    4.616855421) .. (0.8730502552, 4.6297950013).. controls (0.8716979698, 
    4.6304897971) and (0.8370564708, 4.6456126396) .. (0.8207553594, 
    4.6174596765).. controls (0.820279374, 4.6166815369) and (0.814787841, 
    4.6082288667) .. (0.814787841, 4.5974647435).. controls (0.814787841, 
    4.5892235548) and (0.8180586467, 4.5796831031) .. (0.8293818079, 
    4.5713211056).. controls (0.8293818079, 4.5713211056) and (0.8405839752, 
    4.5662677227) .. (0.8494294723, 4.5699690848).. controls (0.8536595777, 
    4.5717140119) and (0.8567038471, 4.5754606177) .. (0.8585393675, 
    4.5810806323).. controls (0.8592341633, 4.5833920323) and (0.8604128821, 
    4.5907570533) .. (0.8574439925, 4.5941109646).. controls (0.8560919717, 
    4.5956820604) and (0.8536596042, 4.5962032896) .. (0.8502150467, 
    4.5956820604).. controls (0.8491273446, 4.5952891542) and (0.84655152, 
    4.5938842167) .. (0.8459018885, 4.5927588115).. controls (0.8448065135, 
    4.5912782031) and (0.8413696025, 4.5866553238) .. (0.83727544, 4.5876599996)..
     controls (0.8334381085, 4.5887098663) and (0.8325694021, 4.5938843754) .. 
    (0.8322143842, 4.597721654).. controls (0.8334380821, 4.6057059588) and 
    (0.8390128265, 4.6194237315) .. (0.8530477552, 4.6184190821).. controls 
    (0.853304401, 4.6184190821) and (0.8637212585, 4.6181994779) .. (0.8704290017,
     4.6086218523).. controls (0.87783924, 4.5980314017) and (0.877574789, 
    4.5810807381) .. (0.8696434538, 4.557815026).. controls (0.8691222246, 
    4.5562439302) and (0.8674680496, 4.5555414615) .. (0.8661233842, 4.556070099) 
    -- cycle;
  
  
  
      \end{scope}
      \begin{scope}[shift={(0.8731702173, -3.5515952167)}]
        \path[fill=c005093,even odd rule] (0.8818895081, 4.5874022425).. controls 
    (0.8803710644, 4.5877951488) and (0.879495029, 4.5893662446) .. (0.879933179, 
    4.5909297733).. controls (0.8800151998, 4.5913300879) and (0.8913922302, 
    4.6340693979) .. (0.8632468342, 4.6460950544).. controls (0.8625443654, 
    4.6462299919) and (0.8403664877, 4.6533315673) .. (0.8292095377, 
    4.6327248912).. controls (0.8292095377, 4.6325502662) and (0.8277366023, 
    4.6277997783) .. (0.8277366023, 4.6224366477).. controls (0.8277366023, 
    4.6187806352) and (0.8284239898, 4.6150717325) .. (0.8305692579, 
    4.6121559448).. controls (0.8330468163, 4.6089758383) and (0.8376622079, 
    4.6073214781) .. (0.8438563154, 4.6076160917).. controls (0.84441538, 
    4.6076160917) and (0.8569697004, 4.6082354813) .. (0.8574456065, 
    4.6183878083).. controls (0.8574456065, 4.6199514958) and (0.8587523835, 
    4.6211751144) .. (0.8601951565, 4.6211751144).. controls (0.861758844, 
    4.6211751144) and (0.863065621, 4.6197852581) .. (0.8629372188, 
    4.6182141094).. controls (0.8624612333, 4.6058334615) and (0.8503903596, 
    4.6021319406) .. (0.8441131729, 4.6019129185).. controls (0.8358795777, 
    4.601746231) and (0.8297836835, 4.603914756) .. (0.8260747808, 4.6089757325)..
     controls (0.818845835, 4.6185237777) and (0.8239446204, 4.6342884994) .. 
    (0.8240276996, 4.6349003483).. controls (0.8377529335, 4.6602131152) and 
    (0.8649843529, 4.6513751323) .. (0.865203375, 4.6513751323).. controls 
    (0.8981453758, 4.6373477177) and (0.8855079763, 4.5898420448) .. 
    (0.8852964154, 4.5893661387).. controls (0.8849413446, 4.5879307742) and 
    (0.8834231654, 4.5870169033) .. (0.8818897198, 4.5874021367) -- cycle;
  
  
  
      \end{scope}
      \begin{scope}[shift={(0.8731702173, -3.5440035271)}]
        \path[fill=c005093,even odd rule] (0.9028508581, 4.6738021152).. controls 
    (0.9014154935, 4.6741950215) and (0.9005394581, 4.6756377944) .. 
    (0.9008037769, 4.6772466727).. controls (0.9011966831, 4.6786820373) and 
    (0.9025939477, 4.6796789873) .. (0.904112365, 4.679293754).. controls 
    (0.9052908192, 4.6790741498) and (0.9335720788, 4.6729788377) .. 
    (0.9405442465, 4.7012526096).. controls (0.9405442465, 4.7013346304) and 
    (0.9447744577, 4.7251831402) .. (0.9363217346, 4.735592219).. controls 
    (0.9339648262, 4.7385079802) and (0.9309132544, 4.7400337529) .. (0.926902171,
     4.7401697223).. controls (0.9266000169, 4.7401697223) and (0.9203228302, 
    4.7410838577) .. (0.9173163963, 4.7376392473).. controls (0.9152314796, 
    4.7351992598) and (0.9151407275, 4.7302667915) .. (0.9171867504, 
    4.7239140133).. controls (0.9196267379, 4.718422401) and (0.9212809923, 
    4.7054827944) .. (0.9097312154, 4.6982916046).. controls (0.9094290613, 
    4.6979894504) and (0.897494104, 4.689272435) .. (0.8801959102, 4.7053544185)..
     controls (0.8791005352, 4.7064873644) and (0.8790552915, 4.7081416981) .. 
    (0.8800133477, 4.7093200729).. controls (0.8811918019, 4.7104456104) and 
    (0.8829820256, 4.7104456104) .. (0.884024431, 4.7094867604).. controls 
    (0.897575966, 4.6967208527) and (0.9061569856, 4.7026958323) .. (0.9065951092,
     4.7029073931).. controls (0.9171024013, 4.7094867604) and (0.9121772883, 
    4.72155774) .. (0.91186762, 4.7220411867).. controls (0.9091709602, 
    4.730758255) and (0.9096014637, 4.7371638175) .. (0.9131368262, 
    4.7412957096).. controls (0.918190315, 4.7471347462) and (0.9271188383, 
    4.7457902396) .. (0.9276023115, 4.7457902396).. controls (0.932784176, 
    4.7454806771) and (0.937271139, 4.7432145208) .. (0.9407156965, 
    4.7391202525).. controls (0.9505204933, 4.7267471717) and (0.9462072294, 
    4.7012531387) .. (0.9460335304, 4.7002484894).. controls (0.9375808073, 
    4.6663471358) and (0.9032414096, 4.6738026708) .. (0.9028485827, 4.6738026708)
     -- cycle;
  
  
  
      \end{scope}
      \begin{scope}[shift={(0.8731702173, -3.5622616292)}]
        \path[fill=c005093,even odd rule] (0.89038684, 4.7030442092).. controls 
    (0.8889514754, 4.7034371154) and (0.8879471171, 4.7047817279) .. 
    (0.8882566796, 4.706345336).. controls (0.8885588338, 4.7079619402) and 
    (0.8900016067, 4.7088832194) .. (0.8915198652, 4.7084905777).. controls 
    (0.8917844485, 4.7084905777) and (0.9153142413, 4.7039130744) .. 
    (0.9225431871, 4.7166336325).. controls (0.9230718246, 4.7176382554) and 
    (0.9356188425, 4.7384641123) .. (0.9163188906, 4.7548935731).. controls 
    (0.9152311885, 4.755898196) and (0.915049949, 4.7577715254) .. (0.9160992865, 
    4.7587308517).. controls (0.9171039094, 4.759992385) and (0.9188412429, 
    4.7600828725) .. (0.9200196177, 4.7591690017).. controls (0.9429831227, 
    4.7395140848) and (0.9276036344, 4.7140200519) .. (0.9273846123, 
    4.713808491).. controls (0.9181916644, 4.6975527292) and (0.8915645798, 
    4.7028253458) .. (0.890386205, 4.7030443679) -- cycle;
  
  
  
      \end{scope}
      \begin{scope}[shift={(0.8731702173, -3.4658826521)}]
        \path[fill=c005093,even odd rule] (1.0006818156, 4.5161181215).. controls 
    (1.0005468781, 4.5164202756) and (1.0005468781, 4.516684859) .. (1.0005468781,
     4.5169944215).. controls (1.0005468781, 4.5182030381) and (1.001332426, 
    4.5192529048) .. (1.0025108802, 4.5196457581).. controls (1.0027304844, 
    4.5197357165) and (1.0245981117, 4.5268369479) .. (1.0304371483, 
    4.5449659863).. controls (1.0313586921, 4.5479648531) and (1.0319630004, 
    4.5510997423) .. (1.0319630004, 4.5544611413).. controls (1.0319630004, 
    4.5649609456) and (1.0270378875, 4.5767750677) .. (1.0174521127, 
    4.589759971).. controls (1.0140151223, 4.5925472771) and (0.9937936002, 
    4.6085840169) .. (0.9937936002, 4.6232157133).. controls (0.9937936002, 
    4.6242203362) and (0.993923246, 4.6250889633) .. (0.994096019, 4.6260936656)..
     controls (0.9955766273, 4.6331187235) and (1.0019443808, 4.6380287023) .. 
    (1.0131390869, 4.6408614108).. controls (1.0145818598, 4.6416998754) and 
    (1.0197637508, 4.6450990833) .. (1.02028498, 4.6487097727).. controls 
    (1.02028498, 4.6488447102) and (1.02028498, 4.6490193352) .. (1.02028498, 
    4.6492384102).. controls (1.02028498, 4.650975929) and (1.0191065258, 
    4.6528492054) .. (1.0166666177, 4.6548962867).. controls (1.0163644635, 
    4.6550259325) and (0.9947077621, 4.665886999) .. (0.9791923044, 
    4.6512780302).. controls (0.9791923044, 4.6512780302) and (0.9687832521, 
    4.6396074181) .. (0.9699616004, 4.6233063331).. controls (0.9708302275, 
    4.6110616546) and (0.9777117754, 4.5995194183) .. (0.9905683029, 
    4.5886193258).. controls (0.9909157008, 4.5883171717) and (1.0070430869, 
    4.5772055977) .. (1.0088258229, 4.5615163617).. controls (1.00969445, 
    4.553894589) and (1.0070430604, 4.5464389481) .. (1.0007659002, 
    4.5396481523).. controls (0.9996781981, 4.5385152065) and (0.9978878685, 
    4.538386619) .. (0.9967548433, 4.5393915065).. controls (0.9956218975, 
    4.5405622877) and (0.9955763892, 4.5422997271) .. (0.9965352392, 
    4.5433950227).. controls (1.001634051, 4.5489697142) and (1.0039001279, 
    4.554854074) .. (1.0032052262, 4.5609122121).. controls (1.0017246179, 
    4.5741162433) and (0.9872137831, 4.5840041194) .. (0.9871307304, 
    4.5840947392).. controls (0.9728314035, 4.5962487979) and (0.9652550333, 
    4.6091885104) .. (0.9643787598, 4.6230421202).. controls (0.963079391, 
    4.6416999548) and (0.9746670033, 4.6546773969) .. (0.9751429094, 
    4.6552061667).. controls (0.9937931504, 4.6727611121) and (1.0187207633, 
    4.659995125) .. (1.0197631687, 4.6594285992).. controls (1.0244238571, 
    4.6556819404) and (1.0263424302, 4.6519807371) .. (1.0258665506, 
    4.6480602735).. controls (1.0249902506, 4.6406953848) and (1.0160617538, 
    4.6361176698) .. (1.0149740252, 4.6356871133).. controls (1.0055620552, 
    4.6331188029) and (1.0005463225, 4.6296668106) .. (0.9995869698, 
    4.6250892015).. controls (0.9974115656, 4.6146725556) and (1.0144905785, 
    4.5991267502) .. (1.0212436183, 4.593891969) -- (1.0215534454, 4.5936803023) 
    -- (1.0217730496, 4.5932799877).. controls (1.0357550617, 4.574501265) and 
    (1.0404610467, 4.5575958188) .. (1.0357550617, 4.5430924981).. controls 
    (1.0290019954, 4.5223044502) and (1.0052528631, 4.5146298931) .. 
    (1.0042104312, 4.5142446333).. controls (1.0027298229, 4.5138517271) and 
    (1.0011587006, 4.5146298667) .. (1.000682821, 4.5161178833) -- cycle;
  
  
  
      \end{scope}
      \begin{scope}[shift={(0.8731702173, -3.4894512063)}]
        \path[fill=c005093,even odd rule] (0.997630085, 4.5521422535).. controls 
    (0.9970181037, 4.5535019473) and (0.99780471, 4.5552846304) .. (0.99928426, 
    4.5556774573).. controls (0.9995038642, 4.5557594781) and (1.0142709215, 
    4.560859216) .. (1.0185010269, 4.5725373952).. controls (1.0207671831, 
    4.5792451383) and (1.019369654, 4.586874584) .. (1.0140972225, 4.595501006).. 
    controls (1.0140152017, 4.5956835685) and (1.0088247117, 4.6042256415) .. 
    (0.9967537321, 4.6147178787).. controls (0.9958395967, 4.6157148288) and 
    (0.9744926958, 4.6364576065) .. (0.9788437158, 4.6575099996).. controls 
    (0.9814194346, 4.6692259083) and (0.9914357923, 4.6784642852) .. 
    (1.0086055838, 4.684749039).. controls (1.0100861921, 4.6853533473) and 
    (1.0116121235, 4.6845664765) .. (1.0123523483, 4.683087191).. controls 
    (1.0128357421, 4.6816065827) and (1.0120049504, 4.6800430275) .. 
    (1.0105695858, 4.6794689346).. controls (0.9952278535, 4.6738413529) and 
    (0.9864276267, 4.666136501) .. (0.9843805454, 4.6562864075).. controls 
    (0.9806339131, 4.638467011) and (1.0003718827, 4.6190386833) .. (1.0005455817,
     4.6188120942).. controls (1.0132284102, 4.6078364102) and (1.0185841854, 
    4.5988171613) .. (1.0188032075, 4.5984243344).. controls (1.0250803677, 
    4.5884457327) and (1.0267346221, 4.5790338685) .. (1.0236828915, 
    4.5705811454).. controls (1.0185840796, 4.5563800054) and (1.0018524381, 
    4.550669371) .. (1.00115751, 4.5504050258).. controls (0.9995864142, 
    4.5500121196) and (0.9981058852, 4.5506696092) .. (0.9976300056, 4.55214228) 
    -- cycle;
  
  
  
      \end{scope}
      \begin{scope}[shift={(0.8731702173, -3.5944262313)}]
        \path[fill=c005093,even odd rule] (1.11571659, 4.6613547394).. controls 
    (1.0900035085, 4.6789096848) and (1.035661505, 4.6849679552) .. (1.0351025463,
     4.685096331).. controls (1.0336145296, 4.685270956) and (1.0324436425, 
    4.686539104) .. (1.032617315, 4.6881026592).. controls (1.0328369192, 
    4.689673755) and (1.0342262463, 4.6908446421) .. (1.0356614785, 
    4.6906709696).. controls (1.0380183869, 4.6904593029) and (1.0919598377, 
    4.6843937035) .. (1.1189797492, 4.6659700517).. controls (1.120112695, 
    4.6649654288) and (1.1205056013, 4.6631827456) .. (1.1196369742, 
    4.6620043973).. controls (1.1187606742, 4.6607354556) and (1.1169781762, 
    4.6604785452) .. (1.1157166429, 4.6613548452) -- cycle;
  
  
  
      \end{scope}
      \begin{scope}[shift={(0.8731702173, -3.5757749583)}]
        \path[fill=c005093,even odd rule] (1.1411121185, 4.6390406277).. controls 
    (1.1405078102, 4.639425861) and (1.09422565, 4.6813494871) .. (1.0322242235, 
    4.6779049296).. controls (1.0306531277, 4.6777699921) and (1.0293915944, 
    4.6789473879) .. (1.0293009746, 4.6805562662).. controls (1.0291713288, 
    4.6820067121) and (1.0303963496, 4.6833888954) .. (1.0319598783, 
    4.6833888954).. controls (1.096235081, 4.6869240727) and (1.1443828713, 
    4.6435729137) .. (1.1448210213, 4.6431725198).. controls (1.1460370463, 
    4.6420848177) and (1.1460370463, 4.6403398112) .. (1.1450326879, 
    4.6391690035).. controls (1.1440354733, 4.6379982223) and (1.142290705, 
    4.6379074702) .. (1.1411123567, 4.6390393577) -- cycle;
  
  
  
      \end{scope}
      \begin{scope}[shift={(0.8731702173, -3.5699363979)}]
        \path[fill=c005093,even odd rule] (1.1593703529, 4.6386483035).. controls 
    (1.1588943675, 4.6390787806) and (1.1153166988, 4.6837821192) .. 
    (1.0242029565, 4.6866601508).. controls (1.022639269, 4.6866601508) and 
    (1.0213703538, 4.6879669279) .. (1.0214144598, 4.6894928594).. controls 
    (1.0214144598, 4.6911094635) and (1.0227225598, 4.6922801654) .. 
    (1.0242937085, 4.6922801654).. controls (1.1178473325, 4.6893266217) and 
    (1.1615987002, 4.644404261) .. (1.163472056, 4.6425687406).. controls 
    (1.1644692706, 4.6413450427) and (1.1644692706, 4.6395623331) .. 
    (1.1633820977, 4.6386483035).. controls (1.1620753206, 4.6374246056) and 
    (1.1604209075, 4.6375153577) .. (1.159370935, 4.6386483035) -- cycle;
  
  
  
      \end{scope}
      \begin{scope}[shift={(0.8731702173, -3.5575097125)}]
        \path[fill=c005093,even odd rule] (1.1054279231, 4.5935895767).. controls 
    (1.1039473148, 4.5947151142) and (1.07784138, 4.613758235) .. (1.0862941296, 
    4.6383306712).. controls (1.0867775233, 4.6398112796) and (1.0884772067, 
    4.6405968275) .. (1.0897839838, 4.6399848462).. controls (1.0913098358, 
    4.63959194) and (1.09205014, 4.637937765) .. (1.0915667463, 4.6365025327).. 
    controls (1.0845114994, 4.6158052104) and (1.108041636, 4.5986353131) .. 
    (1.1083438167, 4.5984163175).. controls (1.1466943771, 4.5798113731) and 
    (1.2063465417, 4.6241294783) .. (1.2069130675, 4.6245222787).. controls 
    (1.2080838488, 4.625473985) and (1.2098740725, 4.6250887517) .. (1.2108787219,
     4.6239028892).. controls (1.2117550219, 4.6227773517) and (1.2114377865, 
    4.6209871015) .. (1.2102670052, 4.6199899927).. controls (1.2076910219, 
    4.6180712344) and (1.147298659, 4.5733151908) .. (1.105428029, 4.5935894973) 
    -- cycle;
  
  
  
      \end{scope}
      \begin{scope}[shift={(0.8731702173, -3.5852245521)}]
        \path[fill=c005093,even odd rule] (1.1258163973, 4.6299676683).. controls 
    (1.1016820846, 4.6327096512) and (1.0997181883, 4.6584227063) .. 
    (1.0997181883, 4.6586417283).. controls (1.0996361675, 4.6602959033) and 
    (1.1007228113, 4.661565189) .. (1.1022486633, 4.6616406746).. controls 
    (1.1038575946, 4.6618232371) and (1.1052551767, 4.6605603808) .. 
    (1.1053004998, 4.6591101996).. controls (1.1054301456, 4.6582415725) and 
    (1.1070829977, 4.6376800344) .. (1.126390649, 4.6355423333).. controls 
    (1.1266023156, 4.6354603125) and (1.1632908429, 4.633502925) .. (1.203212605, 
    4.6662712269).. controls (1.2043833862, 4.6672681769) and (1.20617361, 
    4.666920779) .. (1.2071329363, 4.6658709123).. controls (1.2080846425, 
    4.664609379) and (1.2079184842, 4.6629476633) .. (1.2067326217, 4.661943014)..
     controls (1.1650357171, 4.6277017237) and (1.1274782452, 4.6297941017) .. 
    (1.1258164238, 4.6299678006) -- cycle;
  
  
  
      \end{scope}
      \begin{scope}[shift={(0.8731702173, -3.5770814708)}]
        \path[fill=c005093,even odd rule] (1.1133619835, 4.6447881979).. controls 
    (1.1121835294, 4.6457928208) and (1.1118813752, 4.647447234) .. (1.1127500023,
     4.6487162285).. controls (1.1137546252, 4.649970089) and (1.1155373877, 
    4.650279916) .. (1.1167157625, 4.6493657806).. controls (1.1449517783, 
    4.6284040337) and (1.1937944967, 4.6676231948) .. (1.1942779433, 
    4.6680084546).. controls (1.1954487246, 4.6690130775) and (1.1972767838, 
    4.6687940025) .. (1.1981908133, 4.6676232213).. controls (1.1992861883, 
    4.6664447671) and (1.1990745217, 4.6647073542) .. (1.1978886592, 
    4.6637026783).. controls (1.1956301758, 4.6619577513) and (1.145034884, 
    4.6213032785) .. (1.1133618777, 4.6447879862) -- cycle;
  
  
  
      \end{scope}
      \begin{scope}[shift={(0.8731702173, -3.5545940042)}]
        \path[fill=c005093,even odd rule] (1.1849536033, 4.6045272135).. controls 
    (1.1847710408, 4.6061361448) and (1.1859505533, 4.6074430013) .. 
    (1.1875142408, 4.6074430013).. controls (1.2029843754, 4.6084476242) and 
    (1.2041628296, 4.6167264896) .. (1.2041628296, 4.6169833735) -- (1.2042527879,
     4.6174667673) -- (1.2044723921, 4.6178596735).. controls (1.2061720754, 
    4.6216894908) and (1.2061720754, 4.6254361231) .. (1.2044723921, 
    4.6292734017).. controls (1.1984142275, 4.6434745152) and (1.1727917129, 
    4.6531962594) .. (1.1726179875, 4.6534152815).. controls (1.1611589627, 
    4.656421689) and (1.1441628702, 4.6659997908) .. (1.1434679421, 
    4.6664001848).. controls (1.1421990004, 4.667140489) and (1.141723015, 
    4.6688780077) .. (1.1425085629, 4.670237675).. controls (1.143158115, 
    4.6716201229) and (1.1449485504, 4.6720201729) .. (1.1463006242, 
    4.6712422979).. controls (1.14643027, 4.6712422979) and (1.1632889379, 
    4.6616037388) .. (1.1743628352, 4.6586426544).. controls (1.1757526915, 
    4.6581214252) and (1.2025913898, 4.6478785048) .. (1.2095635575, 
    4.6314943406).. controls (1.2118297138, 4.6262294233) and (1.2118297138, 
    4.6208207579) .. (1.2095635575, 4.615721946) -- (1.20974612, 4.6164622502).. 
    controls (1.2097011408, 4.615941021) and (1.2083033471, 4.6029560119) .. 
    (1.1878250204, 4.6019966592).. controls (1.1863444121, 4.6018220342) and 
    (1.1849546352, 4.6029560383) .. (1.1849546352, 4.6045273988) -- cycle;
  
  
  
      \end{scope}
      \begin{scope}[shift={(0.8731702173, -3.5413907667)}]
        \path[fill=c005093,even odd rule] (1.1786699873, 4.5836940542).. controls 
    (1.1771062998, 4.5839586375) and (1.176192429, 4.5855673042) .. (1.1765398269,
     4.5870025629).. controls (1.1770232206, 4.5884905796) and (1.1784585852, 
    4.5893594712) .. (1.1800674371, 4.5890044004).. controls (1.1947595908, 
    4.5852199062) and (1.20556147, 4.5912327212) .. (1.20556147, 4.5913234733).. 
    controls (1.2229049869, 4.6040817875) and (1.21440702, 4.6202544967) .. 
    (1.2140142196, 4.6209419371) -- (1.2137946154, 4.6212517642) -- (1.2136649696,
     4.6215539183).. controls (1.2093063825, 4.6384215292) and (1.1724893471, 
    4.652713289) .. (1.1586734669, 4.6567696954).. controls (1.1572381023, 
    4.6571549287) and (1.1563242315, 4.6588167767) .. (1.1568454606, 
    4.6601613627).. controls (1.1572383669, 4.6616948877) and (1.158764219, 
    4.6625635148) .. (1.1602448008, 4.6621253648).. controls (1.162420205, 
    4.6614758127) and (1.2131439256, 4.6462624298) .. (1.219111444, 4.6230874698) 
    -- (1.2188918398, 4.6236917781).. controls (1.2188918398, 4.623472174) and 
    (1.230267706, 4.6026091167) .. (1.2086941102, 4.5866177794).. controls 
    (1.2079159706, 4.5862174648) and (1.19575453, 4.5791622444) .. (1.1786678971, 
    4.5836945304) -- cycle;
  
  
  
      \end{scope}
      \begin{scope}[shift={(0.8731702173, -3.5532866979)}]
        \path[fill=c005093,even odd rule] (1.2266522542, 4.6080546385).. controls 
    (1.2252168896, 4.608530624) and (1.2243030188, 4.6099733969) .. (1.2247787396,
     4.6114463058).. controls (1.2281326773, 4.6234265598) and (1.2271356479, 
    4.6343342194) .. (1.2216363892, 4.6437387281).. controls (1.2088780485, 
    4.6654407262) and (1.1763212017, 4.6732361979) .. (1.176154964, 
    4.6732361979).. controls (1.1745383598, 4.6736365125) and (1.1736623244, 
    4.6752454438) .. (1.1739644785, 4.6765976763).. controls (1.1743647931, 
    4.6781613638) and (1.1758377285, 4.6791659867) .. (1.177325851, 
    4.6787656721).. controls (1.178760951, 4.6783804388) and (1.2123981181, 
    4.6705395381) .. (1.2264783171, 4.6465713573).. controls (1.2328385565, 
    4.6357618846) and (1.2340171165, 4.6234266656) .. (1.2300437892, 
    4.6099280208).. controls (1.2296963912, 4.6084400042) and (1.2281327037, 
    4.6076618646) .. (1.2266522277, 4.6080545063) -- cycle;
  
  
  
      \end{scope}
      \begin{scope}[shift={(0.8731702173, -3.6071005667)}]
        \path[fill=c005093,even odd rule] (1.1654271946, 4.643179981).. controls 
    (1.1641656613, 4.6439126123) and (1.1635991883, 4.6455667873) .. 
    (1.1643394925, 4.6470474485).. controls (1.1649512092, 4.6483994694) and 
    (1.1666053842, 4.6488831277) .. (1.1679953992, 4.6482259027).. controls 
    (1.1681700242, 4.6480962569) and (1.189002575, 4.6379829823) .. (1.2076906515,
     4.648618809).. controls (1.2090881806, 4.6494043569) and (1.2108783515, 
    4.6488833923) .. (1.21152793, 4.647614186).. controls (1.2123134779, 
    4.6462243298) and (1.2118300842, 4.6444792969) .. (1.2104780633, 
    4.6437843952).. controls (1.1890933006, 4.6317133892) and (1.1664319763, 
    4.6427345285) .. (1.1654273269, 4.6431800869) -- cycle;
  
  
  
      \end{scope}
      \begin{scope}[shift={(0.8731702173, -3.5261018188)}]
        \path[fill=c005093,even odd rule] (1.1188535958, 4.6250972713).. controls 
    (1.1174560667, 4.6258828192) and (1.1171089333, 4.6275372588) .. 
    (1.1178489729, 4.6287986863).. controls (1.1182418792, 4.6296297425) and 
    (1.1284771531, 4.6464066277) .. (1.1690484938, 4.6501079633).. controls 
    (1.1690484938, 4.6501079633) and (1.2176343281, 4.6660993006) .. 
    (1.2220382383, 4.6988676025).. controls (1.2222049258, 4.7000460567) and 
    (1.2232997717, 4.7090651204) .. (1.2136233506, 4.7113690592).. controls 
    (1.2122334944, 4.7119884488) and (1.2031235462, 4.7152441731) .. (1.197631934,
     4.702085465).. controls (1.1972467006, 4.7010808421) and (1.1867394085, 
    4.6719837398) .. (1.1546660083, 4.6669755742).. controls (1.1540617, 
    4.6668856158) and (1.1025977546, 4.6601695119) .. (1.0962752712, 
    4.6350758729).. controls (1.0958900379, 4.6335047771) and (1.0943189421, 
    4.6327115563) .. (1.0928382808, 4.6331118708).. controls (1.0914859954, 
    4.6335047771) and (1.0905268808, 4.6349851208) .. (1.0907911996, 
    4.6364657821).. controls (1.0980654685, 4.6651021919) and (1.1516143571, 
    4.6722405179) .. (1.1538351373, 4.6725955888).. controls (1.182645246, 
    4.6769012856) and (1.192057216, 4.7029541715) .. (1.1924500429, 
    4.7041401133).. controls (1.2002908377, 4.7228735129) and (1.2151870646, 
    4.7168153748) .. (1.2153683306, 4.716694566).. controls (1.2269558635, 
    4.7139374225) and (1.2285269594, 4.7032638398) .. (1.22752231, 4.6980366256)..
     controls (1.2227332458, 4.6617784431) and (1.1724023256, 4.6452734967) .. 
    (1.1702193015, 4.6447447004).. controls (1.13218595, 4.6410887144) and 
    (1.12277398, 4.6261774327) .. (1.1226833337, 4.626056624).. controls 
    (1.1218977858, 4.6247043385) and (1.1201150233, 4.6241831094) .. 
    (1.1188535165, 4.6250972448) -- cycle;
  
  
  
      \end{scope}
      \begin{scope}[shift={(0.8731702173, -3.5366089521)}]
        \path[fill=c005093,even odd rule] (1.1059937346, 4.640136426).. controls 
    (1.1046417138, 4.6406576552) and (1.1040297325, 4.6424401531) .. 
    (1.1047322013, 4.6438453288).. controls (1.1052534304, 4.6448425433) and 
    (1.11631984, 4.6672846075) .. (1.1545797806, 4.6694979265).. controls 
    (1.1547544056, 4.6694979265) and (1.2090877044, 4.6805265006) .. 
    (1.2115276125, 4.7248823354).. controls (1.2116175708, 4.7264081875) and 
    (1.2130082208, 4.7275789952) .. (1.2144963698, 4.7274583188).. controls 
    (1.2160600573, 4.7273233813) and (1.2172383527, 4.7260155458) .. 
    (1.2171025156, 4.7245350698).. controls (1.2144058558, 4.6758207538) and 
    (1.155841314, 4.66400642) .. (1.1553201113, 4.663832721).. controls 
    (1.1201119013, 4.6618763919) and (1.1101334054, 4.6422667981) .. 
    (1.1097405785, 4.6413905246).. controls (1.1090457827, 4.6401366642) and 
    (1.1073836702, 4.63939636) .. (1.1059938404, 4.6401366642) -- cycle;
  
  
  
      \end{scope}
      \begin{scope}[shift={(0.8731702173, -3.5647468604)}]
        \path[fill=c005093,even odd rule] (1.1137518735, 4.6961166238).. controls 
    (1.1125734194, 4.6967209321) and (1.0871701915, 4.7108841838) .. 
    (1.0840353817, 4.729443805).. controls (1.0828569275, 4.736514186) and 
    (1.0850323317, 4.7432597646) .. (1.0905240233, 4.749446305) -- (1.0906986483, 
    4.7497561321) -- (1.0910084754, 4.7498381529).. controls (1.0913937088, 
    4.7500127779) and (1.099151345, 4.7552466067) .. (1.1075662327, 
    4.7525801358).. controls (1.1133599463, 4.750631215) and (1.1177638565, 
    4.7456985879) .. (1.12080802, 4.7375556125).. controls (1.1207259992, 
    4.7377672792) and (1.1338383523, 4.7151511456) .. (1.1465211808, 
    4.7094480783).. controls (1.1478657933, 4.7086625304) and (1.1485682621, 
    4.7070459263) .. (1.1478657933, 4.705618261).. controls (1.1472993204, 
    4.7042585673) and (1.1456451454, 4.703654259) .. (1.1441644842, 
    4.7043493194).. controls (1.129918021, 4.7108380404) and (1.1164044538, 
    4.7342395896) .. (1.1157170133, 4.7352442654).. controls (1.11314103, 
    4.7419973052) and (1.10996095, 4.7458345837) .. (1.1057761677, 4.7470961171)..
     controls (1.1007226525, 4.7488788796) and (1.0956237612, 4.7461367379) .. 
    (1.0944453865, 4.7455250212).. controls (1.091008396, 4.741521505) and 
    (1.0892183046, 4.7375103158) .. (1.0892183046, 4.7331895906).. controls 
    (1.0892183046, 4.732192376) and (1.0893479504, 4.7313163406) .. (1.0895204588,
     4.7304022846).. controls (1.092179389, 4.7145393496) and (1.1161928665, 
    4.7010785402) .. (1.1164044273, 4.7009123025).. controls (1.1173638065, 
    4.7004212358) and (1.117764121, 4.6994242858) .. (1.117764121, 4.6985479858)..
     controls (1.117764121, 4.6980344296) and (1.1176741627, 4.6975510358) .. 
    (1.1174545585, 4.6970750504).. controls (1.1167142544, 4.6958061088) and 
    (1.115014571, 4.6952018004) .. (1.1137531435, 4.6961156713) -- cycle;
  
  
  
      \end{scope}
      \begin{scope}[shift={(0.8731702173, -3.5727994542)}]
        \path[fill=c005093,even odd rule] (1.1263913898, 4.7094809396).. controls 
    (1.1028537919, 4.7255176) and (1.09283754, 4.7550603665) .. (1.0924447131, 
    4.7563672229).. controls (1.0920518069, 4.757893075) and (1.092792111, 
    4.7593736304) .. (1.0943179631, 4.7598571035).. controls (1.0957985715, 
    4.7604688202) and (1.0973697996, 4.7595549494) .. (1.0978910023, 
    4.7580669327).. controls (1.0979730231, 4.7578923077) and (1.1076957992, 
    4.7290452631) .. (1.129526279, 4.7140510346).. controls (1.130833056, 
    4.7131824075) and (1.1312256977, 4.7113995921) .. (1.1302210748, 
    4.7102287844).. controls (1.1293524477, 4.7088767635) and (1.1276527644, 
    4.708476449) .. (1.1263912575, 4.7094810719) -- cycle;
  
  
  
      \end{scope}
      \begin{scope}[shift={(0.8731702173, -3.4406379625)}]
        \path[fill=c005093,even odd rule] (0.8263403696, 3.7278425094).. controls 
    (0.8239458904, 3.7327675694) and (0.8256831446, 3.7386594698) .. 
    (0.8305251519, 3.7410918108).. controls (0.8336146385, 3.7427536587) and 
    (1.149483694, 3.9054168463) .. (1.483217469, 3.7425724192).. controls 
    (1.4880593969, 3.7404422587) and (1.490098911, 3.7343387975) .. (1.4877496492,
     3.7295421398).. controls (1.4853927408, 3.7247152667) and (1.4793800846, 
    3.7227966406) .. (1.4746287765, 3.7251005265).. controls (1.1499595735, 
    3.8832692104) and (0.8426339404, 3.7255386765) .. (0.8395822098, 
    3.7238389931).. controls (0.8347930927, 3.7212630098) and (0.8289086535, 
    3.7231817681) .. (0.8263403696, 3.7278425094) -- cycle;
  
  
  
      \end{scope}
      \begin{scope}[shift={(0.8731702173, -2.4204917829)}]
        \path[fill=c005093,even odd rule] (0.8237187456, 2.7036102106) -- 
    (0.8187106329, 2.7059517731) -- (0.7441547538, 3.7368319133) -- (0.7597533171,
     3.7288626369).. controls (0.9696503181, 3.6203147579) and (1.1734966952, 
    3.6244088146) .. (1.3075688754, 3.6471532975).. controls (1.453107674, 
    3.6716878983) and (1.5504003631, 3.7227137467) .. (1.5511934781, 3.7231065735)
     -- (1.5669204702, 3.7315592967) -- (1.4952879458, 2.7120628279) -- 
    (1.4908010888, 2.7093963306).. controls (1.181896634, 2.5325320296) and 
    (0.8272539494, 2.7017368283) .. (0.8237187456, 2.7036101313) -- 
    cycle(1.4764564388, 2.7240885108).. controls (1.4772496596, 2.7348073373) and 
    (1.542650294, 3.6657128923) .. (1.5449163973, 3.6981866865).. controls 
    (1.5117552685, 3.6828902773) and (1.4277873135, 3.6475536915) .. 
    (1.3116176883, 3.6279893944).. controls (1.1777193129, 3.6051921271) and 
    (0.9756631596, 3.6010147002) .. (0.7662496052, 3.7040634792).. controls 
    (0.7685157615, 3.6718238412) and (0.8365374906, 2.7301013523) .. 
    (0.8374515202, 2.718891565).. controls (0.879450526, 2.6998635783) and 
    (1.1969286185, 2.5672568458) .. (1.4764564388, 2.7240885108) -- cycle;
  
  
  
      \end{scope}
      \begin{scope}[shift={(0.8731702173, -3.2878783938)}]
        \path[fill=c005093,even odd rule] (1.1595451631, 3.2515922977).. controls 
    (1.033094835, 3.2499379904) and (1.0249895098, 3.4670262606) .. (1.0245967094,
     3.4762645317).. controls (1.0245967094, 3.5139125704) and (1.0038992813, 
    3.5700371513) .. (1.0037709054, 3.5705054902).. controls (1.0018521471, 
    3.5756118692) and (1.0043828867, 3.5813225565) .. (1.0094816456, 
    3.5831505627).. controls (1.0145804575, 3.5850238127) and (1.020155149, 
    3.582448094) .. (1.0221569071, 3.5774021194).. controls (1.0230332071, 
    3.5749094798) and (1.0442062767, 3.5174401542) .. (1.0442062767, 
    3.4765212833).. controls (1.0442882975, 3.4748594354) and (1.0520471773, 
    3.2691095135) .. (1.159854805, 3.2711565683) -- (1.1609425071, 3.2712473204) 
    -- (1.1620302092, 3.2710282454).. controls (1.1628157571, 3.2708544142) and 
    (1.1839512294, 3.2673268833) .. (1.2075265833, 3.285886531).. controls 
    (1.2324465233, 3.3057605494) and (1.2634247338, 3.354255764) .. (1.2693166342,
     3.4725178465).. controls (1.2693166342, 3.4805701756) and (1.2711898842, 
    3.5549145469) .. (1.2957170502, 3.5805747646).. controls (1.2995468675, 
    3.584631171) and (1.3056048998, 3.5847595469) .. (1.3095329304, 
    3.5809674063).. controls (1.3134458004, 3.5772660442) and (1.3136724954, 
    3.5710795038) .. (1.3099181638, 3.5672497394).. controls (1.2935415931, 
    3.5499893546) and (1.2891376829, 3.49264097) .. (1.2888355023, 3.4719058652)..
     controls (1.2838197431, 3.3701186196) and (1.2604256817, 3.3022705629) .. 
    (1.2188949883, 3.2698950085).. controls (1.190832645, 3.2481099048) and 
    (1.164734436, 3.2507235119) .. (1.1595450044, 3.2515922183) -- cycle;
  
  
  
      \end{scope}
      \begin{scope}[shift={(0.8731702173, -3.3454377042)}]
        \path[fill=c005093,even odd rule] (1.1320923669, 3.380006575).. controls 
    (1.0875628067, 3.4253217096) and (1.0929637198, 3.6170064873) .. 
    (1.0936661885, 3.6386632152).. controls (1.0938408135, 3.6440641548) and 
    (1.0983721735, 3.6483773658) .. (1.1038562452, 3.6480752381).. controls 
    (1.1092193758, 3.6478926756) and (1.11344202, 3.6433692531) .. (1.113268321, 
    3.6379683135).. controls (1.1112664835, 3.5812696133) and (1.1133582794, 
    3.4270665837) .. (1.1460820254, 3.3937770792).. controls (1.1506143115, 
    3.38919947) and (1.1541871654, 3.3886329442) .. (1.1564911042, 3.3887160763)..
     controls (1.1568914188, 3.3887160763) and (1.1639919623, 3.389418545) .. 
    (1.1727014635, 3.3996161688).. controls (1.1882622179, 3.4182664627) and 
    (1.2118301106, 3.4733864471) .. (1.1934064323, 3.6337305881).. controls 
    (1.1927947156, 3.6391390948) and (1.19676037, 3.6440641548) .. (1.2020253138, 
    3.6445476279).. controls (1.2074337675, 3.6452426883) and (1.2123136631, 
    3.6414958708) .. (1.2129255121, 3.6360043115).. controls (1.2314473508, 
    3.4736885748) and (1.2084838458, 3.4098968981) .. (1.1859508973, 
    3.3850147406).. controls (1.1718328365, 3.3694161773) and (1.1581076025, 
    3.3692423725) .. (1.1566270471, 3.3692423725).. controls (1.1476909567, 
    3.3690232975) and (1.1394572556, 3.3727700621) .. (1.1320923669, 3.380006575) 
    -- cycle;
  
  
  
      \end{scope}
      \begin{scope}[shift={(0.8731702173, -3.3390324063)}]
        \path[fill=c005093,even odd rule] (0.5508589546, 3.9156977344) -- 
    (0.528764156, 3.910901209).. controls (0.528764156, 3.910901209) and 
    (0.5244585121, 3.8946151788) .. (0.5341726627, 3.8914728019).. controls 
    (0.5341726627, 3.8914728019) and (0.5826226071, 3.8957029073) .. 
    (0.5982287375, 3.9156977344).. controls (0.5982287375, 3.9156977344) and 
    (0.602066069, 3.9226774956) .. (0.602066069, 3.9344387275) -- (0.603063019, 
    4.0980612412).. controls (0.603063019, 4.0980612412) and (0.603063019, 
    4.1723979925) .. (0.5885068875, 4.1739691942).. controls (0.5885068875, 
    4.1739691942) and (0.5750837813, 4.1772324592) .. (0.577704929, 4.1276947598) 
    -- (0.5847602287, 3.9307298248).. controls (0.5847602287, 3.9307298248) and 
    (0.5879403617, 3.9188855402) .. (0.5682023127, 3.9167931358) -- (0.5583974894,
     4.1326122792).. controls (0.5583974894, 4.1326122792) and (0.5545601579, 
    4.1665513888) .. (0.5401399429, 4.1659471333).. controls (0.5401399429, 
    4.1659471333) and (0.5293758727, 4.1654259042) .. (0.533167881, 4.1309126223) 
    -- (0.5508588223, 3.9156977344) -- cycle;
  
  
  
      \end{scope}
      \begin{scope}[shift={(0.8731702173, -3.3293105563)}]
        \path[fill=c005093,even odd rule] (0.6483359525, 3.9091636902).. controls 
    (0.6483359525, 3.9091636902) and (0.630516556, 3.907547086) .. (0.630516556, 
    3.9259406813) -- (0.6374433742, 4.1340549992).. controls (0.6374433742, 
    4.1340549992) and (0.6380553554, 4.1614979) .. (0.6250702404, 4.1631974775).. 
    controls (0.6250702404, 4.1631974775) and (0.6116471606, 4.1658563813) .. 
    (0.6137848617, 4.1330956465) -- (0.6133467117, 3.9225414469).. controls 
    (0.6133467117, 3.9225414469) and (0.6105594056, 3.9027581542) .. 
    (0.6316495548, 3.8952194342).. controls (0.6316495548, 3.8952194342) and 
    (0.6789286385, 3.870374874) .. (0.6869885613, 3.8898183358).. controls 
    (0.6869885613, 3.8898183358) and (0.6886427363, 3.9037174804) .. 
    (0.6661174342, 3.9055077042) -- (0.6821088773, 4.1324837181).. controls 
    (0.6821088773, 4.1324837181) and (0.6848131042, 4.1551376606) .. 
    (0.6714731035, 4.156829671).. controls (0.6714731035, 4.156829671) and 
    (0.6606258483, 4.1605385473) .. (0.6568791631, 4.1265993583) -- (0.6483358996,
     3.9091636902) -- cycle;
  
  
  
      \end{scope}
      \begin{scope}[shift={(0.8731702173, -3.5502434604)}]
        \path[fill=c005093,even odd rule] (1.6809966856, 4.120466105).. controls 
    (1.6809966856, 4.120466105) and (1.699125724, 4.1176335023) .. (1.7018678127, 
    4.1040819673).. controls (1.7018678127, 4.1040819673) and (1.7148073135, 
    4.0832561633) .. (1.7261833915, 4.1040819673).. controls (1.7261833915, 
    4.1040819673) and (1.7223007898, 4.1170291146) .. (1.748754096, 
    4.1215085104).. controls (1.748754096, 4.1215085104) and (1.8007393765, 
    4.1306562675) .. (1.8002181738, 4.1430822121).. controls (1.8002181738, 
    4.1430822121) and (1.8025219008, 4.1578045548) .. (1.7691416442, 
    4.150356481).. controls (1.7691416442, 4.150356481) and (1.7261833915, 
    4.1401588573) .. (1.7137271521, 4.1300972029) -- (1.7092327279, 
    4.1306562675).. controls (1.7092327279, 4.1306562675) and (1.68265094, 
    4.1458695181) .. (1.6566432715, 4.1498806015).. controls (1.6566432715, 
    4.1498806015) and (1.629449661, 4.1571548704) .. (1.624479225, 4.1442605869)..
     controls (1.624479225, 4.1442605869) and (1.62282505, 4.1340629631) .. 
    (1.6447913933, 4.1288735315) -- (1.6713656142, 4.1232989194) -- (1.6809966856,
     4.1204662108) -- cycle;
  
  
  
      \end{scope}
      \begin{scope}[shift={(0.8731702173, -3.5592546396)}]
        \path[fill=c005093,even odd rule] (1.7463594052, 4.1752227371) -- 
    (1.67712136, 4.1752227371).. controls (1.6595588475, 4.1752227371) and 
    (1.6453124108, 4.1895672812) .. (1.6453124108, 4.2072583019) -- (1.6453124108,
     4.208784154).. controls (1.6453124108, 4.2264751746) and (1.6595588475, 
    4.2407593673) .. (1.67712136, 4.2407593673) -- (1.7463594052, 4.2407593673).. 
    controls (1.7639596738, 4.2407593673) and (1.778168275, 4.2264751746) .. 
    (1.778168275, 4.208784154) -- (1.778168275, 4.2072583019).. controls 
    (1.778168275, 4.1895672813) and (1.7639596738, 4.1752227371) .. (1.7463594052,
     4.1752227371) -- cycle(1.7395155075, 4.2245187131) -- (1.6829529354, 
    4.2245187131).. controls (1.6748930127, 4.2245187131) and (1.6683136719, 
    4.2181509067) .. (1.6683136719, 4.2101740631).. controls (1.6683136719, 
    4.20215949) and (1.6748930127, 4.1956632019) .. (1.6829529354, 4.1956632019) 
    -- (1.7395155075, 4.1956632019).. controls (1.7475754302, 4.1956632019) and 
    (1.7541094479, 4.20215949) .. (1.7541094479, 4.2101740631).. controls 
    (1.7541094479, 4.2181509067) and (1.7475754302, 4.2245187131) .. 
    (1.7395155075, 4.2245187131) -- cycle;
  
  
  
      \end{scope}
      \begin{scope}[shift={(0.8731702173, -3.4581255979)}]
        \path[fill=c005093,even odd rule] (1.699125724, 4.1345845627) -- 
    (1.7000848385, 4.2799419894).. controls (1.7000848385, 4.2799419894) and 
    (1.699125724, 4.2961977777) .. (1.7114986725, 4.2961977777).. controls 
    (1.7114986725, 4.2961977777) and (1.7193018171, 4.3001182148) .. 
    (1.7210844473, 4.2776759125) -- (1.7239625848, 4.1334515375) -- (1.6991255917,
     4.1345844833) -- cycle;
  
  
  
      \end{scope}
      \begin{scope}[shift={(0.8731702173, -3.461470725)}]
        \path[fill=c005093,even odd rule] (1.6633884531, 4.2882576583).. controls 
    (1.6633884531, 4.2882576583) and (1.6628219802, 4.3052537508) .. 
    (1.6531983965, 4.3030631596).. controls (1.6531983965, 4.3030631596) and 
    (1.6431293867, 4.3046419283) .. (1.6453124108, 4.2787550685).. controls 
    (1.6453124108, 4.2787550685) and (1.6527151615, 4.1429757438) .. 
    (1.7143313281, 4.1514284669).. controls (1.7143313281, 4.1514284669) and 
    (1.7708485771, 4.1419332325) .. (1.779860259, 4.2832872487).. controls 
    (1.779860259, 4.2832872487) and (1.7815221069, 4.2962345019) .. (1.7718911942,
     4.2977981365).. controls (1.7718911942, 4.2977981365) and (1.7601221042, 
    4.3001550448) .. (1.7577805152, 4.2797143948).. controls (1.7577805152, 
    4.2797143948) and (1.7493657333, 4.1656370681) .. (1.711929176, 
    4.1645946627).. controls (1.711929176, 4.1645946627) and (1.6810341506, 
    4.1497816737) .. (1.6646122304, 4.281497025) -- (1.6633885325, 4.2882576583) 
    -- cycle;
  
  
  
      \end{scope}
    \end{scope}
  
  \end{tikzpicture}

% \end{document}
};
        % To include the PDF is recommended.
        \node[opacity=0.2] at (current page.east) {\includegraphics[width=4cm]{BNU/BNU_icon.pdf}};
        \node[opacity=0.2] at (current page.west) {\includegraphics[width=4cm]{BNU/BNU_icon.pdf}};
    \end{tikzpicture}
}

\begin{document}
    \maketitle
    \tableofcontents

    \section{Section One}
    \begin{frame}{Slide with bullet points}
        This is a bullet list of two points:
        \begin{itemize}
            \item Point one
            \item Point two
        \end{itemize}
    \end{frame}

    \begin{frame}{Slide with two columns}
        \begin{columns}
            \column{.5\textwidth}
            Text goes in first column.
            
            \column{.5\textwidth}
            Text goes in second column
        \end{columns}
    \end{frame}

    \section{Section Two}

    \begin{frame}{Slide with table}
        \begin{table}
            \centering
            \begin{tabular}{ccc}
                \hline
                \textbf{name1} & \textbf{name2} & \textbf{name3} \\
                \hline 
                cell1 & cell2 & cell3 \\ 
                cell4 & cell5 & cell6 \\ 
                cell7 & cell8 & cell9 \\ 
                \hline
            \end{tabular}
            \caption{Caption for table one}
            \label{tab:table1}
        \end{table}
    \end{frame}

    \begin{frame}{Slide with figure}
        \begin{figure}[H]
            \centering
            \includegraphics[width=.5\textwidth]{example-image-a}
            \caption{Caption for figure one.}
            \label{fig:figure1}
        \end{figure}
    \end{frame}

    \begin{frame}{Slide with references}
        This is to reference a figure (Figure \ref{fig:figure1})\\
        This it to reference a table (Table \ref{tab:table1})\\
        This is to cite an article \cite{vaswani2017attention}\footnote{Personally this way of citing is not recommended, as \textbf{printbibliography} seems not to behave well in this theme. Will fix later, and any PR/suggestion is welcome. }\\
        This is to cite an article in the format of footnote \footfullcite{vaswani2017attention}
    \end{frame}

    % \section{References}

    % Adding the option 'allowframebreaks' allows the contents of the slide to be expanded in more than one slide.
    % \begin{frame}[allowframebreaks]{References}
    %     \tiny\bibliography{references}
    %     \bibliographystyle{apalike}
    % \end{frame}
    \printbibliography[title={References},heading=bibnumbered]
\end{document}